\input{preamble}

% OK, start here.
%
\begin{document}

\title{Morphismes de schémas}


\maketitle

\phantomsection
\label{section-phantom}

\tableofcontents

\section{Introduction}
\label{section-introduction}

\noindent
Dans ce chapitre, nous introduisons certains types de morphismes de schémas. Une référence de base est \cite{EGA}.




















\section{Immersions fermées}
\label{section-closed-immersions}

\noindent
Dans cette section, nous élucidons certains des résultats obtenus précédemment sur les immersions fermées des schémas. Rappelons qu'un morphisme des schémas $i : Z \to X$ est défini comme une immersion fermée si (a) $i$ induit un homéomorphisme sur un sous-ensemble fermé de $X$, (b) $i^\sharp : \mathcal{O}_X \to i_*\mathcal{O}_Z$ est surjectif, et (c) le noyau de $i^\sharp$ est généré localement par des sections, voir Schémas, définitions \ref{schemes-definition-immersion} et \ref{schemes-definition-closed-immersion-locally-ringed-spaces}. Il s'avère que, étant donné que $Z$ et $X$ sont des schémas, il existe de nombreuses façons différentes de caractériser une immersion fermée.

\begin{lemma}

\label{lemma-closed-immersion}

Laisser $i : Z \to X$ morphisme des régimes. Les conditions suivantes sont équivalentes:

\begin{enumerate}

\item Le morphisme $i$ est une immersion fermée.
\item Pour chaque ouvert affine $\Spec(R) = U \subset X$, il existe un idéal $I \subset R$ de telle manière que
$i^{-1}(U) = \Spec(R/I)$ en tant que régimes $U = \Spec(R)$.
\item Il existe un recouvrement ouvert affine $X = \bigcup_{j \in J} U_j$,
$U_j = \Spec(R_j)$ et pour chaque $j \in J$ il existe un idéal $I_j \subset R_j$ de telle manière que
$i^{-1}(U_j) = \Spec(R_j/I_j)$ en tant que régimes $U_j = \Spec(R_j)$.
\item Le morphisme $i$ induit un homéomorphisme de $Z$ avec un sous-ensemble fermé de $X$ et $i^\sharp : \mathcal{O}_X \to i_*\mathcal{O}_Z$ C'est surjectif.
\item Le morphisme $i$ induit un homéomorphisme de $Z$ avec un sous-ensemble fermé de $X$, l'application $i^\sharp : \mathcal{O}_X \to i_*\mathcal{O}_Z$ est surjectif, et le noyau $\Ker(i^\sharp)\subset \mathcal{O}_X$ est un faisceau quasi-cohérent d'idéaux.
\item Le morphisme $i$ induit un homéomorphisme de $Z$ avec un sous-ensemble fermé de $X$, l'application $i^\sharp : \mathcal{O}_X \to i_*\mathcal{O}_Z$ est surjectif, et le noyau $\Ker(i^\sharp)\subset \mathcal{O}_X$ est un faisceau d'idées qui est généré localement par des sections.

\end{enumerate}

\end{lemma}

\begin{proof}
Condition (6) est notre définition d'une immersion fermée, voir Schémas, définitions \ref{schemes-definition-closed-immersion-locally-ringed-spaces} et \ref{schemes-definition-immersion}. Donc (6) $\Leftrightarrow$ (1). Nous avons (1) $\Rightarrow$ (2) par Schémas, lemme \ref{schemes-lemma-closed-subspace-scheme}. Trivialement (2) $\Rightarrow$ (3).

\medskip\noindent
Supposons (3). Chacun des morphismes $\Spec(R_j/I_j) \to \Spec(R_j)$ est une immersion fermée, voir Schémas, exemple \ref{schemes-example-closed-immersion-affines}. Par conséquent, $i^{-1}(U_j) \to U_j$ est un homéomorphisme sur son image et $i^\sharp|_{U_j}$ est surjectif. Par conséquent, $i$ est un homéomorphisme sur son image et $i^\sharp$ est surjectif puisque cela peut être vérifié localement. Nous concluons que (3) $\Rightarrow$ (4).

\medskip\noindent
L'implication (4) $\Rightarrow$ (1) est Schémas, lemme \ref{schemes-lemma-characterize-closed-immersions}. L'implication (5) $\Rightarrow$ (6) est triviale. Et l'implication (6) $\Rightarrow$ (5) découle de Schémas, lemme \ref{schemes-lemma-closed-subspace-scheme}.
\end{proof}

\begin{lemma}

\label{lemma-closed-immersion-ideals}
Laissez $X$ être un schéma. Laissez $i : Z \to X$ et $i' : Z' \to X$ être des immersions fermées et considérer les faits d'idéaux $\mathcal{I} = \Ker(i^\sharp)$ et $\mathcal{I}' = \Ker((i')^\sharp)$ de $\mathcal{O}_X$.
\begin{enumerate}
\item Le morphisme $i : Z \to X$ facteurs comme $Z \to Z' \to X$ pour certains $a : Z \to Z'$ si et seulement si $\mathcal{I}' \subset \mathcal{I}$. Si cela se produit, alors $a$ est une immersion fermée.
\item Nous avons $Z \cong Z'$ sur $X$ si et seulement si $\mathcal{I} = \mathcal{I}'$.
\end{enumerate}

\end{lemma}

\begin{proof}
Ceci découle de notre discussion sur les sous-espaces fermés dans Schémas, section \ref{schemes-section-closed-immersion} en particulier Schémas, lemmes \ref{schemes-lemma-closed-immersion} et \ref{schemes-lemma-characterize-closed-subspace}. Il suit également d'une manière simple de la caractérisation (3) dans le lemme \ref{lemma-closed-immersion} ci-dessus.
\end{proof}

\begin{lemma}

\label{lemma-closed-immersion-bijection-ideals}
Laissez $X$ être un schéma. Laissez $\mathcal{I} \subset \mathcal{O}_X$ être un faisceau d'idéaux. Les conditions suivantes sont équivalentes:
\begin{enumerate}
\item $\mathcal{I}$ est généré localement par des sections sous forme de modules $\mathcal{O}_X$,
\item $\mathcal{I}$ est quasi-cohérent comme un faisceau de modules $\mathcal{O}_X$, et
\item il existe une immersion fermée $i : Z \to X$ de schémas dont le faisceau d'idéaux $\Ker(i^\sharp)$ correspondant est égal à $\mathcal{I}$.
\end{enumerate}

\end{lemma}

\begin{proof}
L'équivalence de (1) et (2) est immédiate de Schémas, lemme \ref{schemes-lemma-closed-subspace-scheme}. Si (1) tient, alors il y a un sous-espace fermé $i : Z \to X$ avec $\mathcal{I} = \Ker(i^\sharp)$ par Schémas, définition \ref{schemes-definition-closed-subspace} et exemple \ref{schemes-example-closed-subspace}. Par Schémas, lemme \ref{schemes-lemma-closed-subspace-scheme} il s'agit d'une immersion fermée de schémas et (3) tient. Inversement, si (3) tient, alors (2) tient par Schémas, lemme \ref{schemes-lemma-closed-subspace-scheme} (ce qui s'applique parce qu'une immersion fermée de schémas est a fortori une immersion fermée de locaux).
\end{proof}

\begin{lemma}

\label{lemma-base-change-closed-immersion}

Le changement de base d'une immersion fermée est une immersion fermée.

\end{lemma}

\begin{proof}
Voir Schémas, lemme \ref{schemes-lemma-base-change-immersion}.
\end{proof}

\begin{lemma}

\label{lemma-composition-closed-immersion}

Une composition d'immersions fermées est une immersion fermée.

\end{lemma}

\begin{proof}
Nous l'avons vu dans Schémas, lemme \ref{schemes-lemma-composition-immersion}, mais voici une autre preuve. C'est-à-dire, cela découle de la caractérisation (3) des immersions fermées dans lemme \ref{lemma-closed-immersion}. Puisque si $I \subset R$ est un idéal, et $\overline{J} \subset R/I$ est un idéal, alors $\overline{J} = J/I$ pour un $J \subset R$ idéal qui contient $I$ et $(R/I)/\overline{J} = R/J$.
\end{proof}

\begin{lemma}

\label{lemma-closed-immersion-quasi-compact}

Une fermée d'immersion est quasi compacte.

\end{lemma}

\begin{proof}
Ce lemme est un double de Schémas, lemme \ref{schemes-lemma-closed-immersion-quasi-compact}.
\end{proof}

\begin{lemma}

\label{lemma-closed-immersion-separated}

Une fermée d'immersion est séparée.

\end{lemma}

\begin{proof}
Ce lemme est un cas spécial de Schémas, lemme \ref{schemes-lemma-immersions-monomorphisms}.
\end{proof}





\section{Immersions}
\label{section-immersions}

\noindent
Dans cette section, nous recueillons quelques faits sur les immersions.

\begin{lemma}

\label{lemma-immersion-permanence}
Laissez $Z \to Y \to X$ être des morphismes de schémas.
\begin{enumerate}
\item Si $Z \to X$ est une immersion, alors $Z \to Y$ est une immersion.
\item Si $Z \to X$ est une immersion quasi-compacte et $Y \to X$ est quasi-séparée, alors $Z \to Y$ est une immersion quasi-compacte.
\item Si $Z \to X$ est une immersion fermée et $Y \to X$ est séparé, alors $Z \to Y$ est une immersion fermée.
\end{enumerate}

\end{lemma}

\begin{proof}

Dans chaque cas, la démonstration est de contempler le diagramme commutatif
$$
\xymatrix{
Z \ar[r] \ar[rd] & Y \times_X Z \ar[r] \ar[d] & Z \ar[d] \\
& Y \ar[r] & X
}
$$
où la composition des flèches horizontales supérieures est l'identité. Prouvons (1). La première flèche horizontale est une section de
$Y \times_X Z \to Z$, d'où une immersion par Schémas, lemme \ref{schemes-lemma-section-immersion}La flèche $Y \times_X Z \to Y$ est un changement de base $Z \to X$ d'où une immersion (Schémas, lemme \ref{schemes-lemma-base-change-immersion}Enfin, une composition d'immersions est une immersion (Schémas, lemme \ref{schemes-lemma-composition-immersion}Les deux autres résultats sont prouvés exactement de la même manière.

\end{proof}

\begin{lemma}

\label{lemma-factor-quasi-compact-immersion}

Laisser $h : Z \to X$ être une immersion. Si $h$ est quasi-compact, alors nous pouvons factoriser
$h = i \circ j$ avec $j : Z \to \overline{Z}$ une immersion ouverte et $i : \overline{Z} \to X$ une immersion fermée.

\end{lemma}

\begin{proof}

Notez que $h$ est quasi compacte et quasi séparée (voir Schémas, \ref{schemes-lemma-immersions-monomorphisms}). D'où $h_*\mathcal{O}_Z$ est un faisceau quasi-cohérent de $\mathcal{O}_X$-modules de Schémas, lemme \ref{schemes-lemma-push-forward-quasi-coherent}. Cela implique que
$\mathcal{I} = \Ker(\mathcal{O}_X \to h_*\mathcal{O}_Z)$
est un faisceau quasi-cohérent d'idéaux, voir Schémas, section \ref{schemes-section-quasi-coherent}. Laisse $\overline{Z} \subset X$ soit le sous-régime fermé correspondant à $\mathcal{I}$, voir lemme \ref{lemma-closed-immersion-bijection-ideals}Par Schémas, lemme \ref{schemes-lemma-characterize-closed-subspace}
le morphisme $h$ les facteurs tels que
$h = i \circ j$ où $i : \overline{Z} \to X$ est le morphisme de l'inclusion. $j$ est un ouverte d'immersion, choisissez un sous-système ouvert
$U \subset X$ de telle manière que $h$ induit une immersion fermée de $Z$
dans $U$. Alors il est clair que $\mathcal{I}|_U$ est le faisceau d'idéaux correspondant à l'immersion fermée $Z \to U$. Nous voyons donc que $Z = \overline{Z} \cap U$.

\end{proof}

\begin{lemma}

\label{lemma-factor-reduced-immersion}

Laisser $h : Z \to X$ être une immersion. Si $Z$ est réduit, alors nous pouvons factoriser
$h = i \circ j$ avec $j : Z \to \overline{Z}$ une immersion ouverte et $i : \overline{Z} \to X$ une immersion fermée.

\end{lemma}

\begin{proof}

Laisser $\overline{Z} \subset X$ soit la fermeture de $h(Z)$ avec la structure réduite du sous-système fermé induite, voir Schémas, définition \ref{schemes-definition-reduced-induced-scheme}Par Schémas, lemme \ref{schemes-lemma-map-into-reduction}
le morphisme $h$ les facteurs tels que
$h = i \circ j$ avec $i : \overline{Z} \to X$ le morphisme d'inclusion et $j : Z \to \overline{Z}$. De la définition d'une immersion nous voyons qu'il existe un sous-système ouvert $U \subset X$ de telle manière que
$h$ facteurs par une immersion fermée dans $U$. D'où
$\overline{Z} \cap U$ et $h(Z)$ les sous-systèmes fermés réduits de $U$ avec le même ensemble fermé sous-jacent. D'où l'unicité à Schémas, lemme \ref{schemes-lemma-reduced-closed-subscheme}
nous voyons que $h(Z) \cong \overline{Z} \cap U$Donc... $j$ induit un isomorphisme de $Z$ avec $\overline{Z} \cap U$. Autrement dit $j$ est une immersion ouverte.

\end{proof}



\begin{example}

\label{example-thibaut}

Voici un exemple d'immersion qui n'est pas une composition d'une immersion ouverte suivie d'une immersion fermée. $k$ être un champ. $X = \Spec(k[x_1, x_2, x_3, \ldots])$. Laisse $U = \bigcup_{n = 1}^{\infty} D(x_n)$. Ensuite $U \to X$ est une immersion ouverte. Considérez les idéaux
$$
I_n =
(x_1^n, x_2^n, \ldots, x_{n - 1}^n, x_n - 1, x_{n + 1}, x_{n + 2}, \ldots)
\subset
k[x_1, x_2, x_3, \ldots][1/x_n].
$$
Notez que $I_n k[x_1, x_2, x_3, \ldots][1/x_nx_m] = (1)$
pour tout $m \not = n$. D'où les idéaux quasi-cohérents
$\widetilde I_n$ le $D(x_n)$ d'accord sur $D(x_nx_m)$, à savoir
$\widetilde I_n|_{D(x_nx_m)} = \mathcal{O}_{D(x_n x_m)}$ si
$n \not = m$. D'où ces idéaux collent à un faisceau quasi-cohérent d'idéaux
$\mathcal{I} \subset \mathcal{O}_U$. Laisse $Z \subset U$ soit le sous-régime fermé correspondant à
$\mathcal{I}$. Ainsi $Z \to X$ c'est une immersion.

\medskip\noindent
Nous prétendons que nous ne pouvons pas factoriser $Z \to X$ comme $Z \to \overline{Z} \to X$, où $\overline{Z} \to X$ est fermé et $Z \to \overline{Z}$ est ouvert. À savoir, $\overline{Z}$ devrait être défini par un $I \subset k[x_1, x_2, x_3, \ldots]$ idéal tel que $I_n = I k[x_1, x_2, x_3, \ldots][1/x_n]$. Mais le seul élément $f \in k[x_1, x_2, x_3, \ldots]$ qui finit dans tout $I_n$ est $0$! Donc $I$ n'existe pas.
\end{example}

\begin{lemma}

\label{lemma-check-immersion}
Laissez $f : Y \to X$ être un morphisme de schémas. Si pour tous $y \in Y$ il y a un sous-système ouvert $f(y) \in U \subset X$ tel que $f|_{f^{-1}(U)} : f^{-1}(U) \to U$ est une immersion, alors $f$ est une immersion.
\end{lemma}

\begin{proof}

Cette déclaration fait suite à la discussion des sous-régimes clos à la fin de Schémas, section \ref{schemes-section-immersions}
mais nous donnerons aussi une preuve détaillée. $Z \subset X$ soit la fermeture de $f(Y)$. Depuis la prise de fermetures navettes avec restrictions à l'ouverture, nous voyons à partir de l'hypothèse que
$f(Y) \subset Z$ est ouvert. $Z' = Z \setminus f(Y)$ est fermée. $X' = X \setminus Z'$ est un sous-système ouvert de $X$ et
$f$ les facteurs tels que $f : Y \to X'$ suivi de l'inclusion. Si $y \in Y$ et $U \subset X$ est comme dans l'énoncé du lemme, alors $U' = X' \cap U$ est un quartier ouvert de $f'(y)$ de telle manière que $(f')^{-1}(U') \to U'$ est une immersion (lemme \ref{lemma-immersion-permanence}) avec image fermée. C'est donc une immersion fermée, voir Schémas, lemme \ref{schemes-lemma-immersion-when-closed}. Depuis qu'une immersion fermée est locale sur la cible (par exemple par lemme \ref{lemma-closed-immersion}) nous concluons que $f'$ est une immersion fermée comme souhaité.

\end{proof}








\section{Immersions fermées et faisceaux quasi-cohérents}
\label{section-closed-immersions-quasi-coherent}

\noindent
Le lemme suivant fait finalement pour les faisceaux quasi-cohérents sur les schémas ce que Modules, lemme \ref{modules-lemma-i-star-exact} fait pour les gaufrettes abeliennes. Voir aussi la discussion dans Modules, section \ref{modules-section-closed-immersion}.

\begin{lemma}

\label{lemma-i-star-equivalence}
Laissez $i : Z \to X$ être une immersion fermée de schémas. Laissez $\mathcal{I} \subset \mathcal{O}_X$ être le faisceau quasi-cohérent d'idéaux découpant $Z$. Le functeur $$
i_* :
\QCoh(\mathcal{O}_Z)
\longrightarrow
\QCoh(\mathcal{O}_X)
$$ est exact, entièrement fidèle, avec une image essentielle ceux quasi-cohérents $\mathcal{O}_X$-modules $\mathcal{G}$ tels que $\mathcal{I}\mathcal{G} = 0$.
\end{lemma}

\begin{proof}
Une fermée d'immersion est quasi compacte et séparée, voir lemmes \ref{lemma-closed-immersion-quasi-compact} et \ref{lemma-closed-immersion-separated}. Schémas, lemme \ref{schemes-lemma-push-forward-quasi-coherent} s'applique donc et la poussée d'une faisceau quasi-cohérente sur $Z$ est en effet un faisceau quasi-cohérent sur $X$.

\medskip\noindent

Par modules, lemme \ref{modules-lemma-i-star-equivalence}
le functeur $i_*$ est pleinement fidèle.

\medskip\noindent

Maintenant nous passons à la description de l'image essentielle du functeur $i_*$Nous avons $\mathcal{I}(i_*\mathcal{F}) = 0$
pour toute quasi-cohérence $\mathcal{O}_Z$-module, par exemple par Modules, lemme \ref{modules-lemma-i-star-equivalence}. Ensuite, supposez que $\mathcal{G}$
est un quasi-cohérent $\mathcal{O}_X$-module tel que
$\mathcal{I}\mathcal{G} = 0$. Il suffit de démontrer que l'application canonique
$$
\mathcal{G} \longrightarrow i_* i^*\mathcal{G}
$$
est un isomorphisme\footnote{Cela a été prouvé dans une situation plus générale dans la démonstration de Modules, \ref{modules-lemma-i-star-equivalence}.}. Dans le cas des schémas et des modules quasi-cohérents, en travaillant localement sur un ouvert affine de $X$ et en utilisant lemme \ref{lemma-closed-immersion} et Schémas, lemme \ref{schemes-lemma-widetilde-pullback}
il suffit de prouver l'énoncé algébrique suivante:
$R$, un idéal $I$ et une $R$-module $N$ de telle manière que $IN = 0$ l'application canonique
$$
N \longrightarrow N \otimes_R R/I,\quad
n \longmapsto n \otimes 1
$$
est un isomorphisme de $R$-modules. La démonstration de ce fait algèbre facile est omise.

\end{proof}

\noindent
Laissez $i : Z \to X$ être une immersion fermée. En raison de la lemme ci-dessus, nous souvent, par abus de notation, dénotons $\mathcal{F}$ la gerbe $i_*\mathcal{F}$ sur $X$.

\begin{lemma}

\label{lemma-largest-quasi-coherent-subsheaf}

Laisser $X$ Être un plan. $\mathcal{F}$ être un quasi-cohérent
$\mathcal{O}_X$- Module. $\mathcal{G} \subset \mathcal{F}$
être un $\mathcal{O}_X$- sous-module. Il existe un quasi-cohérent unique
$\mathcal{O}_X$-sous-module $\mathcal{G}' \subset \mathcal{G}$
avec la propriété suivante: Pour chaque quasi-cohérent $\mathcal{O}_X$-module
$\mathcal{H}$ l'application
$$
\Hom_{\mathcal{O}_X}(\mathcal{H}, \mathcal{G}')
\longrightarrow
\Hom_{\mathcal{O}_X}(\mathcal{H}, \mathcal{G})
$$
est bijectif. $\mathcal{G}'$ est le plus grand quasi-cohérent
$\mathcal{O}_X$-sous-module $\mathcal{F}$ figurant dans $\mathcal{G}$.

\end{lemma}

\begin{proof}

Laisser $\mathcal{G}_a$, $a \in A$ être l'ensemble de quasi-cohérents
$\mathcal{O}_X$- sous-modules contenus dans $\mathcal{G}$. Puis l'image $\mathcal{G}'$ de
$$
\bigoplus\nolimits_{a \in A} \mathcal{G}_a \longrightarrow \mathcal{F}
$$
est quasi-cohérent comme l'image d'une application de faisceaux quasi-cohérents sur $X$ est quasi-cohérent et puisque une somme directe de faisceaux quasi-cohérents est quasi-cohérent, voir Schémas, section \ref{schemes-section-quasi-coherent}. Le module $\mathcal{G}'$ est contenu dans $\mathcal{G}$. Il s'agit donc du plus grand quasi-cohérent $\mathcal{O}_X$-module contenu dans $\mathcal{G}$.

\medskip\noindent
Pour prouver la formule, laissez $\mathcal{H}$ être un module $\mathcal{O}_X$ quasi-cohérent et laissez $\alpha : \mathcal{H} \to \mathcal{G}$ être une application $\mathcal{O}_X$-module. L'image de la composition $\mathcal{H} \to \mathcal{G} \to \mathcal{F}$ est quasi-cohérent comme l'image d'une application de faisceaux quasi-cohérents. Par conséquent, elle est contenue dans $\mathcal{G}'$. D'où $\alpha$ facteurs à travers $\mathcal{G}'$ comme souhaité.
\end{proof}

\begin{lemma}

\label{lemma-i-upper-shriek}
Laissez $i : Z \to X$ être une immersion fermée de schémas. Il y a un functeur\footnote{Ceci est probablement une notation non standard.} $i^! : \QCoh(\mathcal{O}_X) \to \QCoh(\mathcal{O}_Z)$ qui est un adjoint droit à $i_*$. (Comparer les modules, lemme \ref{modules-lemma-i-star-right-adjoint}.)
\end{lemma}

\begin{proof}
Compte tenu de la quasi-cohérence $\mathcal{O}_X$-module $\mathcal{G}$ nous considérons la sous-chaîne $\mathcal{H}_Z(\mathcal{G})$ de $\mathcal{G}$ des sections locales anéanties par $\mathcal{I}$. Par lemme \ref{lemma-largest-quasi-coherent-subsheaf} il y a une quasi-cohérente quasi-module canonique $\mathcal{O}_X$-submodule $\mathcal{H}_Z(\mathcal{G})'$. Par construction nous avons $$
\Hom_{\mathcal{O}_X}(i_*\mathcal{F}, \mathcal{H}_Z(\mathcal{G})')
=
\Hom_{\mathcal{O}_X}(i_*\mathcal{F}, \mathcal{G})
$$ pour toute $\mathcal{O}_Z$-module $\mathcal{F}$. Par conséquent, nous pouvons définir $i^!\mathcal{G} = i^*(\mathcal{H}_Z(\mathcal{G})')$. Détails omis.
\end{proof}

\noindent
En utilisant le $1$-to-$1$ correspondant entre des faisceaux quasi-cohérents d'idéaux et des sous-systèmes fermés (voir lemme \ref{lemma-closed-immersion-bijection-ideals}), nous pouvons définir des intersections théoriques et des unions de sous-systèmes fermés.

\begin{definition}

\label{definition-scheme-theoretic-intersection-union}
Laissez $X$ être un schéma. Laissez $Z, Y \subset X$ fermer les sous-systèmes correspondant à des faisceaux d'idéaux quasi-cohérents $\mathcal{I}, \mathcal{J} \subset \mathcal{O}_X$. L'intersection théorique {\it du schéma} de $Z$ et $Y$ est le sous-système fermé de $X$ coupé par $\mathcal{I} + \mathcal{J}$. L'union théorique {\it du schéma} de $Z$ et $Y$ est le sous-système fermé de $X$ coupé par $\mathcal{I} \cap \mathcal{J}$.
\end{definition}

\begin{lemma}

\label{lemma-scheme-theoretic-intersection}
Laissez $X$ être un schéma. Laissez $Z, Y \subset X$ être des sous-systèmes fermés. Laissez $Z \cap Y$ être l'intersection schématique de $Z$ et $Y$. Puis $Z \cap Y \to Z$ et $Z \cap Y \to Y$ sont des immersions fermées et $$
\xymatrix{
Z \cap Y \ar[r] \ar[d] & Z \ar[d] \\
Y \ar[r] & X
}
$$ est un diagramme cartésien de schémas, c'est-à-dire $Z \cap Y = Z \times_X Y$.
\end{lemma}

\begin{proof}
Les morphismes $Z \cap Y \to Z$ et $Z \cap Y \to Y$ sont des immersions fermées par lemme \ref{lemma-closed-immersion-ideals}.Soit $U = \Spec(A)$ un ouvert affine de $X$ et que $Z \cap U$ et $Y \cap U$ correspondent aux idéaux $I \subset A$ et $J \subset A$. Puis $Z \cap Y \cap U$ correspond à $I + J \subset A$. Depuis $A/I \otimes_A A/J = A/(I + J)$ nous voyons que le diagramme est cartésien par notre description des produits fibrés des schémas dans Schémas, section \ref{schemes-section-fibre-products}.
\end{proof}

\begin{lemma}

\label{lemma-scheme-theoretic-union}
Laissez $S$ être un schéma. Laissez $X, Y \subset S$ être des sous-systèmes fermés.Soit $X \cup Y$ l'union schématique de $X$ et $Y$.Soit $X \cap Y$ l'intersection schématique de $X$ et $Y$. Puis $X \to X \cup Y$ et $Y \to X \cup Y$ sont des immersions fermées, il y a une suite exacte court $$
0 \to \mathcal{O}_{X \cup Y} \to \mathcal{O}_X \times \mathcal{O}_Y
\to \mathcal{O}_{X \cap Y} \to 0
$$ des modules $\mathcal{O}_S$, et le schéma $$
\xymatrix{
X \cap Y \ar[r] \ar[d] & X \ar[d] \\
Y \ar[r] & X \cup Y
}
$$ est cocartésien dans la catégorie de schémas, c'est-à-dire $X \cup Y = X \amalg_{X \cap Y} Y$.
\end{lemma}

\begin{proof}

Les morphismes $X \to X \cup Y$ et $Y \to X \cup Y$ sont des immersions fermées par lemme \ref{lemma-closed-immersion-ideals}. Dans la suite exacte courté nous utilisons l'équivalence de lemme \ref{lemma-i-star-equivalence} de penser à des modules quasi-cohérents sur des sous-systèmes fermés de $S$ en tant que modules quasi-cohérents sur $S$. Pour la première carte de la séquence nous utilisons les applications canoniques
$\mathcal{O}_{X \cup Y} \to \mathcal{O}_X$ et
$\mathcal{O}_{X \cup Y} \to \mathcal{O}_Y$
et pour la deuxième carte nous utilisons l'application canonique
$\mathcal{O}_X \to \mathcal{O}_{X \cap Y}$ et le négatif de l'application canonique
$\mathcal{O}_Y \to \mathcal{O}_{X \cap Y}$. Ensuite, pour vérifier l'exactitude, nous pouvons travailler affiner localement. $U = \Spec(A)$ être une affinité d'ouverture $S$ et laisser $X \cap U$ et $Y \cap U$
correspondent aux idéaux $I \subset A$ et $J \subset A$. Ensuite
$(X \cup Y) \cap U$ correspond à $I \cap J \subset A$
et $X \cap Y \cap U$ correspond à $I + J \subset A$. Ainsi, l'exactitude découle de l'exactitude de
$$
0 \to A/I \cap J \to A/I \times A/J \to A/(I + J) \to 0
$$
Pour montrer le diagramme est cocartésien, supposez que nous avons un schéma $T$
et les morphismes des régimes $f : X \to T$, $g : Y \to T$ d'accord comme des morphismes $X \cap Y \to T$Objectif : Montrer qu'il existe un morphisme unique $h : X \cup Y \to T$ d ' accord avec $f$ et $g$. Pour construire $h$ nous pouvons travailler localement sur $X \cup Y$, voir Schémas, section \ref{schemes-section-glueing-schemes}Si $s \in X$, $s \not \in Y$Alors $X \to X \cup Y$ est un isomorphisme dans un quartier de $s$ et il est clair comment construire $h$. De même pour $s \in Y$, $s \not \in X$. Pour $s \in X \cap Y$ Nous pouvons choisir un ouvert affine
$V = \Spec(B) \subset T$ contenant $f(s) = g(s)$. Ensuite, nous pouvons choisir un ouvert affine $U = \Spec(A) \subset S$
contenant $s$ de telle manière que $f(X \cap U)$ et $g(Y \cap U)$
sont contenues dans $V$Les morphismes $f|_{X \cap U}$
et $g|_{Y \cap V}$ dans $V$ correspondent à des homomorphismes d'anneaux
$$
B \to A/I
\quad\text{and}\quad
B \to A/J
$$
qui conviennent comme des cartes dans $A/(I + J)$. Par la suite exacte courtoisie exposée ci-dessus, il y a un soulèvement unique de ces homomorphismes d'anneaux à un homomorphisme d'anneaux $B \to A/I \cap J$ comme souhaité.

\end{proof}









\section{Supports des modules}
\label{section-support}

\noindent
Dans cette section, nous recueillons des résultats élémentaires sur les supports de modules quasi-cohérents sur les schémas. Rappelons que le support d'un faisceau de modules a été défini dans Modules, section \ref{modules-section-support}. D'autre part, le support d'un module a été défini dans Algèbre, section \ref{algebra-section-support}.

\begin{lemma}

\label{lemma-support-affine-open}
Laisser $X$ être un schéma. Laisser $\mathcal{F}$ être un faisceau quasi-cohérent sur $X$. Laisser $\Spec(A) = U \subset X$ être une affine ouverte, et définir $M = \Gamma(U, \mathcal{F})$. Laisser $x \in U$, et laisser $\mathfrak p \subset A$ être le premier correspondant. Les suivants sont équivalents
\begin{enumerate}
\item $\mathfrak p$ est dans le support de $M$, et
\item $x$ est dans le support de $\mathcal{F}$.
\end{enumerate}

\end{lemma}

\begin{proof}
Ceci découle de l'égalité $\mathcal{F}_x = M_{\mathfrak p}$, voir Schémas, lemme \ref{schemes-lemma-spec-sheaves} et les définitions.
\end{proof}

\begin{lemma}

\label{lemma-support-closed-specialization}
Laissez $X$ être un schéma. Laissez $\mathcal{F}$ être un faisceau quasi-cohérent sur $X$. Le support de $\mathcal{F}$ est fermé sous spécialisation.
\end{lemma}

\begin{proof}
Si $x' \leadsto x$ est une spécialisation et $\mathcal{F}_x = 0$ alors $\mathcal{F}_{x'}$ est zéro, comme $\mathcal{F}_{x'}$ est une localisation du module $\mathcal{F}_x$. D'où le complément de $\text{Supp}(\mathcal{F})$ est fermé sous générisation.
\end{proof}

\noindent
Pour les modules de type fini quasi-cohérents, le support est fermé, peut être vérifié sur les fibres et se déplace avec le changement de base.

\begin{lemma}

\label{lemma-support-finite-type}
Laissez $\mathcal{F}$ être un module de type fini quasi-cohérent sur un schéma $X$. Ensuite,
\begin{enumerate}
\item Le support de $\mathcal{F}$ est fermé.
\item Pour $x \in X$ nous avons $$
x \in \text{Supp}(\mathcal{F})
\Leftrightarrow
\mathcal{F}_x \not = 0
\Leftrightarrow
\mathcal{F}_x \otimes_{\mathcal{O}_{X, x}} \kappa(x) \not = 0.
$$
\item Pour tout morphisme de schémas $f : Y \to X$ le retrait $f^*\mathcal{F}$ est de type fini aussi bien et nous avons $\text{Supp}(f^*\mathcal{F}) = f^{-1}(\text{Supp}(\mathcal{F}))$.
\end{enumerate}

\end{lemma}

\begin{proof}

La partie 1 est une reformulation des modules, \ref{modules-lemma-support-finite-type-closed}. Vous pouvez aussi combiner lemme \ref{lemma-support-affine-open}, Propriétés, lemme \ref{properties-lemma-finite-type-module}, et Algèbre, lemme \ref{algebra-lemma-support-closed}
La première équivalence en (2) est la définition du support, et la seconde équivalence découle du lemme de Nakayama, voir Algèbre, lemme \ref{algebra-lemma-NAK}. Laisse $f : Y \to X$ morphisme des schémas. Notez que
$f^*\mathcal{F}$ est de type fini par Modules, lemme \ref{modules-lemma-pullback-finite-type}. Pour l'affirmation finale, laisser $y \in Y$ avec image $x \in X$. Rappelez-vous que
$$
(f^*\mathcal{F})_y =
\mathcal{F}_x \otimes_{\mathcal{O}_{X, x}} \mathcal{O}_{Y, y},
$$
Voir Faisceaux, lemme \ref{sheaves-lemma-stalk-pullback-modules}. D'où $(f^*\mathcal{F})_y \otimes \kappa(y)$ est non zéro si et seulement si $\mathcal{F}_x \otimes \kappa(x)$ est non zéro. Par (2) cela implique $x \in \text{Supp}(\mathcal{F})$ si et seulement si $y \in \text{Supp}(f^*\mathcal{F})$, qui est le contenu de l'affirmation (3).

\end{proof}

\begin{lemma}

\label{lemma-scheme-theoretic-support}

Laisser $\mathcal{F}$ être un module de type fini quasi-cohérent sur un schéma $X$. Il existe un sous-système fermé le plus petit
$i : Z \to X$ de telle sorte qu'il existe une quasi-cohérence
$\mathcal{O}_Z$-module $\mathcal{G}$ avec
$i_*\mathcal{G} \cong \mathcal{F}$. En outre:

\begin{enumerate}

\item Si $\Spec(A) \subset X$ est toute ouvert affine, et
$\mathcal{F}|_{\Spec(A)} = \widetilde{M}$ alors
$Z \cap \Spec(A) = \Spec(A/I)$ où $I = \text{Ann}_A(M)$.
\item Le faisceau quasi-cohérent $\mathcal{G}$ est unique jusqu'à un isomorphisme unique.
\item Le faisceau quasi-cohérent $\mathcal{G}$ est de type fini.
\item L'appui de $\mathcal{G}$ et des $\mathcal{F}$ est $Z$.

\end{enumerate}

\end{lemma}

\begin{proof}

Supposons que $i' : Z' \to X$ est un sous-système fermé qui satisfait à la description sur les affines ouvertes du lemme. \ref{lemma-i-star-equivalence}
nous voyons que $\mathcal{F} \cong i'_*\mathcal{G}'$ pour un faisceau quasi-cohérent unique $\mathcal{G}'$ le $Z'$. En outre, il est clair que $Z'$ est le plus petit sous-système fermé avec cette propriété (par le même lemme). \ref{properties-lemma-finite-type-module}
et d ' Algèbre, lemme \ref{algebra-lemma-finite-over-subring}
il en résulte que $\mathcal{G}'$ est de type fini.
$\text{Supp}(\mathcal{G}') = Z$ par Algèbre, lemme \ref{algebra-lemma-support-closed}. Donc, pour prouver le lemme il suffit de montrer que la caractérisation dans (1) ne définit en fait un sous-système fermé. Et, pour ce faire, il suffit de prouver que la règle donnée produit un faisceau quasi-cohérent d'idéaux, voir lemme \ref{lemma-closed-immersion-bijection-ideals}. Cela revient au fait de l'algèbre suivant: Si $A$ est une bague, $f \in A$, $M$ est un fini $A$- Module, alors
$\text{Ann}_A(M)_f = \text{Ann}_{A_f}(M_f)$Nous omettons la démonstration.

\end{proof}

\begin{definition}

\label{definition-scheme-theoretic-support}
Laissez $X$ être un schéma. Laissez $\mathcal{F}$ être un quasi-cohérent $\mathcal{O}_X$-module de type fini. Le {\it support schématique de $\mathcal{F}$} est le sous-système fermé $Z \subset X$ construit en lemme \ref{lemma-scheme-theoretic-support}.
\end{definition}

\noindent
Dans cette situation, nous considérons souvent $\mathcal{F}$ comme une fourchette quasi-cohérente de type fini sur $Z$ (par l'équivalence des catégories de lemme \ref{lemma-i-star-equivalence}).













\section{Image schématique}
\label{section-scheme-theoretic-image}

\noindent
Attention : Une partie du matériel de cette section est ultra-générale et se comporte différemment de ce que vous pourriez attendre.

\begin{lemma}

\label{lemma-scheme-theoretic-image}
Laissez $f : X \to Y$ être un morphisme de schémas. Il existe un sous-système fermé $Z \subset Y$ tel que $f$ facteurs par $Z$ et tel que pour tout autre sous-système fermé $Z' \subset Y$ tel que $f$ facteurs par $Z'$ nous avons $Z \subset Z'$.
\end{lemma}

\begin{proof}
Laissez $\mathcal{I} = \Ker(\mathcal{O}_Y \to f_*\mathcal{O}_X)$. Si $\mathcal{I}$ est quasi-cohérent, alors nous prenons juste $Z$ pour être le sous-système fermé déterminé par $\mathcal{I}$, voir lemme \ref{lemma-closed-immersion-bijection-ideals}. Cela fonctionne par lemme \ref{lemma-closed-immersion-ideals}. En général, le même lemme nous oblige à montrer qu'il existe un faisceau quasi-cohérent plus grand des idéaux $\mathcal{I}'$ contenus dans $\mathcal{I}$. Ceci suit de lemme \ref{lemma-largest-quasi-coherent-subsheaf}.
\end{proof}

\begin{definition}

\label{definition-scheme-theoretic-image}
Laissez $f : X \to Y$ être un morphisme de schémas. L'image {\it schématique} de $f$ est le plus petit sous-système fermé $Z \subset Y$ par lequel $f$ facteurs, voir lemme \ref{lemma-scheme-theoretic-image} ci-dessus.
\end{definition}

\noindent

Pour un morphisme $f : X \to Y$ de schémas avec image schématique $Z$
nous dénotons souvent $f : X \to Z$ de la factorisation $f$
à travers son image schématique. $f$ n'est pas quasi-compact, alors (en général)

\begin{enumerate}

\item l'inclusion théorique de l'ensemble $\overline{f(X)} \subset Z$
n'est pas une égalité, c'est-à-dire, $f(X) \subset Z$ n'est pas un sous-ensemble dense, et
\item la construction de l'image schématique ne se déplace pas avec restriction pour ouvrir des sous-systèmes à $Y$.

\end{enumerate}

Dans l'exemple, section \ref{examples-section-scheme-theoretic-image}
le lecteur trouve un exemple pour les deux phénomènes. Ces phénomènes peuvent survenir même pour les immersions, voir exemple, section \ref{examples-section-strange-immersion}. Cependant, le prochain lemme montre que les deux catastrophes sont évitées lorsque le morphisme est quasi-compact.

\begin{lemma}

\label{lemma-quasi-compact-scheme-theoretic-image}

Laisser $f : X \to Y$ être un morphisme de schémas. $Z \subset Y$ être l'image schématique de $f$Si $f$ est quasi-compact alors

\begin{enumerate}

\item le faisceau d'idées
$\mathcal{I} = \Ker(\mathcal{O}_Y \to f_*\mathcal{O}_X)$
est quasi-cohérent,
\item l'image schématique $Z$ est le sous-système fermé déterminé par: $\mathcal{I}$,
\item pour toute ouverture $U \subset Y$ l'image schématique de
$f|_{f^{-1}(U)} : f^{-1}(U) \to U$ est égal à $Z \cap U$,
\item l'image $f(X) \subset Z$ est un sous-ensemble dense de $Z$, en d'autres termes le morphisme $X \to Z$ est dominante (voir définition) \ref{definition-dominant}).

\end{enumerate}

\end{lemma}

\begin{proof}

La partie 4) suit de la partie 3). Pour montrer (3) il suffit de prouver (1) depuis la formation de $\mathcal{I}$ commutations avec restriction aux sous-régimes ouverts de $Y$. Et si (1) tient alors dans la démonstration de lemme \ref{lemma-scheme-theoretic-image}
nous avons montré (2). Il suffit donc de prouver que $\mathcal{I}$ Comme la propriété d'être quasi-cohérent est locale, nous pouvons supposer $Y$ C'est affine. $f$ est quasi-compact, nous pouvons trouver un recouvrement fini ouvert affine
$X = \bigcup_{i = 1, \ldots, n} U_i$. Dénotation $f'$ la composition
$$
X' = \coprod U_i \longrightarrow X \longrightarrow Y.
$$
Alors $f_*\mathcal{O}_X$ est un sous-chapeau de $f'_*\mathcal{O}_{X'}$, et donc $\mathcal{I} = \Ker(\mathcal{O}_Y \to f'_*\mathcal{O}_{X'})$Par Schémas, lemme \ref{schemes-lemma-push-forward-quasi-coherent}
les gerbes $f'_*\mathcal{O}_{X'}$ est quasi-cohérent sur $Y$. C'est pourquoi nous avons gagné.

\end{proof}

\begin{example}

\label{example-scheme-theoretic-image}
Si $A \to B$ est un homomorphisme d'anneaux avec le noyau $I$, alors l'image schématique de $\Spec(B) \to \Spec(A)$ est le sous-système fermé $\Spec(A/I)$ de $\Spec(A)$. Ceci suit de lemme \ref{lemma-quasi-compact-scheme-theoretic-image}.
\end{example}

\noindent
Si le morphisme est quasi-compact, alors l'image schématique ajoute seulement des points qui sont des spécialisations de points dans l'image.

\begin{lemma}

\label{lemma-reach-points-scheme-theoretic-image}

Laisser $f : X \to Y$ être un morphisme quasi-compact. $Z$ être l'image schématique de $f$. Laisse $z \in Z$\footnote{Par lemme \ref{lemma-quasi-compact-scheme-theoretic-image} en théorie
$Z$ est d'accord avec la fermeture $f(X)$ en $Y$.}. Il existe un anneau de valorisation $A$ avec champ de fractions $K$ et un diagramme commutatif
$$
\xymatrix{
\Spec(K) \ar[rr] \ar[d] & & X \ar[d] \ar[ld] \\
\Spec(A) \ar[r] & Z \ar[r] & Y
}
$$
tel que le point fermé de $\Spec(A)$ cartes à $z$. En particulier tout point de $Z$ est la spécialisation d'un point de $f(X)$.

\end{lemma}

\begin{proof}

Laisser $z \in \Spec(R) = V \subset Y$ être un quartier affiné ouvert de $z$. Par lemme \ref{lemma-quasi-compact-scheme-theoretic-image}
l'intersection $Z \cap V$ est l'image schématique de
$f^{-1}(V) \to V$. Par conséquent, nous pouvons remplacer $Y$ par $V$
et supposer $Y = \Spec(R)$ Dans ce cas, $X$ est quasi-compact comme $f$ C'est quasi-compact. $X = U_1 \cup \ldots \cup U_n$
est un recouvrement fini ouvert affine. $U_i = \Spec(A_i)$. Laisse $I = \Ker(R \to A_1 \times \ldots \times A_n)$. Par lemme \ref{lemma-quasi-compact-scheme-theoretic-image}
Encore une fois nous voyons que $Z$ correspond au sous-système fermé
$\Spec(R/I)$ de $Y$Si $\mathfrak p \subset R$ est le premier correspondant à $z$, alors nous voyons que
$I_{\mathfrak p} \subset R_{\mathfrak p}$ n'est donc pas une égalité. (comme la localisation est exacte, voir Algèbre, proposition \ref{algebra-proposition-localization-exact}) nous voyons que
$R_{\mathfrak p} \to
(A_1)_{\mathfrak p} \times \ldots \times (A_n)_{\mathfrak p}$
n'est pas zéro. D'où un des anneaux $(A_i)_{\mathfrak p}$ n'est pas zéro. Il existe donc un $i$ et une première $\mathfrak q_i \subset A_i$
couché au-dessus d'un $\mathfrak p_i \subset \mathfrak p$. Par Algèbre, lemme \ref{algebra-lemma-dominate} nous pouvons choisir un anneau de valorisation
$A \subset K = \kappa(\mathfrak q_i)$ dominant l'Anneau local
$R_{\mathfrak p}/\mathfrak p_iR_{\mathfrak p} \subset \kappa(\mathfrak q_i)$. Ceci donne le diagramme souhaité. Certains détails omis.

\end{proof}

\begin{lemma}

\label{lemma-factor-factor}
Laissez $$
\xymatrix{
X_1 \ar[d] \ar[r]_{f_1} & Y_1 \ar[d] \\
X_2 \ar[r]^{f_2} & Y_2
}
$$ être un diagramme commutatif des schémas. Laissez $Z_i \subset Y_i$, $i = 1, 2$ être l'image schématique de $f_i$. Puis le morphisme $Y_1 \to Y_2$ induit un morphisme $Z_1 \to Z_2$ et un diagramme commutatif $$
\xymatrix{
X_1 \ar[r] \ar[d] & Z_1 \ar[d] \ar[r] & Y_1 \ar[d] \\
X_2 \ar[r] & Z_2 \ar[r] & Y_2
}
$$
\end{lemma}

\begin{proof}
L'image inverse schématique de $Z_2$ dans $Y_1$ est un sous-système fermé de $Y_1$ par lequel $f_1$ facteurs. Par conséquent $Z_1$ est contenu dans ceci. Cela prouve la lemme.
\end{proof}

\begin{lemma}

\label{lemma-scheme-theoretic-image-reduced}
Laissez $f : X \to Y$ être un morphisme de schémas. Si $X$ est réduit, alors l'image schématique de $f$ est la structure de schéma induite réduite sur $\overline{f(X)}$.
\end{lemma}

\begin{proof}
Ceci est vrai parce que la structure de schéma induite réduite sur $\overline{f(X)}$ est clairement le plus petit sous-système fermé de $Y$ par lequel $f$ facteurs, voir Schémas, lemme \ref{schemes-lemma-map-into-reduction}.
\end{proof}

\begin{lemma}

\label{lemma-scheme-theoretic-image-of-partial-section}
Laissez $f : X \to Y$ être un morphisme séparé de schémas. Laissez $V \subset Y$ être un rétrocompact ouvert. Laissez $s : V \to X$ être un morphisme tel que $f \circ s = \text{id}_V$. Laissez $Y'$ être l'image schématique de $s$. Alors $Y' \to Y$ est un isomorphisme sur $V$.
\end{lemma}

\begin{proof}

L'hypothèse selon laquelle $V$ est rétrocompact dans $Y$
(Topologie, définition) \ref{topology-definition-quasi-compact}) signifie que $V \to Y$ est un morphisme quasi-compact. Par Schémas, \ref{schemes-lemma-quasi-compact-permanence}
le morphisme $s : V \to X$ est quasi-compact. D'où la construction de l'image schématique $Y'$
de $s$ commutations avec restriction à ouvrir par lemme \ref{lemma-quasi-compact-scheme-theoretic-image}. En particulier, nous voyons que $Y' \cap f^{-1}(V)$ est l'image schématique d'une section du morphisme séparé $f^{-1}(V) \to V$. Comme une section d'un morphisme séparé est une immersion fermée (Schémas, lemme \ref{schemes-lemma-section-immersion}), nous concluons que
$Y' \cap f^{-1}(V) \to V$ est un isomorphisme tel que souhaité.

\end{proof}







\section{Adhérence schématique et densité}
\label{section-scheme-theoretic-closure}

\noindent
Nous prenons la définition suivante de \cite[IV, définition 11.10.2]{EGA}.

\begin{definition}

\label{definition-scheme-theoretically-dense}

Laisser $X$ Être un plan. $U \subset X$ être un sous-système ouvert.

\begin{enumerate}

\item L'image schématique du morphisme $U \to X$
est appelé {\it Adhérence schématique de $U$ en $X$}.
\item Nous disons $U$ est {\it schématiquement dense $X$} si pour chaque ouverture $V \subset X$ l'adhérence schématique de $U \cap V$ en $V$ est égal à $V$.

\end{enumerate}

\end{definition}

\noindent

Avec cette définition, il est {\bf pas} dans le cas où $U$ est schématiquement dense en $X$ si et seulement si l'adhérence schématique de $U$ est $X$, voir exemple \ref{example-scheme-theretically-dense-not-dense}. C'est un peu inélégant; mais voir lemmes \ref{lemma-scheme-theoretically-dense-quasi-compact} et
\ref{lemma-reduced-scheme-theoretically-dense}
ci-dessous. D'autre part, avec cette définition $U$ est schématiquement dense en $X$ si et seulement si pour chaque $V \subset X$ ouvrir l'homomorphisme d'anneaux
$\mathcal{O}_X(V) \to \mathcal{O}_X(U \cap V)$ est injectable, voir lemme \ref{lemma-characterize-scheme-theoretically-dense}
ci-dessous. En particulier, nous voyons que le schématiquement dense implique dense qui est agréable.

\begin{example}

\label{example-scheme-theretically-dense-not-dense}

Voici un exemple où l'adhérence schématique $X$ ne signifie pas dense pour les espaces topologiques sous-jacents. $k$ être un champ. $A = k[x, z_1, z_2, \ldots]/(x^n z_n)$
Définir $I = (z_1, z_2, \ldots) \subset A$. Considérer le régime affiné $X = \Spec(A)$ et le sous-système ouvert $U = X \setminus V(I)$Depuis $A \to \prod_n A_{z_n}$ nous voyons que l'adhérence schématique de $U$ est $X$. Considérez le morphisme
$X \to \Spec(k[x])$. Ce morphisme est surjectif (ensemble $z_n = 0$ Mais la restriction de ce morphisme à $U$ n'est pas surjectif parce qu'il cartographie au point $x = 0$. D'où $U$ ne peut être topologiquement dense dans $X$.

\end{example}

\begin{lemma}

\label{lemma-scheme-theoretically-dense-quasi-compact}
Laissez $X$ être un schéma. Laissez $U \subset X$ être un sous-système ouvert. Si le morphisme d'inclusion $U \to X$ est quasi-compact, alors $U$ est dense schématiquement dans $X$ si et seulement si la adhérence schématique de $U$ dans $X$ est $X$.
\end{lemma}

\begin{proof}
Suit de lemme \ref{lemma-quasi-compact-scheme-theoretic-image} partie (3).
\end{proof}

\begin{example}

\label{example-scheme-theoretic-closure}

Laisser $A$ être une bague et $X = \Spec(A)$. Laisse $f_1, \ldots, f_n \in A$ et laisser $U = D(f_1) \cup \ldots \cup D(f_n)$. Laisse $I = \Ker(A \to \prod A_{f_i})$. Puis l'adhérence schématique de $U$ en $X$
est le sous-système fermé $\Spec(A/I)$ de $X$. Notez que $U \to X$ est quasi-compact. \ref{lemma-scheme-theoretically-dense-quasi-compact}
nous voyons $U$ est schématiquement dense en $X$ si et seulement si $I = 0$.

\end{example}

\begin{lemma}

\label{lemma-characterize-scheme-theoretically-dense}
Laissez $j : U \to X$ être une immersion ouverte de schémas. Alors $U$ est dense schématiquement dans $X$ si et seulement si $\mathcal{O}_X \to j_*\mathcal{O}_U$ est injectable.
\end{lemma}

\begin{proof}

Si $\mathcal{O}_X \to j_*\mathcal{O}_U$ est injectable, alors il en est de même lorsqu'il est limité à tout $V$ de $X$. D'où l'adhérence schématique $U \cap V$ en $V$
est égal à $V$, voir la démonstration de lemme \ref{lemma-scheme-theoretic-image}. Inversement, supposer que l'adhérence schématique de $U \cap V$ est égal à $V$ pour tous ouverts $V$. Supposons que $\mathcal{O}_X \to j_*\mathcal{O}_U$ Alors on peut trouver un ouvert affine, disons $\Spec(A) = V \subset X$
et un élément non zéro $f \in A$ de telle manière que $f$ Cartes à zéro dans
$\Gamma(V \cap U, \mathcal{O}_X)$. Dans ce cas, l'adhérence schématique de $V \cap U$ en $V$ est clairement contenu dans $\Spec(A/(f))$
une contradiction.

\end{proof}

\begin{lemma}

\label{lemma-intersection-scheme-theoretically-dense}
Laissez $X$ être un schéma. Si $U$, $V$ sont des sous-systèmes ouverts denses de $X$, alors $U \cap V$.
\end{lemma}

\begin{proof}
Laissez $W \subset X$ être ouvert. Considérez l'application $\mathcal{O}_X(W) \to \mathcal{O}_X(W \cap V)
\to \mathcal{O}_X(W \cap V \cap U)$. Par lemme \ref{lemma-characterize-scheme-theoretically-dense} les deux cartes sont injectables. Par conséquent, le composite est injectable. Par lemme \ref{lemma-characterize-scheme-theoretically-dense} $U \cap V$ est dense en $X$.
\end{proof}

\begin{lemma}

\label{lemma-quasi-compact-immersion}
Laissez $h : Z \to X$ être une immersion. Supposons que $h$ soit quasi-compact ou $Z$ est réduit.Soit $\overline{Z} \subset X$ l'image schématique de $h$. Ensuite, le morphisme $Z \to \overline{Z}$ est une ouverture d'immersion qui identifie $Z$ avec un sous-système ouvert dense schématiquement de $\overline{Z}$. De plus, $Z$ est topologiquement dense dans $\overline{Z}$.
\end{lemma}

\begin{proof}

Par moi \ref{lemma-factor-quasi-compact-immersion} ou laisser-moi \ref{lemma-factor-reduced-immersion} nous pouvons factoriser
$Z \to X$ comme $Z \to \overline{Z}_1 \to X$ avec $Z \to \overline{Z}_1$
ouvert et $\overline{Z}_1 \to X$ fermé. D'un autre côté, laisser
$Z \to \overline{Z} \subset X$ être l'image schématique de
$Z \to X$. Nous concluons que $\overline{Z} \subset \overline{Z}_1$Depuis $Z$ est un sous-système ouvert de $\overline{Z}_1$ il en résulte que $Z$ est un sous-système ouvert de $\overline{Z}$ Dans le cas où $Z$ est réduit nous savons que $Z \subset \overline{Z}_1$
est topologiquement dense par la construction de $\overline{Z}_1$ dans la démonstration de lemme \ref{lemma-factor-reduced-immersion}. D'où $\overline{Z}_1$ et $\overline{Z}$ ont les mêmes espaces topologiques sous-jacents. $\overline{Z} \subset \overline{Z}_1$
est une immersion fermée dans un schéma réduit qui induit un bijection sur les espaces topologiques sous-jacents, et donc c'est un isomorphisme. $Z \to X$ est quasi-compact nous argumentons comme suit: L'affirmation que $Z$ est schématiquement dense en
$\overline{Z}$ de lemme \ref{lemma-quasi-compact-scheme-theoretic-image} Partie 3). La dernière affirmation découle de \ref{lemma-quasi-compact-scheme-theoretic-image} partie 4).

\end{proof}

\begin{lemma}

\label{lemma-reduced-scheme-theoretically-dense}
Laissez $X$ être un schéma réduit et laissez $U \subset X$ être un sous-système ouvert. Ensuite, les suivants sont équivalents
\begin{enumerate}
\item $U$ est topologiquement dense dans $X$,
\item l'adhérence schématique de $U$ dans $X$ est $X$, et
\item $U$ est dense dans $X$.
\end{enumerate}

\end{lemma}

\begin{proof}

Cela s'explique par le fait qu'il s'agit d'une \ref{lemma-quasi-compact-immersion}
et le fait qu'un sous-système fermé $Z$ de $X$ dont l'espace topologique sous-jacent est égal à $X$ doit être égal à $X$
dans le cadre d'un programme.

\end{proof}

\begin{lemma}

\label{lemma-reduced-subscheme-closure}
Laissez $X$ être un schéma et laissez $U \subset X$ être un sous-système ouvert réduit. Ensuite, les éléments suivants sont équivalents
\begin{enumerate}
\item l'adhérence schématique de $U$ dans $X$ est $X$, et
\item $U$ est dense en $X$.
\end{enumerate}
Si cela tient alors $X$ est un schéma réduit.
\end{lemma}

\begin{proof}
Ceci découle de \ref{lemma-quasi-compact-immersion} et du fait que l'adhérence schématique de $U$ dans $X$ est réduite par \ref{lemma-scheme-theoretic-image-reduced}.
\end{proof}

\begin{lemma}

\label{lemma-equality-of-morphisms}
Laissez $S$ être un schéma. Laissez $X$, $Y$ être des schémas sur $S$. Laissez $f, g : X \to Y$ être des morphismes de schémas sur $S$. Laissez $U \subset X$ être un sous-système ouvert tel que $f|_U = g|_U$. Si l'adhérence schématique de $U$ dans $X$ est $X$ et $Y \to S$ est séparé, alors $f = g$.
\end{lemma}

\begin{proof}
Suit des définitions et Schémas, lemme \ref{schemes-lemma-where-are-they-equal}.
\end{proof}




\section{Morphismes dominants}
\label{section-dominant}

\noindent
La définition d'un morphisme des schémas étant dominant est un peu différente de ce que vous pourriez attendre si vous êtes habitué à la notion d'un morphisme dominant des variétés.

\begin{definition}

\label{definition-dominant}
Un morphisme $f : X \to S$ des schémas est appelé {\it dominant} si l'image de $f$ est un sous-ensemble dense de $S$.
\end{definition}

\noindent
Par exemple, si $k$ est un champ infini et $\lambda_1, \lambda_2, \ldots$ est une collection comptable d'éléments distincts de $k$, alors le morphisme $$
\coprod\nolimits_{i = 1, 2, \ldots } \Spec(k)
\longrightarrow
\Spec(k[x])
$$ avec $i$th mapping de facteur au point $x = \lambda_i$ est dominant.

\begin{lemma}

\label{lemma-generic-points-in-image-dominant}
Laissez $f : X \to S$ être un morphisme de schémas. Si chaque point général de chaque composite irréductible de $S$ est à l'image de $f$, alors $f$ est dominant.
\end{lemma}

\begin{proof}
C'est un fait topologique qui découle directement du fait que l'espace topologique sous-jacent à un schéma est sobre, voir Schémas, lemme \ref{schemes-lemma-scheme-sober}, et que chaque point de $S$ est contenu dans un composante irréductible de $S$, voir Topologie, lemme \ref{topology-lemma-irreducible}.
\end{proof}

\noindent
L'attente que les morphismes sont dominants seulement si les points génériques de la cible sont dans l'image tient si le morphisme est quasi-compact.

\begin{lemma}

\label{lemma-quasi-compact-dominant}

\begin{slogan}
Morphismes dont l'image contient les points génériques sont dominantes
\end{slogan}
Laissez $f : X \to S$ être un morphisme quasi-compact de schémas. Alors $f$ est dominante si et seulement si pour chaque composant irréductible $Z \subset S$ le point général de $Z$ est dans l'image de $f$.
\end{lemma}

\begin{proof}

Laisser $V \subset S$ être un ouvert affine. $f$ est quasi-compact nous pouvons choisir finiment de nombreux ouverts affines $U_i \subset f^{-1}(V)$, $i = 1, \ldots, n$ couvrent
$f^{-1}(V)$. Considérez le morphisme des affines
$$
f' :
\coprod\nolimits_{i = 1, \ldots, n} U_i
\longrightarrow
V.
$$
Une union disjointe des affines est affine, voir Schémas, lemme \ref{schemes-lemma-disjoint-union-affines}. Points génériques de composantes irréductibles de $V$
sont exactement les points génériques des composantes irréductibles de
$S$ qui se rencontrent $V$. Aussi, $f$ est dominante si et seulement si $f'$ est dominant peu importe les choix de $V, n, U_i$ Nous avons donc réduit le lemme au cas d'un morphisme des schémas affinés. \ref{algebra-lemma-image-dense-generic-points}.

\end{proof}

\begin{lemma}

\label{lemma-base-change-dominant}
Laissez $f : X \to S$ être un morphisme quasi-compact dominant des schémas. Laissez $g : S' \to S$ être un morphisme des schémas, et laissez $f' : X' \to S'$ être le changement de base de $f$ par $g$. Si les générisations remontent le long de $g$, alors $f'$ est dominante.
\end{lemma}

\noindent
Nous verrons que les générisations remontent le long de $g$ si $g$ est ouvert (lemme \ref{lemma-open-generizing}) ou plat (lemme \ref{lemma-generalizations-lift-flat}).

\begin{proof}

Observer que $f'$ est quasi-compact par Schémas, lemme
\ref{schemes-lemma-quasi-compact-preserved-base-change}. Laisse $\eta' \in S'$ être le point général d'un composante irréductible de $S'$. Si les générisations se lèvent le long $g$Alors $\eta = g(\eta')$ est le point général d'un composante irréductible de $S$. Par lemme \ref{lemma-quasi-compact-dominant}
nous voyons que $\eta$ est à l'image de $f$. D'où $\eta'$ est à l'image de $f'$ par Schémas, lemme \ref{schemes-lemma-points-fibre-product}Il s'ensuit que $f'$ est dominant par me \ref{lemma-quasi-compact-dominant}.

\end{proof}

\noindent
Une partie de cela s'applique aux morphismes arbitraires (pas nécessairement quasi-compacts) dominants.

\begin{lemma}

\label{lemma-open-base-change-dominant}
Laisser $f : X \to S$ être un morphisme dominant des schémas. Laisser $g : S' \to S$ être un morphisme ouvert des schémas, et laisser $f' : X' \to S'$ être le changement de base de $f$ par $g$. Alors $f'$ est dominant.
\end{lemma}

\begin{proof}

Si $f'$ n'est pas dominant, alors $f'(X')$ est contenu dans un sous-ensemble fermé $B$ non égal à $S'$. Ensuite $S'-B$ est ouvert en $S'$ et non vide. Depuis $g$ est ouvert, $g(S'-B)$ est non vide et ouvert dans $S$Depuis $f(X)$
est dense en $S$, il y a un point $x\in X$ avec $f(x) \in g(S'-B)$. Laisse $W = \Spec(\kappa (f(x)))$, et écrire $X_W$, $S_W$, $S'_W$ pour les retraits de ces régimes $S$ à $W$. Ensuite $X_W$ et $(S'-B)_W$ ne sont pas vides, et donc $X_W \times_W (S'-B)_W$ est non vide, comme $W$ est le spectre d'un champ. Cela contredit que $f'(X \times_S S')$ est contenu dans $B \subset S'$. Donc en fait $f'$ est dominant.

\end{proof}

\begin{lemma}

\label{lemma-quasi-compact-generic-point-not-in-image}
Laissez $f : X \to S$ être un morphisme quasi-compact de schémas. Laissez $\eta \in S$ être un point général d'un composant irréductible de $S$. Si $\eta \not \in f(X)$ alors il existe un quartier ouvert $V \subset S$ de $\eta$ tel que $f^{-1}(V) = \emptyset$.
\end{lemma}

\begin{proof}

Laisser $Z \subset S$ être l'image schématique de $f$. Nous devons montrer que $\eta \not \in Z$. Cela résulte de l'expérience acquise \ref{lemma-reach-points-scheme-theoretic-image}
mais peut aussi être vu comme suit. \ref{lemma-quasi-compact-scheme-theoretic-image}
le morphisme $X \to Z$ est dominant, qui par lemme \ref{lemma-quasi-compact-dominant}
désigne tous les points génériques de tous les composantes irréductibles de $Z$
sont à l'image de $X \to Z$. Par hypothèse, nous voyons que
$\eta \not \in Z$ depuis $\eta$ serait le point général de certains composantes irréductibles de $Z$ s'il était dans $Z$.

\end{proof}

\noindent
Il y a un autre cas où dominante est la même que d'avoir tous les points génériques de composantes irréductibles dans l'image.

\begin{lemma}

\label{lemma-dominant-finite-number-irreducible-components}

Laisser $f : X \to S$ être un morphisme de schémas. Supposons que $X$ a finiment beaucoup de composantes irréductibles. $f$ est dominante (si et) seulement si pour chaque Composante irréductible $Z \subset S$ le point général de $Z$ est à l'image de $f$. Si oui, alors $S$ a finiment beaucoup de composantes irréductibles ainsi.

\end{lemma}

\begin{proof}
Supposons que $f$ est dominante. Écrivons $X = Z_1 \cup Z_2 \cup \ldots \cup Z_n$ est la décomposition de $X$ en composantes irréductibles. Que $\xi_i \in Z_i$ soit son point général, donc $Z_i = \overline{\{\xi_i\}}$. Notez que $f(Z_i)$ est un sous-ensemble irréductible de $S$. Par conséquent $$
S = \overline{f(X)} = \bigcup \overline{f(Z_i)} =
\bigcup \overline{\{f(\xi_i)\}}
$$ est une union finie de sous-ensembles irréductibles dont les points génériques sont à l'image de $f$. La lemme suit.
\end{proof}

\begin{lemma}

\label{lemma-dominant-between-integral}

Laisser $f : X \to Y$ être un morphisme de schémas intégrés.

\begin{enumerate}

\item  $f$ est dominant,
\item  $f$ cartographie le point général de $X$ au point général de $Y$,
\item pour certains ouverts non vides affines $U \subset X$ et $V \subset Y$
avec $f(U) \subset V$ l'homomorphisme d'anneaux $\mathcal{O}_Y(V) \to \mathcal{O}_X(U)$
est injectable,
\item pour tous les ouverts non vides affinés $U \subset X$ et $V \subset Y$
avec $f(U) \subset V$ l'homomorphisme d'anneaux $\mathcal{O}_Y(V) \to \mathcal{O}_X(U)$
est injectable,
\item pour certains $x \in X$ avec image $y = f(x) \in Y$ l'application locale de l'anneau $\mathcal{O}_{Y, y} \to \mathcal{O}_{X, x}$ est injectable, et
\item pour tous $x \in X$ avec image $y = f(x) \in Y$ l'application locale de l'anneau $\mathcal{O}_{Y, y} \to \mathcal{O}_{X, x}$ est injectable.

\end{enumerate}

\end{lemma}

\begin{proof}

L'équivalence des points 1) et 2) découle de l'équivalence des points 2) et 2) de l'équivalence des points 2) et 2) de l'équivalence des points 2) et 2) de l'équivalence des points 2) et 2) de l'équivalence des points 2) et 2) de l'équivalence des points 2) et 2) de l'équivalence des points 2) et 2) de l'équivalence des points 2) à 2) de l'équivalence des points 2) et 2) de l'équivalence des points 2) à 2) de l'équivalence des points 2) et 2) de l'équivalence des points 2) à 2) de l'équivalence des points 2) et 2) de l'équivalence des points 2) à 2) de l'équivalence des points 2) et 2) de l'équivalence des points 2) de l'équivalence des points 2) à 2) de l'équivalence des points 2) et 2) de l'équivalence des points 2) de l'équivalence entre les points 2) de l'équivalence des points 2) de l'équivalence des points 2) de l'équivalence des points 2) de l'équi \ref{lemma-dominant-finite-number-irreducible-components}. Laisse $U \subset X$ et $V \subset Y$ être non vide ouverts affine avec
$f(U) \subset V$. Rappelez-vous que les anneaux $A = \mathcal{O}_X(U)$
et $B = \mathcal{O}_Y(V)$ sont des domaines intégrés. Le morphisme $f|_U : U \to V$ correspond à un homomorphisme d'anneaux
$\varphi : B \to A$. Les points génériques de $X$ et $Y$ correspondent aux idéaux principaux $(0) \subset A$ et $(0) \subset B$. Ainsi (2) est équivalent à la condition $(0) = \varphi^{-1}((0))$, c'est-à-dire à la condition que $\varphi$ De cette façon, nous voyons que (2), (3) et (4) sont équivalents. $x$ et $y$ comme dans (5) les anneaux locaux
$\mathcal{O}_{X, x}$ et $\mathcal{O}_{Y, y}$ sont les domaines et les idéaux principaux $(0) \subset \mathcal{O}_{X, x}$
et $(0) \subset \mathcal{O}_{Y, y}$ correspondent aux points génériques de $X$ et $Y$ (par l'identification du spectre de l'anneau local à $x$
avec l'ensemble de points spécialisés à $x$, voir Schémas, lemme \ref{schemes-lemma-specialize-points}Ainsi, nous pouvons argumenter de la même manière que ci-dessus pour voir que (2), (5), et (6) sont équivalents.

\end{proof}







\section{Morphismes surjectifs}
\label{section-surjective}

\begin{definition}

\label{definition-surjective}

Un morphisme des schémas est dit être {\it surjectif} s'il est surjectif sur les espaces topologiques sous-jacents.

\end{definition}

\begin{lemma}

\label{lemma-composition-surjective}

La composition des morphismes surjectifs est surjective.

\end{lemma}

\begin{proof}
Démonstration omise.
\end{proof}

\begin{lemma}

\label{lemma-when-point-maps-to-pair}
Laissez $X$ et $Y$ être des schémas sur un schéma de base $S$. Compte tenu des points $x \in X$ et $y \in Y$, il ya un point de $X \times_S Y$ cartographie à $x$ et $y$ sous les projections si et seulement si $x$ et $y$ se trouvent au-dessus du même point de $S$.
\end{lemma}

\begin{proof}
L'état est évidemment nécessaire, et l'inverse suit de la démonstration de Schémas, lemme \ref{schemes-lemma-points-fibre-product}.
\end{proof}

\begin{lemma}

\label{lemma-base-change-surjective}
Le changement de base d'un morphisme surjectif est surjectif.
\end{lemma}

\begin{proof}

Laisser $f: X \to Y$ être un morphisme des régimes sur un régime de base $S$Si $S' \to S$ est un morphisme de schémas, let $p: X_{S'} \to X$
et $q: Y_{S'} \to Y$ les projections canoniques. Le carré commutatif
$$
\xymatrix{
X_{S'} \ar[d]_{f_{S'}} \ar[r]_p & X \ar[d]^{f} \\
Y_{S'} \ar[r]^{q} & Y.
}
$$
identifier $X_{S'}$ comme produit fibré de $X \to Y$ et
$Y_{S'} \to Y$. Laisse $Z$ être un sous-ensemble de l'espace topologique sous-jacent de $X$. Ensuite $q^{-1}(f(Z)) = f_{S'}(p^{-1}(Z))$Parce que
$y' \in q^{-1}(f(Z))$ si et seulement si $q(y') = f(x)$ pour certains $x \in Z$, si et seulement si, par lemme \ref{lemma-when-point-maps-to-pair}, il existe
$x' \in X_{S'}$ de telle manière que $f_{S'}(x') = y'$ et $p(x') = x$. En particulier prendre $Z = X$ nous voyons que si $f$ est surjectif ainsi que le changement de base
$f_{S'}: X_{S'} \to Y_{S'}$.

\end{proof}

\begin{example}

\label{example-injective-not-preserved-base-change}
La bijectivité n'est pas stable sous le changement de base, et donc pas non plus l'injectivité. Par exemple, considérer le bijection $\Spec(\mathbf{C}) \to \Spec(\mathbf{R})$. Le changement de base $\Spec(\mathbf{C} \otimes_{\mathbf{R}} \mathbf{C}) \to
\Spec(\mathbf{C})$ n'est pas injectable, car il y a un isomorphisme $\mathbf{C} \otimes_{\mathbf{R}} \mathbf{C} \cong \mathbf{C} \times \mathbf{C}$ (la décomposition vient de l'idempotent $\frac{1 \otimes 1 + i \otimes i}{2}$) et donc $\Spec(\mathbf{C} \otimes_{\mathbf{R}} \mathbf{C})$ a deux points.
\end{example}

\begin{lemma}

\label{lemma-surjection-from-quasi-compact}
Laissez $$
\xymatrix{
X \ar[rr]_f \ar[rd]_p & &
Y \ar[dl]^q \\
& Z
}
$$ être un diagramme commutatif des morphismes des schémas. Si $f$ est surjectif et $p$ est quasi-compact, alors $q$ est quasi-compact.
\end{lemma}

\begin{proof}
Laissez $W \subset Z$ être un quasi-compact ouvert. Par hypothèse $p^{-1}(W)$ est quasi-compact. Par conséquent, par Topologie, lemme \ref{topology-lemma-image-quasi-compact} l'image inverse $q^{-1}(W) = f(p^{-1}(W))$ est quasi-compact aussi. Cela prouve la lemme.
\end{proof}




\section{Morphismes radiciels et universellement injectifs}
\label{section-radicial}

\noindent
Dans cette section, nous définissons ce que cela signifie pour un morphisme de schémas d'être {\it radicial} et ce que cela signifie pour un morphisme de schémas d'être {\it universellement injectable}. Nous montrons ensuite que ces notions sont d'accord. La raison pour introduire les deux est que dans le cas des espaces algébriques il y a des notions correspondantes qui ne peuvent pas toujours être d'accord.

\begin{definition}

\label{definition-universally-injective}
Laissez $f : X \to S$ être un morphisme.
\begin{enumerate}
\item Nous disons que $f$ est {\it universellement injectable} si et seulement si pour tout morphisme des schémas $S' \to S$ le changement de base $f' : X_{S'} \to S'$ est injectable (sur les espaces topologiques sous-jacents).
\item Nous disons que $f$ est {\it radical} si $f$ est injectable comme une application des espaces topologiques, et pour chaque $x \in X$ l'extension de champ $\kappa(x)/\kappa(f(x))$ est purement inséparable.
\end{enumerate}

\end{definition}

\begin{lemma}

\label{lemma-universally-injective}

\begin{reference}
\cite[proposition 3.7.1]{EGA1-seconde}
\end{reference}
Laissez $f : X \to S$ être un morphisme de schémas. Les conditions suivantes sont équivalentes:
\begin{enumerate}
\item Pour chaque champ $K$ l'application induite $\Mor(\Spec(K), X) \to \Mor(\Spec(K), S)$ est injectable.
\item Le morphisme $f$ est universellement injectable.
\item Le morphisme $f$ est radial.
\item Le morphisme diagonal $\Delta_{X/S} : X \longrightarrow X \times_S X$ est surjectif.
\end{enumerate}

\end{lemma}

\begin{proof}
Laissez $K$ être un champ, et laissez $s : \Spec(K) \to S$ être un morphisme. Donner un morphisme $x : \Spec(K) \to X$ tel que $f \circ x = s$ est le même que donner une section de la projection $X_K = \Spec(K) \times_S X \to \Spec(K)$, qui à son tour est le même que donner un point $x \in X_K$ dont le corps résiduel est $K$.

\medskip\noindent

À l'inverse, supposez que (1) tient. Supposons que $x, x' \in X_{S'}$
carte au même point $s' \in S'$. Choisissez un diagramme commutatif
$$
\xymatrix{
K & \kappa(x) \ar[l] \\
\kappa(x') \ar[u] & \kappa(s') \ar[l] \ar[u]
}
$$
des champs. Par Schémas, me laisse \ref{schemes-lemma-characterize-points}
nous obtenons deux morphismes $a, a' : \Spec(K) \to X_{S'}$. Un correspondant au point $x$ et de l ' intégration $\kappa(x) \subset K$ et l'autre correspondant au point $x'$ et de l ' intégration
$\kappa(x') \subset K$. Aussi nous avons $f' \circ a = f' \circ a'$. Condition (1) implique maintenant que les compositions de $a$ et $a'$ avec
$X_{S'} \to X$ sont égaux. $X_{S'}$ est le produit fibré de $S'$ et $X$ terminé $S$ nous voyons que $a = a'$. D'où $x = x'$. Ainsi 1) implique (2).

\medskip\noindent
S'il y a deux points différents de la cartographie $x, x' \in X$ au même point de $s$ (2) est violé. Si pour certains $s = f(x)$, $x \in X$ l'extension de champ $\kappa(x)/\kappa(s)$ n'est pas purement inséparable, alors nous pouvons trouver une extension de champ $K/\kappa(s)$ telle que $\kappa(x)$ a deux homomorphismes $\kappa(s)$ en $K$. Par Schémas, lemme \ref{schemes-lemma-characterize-points} cela implique que l'application $\Mor(\Spec(K), X) \to \Mor(\Spec(K), S)$ n'est pas injectable, et donc (1) est violée. Ainsi, nous voyons que les conditions équivalentes (1) et (2) impliquent que $f$ est radial, c'est-à-dire qu'ils impliquent (3).

\medskip\noindent
Assumme (3). Par Schémas, lemme \ref{schemes-lemma-characterize-points} un morphisme $\Spec(K) \to X$ est donné par une paire $(x, \kappa(x) \to K)$. La propriété (3) dit exactement que l'association à la paire $(x, \kappa(x) \to K)$ la paire $(s, \kappa(s) \to \kappa(x) \to K)$ est injectable. En d'autres termes (1) tient. À ce moment nous savons que (1), (2) et (3) sont tous équivalents.

\medskip\noindent

Enfin, nous prouvons l'équivalence de (4) avec (1), (2) et (3). $X \times_S X$ est administrée par un
$(x_1, x_2, s, \mathfrak p)$où: $x_1, x_2 \in X$,
$f(x_1) = f(x_2) = s$ et
$\mathfrak p \subset \kappa(x_1) \otimes_{\kappa(s)} \kappa(x_2)$
est un idéal idéal, voir Schémas, lemme \ref{schemes-lemma-points-fibre-product}Si $f$ est universellement injectable, puis en prenant $S'=X$ dans la définition de l'injection universelle,
$\Delta_{X/S}$ doit être surjectif puisqu'il s'agit d'une section du morphisme injectable
$X \times_S X  \longrightarrow X$. Inversement, si
$\Delta_{X/S}$ est surjectif, alors toujours $x_1 = x_2 = x$ et il y a exactement un tel idéal $\mathfrak p$, ce qui signifie que
$\kappa(s) \subset \kappa(x)$ est purement inséparable. $f$ est radial. $\Delta_{X/S}$ est surjectif, puis pour n'importe quel $S' \to S$ le changement de base $\Delta_{X_{S'}/S'}$ est surjectif ce qui implique que $f$
est universellement injectable. Ceci termine la démonstration de la lemme.

\end{proof}

\begin{lemma}

\label{lemma-universally-injective-separated}
Un morphisme universel injectif est séparé.
\end{lemma}

\begin{proof}
Combinez lemme \ref{lemma-universally-injective} avec la remarque que $X \to S$ est séparé si et seulement si l'image de $\Delta_{X/S}$ est fermée dans $X \times_S X$, voir Schémas, définition \ref{schemes-definition-separated} et la discussion qui suit.
\end{proof}

\begin{lemma}

\label{lemma-base-change-universally-injective}

Un changement de base d'un morphisme universellement injectif est universellement injectable.

\end{lemma}

\begin{proof}
Ceci est formel.
\end{proof}

\begin{lemma}

\label{lemma-composition-universally-injective}

Une composition de morphismes radiciels est radicial, et donc la même chose vaut pour l'équivalent condition d'être universellement injectable.

\end{lemma}

\begin{proof}
Démonstration omise.
\end{proof}









\section{Morphismes affines}
\label{section-affine}

\begin{definition}

\label{definition-affine}
Un morphisme des schémas $f : X \to S$ est appelé {\it affine} si l'image inverse de chaque ouvert affine de $S$ est une affine ouverte de $X$.
\end{definition}

\begin{lemma}

\label{lemma-affine-separated}
Un affine morphisme est séparé et quasi compact.
\end{lemma}

\begin{proof}
Soit $f : X \to S$ un morphisme affine. La quasi-compacité résulte immédiatement de Schémas, lemme \ref{schemes-lemma-quasi-compact-affine}. Nous allons montrer que $f$ est séparé en utilisant Schémas, lemme \ref{schemes-lemma-characterize-separated}. Soient $x_1, x_2 \in X$ des points de $X$ qui s'envoient sur le même point $s \in S$. Choisissons un ouvert affine $W \subset S$ contenant $s$. Par hypothèse, $f^{-1}(W)$ est affine. Appliquons le lemme cité avec $U = V = f^{-1}(W)$.
\end{proof}

\begin{lemma}

\label{lemma-characterize-affine}

\begin{reference}

\cite[II, corollaire 1.3.2]{EGA}

\end{reference}

Laisser $f : X \to S$ morphisme des régimes. Les suivants sont équivalents

\begin{enumerate}

\item Le morphisme $f$ C'est affiné.
\item Il existe un recouvrement ouvert affine $S = \bigcup W_j$
de telle sorte que chacun $f^{-1}(W_j)$ C'est affiné.
\item Il existe un faisceau quasi-cohérent de $\mathcal{O}_S$-Algèbres
$\mathcal{A}$ et un isomorphisme
$X \cong \underline{\Spec}_S(\mathcal{A})$
des régimes $S$. Voir Constructions, section \ref{constructions-section-spec} pour la notation.

\end{enumerate}

En outre, en l'espèce, $X = \underline{\Spec}_S(f_*\mathcal{O}_X)$.

\end{lemma}

\begin{proof}
Il est évident que (1) implique (2).

\medskip\noindent

Présumer $S = \bigcup_{j \in J} W_j$ est un recouvrement ouvert affine de telle sorte que chacun $f^{-1}(W_j)$ Par Schémas, lemme \ref{schemes-lemma-quasi-compact-affine} nous voyons que $f$ est quasi compact. Par Schémas, laisse \ref{schemes-lemma-characterize-quasi-separated}
nous voyons le morphisme $f$ est quasi-séparé. Ainsi par Schémas, lemme \ref{schemes-lemma-push-forward-quasi-coherent} les gerbes $\mathcal{A} = f_*\mathcal{O}_X$ est un faisceau quasi-cohérent de $\mathcal{O}_S$- des algèbres. Ainsi nous avons le schéma
$g : Y = \underline{\Spec}_S(\mathcal{A}) \to S$ terminé $S$. La carte d'identité
$\text{id} : \mathcal{A} = f_*\mathcal{O}_X \to f_*\mathcal{O}_X$
fournit, par la définition du spectre relatif, un morphisme $can : X \to Y$ terminé $S$, voir Constructions, lemme \ref{constructions-lemma-canonical-morphism}. Par hypothèse et le lemme vient de citer la restriction $can|_{f^{-1}(W_j)} : f^{-1}(W_j) \to g^{-1}(W_j)$
est un isomorphisme. $can$ Nous avons montré que (2) implique (3).

\medskip\noindent

Supposons que (3). Par Constructions, lemme \ref{constructions-lemma-spec-properties}
nous voyons que l'image inverse de chaque ouvert affine est affine, et donc le morphisme est affine par définition.

\end{proof}

\begin{remark}

\label{remark-direct-argument}

Nous pouvons également dire directement que (2) implique (1) dans le lemme \ref{lemma-characterize-affine} ci-dessus, comme suit : $S = \bigcup W_j$ est un recouvrement ouvert affine de telle sorte que chacun $f^{-1}(W_j)$ est affine. D'abord argumenter que $\mathcal{A} = f_*\mathcal{O}_X$ est quasi-cohérent comme dans la démonstration ci-dessus. $\Spec(R) = V \subset S$ Nous devons montrer que $f^{-1}(V)$ C'est affiné.
$A = \mathcal{A}(V) = f_*\mathcal{O}_X(V) = \mathcal{O}_X(f^{-1}(V))$Par Schémas, lemme \ref{schemes-lemma-morphism-into-affine} il y a un morphisme canonique $\psi : f^{-1}(V) \to \Spec(A)$ terminé
$\Spec(R) = V$Par Schémas, lemme \ref{schemes-lemma-good-subcover} il existe un entier $n \geq 0$, une ouverture de recouvrement standard
$V = \bigcup_{i = 1, \ldots, n} D(h_i)$, $h_i \in R$, et une application
$a : \{1, \ldots, n\} \to J$ de telle sorte que chacun $D(h_i)$ est également un principal ouvert du régime affine $W_{a(i)}$. L'image inverse d'un principal d'ouverture sous un morphisme de schémas d'affine est principal d'ouverture, voir Algèbre, lemme \ref{algebra-lemma-spec-functorial}. Nous voyons donc que $f^{-1}(D(h_i))$ est un principal ouvert de $f^{-1}(W_{a(i)})$, en particulier que $f^{-1}(D(h_i))$ C'est affine. $\mathcal{A}$
est quasi-cohérent que nous avons
$A_{h_i} = \mathcal{A}(D(h_i)) = \mathcal{O}_X(f^{-1}(D(h_i)))$Donc... $f^{-1}(D(h_i))$ est le spectre de $A_{h_i}$Il s'ensuit que le morphisme $\psi$ induit un isomorphisme de l'ouvert
$f^{-1}(D(h_i))$ avec l'ouverture $\Spec(A_{h_i})$ de
$\Spec(A)$Depuis $f^{-1}(V) = \bigcup f^{-1}(D(h_i))$
et $\Spec(A) = \bigcup \Spec(A_{h_i})$ Nous avons gagné.

\end{remark}

\begin{lemma}

\label{lemma-affine-equivalence-algebras}

Laisser $S$ Il y a une anti-équivalence de catégories
$$
\begin{matrix}
\text{Schemes affine} \\
\text{over }S
\end{matrix}
\longleftrightarrow
\begin{matrix}
\text{quasi-coherent sheaves} \\
\text{of }\mathcal{O}_S\text{-algebras}
\end{matrix}
$$
qui s ' associe à $f : X \to S$ les gerbes $f_*\mathcal{O}_X$En outre, cette équivalence est compatible avec un changement de base arbitraire.

\end{lemma}

\begin{proof}
Le functeur de droite à gauche est donné par $\underline{\Spec}_S$. Les deux functeurs sont mutuellement inversés par lemme \ref{lemma-characterize-affine} et Constructions, lemme \ref{constructions-lemma-spec-properties} partie (3). L'énoncé final est Constructions, lemme \ref{constructions-lemma-spec-properties} partie (2).
\end{proof}

\begin{lemma}

\label{lemma-quasi-coherence-scalar-restriction}

\begin{reference}
\cite[I, proposition 9.6.1]{EGA}
\end{reference}
Laisser $S$ être un schéma et laisser $\mathcal{A}$ être un $\mathcal{O}_S$-algèbre quasi-cohérent. Un module $\mathcal{A}$ est quasi-cohérent comme un module $\mathcal{O}_S$ si et seulement si il est quasi-cohérent comme un module $\mathcal{A}$.
\end{lemma}

\begin{proof}
Laissez $\mathcal{F}$ être un module $\mathcal{A}$. Si $\mathcal{F}$ est quasi-cohérent comme un module $\mathcal{A}$, alors pour chaque $s \in S$ il existe un quartier ouvert $U$ de $s$ et une suite exacte $$
\bigoplus\nolimits_J\mathcal{A}|_U\to
\bigoplus\nolimits_I\mathcal{A}|_U\to
\mathcal{F}|_U
$$ de $\mathcal{A}|_U$-modules. Ensuite, c'est aussi une suite exacte de $\mathcal{O}_U$-modules. D'où $\mathcal{F}|_U$ est quasi-cohérent comme le cokernel d'un morphisme de $\mathcal{O}_U$-modules quasi-cohérent sur un schéma. Il s'ensuit que $\mathcal{F}$ est quasi-cohérent comme un module $\mathcal{O}_X$.

\medskip\noindent

À l'inverse, supposer $\mathcal{F}$ est quasi-cohérent comme un $\mathcal{O}_X$- Module. Choisissez une affine ouverte $\Spec(R) = U \subset S$. Nous avons des isomorphismes
$\mathcal{O}_U$-modules $\mathcal{A}|_U \cong \widetilde{A}$ et
$\mathcal{F}|_U \cong \widetilde{M}$, pour certains $R$-Algèbre $A$
et certains $R$-module $M$. Les $\mathcal{A}$-module structure sur $\mathcal{F}$
traduit en un $A$-module structure sur $M$ compatible avec le $R$- structure du module (détails omis). Choisissez une suite exacte
$$
\bigoplus\nolimits_J A \to
\bigoplus\nolimits_I A \to
M \to 0
$$
de $A$Depuis le functeur. $\widetilde{\ }$ est exact, cela produit une suite exacte
$$
\bigoplus\nolimits_J \widetilde{A}\to
\bigoplus\nolimits_I \widetilde{A}\to
\widetilde{M} \to 0
$$
de $\widetilde{A}$- des modules.
$\mathcal{F}$ est quasi-cohérent comme un $\mathcal{A}$- Module.

\end{proof}

\begin{lemma}

\label{lemma-affine-equivalence-modules}
Laissez $f : X \to S$ être un morphisme affine de schémas. Laissez $\mathcal{A} = f_*\mathcal{O}_X$. Le functeur $\mathcal{F} \mapsto f_*\mathcal{F}$ induit une équivalence des catégories $$
\left\{
\begin{matrix}
\text{category of quasi-coherent}\\
\mathcal{O}_X\text{-modules}
\end{matrix}
\right\}
\longrightarrow
\left\{
\begin{matrix}
\text{category of quasi-coherent}\\
\mathcal{A}\text{-modules}
\end{matrix}
\right\}
$$ En outre, un module $\mathcal{A}$ est quasi-cohérent comme un module $\mathcal{O}_S$ si et seulement s'il est quasi-cohérent comme un module $\mathcal{A}$.
\end{lemma}

\begin{proof}
L'instruction finale est lemme \ref{lemma-quasi-coherence-scalar-restriction}. Par lemme \ref{lemma-affine-separated} et Schémas, lemme \ref{schemes-lemma-push-forward-quasi-coherent} la poussée $f_*\mathcal{F}$ d'un module $\mathcal{O}_X$ quasi-cohérent $\mathcal{F}$ est un module $\mathcal{O}_S$ quasi-cohérent.

\medskip\noindent

Nous allons construire un quasi-inverse $h$ à ce functeur. $\mathcal{G}$
être un quasi-cohérent $\mathcal{A}$- Module.
$$
h(\mathcal{G}) = f^*\mathcal{G} \otimes_{f^*\mathcal{A}} \mathcal{O}_X =
f^*\mathcal{G} \otimes_{f^*f_*\mathcal{O}_X} \mathcal{O}_X
$$
Élucidation : le retrait $f^*\mathcal{A} = f^*f_*\mathcal{O}_X$
est un $\mathcal{O}_X$-l'algèbre, l'application d'adjonction
$f^*f_*\mathcal{O}_X \to \mathcal{O}_X$ est un homomorphisme de l'algèbre, et le retrait $f^*\mathcal{G}$ est un $f^*\mathcal{A}$- Module. Observez que
$h(\mathcal{G})$ est quasi-cohérente car la quasi-cohérence est préservée par des retraits et le changement d'anneaux. Observez qu'il y a une application functoire
$$
h(f_*\mathcal{F}) =
f^*f_*\mathcal{F} \otimes_{f^*f_*\mathcal{O}_X} \mathcal{O}_X \to \mathcal{F}
$$
provenant de l'application d'adjonction $f^*f_*\mathcal{F} \to \mathcal{F}$ et une application functoire
$$
\mathcal{G} \to f_*h(\mathcal{G}) \quad\text{adjoint to the map}\quad
f^*\mathcal{G} \to f^*\mathcal{G} \otimes_{f^*f_*\mathcal{O}_X} \mathcal{O}_X
$$
qui envoie une section locale $s$ de $f^*\mathcal{F}$ à $s \otimes 1$. Pour terminer la démonstration il suffit de montrer que ces cartes sont des isomorphismes pour $\mathcal{F}$ et $\mathcal{G}$ comme ci-dessus. Ceci peut être vérifié sur les membres d'une couverture d'affinage, c'est-à-dire lorsque $X$ et $S$ Ils sont affinés.

\medskip\noindent
L'observation clé de l'algèbre qui rend ce travail est la suivante : Que $R \to A$ soit un homomorphisme d'anneaux.Soit $N$ un module $A$. Puis $$
(N \otimes_R A) \otimes_{(A \otimes_R A)} A = N
$$ Autrement dit, le côté gauche de cette égalité est l'effet d'appliquer $h$ à la quasi-cohérente $\widetilde{A}$-module $\widetilde{N}$ sur $\Spec(R)$. Nous omettons les détails.
\end{proof}

\begin{lemma}

\label{lemma-composition-affine}

La composition des morphismes affines est affine.

\end{lemma}

\begin{proof}
Laisser $f : X \to Y$ et $g : Y \to Z$ être des morphismes affines. Laisser $U \subset Z$ être ouvert affine. Alors $g^{-1}(U)$ est affine par hypothèse sur $g$. Sur quoi $f^{-1}(g^{-1}(U))$ est affine par hypothèse sur $f$. Par conséquent $(g \circ f)^{-1}(U)$ est affine.
\end{proof}

\begin{lemma}

\label{lemma-base-change-affine}

Le changement de base d'une affine morphisme est affine.

\end{lemma}

\begin{proof}

Laisser $f : X \to S$ être un morphisme affine. $S' \to S$ être tout morphisme. $f' : X_{S'} = S' \times_S X \to S'$ le changement de base de $f$. Pour chaque $s' \in S'$ il existe un quartier affiné ouvert
$s' \in V \subset S'$ qui cartographie dans une affine ouverte $U \subset S$. Par hypothèse $f^{-1}(U)$ est affine. Par le matériel dans Schémas, section \ref{schemes-section-fibre-products} nous voyons que $f^{-1}(U)_V = V \times_U f^{-1}(U)$ est affine et égale à $(f')^{-1}(V)$. Cela prouve que $S'$ a un recouvrement ouvert par affines dont l'image inverse sous $f'$ Nous concluons par lemme \ref{lemma-characterize-affine} ci-dessus.

\end{proof}

\begin{lemma}

\label{lemma-closed-immersion-affine}

Une fermée d'immersion est affine.

\end{lemma}

\begin{proof}
La première indication est Schémas, lemme \ref{schemes-lemma-closed-immersion-affine-case}. Voir Schémas, lemme \ref{schemes-lemma-closed-subspace-scheme} pour un énoncé complète.
\end{proof}

\begin{lemma}

\label{lemma-affine-s-open}
Laissez $X$ être un schéma. Laissez $\mathcal{L}$ être un inversible $\mathcal{O}_X$-module. Laissez $s \in \Gamma(X, \mathcal{L})$. Le morphisme d'inclusion $j : X_s \to X$ est affine.
\end{lemma}

\begin{proof}

Cela découle de Propriétés, me laisse \ref{properties-lemma-affine-cap-s-open}
et la définition.

\end{proof}

\begin{lemma}

\label{lemma-affine-permanence}

Supposons $g : X \to Y$ est un morphisme de schémas sur $S$.

\begin{enumerate}

\item Si $X$ C'est fini $S$ et $\Delta : Y \to Y \times_S Y$ est affine, alors $g$ C'est affiné.
\item Si $X$ C'est fini $S$ et $Y$ est séparé sur $S$Alors $g$ C'est affiné.
\item Un morphisme d'un schéma affiné à un schéma avec diagonale affiné est affiné.
\item Un morphisme d'un schéma affiné à un schéma séparé est affiné.

\end{enumerate}

\end{lemma}

\begin{proof}

(1). Le changement de base $X \times_S Y \to Y$ est affiné par Lemme \ref{lemma-base-change-affine}Le morphisme $(1, g) : X \to X \times_S Y$ est le changement de base de
$Y \to Y \times_S Y$ par le morphisme $X \times_S Y \to Y \times_S Y$. Il est donc affiné par lemme \ref{lemma-base-change-affine}. La composition des morphismes affines est affine (voir lemme \ref{lemma-composition-affine}) et (1) suivant. La partie (2) suit de (1) comme une immersion fermée est affine (voir lemme \ref{lemma-closed-immersion-affine}) et $Y/S$ moyens séparés $\Delta$ Les parties (3) et (4) sont des cas spéciaux de (1) et (2).

\end{proof}

\begin{lemma}

\label{lemma-morphism-affines-affine}
Un morphisme entre les schémas affinés est affiné.
\end{lemma}

\begin{proof}
Immédiate de lemme \ref{lemma-affine-permanence} avec $S = \Spec(\mathbf{Z})$. Elle découle également directement de l'équivalence de (1) et (2) dans le lemme \ref{lemma-characterize-affine}.
\end{proof}

\begin{lemma}

\label{lemma-Artinian-affine}
Laissez $S$ être un schéma. Laissez $A$ être un anneau artinien. Tout morphisme $\Spec(A) \to S$ est affine.
\end{lemma}

\begin{proof}
Démonstration omise.
\end{proof}

\begin{lemma}

\label{lemma-get-affine}
Laissez $j : Y \to X$ être une immersion de schémas. Supposons qu'il existe un $U \subset X$ ouvert avec complément $Z = X \setminus U$ tel que
\begin{enumerate}
\item $U \to X$ est affine,
\item $j^{-1}(U) \to U$ est affine, et
\item $j(Y) \cap Z$ est fermé.
\end{enumerate}
$j$ alors affine. En particulier, si $X$ est affine, ainsi que $Y$.
\end{lemma}

\begin{proof}

Par Schémas, définition \ref{schemes-definition-immersion} il existe un sous-système ouvert $W \subset X$ de telle manière que $j$ facteurs comme une immersion fermée
$i : Y \to W$ suivi du morphisme d'inclusion $W \to X$. Comme une immersion fermée est affine (lemme \ref{lemma-closed-immersion-affine}), nous voyons cela pour chaque $x \in W$ il y a un quartier affiné d'ouverture de $x$ en $X$ dont l'image inverse sous $j$ est affine. Si $x \in U$, alors la même chose est vraie par hypothèse (2). $x \in Z$ et $x \not \in W$. Ensuite $x \not \in j(Y) \cap Z$. Par hypothèse (3) nous pouvons trouver un quartier affiné ouvert
$V \subset X$ de $x$ qui ne se rencontre pas $j(Y) \cap Z$. Ensuite $j^{-1}(V) = j^{-1}(V \cap U)$ qui est affiné par les hypothèses (1) et (2). Il s'ensuit que $j$ est affiné par Lemme \ref{lemma-characterize-affine}.

\end{proof}







\section{Familles de modules inversibles amples}
\label{section-families-ample-invertible-modules}

\noindent
Une courte section sur la notion d'une famille de modules inversibles.

\begin{definition}

\label{definition-family-ample-invertible-modules}

\begin{reference}

\cite[II définition 2.2.4]{SGA6}

\end{reference}

Laisser $X$ Être un plan. $\{\mathcal{L}_i\}_{i \in I}$
être une famille d'inversibles $\mathcal{O}_X$- des modules.
$\{\mathcal{L}_i\}_{i \in I}$ est un {\it ample famille de modules inversibles sur $X$} si \begin{enumerate}

\item  $X$ est quasi compacte, et
\item pour chaque $x \in X$ il existe un $i \in I$, un $n \geq 1$, $s \in \Gamma(X, \mathcal{L}_i^{\otimes n})$ de telle manière que $x \in X_s$ et $X_s$ C'est affiné.

\end{enumerate}

\end{definition}

\noindent
Si $\{\mathcal{L}_i\}_{i \in I}$ est une ample famille de modules inversibles sur un schéma $X$, alors il existe un sous-ensemble fini $I' \subset I$ de telle sorte que $\{\mathcal{L}_i\}_{i \in I'}$ est une ample famille de modules inversibles sur $X$ (suivant immédiatement de quasi-compactness). Un schéma ayant une ample famille de modules inversibles a une diagonale affinée par le prochain lemme et est donc a fortiori quasi-séparé.

\begin{lemma}

\label{lemma-affine-diagonal}
Laissez $X$ être un schéma tel que pour chaque point $x \in X$ il existe un inversible $\mathcal{O}_X$-module $\mathcal{L}$ et une section globale $s \in \Gamma(X, \mathcal{L})$ telle que $x \in X_s$ et $X_s$ est affine. Puis la diagonale de $X$ est un morphisme affine.
\end{lemma}

\begin{proof}

Données inversibles $\mathcal{O}_X$-modules $\mathcal{L}$, $\mathcal{M}$
et sections mondiales $s \in \Gamma(X, \mathcal{L})$,
$t \in \Gamma(X, \mathcal{M})$ de telle manière que $X_s$ et $X_t$ Nous devons prouver $X_s \cap X_t$ c'est affine. A savoir, puis lemme \ref{lemma-characterize-affine}
appliqué à $\Delta : X \to X \times X$ et le fait que
$\Delta^{-1}(X_s \times X_t) = X_s \cap X_t$ montre que $\Delta$
Le fait que $X_s \cap X_t$ est affine suit de Propriétés, lemme \ref{properties-lemma-affine-cap-s-open}.

\end{proof}

\begin{remark}

\label{remark-affine-s-opens-cover-family}

Dans Propriétés, lemme \ref{properties-lemma-affine-s-opens-cover-quasi-separated}
nous voyons qu'un schéma qui a un module ample inversible est séparé. C'est mal pour les schémas ayant une ample famille de modules inversibles. $X$ être comme dans Schémas, exemple \ref{schemes-example-affine-space-zero-doubled}
avec $n = 1$, c'est-à-dire que la ligne affine avec zéro a doublé. Nous utilisons la notation de cet exemple sauf que nous écrivons $x$ pour $x_1$ et $y$
pour $y_1$. Il y a, pour chaque entier $n$, une gerbe inversible $\mathcal{L}_n$ le $X$ qui est trivial sur $X_1$ et
$X_2$ et dont la fonction de transition $U_{12} \to U_{21}$ est
$f(x) \mapsto y^n f(y)$. Les sections mondiales
$\mathcal{L}_n$ sont paires $(f(x), g(y)) \in k[x] \oplus
k[y]$ de telle manière que $y^n f(y) = g(y)$. Les sections $s = (1,
y)$ de $\mathcal{L}_1$ et $t = (x, 1)$ de
$\mathcal{L}_{-1}$ déterminer un couvercle affiné ouvert parce que
$X_s = X_1$ et $X_t = X_2$. Par conséquent $X$ dispose d'une ample famille de modules inversibles mais il n'est pas séparé.

\end{remark}









\section{Morphismes quasi-affines}
\label{section-quasi-affine}

\noindent
Rappelez-vous qu'un schéma $X$ est appelé {\it quasi-affine} s'il est quasi-compact et isomorphique à un sous-système ouvert d'un schéma affine, voir Propriétés, définition \ref{properties-definition-quasi-affine}.

\begin{definition}

\label{definition-quasi-affine}
Un morphisme de schémas $f : X \to S$ est appelé {\it quasi-affine} si l'image inverse de chaque ouvert affine de $S$ est un schéma quasi-affine.
\end{definition}

\begin{lemma}

\label{lemma-quasi-affine-separated}
Un morphisme quasi-affine est séparé et quasi-compact.
\end{lemma}

\begin{proof}
Laisser $f : X \to S$ être quasi-affinée. La quasi-compactivité est immédiate de Schémas, lemme \ref{schemes-lemma-quasi-compact-affine}. Laisser $U \subset S$ être une affine ouverte. Si nous pouvons montrer que $f^{-1}(U)$ est un schéma séparé, alors $f$ est séparé (Schémas, lemme \ref{schemes-lemma-characterize-separated} montre qu'être séparé est local sur la base). Par hypothèse $f^{-1}(U)$ est isomorphe à un sous-système ouvert d'un schéma affiné. Un schéma affiné est séparé et donc chaque sous-système ouvert d'un schéma affiné est séparé comme désiré.
\end{proof}

\begin{lemma}

\label{lemma-characterize-quasi-affine}

Laisser $f : X \to S$ morphisme des régimes. Les suivants sont équivalents

\begin{enumerate}

\item Le morphisme $f$ C'est presque affiné.
\item Il existe un recouvrement ouvert affine $S = \bigcup W_j$
de telle sorte que chacun $f^{-1}(W_j)$ C'est presque affiné.
\item Il existe un faisceau quasi-cohérent de $\mathcal{O}_S$-Algèbres
$\mathcal{A}$ et une immersion quasi compacte ouverte
$$
\xymatrix{
X \ar[rr] \ar[rd] & & \underline{\Spec}_S(\mathcal{A}) \ar[dl] \\
& S &
}
$$
terminé $S$.
\item Même que dans (3) mais avec $\mathcal{A} = f_*\mathcal{O}_X$
et la flèche horizontale le morphisme canonique de Constructions, lemme \ref{constructions-lemma-canonical-morphism}.

\end{enumerate}

\end{lemma}

\begin{proof}
Il est évident que (1) implique (2) et que (4) implique (3).

\medskip\noindent

Présumer $S = \bigcup_{j \in J} W_j$ est un recouvrement ouvert affine de telle sorte que chacun $f^{-1}(W_j)$ Par Schémas, lemme \ref{schemes-lemma-quasi-compact-affine} nous voyons que $f$ est quasi compact. Par Schémas, laisse \ref{schemes-lemma-characterize-quasi-separated}
nous voyons le morphisme $f$ est quasi-séparé. Ainsi par Schémas, lemme \ref{schemes-lemma-push-forward-quasi-coherent} les gerbes $\mathcal{A} = f_*\mathcal{O}_X$ est un faisceau quasi-cohérent de $\mathcal{O}_X$- des algèbres. Ainsi nous avons le schéma
$g : Y = \underline{\Spec}_S(\mathcal{A}) \to S$ terminé $S$. La carte d'identité
$\text{id} : \mathcal{A} = f_*\mathcal{O}_X \to f_*\mathcal{O}_X$
fournit, par la définition du spectre relatif, un morphisme $can : X \to Y$ terminé $S$, voir Constructions, lemme \ref{constructions-lemma-canonical-morphism}. Par hypothèse, le lemme vient de citer, et Propriétés, lemme \ref{properties-lemma-quasi-affine}
la restriction $can|_{f^{-1}(W_j)} : f^{-1}(W_j) \to g^{-1}(W_j)$
est un ouverte d'immersion quasi-compact. $can$ Nous avons montré que (2) implique (4).

\medskip\noindent

Supposons que (3). Choisissez n'importe quelle ouvert affine $U \subset S$. Par Constructions, lemme \ref{constructions-lemma-spec-properties}
nous voyons que l'image inverse de $U$ dans le spectre relatif est affine. $f^{-1}(U)$ est quasi-affinée (notez que la quasi-compactitude est également encodée dans (3)). Ainsi (3) implique (1).

\end{proof}

\begin{lemma}

\label{lemma-composition-quasi-affine}

La composition des morphismes quasi-affines est quasi-affinée.

\end{lemma}

\begin{proof}

Laisser $f : X \to Y$ et $g : Y \to Z$ être morphismes quasi-affines. $U \subset Z$ être ouvert affine. Alors $g^{-1}(U)$ est quasi-affiner par hypothèse sur $g$. Laisse $j : g^{-1}(U) \to V$ être une quasi-immersion compacte ouverte dans un schéma d'affines $V$. Par lemme \ref{lemma-characterize-quasi-affine} ci-dessus nous voyons que $f^{-1}(g^{-1}(U))$
est un sous-système ouvert quasi-compact du spectre relatif
$\underline{\Spec}_{g^{-1}(U)}(\mathcal{A})$ pour certains faisceau quasi-cohérent de $\mathcal{O}_{g^{-1}(U)}$-Algèbres
$\mathcal{A}$Par Schémas, lemme \ref{schemes-lemma-push-forward-quasi-coherent}
les gerbes $\mathcal{A}' = j_*\mathcal{A}$
est un faisceau quasi-cohérent de $\mathcal{O}_V$-algèbres avec la propriété qui $j^*\mathcal{A}' = \mathcal{A}$. Nous obtenons donc un diagramme commutatif
$$
\xymatrix{
f^{-1}(g^{-1}(U)) \ar[r] &
\underline{\Spec}_{g^{-1}(U)}(\mathcal{A})
\ar[r] \ar[d] &
\underline{\Spec}_V(\mathcal{A}') \ar[d] \\
& g^{-1}(U) \ar[r]^j & V
}
$$
avec le carré étant un carré en fibre, voir Constructions, lemme \ref{constructions-lemma-spec-properties}. Notez que le coin supérieur droit est un schéma affiné. $(g \circ f)^{-1}(U)$ C'est presque affiné.

\end{proof}

\begin{lemma}

\label{lemma-base-change-quasi-affine}
Le changement de base d'un morphisme quasi-affine est quasi-affiné.
\end{lemma}

\begin{proof}

Laisser $f : X \to S$ être un morphisme quasi-affiné. Par lemme \ref{lemma-characterize-quasi-affine} ci-dessus nous pouvons trouver un faisceau quasi-cohérent de $\mathcal{O}_S$-Algèbres $\mathcal{A}$ et une immersion quasi compacte ouverte $X \to \underline{\Spec}_S(\mathcal{A})$
terminé $S$. Laisse $g : S' \to S$ être tout morphisme. $f' : X_{S'} = S' \times_S X \to S'$ le changement de base de $f$. Puisque le changement de base d'un ouverte d'immersion quasi-compact est un ouverte d'immersion quasi-compact, nous voyons que
$X_{S'} \to \underline{\Spec}_{S'}(g^*\mathcal{A})$
est un ouverte d'immersion quasi-compact (nous avons utilisé Schémas, lemmes
\ref{schemes-lemma-quasi-compact-preserved-base-change} et
\ref{schemes-lemma-base-change-immersion} et des travaux de construction \ref{constructions-lemma-spec-properties}). Par lemme \ref{lemma-characterize-quasi-affine} Nous concluons à nouveau que $X_{S'} \to S'$ C'est presque affiné.

\end{proof}

\begin{lemma}

\label{lemma-quasi-compact-immersion-quasi-affine}

Une immersion quasi compacte est quasi-affinée.

\end{lemma}

\begin{proof}

Laisser $X \to S$ Nous devons montrer l'image inverse de chaque ouvert affine est quasi-affinée. $S$ c'est un schéma affiné, nous devons montrer
$X$ C'est presque affiné. \ref{lemma-quasi-compact-immersion}
le morphisme $X \to S$ les facteurs tels que $X \to Z \to S$ où $Z$ est un sous-système fermé de $S$ et $X \subset Z$ est un quasi-compact ouvert. $S$ est affine lemme \ref{lemma-closed-immersion} implique
$Z$ C'est là qu'on gagne.

\end{proof}

\begin{lemma}

\label{lemma-affine-quasi-affine}
Laissez $S$ être un schéma. Laissez $X$ être un schéma d'affine. Un morphisme $f : X \to S$ est quasi-affine si et seulement s'il est quasi-compact. En particulier tout morphisme d'un schéma d'affine à un schéma quasi-séparé est quasi-affine.
\end{lemma}

\begin{proof}
Laissez $V \subset S$ être un ouvert affine. Alors $f^{-1}(V)$ est un sous-système ouvert du schéma d'affinage $X$, donc quasi-affiner si et seulement si elle est quasi-compact. Cela prouve la première affirmation. La quasi-compactitude de n'importe quel $f : X \to S$ où $X$ est affiné et $S$ quasi-séparé suit de Schémas, lemme \ref{schemes-lemma-quasi-compact-permanence} appliqué à $X \to S \to \Spec(\mathbf{Z})$.
\end{proof}

\begin{lemma}

\label{lemma-quasi-affine-permanence}
Supposons que $g : X \to Y$ est un morphisme de schémas sur $S$. Si $X$ est quasi-affiné sur $S$ et $Y$ est quasi-séparé sur $S$, alors $g$ est quasi-affiné. En particulier, tout morphisme d'un schéma quasi-affiné à un schéma quasi-affiné est quasi-affiné.
\end{lemma}

\begin{proof}

Le changement de base $X \times_S Y \to Y$ est quasi-affinée par lemme \ref{lemma-base-change-quasi-affine}Le morphisme $X \to X \times_S Y$ est une immersion quasi-compacte comme $Y \to S$ est quasi-séparée, voir Schémas, lemme \ref{schemes-lemma-section-immersion}. Une quasi-compacte d'immersion est quasi-affinée par lemme \ref{lemma-quasi-compact-immersion-quasi-affine}
et la composition des morphismes quasi-affines est quasi-affinée (voir lemme \ref{lemma-composition-quasi-affine}Ainsi, nous gagnons.

\end{proof}











\section{Types de morphismes définis par des propriétés d'homomorphismes d'anneaux}
\label{section-properties-ring-maps}

\noindent
Dans cette section, nous étudions quelles propriétés des homomorphismes d'anneaux permettent de définir les propriétés locales des morphismes des schémas.

\begin{definition}

\label{definition-property-local}

Laisser $P$ être une propriété d'homomorphismes d'anneaux.

\begin{enumerate}

\item Nous disons que $P$ est {\it local} si les conditions suivantes sont maintenues:

\begin{enumerate}

\item Pour tout homomorphisme d'anneaux $R \to A$, et tout $f \in R$ Nous avons
$P(R \to A) \Rightarrow P(R_f \to A_f)$.
\item Pour tous les anneaux $R$, $A$, tout $f \in R$, $a\in A$, et tout homomorphisme d'anneaux
$R_f \to A$ Nous avons $P(R_f \to A) \Rightarrow P(R \to A_a)$.
\item Pour tout homomorphisme d'anneaux $R \to A$, $a_i \in A$ de telle manière que
$(a_1, \ldots, a_n) = A$ alors
$\forall i, P(R \to A_{a_i}) \Rightarrow P(R \to A)$.

\end{enumerate}

\item Nous disons que $P$ est {\it stable sous changement de base} si pour les homomorphismes d'anneaux $R \to A$, $R \to R'$ Nous avons
$P(R \to A) \Rightarrow P(R' \to R' \otimes_R A)$.
\item Nous disons que $P$ est {\it stable sous composition} si pour les homomorphismes d'anneaux $A \to B$, $B \to C$ Nous avons
$P(A \to B) \wedge P(B \to C) \Rightarrow P(A \to C)$.

\end{enumerate}

\end{definition}

\begin{definition}

\label{definition-locally-P}
Laissez $P$ être une propriété d'homomorphismes d'anneaux. Laissez $f : X \to S$ être un morphisme de schémas. Nous disons $f$ est {\it localement de type $P$} si pour n'importe quel $x \in X$ il existe un quartier affine ouvert $U$ de $x$ dans $X$ qui cartographie dans une affine ouvert $V \subset S$ de telle sorte que l'homomorphisme d'anneaux induit $\mathcal{O}_S(V) \to \mathcal{O}_X(U)$ possède la propriété $P$.
\end{definition}

\noindent
Ce n'est pas une définition « bon » à moins que la propriété $P$ soit une propriété locale. Même si $P$ est une propriété locale, nous n'utiliserons pas automatiquement cette définition pour dire qu'un morphisme est « localement de type $P$ » à moins que nous n'expliquions également explicitement la définition ailleurs.

\begin{lemma}

\label{lemma-locally-P}
Laissez $f : X \to S$ être un morphisme de schémas. Laissez $P$ être une propriété d'homomorphismes d'anneaux. Laissez $U$ être une affine ouverte de $X$, et $V$ une affine ouverte de $S$ telle que $f(U) \subset V$. Si $f$ est localement de type $P$ et $P$ est locale, alors $P(\mathcal{O}_S(V) \to \mathcal{O}_X(U))$ tient.
\end{lemma}

\begin{proof}

Comme $f$ est localement de type $P$ pour chaque $u \in U$ il existe un ouvert affine $U_u \subset X$ cartographie dans un ouvert affine $V_u \subset S$
de telle manière que $P(\mathcal{O}_S(V_u) \to \mathcal{O}_X(U_u))$ s'arrête. Choisissez un quartier ouvert $U'_u \subset U \cap U_u$ de $u$
qui est un ouvert affine standard dans les deux $U$ et $U_u$, voir Schémas, lemme \ref{schemes-lemma-standard-open-two-affines}. Par définition \ref{definition-property-local} (1) b) nous voyons que $P(\mathcal{O}_S(V_u) \to \mathcal{O}_X(U'_u))$ Nous pouvons donc supposer que $U_u \subset U$ est un ouvert affine standard. Choisissez un quartier ouvert $V'_u \subset V \cap V_u$
de $f(u)$ qui est un ouvert affine standard dans les deux $V$ et $V_u$, voir Schémas, lemme \ref{schemes-lemma-standard-open-two-affines}. Ensuite $U'_u = f^{-1}(V'_u) \cap U_u$ est une affinité d'ouverture standard de $U_u$ (d'où $U$) et nous avons
$P(\mathcal{O}_S(V'_u) \to \mathcal{O}_X(U'_u))$ par définition \ref{definition-property-local} 1) a). Par conséquent, nous pouvons supposer que les deux $U_u \subset U$ et $V_u \subset V$
sont affines d'ouverture standard. Application de la définition \ref{definition-property-local} (1) b) une fois de plus, nous concluons que $P(\mathcal{O}_S(V) \to \mathcal{O}_X(U_u))$
Parce que $U$ est quasi-compact nous pouvons choisir un nombre fini de points $u_1, \ldots, u_n \in U$ de telle manière que
$$
U = U_{u_1} \cup \ldots \cup U_{u_n}.
$$
Par définition \ref{definition-property-local} (1) c) nous concluons que $P(\mathcal{O}_S(V) \to \mathcal{O}_X(U))$ Attendez.

\end{proof}

\begin{lemma}

\label{lemma-locally-P-characterize}

Laisser $P$ être une propriété locale des homomorphismes d'anneaux. $f : X \to S$ morphisme des régimes. Les suivants sont équivalents

\begin{enumerate}

\item Le morphisme $f$ est localement de type $P$.
\item Pour tous les ouverts $U \subset X$, $V \subset S$
avec $f(U) \subset V$ Nous avons $P(\mathcal{O}_S(V) \to \mathcal{O}_X(U))$.
\item Il existe une ouverture de recouvrement $S = \bigcup_{j \in J} V_j$
et couvertures ouvertes $f^{-1}(V_j) = \bigcup_{i \in I_j} U_i$ de sorte que chacun des morphismes $U_i \to V_j$, $j\in J, i\in I_j$
est localement de type $P$.
\item Il existe un recouvrement ouvert affine $S = \bigcup_{j \in J} V_j$
et revêtements pour affines d'ouverture $f^{-1}(V_j) = \bigcup_{i \in I_j} U_i$ de telle manière que $P(\mathcal{O}_S(V_j) \to \mathcal{O}_X(U_i))$ pour tous
$j\in J, i\in I_j$.

\end{enumerate}

En outre, si $f$ est localement de type $P$ puis pour les sous-systèmes ouverts $U \subset X$, $V \subset S$ avec $f(U) \subset V$
la restriction $f|_U : U \to V$ est localement de type $P$.

\end{lemma}

\begin{proof}
Ceci découle de \ref{lemma-locally-P} ci-dessus.
\end{proof}

\begin{lemma}

\label{lemma-composition-type-P}
Laissez $P$ être une propriété d'homomorphismes d'anneaux. Supposons que $P$ est local et stable sous la composition. La composition des morphismes localement de type $P$ est localement de type $P$.
\end{lemma}

\begin{proof}

Laisser $f : X \to Y$ et $g : Y \to Z$ morphismes locaux de type $P$. Laisse $x \in X$. Choisissez un quartier affiné d'ouverture $W \subset Z$ de
$g(f(x))$. Choisissez un quartier affiné d'ouverture $V \subset g^{-1}(W)$
de $f(x)$. Choisissez un quartier affiné d'ouverture $U \subset f^{-1}(V)$
de $x$. Par lemme \ref{lemma-locally-P-characterize} les homomorphismes d'anneaux
$\mathcal{O}_Z(W) \to \mathcal{O}_Y(V)$ et
$\mathcal{O}_Y(V) \to \mathcal{O}_X(U)$ satisfaire $P$. D'où $\mathcal{O}_Z(W) \to \mathcal{O}_X(U)$ satisfait $P$
comme $P$ est supposé stable sous la composition.

\end{proof}

\begin{lemma}

\label{lemma-base-change-type-P}
Laissez $P$ être une propriété d'homomorphismes d'anneaux. Supposons que $P$ est local et stable sous changement de base. Le changement de base d'un morphisme localement de type $P$ est localement de type $P$.
\end{lemma}

\begin{proof}

Laisser $f : X \to S$ être un morphisme local de type $P$. Laisse $S' \to S$ être tout morphisme.
$f' : X_{S'} = S' \times_S X \to S'$ le changement de base de $f$. Pour chaque $s' \in S'$ il existe un quartier affiné ouvert
$s' \in V' \subset S'$ qui cartographie dans une affine ouverte $V \subset S$. Par lemme \ref{lemma-locally-P-characterize} l'ouverture $f^{-1}(V)$ est une union d'affines $U_i$ tels que les homomorphismes d'anneaux
$\mathcal{O}_S(V) \to \mathcal{O}_X(U_i)$ tous satisfaits $P$. Par le matériel de Schémas, section \ref{schemes-section-fibre-products}
nous voyons que $f^{-1}(U)_{V'} = V' \times_V f^{-1}(V)$ est l'union des ouverts affines $V' \times_V U_i$Depuis $\mathcal{O}_{X_{S'}}(V' \times_V U_i) =
\mathcal{O}_{S'}(V') \otimes_{\mathcal{O}_S(V)} \mathcal{O}_X(U_i)$
nous voyons que les homomorphismes d'anneaux
$\mathcal{O}_{S'}(V') \to \mathcal{O}_{X_{S'}}(V' \times_V U_i)$
satisfaire $P$ comme $P$ est supposé stable sous le changement de base.

\end{proof}

\begin{lemma}

\label{lemma-properties-local}

Les propriétés suivantes d'un homomorphisme d'anneaux $R \to A$ Ils sont locaux.

\begin{enumerate}

\item (Isomorphisme sur les anneaux locaux.) $\mathfrak q$ de $A$ couché sur $\mathfrak p \subset R$
l'homomorphisme d'anneaux $R \to A$ induit un isomorphisme
$R_{\mathfrak p} \to A_{\mathfrak q}$.
\item (Immersion ouverte.) Pour chaque prime $\mathfrak q$ de $A$ il existe un $f \in R$,
$\varphi(f) \not \in \mathfrak q$ tel que l'homomorphisme d'anneaux $\varphi : R \to A$
induit un isomorphisme $R_f \to A_f$.
\item (Fibres réduites.) $\mathfrak p$ de $R$ l'anneau en fibres
$A \otimes_R \kappa(\mathfrak p)$ est réduit.
\item (Fibres de dimension au plus $n$.) Pour chaque première $\mathfrak p$ de $R$ l'anneau en fibres
$A \otimes_R \kappa(\mathfrak p)$ a la dimension Krull au plus $n$.
\item (localement noethérien sur la cible.) L'homomorphisme d'anneaux $R \to A$ a la propriété qui $A$ C'est Noétherien.
\item Ajoutez-en plus ici au besoin\footnote{Mais seulement les propriétés qui ne sont pas déjà traitées séparément ailleurs.}.

\end{enumerate}

\end{lemma}

\begin{proof}
Démonstration omise.
\end{proof}

\begin{lemma}

\label{lemma-properties-base-change}

Les propriétés suivantes des homomorphismes d'anneaux sont stables sous le changement de base.

\begin{enumerate}

\item (Isomorphisme sur les anneaux locaux.) $\mathfrak q$ de $A$ couché sur $\mathfrak p \subset R$
l'homomorphisme d'anneaux $R \to A$ induit un isomorphisme
$R_{\mathfrak p} \to A_{\mathfrak q}$.
\item (Immersion ouverte.) Pour chaque prime $\mathfrak q$ de $A$ il existe un $f \in R$,
$\varphi(f) \not \in \mathfrak q$ tel que l'homomorphisme d'anneaux $\varphi : R \to A$
induit un isomorphisme $R_f \to A_f$.
\item Ajoutez-en plus ici au besoin\footnote{Mais seulement les propriétés qui ne sont pas déjà traitées séparément ailleurs.}.

\end{enumerate}

\end{lemma}

\begin{proof}
Démonstration omise.
\end{proof}

\begin{lemma}

\label{lemma-properties-composition}

Les propriétés suivantes des homomorphismes d'anneaux sont stables sous composition.

\begin{enumerate}

\item (Isomorphisme sur les anneaux locaux.) $\mathfrak q$ de $A$ couché sur $\mathfrak p \subset R$
l'homomorphisme d'anneaux $R \to A$ induit un isomorphisme
$R_{\mathfrak p} \to A_{\mathfrak q}$.
\item (Immersion ouverte.) Pour chaque prime $\mathfrak q$ de $A$ il existe un $f \in R$,
$\varphi(f) \not \in \mathfrak q$ tel que l'homomorphisme d'anneaux $\varphi : R \to A$
induit un isomorphisme $R_f \to A_f$.
\item (localement noethérien sur la cible.) L'homomorphisme d'anneaux $R \to A$ a la propriété qui $A$ C'est Noétherien.
\item Ajoutez-en plus ici au besoin\footnote{Mais seulement les propriétés qui ne sont pas déjà traitées séparément ailleurs.}.

\end{enumerate}

\end{lemma}

\begin{proof}
Démonstration omise.
\end{proof}








\section{Morphismes de type fini}
\label{section-finite-type}

\noindent
Rappelez-vous qu'un homomorphisme d'anneaux $R \to A$ est dit de type fini si $A$ est isomorphe à un quotient de $R[x_1, \ldots, x_n]$ comme algèbre $R$, voir Algèbre, définition \ref{algebra-definition-finite-type}.

\begin{definition}

\label{definition-finite-type}

Laisser $f : X \to S$ être un morphisme de schémas.

\begin{enumerate}

\item Nous disons que $f$ est de {\it type fini à $x \in X$} s'il existe un quartier affiné d'ouverture $\Spec(A) = U \subset X$
de $x$ et un ouvert affine $\Spec(R) = V \subset S$
avec $f(U) \subset V$ tel que l'homomorphisme d'anneaux induit
$R \to A$ est de type fini.
\item Nous disons que $f$ est {\it localement de type fini} s'il est de type fini à chaque point de $X$.
\item Nous disons que $f$ est de {\it type fini} s'il s'agit de localisation de type fini et quasi-compact.

\end{enumerate}

\end{definition}

\begin{lemma}

\label{lemma-locally-finite-type-characterize}

Laisser $f : X \to S$ morphisme des régimes. Les suivants sont équivalents

\begin{enumerate}

\item Le morphisme $f$ est localement de type fini.
\item Pour tous les ouverts affines $U \subset X$, $V \subset S$
avec $f(U) \subset V$ l'homomorphisme d'anneaux
$\mathcal{O}_S(V) \to \mathcal{O}_X(U)$ est de type fini.
\item Il existe un recouvrement ouvert $S = \bigcup_{j \in J} V_j$
et couvertures ouvertes $f^{-1}(V_j) = \bigcup_{i \in I_j} U_i$ de sorte que chacun des morphismes $U_i \to V_j$, $j\in J, i\in I_j$
est localement de type fini.
\item Il existe un recouvrement ouvert affine $S = \bigcup_{j \in J} V_j$
et revêtements pour affines d'ouverture $f^{-1}(V_j) = \bigcup_{i \in I_j} U_i$ tel que l'homomorphisme d'anneaux $\mathcal{O}_S(V_j) \to \mathcal{O}_X(U_i)$ est de type fini, pour tous $j\in J, i\in I_j$.

\end{enumerate}

En outre, si $f$ est localement de type fini puis pour tout sous-système ouvert $U \subset X$, $V \subset S$ avec $f(U) \subset V$
la restriction $f|_U : U \to V$ est localement de type fini.

\end{lemma}

\begin{proof}

Cela s'explique par le fait qu'il s'agit d'une \ref{lemma-locally-P} si nous montrons que la propriété ``$R \to A$ est de type fini'' est local. Nous vérifions les conditions (a), (b) et (c) de la définition
\ref{definition-property-local}. Par Algèbre, lemme \ref{algebra-lemma-base-change-finiteness}
être de type fini est stable sous changement de base et donc nous concluons (a) tient. Par Algèbre, lemme
\ref{algebra-lemma-compose-finite-type} être de type fini est stable sous composition et trivialement pour n'importe quel anneau
$R$ l'homomorphisme d'anneaux $R \to R_f$ est de type fini. Nous concluons (b) détient. Enfin, la propriété (c) est vraie selon Algèbre, lemme \ref{algebra-lemma-cover-upstairs}.

\end{proof}

\begin{lemma}

\label{lemma-composition-finite-type}
La composition de deux morphismes qui sont localement de type fini est localement de type fini. Il en va de même pour les morphismes de type fini.
\end{lemma}

\begin{proof}

Dans la démonstration de lemme \ref{lemma-locally-finite-type-characterize}
Nous avons vu qu'être de type fini est une propriété locale des homomorphismes d'anneaux. \ref{lemma-composition-type-P} combiné avec le fait qu'être de type fini est une propriété d'homomorphismes d'anneaux qui est stable sous composition, voir Algèbre, lemme \ref{algebra-lemma-compose-finite-type}. Par ce qui précède et le fait que les compositions des morphismes quasi-compacts sont quasi-compacts, voir Schémas, lemme \ref{schemes-lemma-composition-quasi-compact}
nous voyons que la composition des morphismes de type fini est de type fini.

\end{proof}

\begin{lemma}

\label{lemma-base-change-finite-type}

Le changement de base d'un morphisme qui est localement de type fini est localement de type fini. Il en est de même pour les morphismes de type fini.

\end{lemma}

\begin{proof}

Dans la démonstration de lemme \ref{lemma-locally-finite-type-characterize}
Nous avons vu qu'être de type fini est une propriété locale des homomorphismes d'anneaux. \ref{lemma-base-change-type-P} combiné avec le fait qu'être de type fini est une propriété d'homomorphismes d'anneaux stable sous changement de base, voir Algèbre, lemme \ref{algebra-lemma-base-change-finiteness}. Par ce qui précède et le fait qu'un changement de base d'un morphisme quasi-compact est quasi-compact, voir Schémas, lemme \ref{schemes-lemma-quasi-compact-preserved-base-change}
nous voyons que le changement de base d'un morphisme de type fini est un morphisme de type fini.

\end{proof}

\begin{lemma}

\label{lemma-immersion-locally-finite-type}
Une immersion fermée est de type fini. Une immersion est localement de type fini.
\end{lemma}

\begin{proof}

C'est vrai parce qu'une immersion ouverte est un isomorphisme local, et une immersion fermée est évidemment de type fini.

\end{proof}

\begin{lemma}

\label{lemma-finite-type-noetherian}
Laissez $f : X \to S$ être un morphisme. Si $S$ est (localement) Noétherien et $f$ (localement) de type fini alors $X$ est (localement) Noétherien.
\end{lemma}

\begin{proof}

Cela s'explique immédiatement par le fait qu'un anneau de type fini sur un anneau Noétherien est Noétherien, voir Algèbre, lemme \ref{algebra-lemma-Noetherian-permanence}. (Également: utiliser le fait que la source d'un morphisme quasi-compact avec la cible quasi-compact est quasi-compact.)

\end{proof}

\begin{lemma}

\label{lemma-finite-type-Noetherian-quasi-separated}
Laissez $f : X \to S$ être localement de type fini avec $S$ localement noethérien. Puis $f$ est quasi-séparé.
\end{lemma}

\begin{proof}

En fait, il est vrai que $X$ est quasi-séparée, voir Propriétés, lemme \ref{properties-lemma-locally-Noetherian-quasi-separated}
et lemme \ref{lemma-finite-type-noetherian} ci-dessus. Appliquer ensuite Schémas, lemme \ref{schemes-lemma-compose-after-separated}
de conclure que $f$ est quasi-séparée.

\end{proof}

\begin{lemma}

\label{lemma-permanence-finite-type}
Laissez $X \to Y$ être un morphisme de schémas sur un schéma de base $S$. Si $X$ est localement de type fini sur $S$, alors $X \to Y$ est localment de type fini.
\end{lemma}

\begin{proof}
Via lemme \ref{lemma-locally-finite-type-characterize} ceci se traduit par le fait d'algèbre suivant: Compte tenu des homomorphismes d'anneaux $A \to B \to C$ de telle sorte que $A \to C$ est de type fini, alors $B \to C$ est de type fini. (Voir Algèbre, lemme \ref{algebra-lemma-compose-finite-type}).
\end{proof}







\section{Points de type fini et schémas de Jacobson}
\label{section-points-finite-type}

\noindent
Laissez $S$ être un schéma. Un point de fin de type $s$ de $S$ est un point tel que le morphisme $\Spec(\kappa(s)) \to S$ est de type fini. La raison pour étudier ceci est que les points de fin de type peuvent remplacer les points fermés dans un certain sens et dans certaines situations. Il y en a toujours assez par exemple. De plus, un schéma est Jacobson si et seulement si tous les points de fin de type sont des points fermés.

\begin{lemma}

\label{lemma-point-finite-type}

Laisser $S$ Être un plan. $k$ être un champ. $f : \Spec(k) \to S$ être un morphisme. Les conditions suivantes sont équivalentes:

\begin{enumerate}

\item Le morphisme $f$ est de type fini.
\item Le morphisme $f$ est localement de type fini.
\item Il existe un ouvert affine $U = \Spec(R)$ de $S$
de telle manière que $f$ correspond à un homomorphisme fini d'anneaux $R \to k$.
\item Il existe un ouvert affine $U = \Spec(R)$ de $S$
de telle sorte que l'image $f$ se compose d'un point fermé $u$ en $U$
et l ' extension sur le terrain $k/\kappa(u)$ C'est fini.

\end{enumerate}

\end{lemma}

\begin{proof}
L'équivalence de (1) et (2) est évidente car $\Spec(k)$ est un singleton et donc tout morphisme de celui-ci est quasi-compact.

\medskip\noindent
Supposons que $f$ soit localement de type fini. Choisissez n'importe quel ouvert affine $\Spec(R) = U \subset S$ tel que l'image de $f$ soit contenue dans $U$, et l'homomorphisme d'anneaux $R \to k$ est de type fini. Laissez $\mathfrak p \subset R$ être le noyau. Alors $R/\mathfrak p \subset k$ est de type fini. Par Algèbre, lemme \ref{algebra-lemma-field-finite-type-over-domain} existe un $\overline{f} \in R/\mathfrak p$ tel que $(R/\mathfrak p)_{\overline{f}}$ est un champ et $(R/\mathfrak p)_{\overline{f}} \to k$ est une extension de champ fini. Si $f \in R$ est un ascenseur de $\overline{f}$, alors nous voyons que $k$ est un $R_f$-module fini. Ainsi (2) $\Rightarrow$ (3).

\medskip\noindent
Supposons que $\Spec(R) = U \subset S$ est une affine ouverte telle que $f$ correspond à un homomorphisme fini d'anneaux $R \to k$. Alors $f$ est localement de type fini par lemme \ref{lemma-locally-finite-type-characterize}. Ainsi (3) $\Rightarrow$ (2).

\medskip\noindent
Supposons que $R \to k$ soit fini. L'image de $R \to k$ est un champ sur lequel $k$ est fini par Algèbre, lemme \ref{algebra-lemma-integral-under-field}. Ainsi, le noyau de $R \to k$ est un idéal maximal. Ainsi (3) $\Rightarrow$ (4).

\medskip\noindent

Implications (4) $\Rightarrow$ 3) est immédiate.

\end{proof}

\begin{lemma}

\label{lemma-artinian-finite-type}
Laissez $S$ être un schéma. Laissez $A$ être un anneau artinien local avec corps résiduel $\kappa$. Laissez $f : \Spec(A) \to S$ être un morphisme de schémas. Alors $f$ est de type fini si et seulement si la composition $\Spec(\kappa) \to \Spec(A) \to S$ est de type fini.
\end{lemma}

\begin{proof}
Puisque le morphisme $\Spec(\kappa) \to \Spec(A)$ est de type fini, il est clair que si $f$ est de type fini ainsi que la composition $\Spec(\kappa) \to S$ (voir lemme \ref{lemma-composition-finite-type}). Pour l'inverse, notez que $\Spec(A) \to S$ mape dans un ouvert affine $U = \Spec(B)$ de $S$ comme $\Spec(A)$ n'a qu'un point. Pour terminer appliquer Algèbre, lemme \ref{algebra-lemma-essentially-of-finite-type-into-artinian-local} à $B \to A$.
\end{proof}

\noindent
Rappelez-vous que compte tenu d'un point $s$ d'un schéma $S$ il y a un morphisme canonique $\Spec(\kappa(s)) \to S$, voir Schémas, section \ref{schemes-section-points}.

\begin{definition}

\label{definition-finite-type-point}
Laissez $S$ être un schéma. Disons qu'un point $s$ de $S$ est un {\it type fini point} si le morphisme canonique $\Spec(\kappa(s)) \to S$ est de type fini. Nous dénotons $S_{\text{ft-pts}}$ le jeu de type fini points de $S$.
\end{definition}

\noindent
Nous pouvons décrire les points de type fini comme suit.

\begin{lemma}

\label{lemma-identify-finite-type-points}
Laissez $S$ être un schéma. Nous avons $$
S_{\text{ft-pts}} = \bigcup\nolimits_{U \subset S\text{ open}} U_0
$$ où $U_0$ est l'ensemble des points fermés de $U$. Ici nous pouvons laisser $U$ gamme sur toutes les ouvertures ou sur tous les ouverts affines de $S$.
\end{lemma}

\begin{proof}

Immédiatement de lemme \ref{lemma-point-finite-type}.

\end{proof}

\begin{lemma}

\label{lemma-finite-type-points-morphism}
Laissez $f : T \to S$ être un morphisme de schémas. Si $f$ est localement de type fini, alors $f(T_{\text{ft-pts}}) \subset S_{\text{ft-pts}}$.
\end{lemma}

\begin{proof}

Si $T$ est le spectre d'un champ ceci est lemme \ref{lemma-point-finite-type}. En général, il s'ensuit que la composition des morphismes localement de type fini est localement de type fini (lemme \ref{lemma-composition-finite-type}).

\end{proof}

\begin{lemma}

\label{lemma-finite-type-points-surjective-morphism}
Laissez $f : T \to S$ être un morphisme de schémas. Si $f$ est localement de type fini et surjectif, alors $f(T_{\text{ft-pts}}) = S_{\text{ft-pts}}$.
\end{lemma}

\begin{proof}
Nous avons $f(T_{\text{ft-pts}}) \subset S_{\text{ft-pts}}$ par lemme \ref{lemma-finite-type-points-morphism}. Laissez $s \in S$ être un point fini de type. Comme $f$ est surjectif, le schéma $T_s = \Spec(\kappa(s)) \times_S T$ n'est pas vide, donc a un point fini de type $t \in T_s$ par lemme \ref{lemma-identify-finite-type-points}. Maintenant $T_s \to T$ est un morphisme de type fini comme un changement de base de $s \to S$ (lemme \ref{lemma-base-change-finite-type}). D'où l'image de $t$ dans $T$ est un point fini de type par lemme \ref{lemma-finite-type-points-morphism} qui cartographie $s$ par construction.
\end{proof}

\begin{lemma}

\label{lemma-enough-finite-type-points}
Laissez $S$ être un schéma. Pour tout sous-ensemble fermé localement $T \subset S$ nous avons $$
T \not = \emptyset
\Rightarrow
T \cap S_{\text{ft-pts}} \not = \emptyset.
$$ En particulier, pour tout sous-ensemble fermé $T \subset S$ nous voyons que $T \cap S_{\text{ft-pts}}$ est dense en $T$.
\end{lemma}

\begin{proof}
Notez que $T$ comporte une structure de schéma (voir Schémas, lemme \ref{schemes-lemma-reduced-closed-subscheme}) telle que $T \to S$ est une immersion fermée locale. Toute fermée d'immersion locale est localement de type fini, voir lemme \ref{lemma-immersion-locally-finite-type}. Par conséquent par lemme \ref{lemma-finite-type-points-morphism} nous voyons $T_{\text{ft-pts}} \subset S_{\text{ft-pts}}$. Enfin, toute ouvert affine non vide de $T$ a au moins un point fermé qui est un point de type fini de $T$ par lemme \ref{lemma-identify-finite-type-points}.
\end{proof}

\noindent
Il s'ensuit que la plupart des matériaux de Topologie, la section \ref{topology-section-space-jacobson} passe par l'ensemble des points fermés remplacés par l'ensemble des points de type fini. En fait, si $S$ est Jacobson, nous récupérons les points fermés comme les points de type fini.

\begin{lemma}

\label{lemma-jacobson-finite-type-points}
Laissez $S$ être un schéma. Les conditions suivantes sont équivalentes:
\begin{enumerate}
\item le schéma $S$ est Jacobson,
\item $S_{\text{ft-pts}}$ est l'ensemble des points fermés de $S$,
\item pour tous les $T \to S$ localement de type finis carte de points fermés à des points fermés, et
\item pour tous $T \to S$ localement de type finis carte $t \in T$ $s \in S$ avec $\kappa(s) \subset \kappa(t)$ fini.
\end{enumerate}

\end{lemma}

\begin{proof}

Nous avons trivialement (4) $\Rightarrow$ (3) $\Rightarrow$ (2). \ref{lemma-enough-finite-type-points} montre que (2) implique (1). Il suffit donc de montrer que (1) implique (4). Supposons que $T \to S$ est localement de type fini. Choisissez $t \in T$ fermé et laisser $s \in S$ soit l'image. Choisissez les quartiers affinés d'ouverture $\Spec(R) = U \subset S$ de $s$ et
$\Spec(A) = V \subset T$ de $t$ avec $V$ cartographie dans $U$. L'homomorphisme d'anneaux induit $R \to A$ est de type fini (voir lemme \ref{lemma-locally-finite-type-characterize}) et $R$ est Jacobson par Propriétés, lemme \ref{properties-lemma-locally-jacobson}. Ainsi, le résultat suit d'Algèbre, proposition \ref{algebra-proposition-Jacobson-permanence}.

\end{proof}

\begin{lemma}

\label{lemma-Jacobson-universally-Jacobson}
Laissez $S$ être un schéma de Jacobson. Tout schéma de type de fini sur $S$ est Jacobson.
\end{lemma}

\begin{proof}

C'est clair d'Algèbre, proposition \ref{algebra-proposition-Jacobson-permanence}
(et Propriétés, \ref{properties-lemma-locally-jacobson} et lemme \ref{lemma-locally-finite-type-characterize}).

\end{proof}

\begin{lemma}

\label{lemma-ubiquity-Jacobson-schemes}
Les types de schémas suivants sont Jacobson.
\begin{enumerate}
\item Tout schéma local de type fini sur un champ.
\item Tout schéma local de type fini sur $\mathbf{Z}$.
\item Tout schéma local de type fini sur un $1$-dimensionnel domaine Noéthérien avec infiniment de primes.
\item Un schéma de la forme $\Spec(R) \setminus \{\mathfrak m\}$ où $(R, \mathfrak m)$ est un local anneau Noéthérien. Aussi tout schéma local de type fini sur lui.
\end{enumerate}

\end{lemma}

\begin{proof}
Nous allons utiliser lemme \ref{lemma-Jacobson-universally-Jacobson} sans mention. Le spectre d'un champ est clairement Jacobson. Le spectre de $\mathbf{Z}$ est Jacobson, voir Algèbre, lemme \ref{algebra-lemma-pid-jacobson}. Pour (3) voir Algèbre, lemme \ref{algebra-lemma-noetherian-dim-1-Jacobson}. Pour (4) voir Propriétés, lemme \ref{properties-lemma-complement-closed-point-Jacobson}.
\end{proof}







\section{Schémas universellement caténaires}
\label{section-universally-catenary}

\noindent
Rappelez qu'un espace topologique $X$ est appelé {\it caténaire} si pour chaque paire de sous-ensembles fermés irréductibles $T \subset T'$ il existe une chaîne maximale de sous-ensembles fermés irréductibles $$
T = T_0 \subset T_1 \subset \ldots \subset T_e = T'
$$ et chaque chaîne a la même longueur. Voir Topologie, définition \ref{topology-definition-catenary}. Rappelez qu'un schéma est caténaire si son espace topologique sous-jacent est caténaire. Voir Propriétés, définition \ref{properties-definition-catenary}.

\begin{definition}

\label{definition-universally-catenary}
Laissez $S$ être un schéma. Supposons que $S$ est localement noethérien. Nous disons que $S$ est {\it universellement caténaire} si pour chaque morphisme $X \to S$ localement de type fini le schéma $X$ est caténaire.
\end{definition}

\noindent
C'est une notion "mieux" que caténaire car il existe des schémas noéthériens qui sont caténaires mais pas universellement caténaire. Voir exemple, section \ref{examples-section-non-catenary-Noetherian-local}. Beaucoup de schémas sont universellement caténaire, voir lemme \ref{lemma-ubiquity-uc} ci-dessous.

\medskip\noindent

Rappelez-vous qu'une bague $A$ est appelé {\it caténaire} si pour une paire d'idéaux principaux $\mathfrak p \subset \mathfrak q$
il existe une chaîne maximale de primes
$$
\mathfrak p =
\mathfrak p_0 \subset \ldots \subset \mathfrak p_e
= \mathfrak q
$$
et tous ont la même longueur. Voir Algèbre, définition \ref{algebra-definition-catenary}. Nous avons vu la relation entre les schémas caténaires et les anneaux caténaires dans Propriétés, section \ref{properties-section-catenary}. Rappelez-vous qu'une bague $A$ est appelé {\it universalement caténaire} si
$A$ est noethérien et pour chaque type de fini homomorphisme d'anneaux $A \to B$
la bague $B$ Voir Algèbre, définition \ref{algebra-definition-universally-catenary}. Beaucoup d'anneaux intéressants qui viennent dans la géométrie algébrique satisfaire cette propriété.

\begin{lemma}

\label{lemma-universally-catenary-local}

Laisser $S$ soit un régime localement noethérien.

\begin{enumerate}

\item  $S$ est universellement caténaire,
\item il existe un recouvrement ouvert de $S$ dont tous les membres sont des schémas universels caténaires,
\item pour chaque ouvert affine $\Spec(R) = U \subset S$ la bague
$R$ est universellement caténaire, et
\item il existe un recouvrement ouvert affine $S = \bigcup U_i$ de telle sorte que chacun $U_i$ est le spectre d'un anneau universel caténaire.

\end{enumerate}

En outre, dans ce cas, tout régime de localisation de type fini sur $S$
est universellement caténaire aussi bien.

\end{lemma}

\begin{proof}
Par lemme \ref{lemma-immersion-locally-finite-type} une immersion ouverte est localement de type fini. Une composition de morphismes localement de type fini est localment de type fini (lemme \ref{lemma-composition-finite-type}). Ainsi, il est clair que si $S$ est universellement caténaire alors tout schéma ouvert et tout schéma local de type fini sur $S$ est aussi universellement caténaire. Ceci prouve l'énoncé final de la lemme et que (1) implique (2).

\medskip\noindent
Si $\Spec(R)$ est un schéma universel caténaire, alors chaque schéma $\Spec(A)$ avec $A$ un type fini $R$-algèbre est caténaire. Ainsi, tous ces anneaux $A$ sont caténaires par Algèbre, lemme \ref{algebra-lemma-catenary}. Ainsi $R$ est universellement caténaire. Combiné avec les remarques ci-dessus, nous concluons que (1) implique (3), et (2) implique (4). Bien sûr (3) implique (4) trivialement.

\medskip\noindent
Pour terminer la démonstration nous montrons que (4) implique (1). Supposons que (4) et laisser $X \to S$ être un morphisme localement de type fini. Nous pouvons trouver un recouvrement ouvert affine $X = \bigcup V_j$ de telle sorte que chaque $V_j \to S$ cartographie dans l'un des $U_i$. Par lemme \ref{lemma-locally-finite-type-characterize} l'homomorphisme d'anneaux $\mathcal{O}(U_i) \to \mathcal{O}(V_j)$ induit est de type fini. Par conséquent, $\mathcal{O}(V_j)$ est caténaire. Par conséquent, $X$ est caténaire par Propriétés, lemme \ref{properties-lemma-catenary-local}.
\end{proof}

\begin{lemma}

\label{lemma-universally-catenary-local-rings-universally-catenary}
Laissez $S$ être un schéma de localisation noéthérien. Les conditions suivantes sont équivalentes:
\begin{enumerate}
\item $S$ est universellement caténaire, et
\item tous les anneaux locaux $\mathcal{O}_{S, s}$ de $S$ sont universellement caténaire.
\end{enumerate}

\end{lemma}

\begin{proof}
Supposons que tous les anneaux locaux de $S$ sont universels caténaires. Laissez $f : X \to S$ être localement de type fini. Nous savons que $X$ est caténaire si et seulement si $\mathcal{O}_{X, x}$ est caténaire pour tous $x \in X$. Si $f(x) = s$, alors $\mathcal{O}_{X, x}$ est essentiellement de type fini sur $\mathcal{O}_{S, s}$. Par conséquent, $\mathcal{O}_{X, x}$ est caténaire par l'hypothèse que $\mathcal{O}_{S, s}$ est universalement caténaire.

\medskip\noindent

À l'inverse, supposer que $S$ est l'universellement caténaire. $s \in S$. Nous pouvons remplacer $S$ par un quartier affiné d'ouverture de $s$ par lemme \ref{lemma-universally-catenary-local}. Dites $S = \Spec(R)$
et $s$ correspond à l'idéal premier $\mathfrak p$. Tout type de fini $R_{\mathfrak p}$-Algèbre $A'$ est de la forme
$A_{\mathfrak p}$ pour certains types de fini $R$-Algèbre $A$. Par hypothèse (et \ref{lemma-universally-catenary-local} si vous voulez) la bague $A$ est caténaire, et donc $A'$ (une localisation de $A$) est caténaire. $R_{\mathfrak p}$ c'est l'universellement caténaire.

\end{proof}

\begin{lemma}

\label{lemma-catenary-check-irreducible}
Laissez $S$ être un schéma localement noethérien. Alors $S$ est universellement caténaire si et seulement si les composantes irréductibles de $S$ sont universellement caténaire.
\end{lemma}

\begin{proof}
Démonstration omise. Pour l'affaire affine, voir Algèbre, lemme \ref{algebra-lemma-catenary-check-irreducible}.
\end{proof}

\begin{lemma}

\label{lemma-ubiquity-uc}
Les types de schémas suivants sont universels caténaires.
\begin{enumerate}
\item Tout schéma local de type fini sur un champ.
\item Tout schéma local de type fini sur un domaine de Cohen--Macaulay.
\item Tout schéma local de type fini sur $\mathbf{Z}$.
\item Tout schéma local de type fini sur un domaine de $1$-dimensionnel noethérien.
\item Et ainsi de suite.
\end{enumerate}

\end{lemma}

\begin{proof}
Tout cela découle du fait qu'une bague de Cohen--Macaulay est universellement caténaire, voir Algèbre, lemme \ref{algebra-lemma-CM-ring-catenary}. Aussi, utilisez la dernière affirmation de lemme \ref{lemma-universally-catenary-local}. Quelques détails omis.
\end{proof}























\section{Schémas de Nagata, reprise}
\label{section-nagata}

\noindent
Voir Propriétés, section \ref{properties-section-nagata} pour les définitions et les propriétés de base des schémas Nagata et universellement japonais.

\begin{lemma}

\label{lemma-finite-type-nagata}
Laissez $f : X \to S$ être un morphisme. Si $S$ est Nagata et $f$ localement de type fini alors $X$ est Nagata. Si $S$ est universellement japonais et $f$ localement de type fini alors $X$ est universellement japonais.
\end{lemma}

\begin{proof}
Pour ``universellement japonais'' cela découle d'Algèbre, lemme \ref{algebra-lemma-universally-japanese}. Pour ``Nagata'' cela découle d'Algèbre, proposition \ref{algebra-proposition-nagata-universally-japanese}.
\end{proof}

\begin{lemma}

\label{lemma-ubiquity-nagata}
Les types de schémas suivants sont Nagata.
\begin{enumerate}
\item Tout schéma local de type fini sur un champ.
\item Tout schéma local de type fini sur un anneau local complet noethérien.
\item Tout schéma local de type fini sur $\mathbf{Z}$.
\item Tout schéma local de type fini sur un anneau Dedekind de zéro caractéristique.
\item Et ainsi de suite.
\end{enumerate}

\end{lemma}

\begin{proof}
Par lemme \ref{lemma-finite-type-nagata} nous n'avons qu'à montrer que les anneaux mentionnés ci-dessus sont des anneaux Nagata. Pour cela voir Algèbre, proposition \ref{algebra-proposition-ubiquity-nagata}.
\end{proof}


\section{Le lieu singulier, reprise}
\label{section-singular-locus}

\noindent
Nous cherchons un critère qui implique l'ouverture du locus régulier pour n'importe quel schéma localement de type fini sur la base. Voici la définition.

\begin{definition}

\label{definition-J}
Laissez $X$ être un schéma de localisation noethérien. Nous disons que $X$ est {\it J-2} si pour chaque morphisme $Y \to X$ qui est localement de type fini le locus régulier $\text{Reg}(Y)$ est ouvert dans $Y$.
\end{definition}

\noindent
C'est l'analogue de la notion correspondante pour les anneaux Noétheriens, voir Compléments sur l'algèbre, définition \ref{more-algebra-definition-J}.

\begin{lemma}

\label{lemma-J}

Laisser $X$ soit un régime localement noethérien.

\begin{enumerate}

\item  $X$ est J-2,
\item il existe un recouvrement ouvert de $X$ dont tous les membres sont des régimes J-2,
\item pour chaque ouvert affine $\Spec(R) = U \subset X$ la bague
$R$ est J-2, et
\item il existe un recouvrement ouvert affine $S = \bigcup U_i$ de telle sorte que chacun $\mathcal{O}(U_i)$ est J-2 pour tous $i$.

\end{enumerate}

En outre, dans ce cas, tout régime de localisation de type fini sur $X$
C'est J-2 aussi.

\end{lemma}

\begin{proof}
Par lemme \ref{lemma-immersion-locally-finite-type} une immersion ouverte est localement de type fini. Une composition de morphismes localement de type fini est localment de type fini (lemme \ref{lemma-composition-finite-type}). Ainsi, il est clair que si $X$ est J-2 alors tout schéma ouvert localement de type fini sur $X$ est aussi J-2. Ceci prouve l'énoncé finale de la lemme.

\medskip\noindent
Si $\Spec(R)$ est J-2, alors pour chaque type fini $R$-algèbre $A$ le locus régulier du schéma $\Spec(A)$ est ouvert. Par conséquent, $R$ est J-2, par définition (voir Compléments sur l'algèbre, définition \ref{more-algebra-definition-J}). Combiné avec les remarques ci-dessus nous concluons que (1) implique (3), et (2) implique (4). Bien sûr (1) $\Rightarrow$ (2) et (3) $\Rightarrow$ (4) trivialement.

\medskip\noindent
Pour terminer la démonstration nous montrons que (4) implique (1). Supposons que (4) et laisser $Y \to X$ être un morphisme localement de type fini. Nous pouvons trouver un recouvrement ouvert affine $Y = \bigcup V_j$ de telle sorte que chaque $V_j \to X$ cartographie dans l'un des $U_i$. Par lemme \ref{lemma-locally-finite-type-characterize} l'homomorphisme d'anneaux $\mathcal{O}(U_i) \to \mathcal{O}(V_j)$ induit est de type fini. D'où le locus régulier de $V_j = \Spec(\mathcal{O}(V_j))$ est ouvert. Depuis $\text{Reg}(Y) \cap V_j = \text{Reg}(V_j)$ nous concluons que $\text{Reg}(Y)$ est ouvert comme désiré.
\end{proof}

\begin{lemma}

\label{lemma-ubiquity-J-2}

Les types de régimes suivants sont J-2.

\begin{enumerate}

\item Tout système de localisation de type fini sur un champ.
\item Tout schéma localement de type fini sur un anneau complet noethérien local.
\item Tout système de localisation de type fini sur $\mathbf{Z}$.
\item Tout schéma localement de type fini sur un anneau noétherien local de dimension $1$.
\item Tout schéma localement de type fini sur un anneau de dimension Nagata $1$.
\item Tout schéma localement de type fini sur un anneau Dedekind de zéro caractéristique.
\item Et ainsi de suite.

\end{enumerate}

\end{lemma}

\begin{proof}
Par lemme \ref{lemma-J} nous n'avons qu'à montrer que les anneaux mentionnés ci-dessus sont J-2. Pour cela voir Compléments sur l'algèbre, proposition \ref{more-algebra-proposition-ubiquity-J-2}.
\end{proof}




\section{Schémas excellents}
\label{section-excellent}

\noindent
Nous nous rappelons qu'un anneau est quasi-excellent s'il s'agit d'un anneau G et J-2. Un anneau est excellent s'il s'agit d'un anneau quasi-excellent et universel caténaire.

\begin{definition}

\label{definition-excellent}

Laisser $X$ être un régime.

\begin{enumerate}

\item Nous disons $X$ est {\it quasi-excellent} si pour chaque $x \in X$ il existe un quartier affiné d'ouverture $x \in U \subset X$ de telle sorte que l'anneau
$\mathcal{O}_X(U)$ est quasi-excellent (voir Compléments sur l'algèbre, définition \ref{more-algebra-definition-excellent}).
\item Nous disons $X$ est {\it excellent} si pour chaque $x \in X$ il existe un quartier affiné d'ouverture $x \in U \subset X$ de telle sorte que l'anneau
$\mathcal{O}_X(U)$ est excellent (voir Compléments sur l'algèbre, définition \ref{more-algebra-definition-excellent}).

\end{enumerate}

\end{definition}

\noindent
Quasi-schémas excellents sont localement noethérien (les anneaux G sont noethériens).

\begin{lemma}

\label{lemma-characterize-quasi-excellent}
Laissez $X$ être un schéma. Les conditions suivantes sont équivalentes
\begin{enumerate}
\item $X$ est quasi-excellent, et
\item $X$ est un schéma G et J-2.
\end{enumerate}

\end{lemma}

\begin{proof}

Implications (1) $\Rightarrow$ (2) suit sur la combinaison de la définition de quasi-schémas excellents (définition \ref{definition-excellent}), la définition des anneaux quasi-excellents (Compléments sur l'algèbre, définition \ref{more-algebra-definition-excellent}), la définition des régimes G (Propriétés, définition) \ref{properties-definition-G-scheme}), et lemme \ref{lemma-J}. Inversement, si $X$ est un schéma G et J-2, puis pour n'importe quel $x \in X$
nous pouvons trouver un quartier affiné ouvert $x \in U \subset X$
de telle manière que $\mathcal{O}_X(U)$ C'est un anneau G. \ref{lemma-J}
la bague $\mathcal{O}_X(U)$ est aussi J-2, d'où quasi-excellent.

\end{proof}

\begin{lemma}

\label{lemma-quasi-excellent-nagata}
Un quasi-schéma excellent est Nagata.
\end{lemma}

\begin{proof}
Voir Compléments sur l'algèbre, lemme \ref{more-algebra-lemma-quasi-excellent-nagata}.
\end{proof}

\begin{lemma}

\label{lemma-characterize-excellent}
Laissez $X$ être un schéma. Les conditions suivantes sont équivalentes
\begin{enumerate}
\item $X$ est excellent, et
\item $X$ est quasi-excellent et universel caténaire.
\end{enumerate}

\end{lemma}

\begin{proof}
Supposons (1). Comme d'excellents anneaux sont quasi-excellents, il est clair que $X$ est quasi-excellent. Puisque d'excellents anneaux sont universellement caténaires, il s'ensuit également que $X$ est universellement caténaire par lemme \ref{lemma-universally-catenary-local}. Supposons que (2). Laisser $x \in X$. Puisque $X$ est quasi-excellent, il existe un quartier affiné ouvert $x \in U \subset X$ de telle sorte que $\mathcal{O}_X(U)$ est quasi-excellent. Puisque $X$ est une catégorie universelle, l'anneau $\mathcal{O}_X(U)$ est également universellement caténaire par lemme \ref{lemma-universally-catenary-local}, où Excellent.
\end{proof}

\begin{lemma}

\label{lemma-locally-excellent}

Laisser $X$ Les conditions suivantes sont équivalentes:

\begin{enumerate}

\item Le régime $X$ est (quasi-) excellent.
\item Pour chaque ouvert affine $U \subset X$ la bague $\mathcal{O}_X(U)$
est (quasi-) excellent.
\item Il existe un recouvrement ouvert affine $X = \bigcup U_i$ de telle sorte que chacun $\mathcal{O}_X(U_i)$ est (quasi-) excellent.
\item Il existe une ouverture de recouvrement $X = \bigcup X_j$
de telle sorte que chaque sous-système ouvert $X_j$ est (quasi-) excellent.

\end{enumerate}

En outre, si $X$ est (quasi-)excellent alors chaque sous-système ouvert est (quasi-) excellent.

\end{lemma}

\begin{proof}
Par lemme \ref{lemma-characterize-quasi-excellent} être quasi-excellent est le même qu'être un schéma G et J-2. Ainsi, pour le cas quasi-excellent, le lemme suit de lemme \ref{lemma-J} et Propriétés, lemme \ref{properties-lemma-locally-G-scheme}. Pour l'excellent cas, utilisez lemme \ref{lemma-characterize-excellent}, le cas quasi-excellent, et lemme \ref{lemma-universally-catenary-local}.
\end{proof}

\begin{lemma}

\label{lemma-finite-type-excellent}
Laissez $f : X \to S$ être un morphisme. Si $S$ est (quasi-)excellent et $f$ localement de type fini alors $X$ est (quasi-)excellent.
\end{lemma}

\begin{proof}
Voir Compléments sur l'algèbre, lemme \ref{more-algebra-lemma-finite-type-over-excellent}.
\end{proof}

\begin{lemma}

\label{lemma-ubiquity-excellent}
Les types de schémas suivants sont excellents.
\begin{enumerate}
\item Tout schéma local de type fini sur un champ.
\item Tout schéma local de type fini sur un anneau complet Noétherien.
\item Tout schéma local de type fini sur $\mathbf{Z}$.
\item Tout schéma local de type fini sur un anneau Dedekind de zéro caractéristique.
\end{enumerate}

\end{lemma}

\begin{proof}
Par lemmes \ref{lemma-locally-excellent} et \ref{lemma-finite-type-excellent} nous n'avons qu'à montrer que les anneaux mentionnés ci-dessus sont excellents. Pour cela voir Compléments sur l'algèbre, proposition \ref{more-algebra-proposition-ubiquity-excellent}.
\end{proof}






\section{Morphismes quasi-finis}
\label{section-quasi-finite}

\noindent
Un traitement solide des morphismes quasi-finis est à la base de nombreux développements plus loin dans la route. Il conduira à différentes versions du Théorème principal de Zariski, le comportement des dimensions des fibres, la descente des morphismes \'etale, etc. Avant de lire cette section, il peut être une bonne idée de jeter un coup d'oeil aux résultats algèbres dans la section \ref{algebra-section-quasi-finite}.

\medskip\noindent
Rappelons qu'un type fini homomorphisme d'anneaux $R \to A$ est quasi-fini à un $\mathfrak q$ initial si $\mathfrak q$ définit un point isolé de sa fibre, voir Algèbre, définition \ref{algebra-definition-quasi-finite}.

\begin{definition}

\label{definition-quasi-finite}

\begin{reference}

\cite[II définition 6.2.3]{EGA}

\end{reference}

Laisser $f : X \to S$ être un morphisme de schémas.

\begin{enumerate}

\item Nous disons que $f$ est {\it quasi-fini à un point $x \in X$} s'il existe un quartier affiné $\Spec(A) = U \subset X$
de $x$ et un ouvert affine $\Spec(R) = V \subset S$ de telle manière que
$f(U) \subset V$, l'homomorphisme d'anneaux $R \to A$ est de type fini, et $R \to A$ est quasi-finie au sommet de $A$ correspondant à $x$
(voir ci-dessus).
\item Nous disons $f$ est {\it localement quasi-fini} si $f$ est quasi-fini à chaque point $x$ de $X$.
\item Nous disons que $f$ est {\it quasi-fini} si $f$ est de type fini et chaque point $x$ est un point isolé de sa fibre.

\end{enumerate}

\end{definition}

\noindent
Triviellement, un morphisme quasi-fini local est localement de type fini. Nous verrons ci-dessous qu'un morphisme $f$ qui est localement de type fini est quasi-fini à $x$ si et seulement si $x$ est isolé dans sa fibre. De plus, l'ensemble de points où un morphisme est quasi-fini est ouvert; nous verrons cela dans la section \ref{section-Zariski} sur le théorème principal de Zariski.

\begin{lemma}

\label{lemma-algebraic-residue-field-extension-closed-point-fibre}
Laissez $f : X \to S$ être un morphisme de schémas. Laissez $x \in X$ être un point. Posons $s = f(x)$. Si $\kappa(x)/\kappa(s)$ est une extension de champ algébrique, alors
\begin{enumerate}
\item $x$ est un point fermé de sa fibre, et
\item si en plus $s$ est un point fermé de $S$, alors $x$ est un point fermé de $X$.
\end{enumerate}

\end{lemma}

\begin{proof}

La deuxième affirmation suit de la première par topologie élémentaire. Selon Schémas, lemme \ref{schemes-lemma-fibre-topological}
pour prouver la première déclaration que nous pouvons remplacer $X$ par $X_s$ et $S$ par $\Spec(\kappa(s))$. Ainsi, nous pouvons supposer que $S = \Spec(k)$ est le spectre d'un champ. Dans ce cas, laissez $\Spec(A) = U \subset X$ soit toute ouvert affine contenant: $x$. Le point $x$ correspond à un idéal premier
$\mathfrak q \subset A$ de telle manière que $\kappa(\mathfrak q)/k$
est une extension de champ algébrique. \ref{algebra-lemma-finite-residue-extension-closed}
nous voyons que $\mathfrak q$ est un idéal maximal, c'est-à-dire, $x \in U$ est un point fermé. Puisque les ouverts affines forment une base de la topologie de $X$ nous concluons que $\{x\}$ C'est fermé.

\end{proof}

\noindent
Le lemme suivant est une version du Hilbert Nullstellinsatz.

\begin{lemma}

\label{lemma-closed-point-fibre-locally-finite-type}
Laissez $f : X \to S$ être un morphisme de schémas. Laissez $x \in X$ être un point. Posons $s = f(x)$. Supposons que $f$ est un emplacement de type fini. Alors $x$ est un point fermé de sa fibre si et seulement si $\kappa(x)/\kappa(s)$ est une extension de champ fini.
\end{lemma}

\begin{proof}

Si l'extension est finie, alors $x$ est un point fermé de la fibre par lemme \ref{lemma-algebraic-residue-field-extension-closed-point-fibre}
ci-dessus. Pour l'inverse, supposez $x$ est un point fermé de sa fibre. Choisissez ouverts affines $\Spec(A) = U \subset X$
et $\Spec(R) = V \subset S$ de telle manière que $f(U) \subset V$. Par lemme \ref{lemma-locally-finite-type-characterize} l'homomorphisme d'anneaux
$R \to A$ est de type fini. $\mathfrak q \subset A$, rép.\ $\mathfrak p \subset R$ être l'idéal premier correspondant à $x$, rép.\ $s$. Considérez la bague en fibres
$\overline{A} = A \otimes_R \kappa(\mathfrak p)$. Laisse $\overline{\mathfrak q}$ être le meilleur de $\overline{A}$
correspondant à $\mathfrak q$. L'hypothèse que $x$
est un point fermé de sa fibre implique que $\overline{\mathfrak q}$
est un idéal maximal de $\overline{A}$Depuis $\overline{A}$
est un algèbre de type fini sur le champ $\kappa(\mathfrak p)$
nous voyons par le Hilbert Nullstellinsatz, voir Algèbre, théorème \ref{algebra-theorem-nullstellensatz}, que $\kappa(\overline{\mathfrak q})$ est une extension finie de $\kappa(\mathfrak p)$Depuis $\kappa(s) = \kappa(\mathfrak p)$ et
$\kappa(x) = \kappa(\mathfrak q) = \kappa(\overline{\mathfrak q})$
Nous avons gagné.

\end{proof}

\begin{lemma}

\label{lemma-base-change-closed-point-fibre-locally-finite-type}
Laissez $f : X \to S$ être un morphisme de schémas qui est localement de type fini. Laissez $g : S' \to S$ être n'importe quel morphisme. Notez $f' : X' \to S'$ le changement de base. Si $x' \in X'$ mappe à un point $x \in X$ qui est fermé dans $X_{f(x)}$ alors $x'$ est fermé dans $X'_{f'(x')}$.
\end{lemma}

\begin{proof}

Le corps résiduel $\kappa(x')$ est un quotient de
$\kappa(f'(x')) \otimes_{\kappa(f(x))} \kappa(x)$, voir Schémas, lemme \ref{schemes-lemma-points-fibre-product}. C'est donc une extension finie de $\kappa(f'(x'))$ comme
$\kappa(x)$ est une extension finie de $\kappa(f(x))$ par lemme \ref{lemma-closed-point-fibre-locally-finite-type}. Ainsi nous voyons que $x'$ est fermé dans sa fibre en appliquant une fois de plus ce lemme.

\end{proof}

\begin{lemma}

\label{lemma-residue-field-quasi-finite}
Laissez $f : X \to S$ être un morphisme de schémas. Laissez $x \in X$ être un point. Posons $s = f(x)$. Si $f$ est quasi-fini à $x$, alors l'extension du corps résiduel $\kappa(x)/\kappa(s)$ est fini.
\end{lemma}

\begin{proof}
Ceci est clair d'Algèbre, définition \ref{algebra-definition-quasi-finite}.
\end{proof}

\begin{lemma}

\label{lemma-quasi-finite-at-point-characterize}

Laisser $f : X \to S$ être un morphisme de schémas. $x \in X$ Soyez un point. $s = f(x)$. Laisse $X_s$ soit la fibre de $f$ au $s$Supposons que $f$ est localement de type fini. Les conditions suivantes sont équivalentes:

\begin{enumerate}

\item Le morphisme $f$ est quasi-fini à $x$.
\item Le point $x$ est isolé dans $X_s$.
\item Le point $x$ est fermé en $X_s$
et il n'y a pas de raison $x' \in X_s$, $x' \not = x$
qui se spécialise dans $x$.
\item Pour n'importe quelle paire d'ouverts affines
$\Spec(A) = U \subset X$, $\Spec(R) = V \subset S$ avec
$f(U) \subset V$ et $x \in U$ correspondant à $\mathfrak q \subset A$
l'homomorphisme d'anneaux $R \to A$ est quasi-fini à $\mathfrak q$.

\end{enumerate}

\end{lemma}

\begin{proof}
Supposons que $f$ est quasi-fini à $x$. Par hypothèse, il existe des ouvertures $U \subset X$, $V \subset S$ tel que $f(U) \subset V$, $x \in U$ et $x$ un point isolé de $U_s$. Par conséquent, $\{x\} \subset U_s$ est un sous-ensemble ouvert. Puisque $U_s = U \cap X_s \subset X_s$ est également ouvert, nous concluons que $\{x\} \subset X_s$ est un sous-ensemble ouvert également. Ainsi nous concluons que $x$ est un point isolé de $X_s$.

\medskip\noindent
Notez que $X_s$ est un schéma de Jacobson par lemme \ref{lemma-ubiquity-Jacobson-schemes} (et lemme \ref{lemma-base-change-finite-type}). Si $x$ est isolé dans $X_s$, c'est-à-dire $\{x\} \subset X_s$ est ouvert, alors $\{x\}$ contient un point fermé (par la propriété Jacobson), donc $x$ est fermé dans $X_s$. Il est clair qu'il n'y a pas de point $x' \in X_s$, distinct de $x$, spécialisé dans $x$.

\medskip\noindent

Supposons que $x$ est fermé en $X_s$ et qu'il n'y a pas de raison $x' \in X_s$, distinctes des $x$, spécialisée dans $x$. Considérez une paire d'ouverts affines
$\Spec(A) = U \subset X$, $\Spec(R) = V \subset S$ avec
$f(U) \subset V$ et $x \in U$. Laisse $\mathfrak q \subset A$ correspond à
$x$ et $\mathfrak p \subset R$ correspond à $s$. Par lemme \ref{lemma-locally-finite-type-characterize} l'homomorphisme d'anneaux
$R \to A$ est de type fini. Considérez l'anneau de fibre
$\overline{A} = A \otimes_R \kappa(\mathfrak p)$. Laisse $\overline{\mathfrak q}$ être le meilleur de $\overline{A}$ correspondant à $\mathfrak q$Depuis $\Spec(\overline{A})$ est un sous-système ouvert de la fibre $X_s$ nous voyons que $\overline{q}$ est un idéal maximal de $\overline{A}$ et qu'il n'y a aucune raison de $\Spec(\overline{A})$
spécialisée dans $\overline{\mathfrak q}$. Cela implique que $\dim(\overline{A}_{\overline{q}}) = 0$. D'où par Algèbre, définition \ref{algebra-definition-quasi-finite}
nous voyons que $R \to A$ est quasi-fini à $\mathfrak q$, i.e.,
$X \to S$ est quasi-fini à $x$ par définition.

\medskip\noindent

Il est clair que (4) implique (1). Et il est aussi clair que (2) implique (4) puisque si $x$ est un point isolé de $X_s$
alors c'est aussi un point isolé de $U_s$ pour toute ouverture $U$
qui le contient.

\end{proof}

\begin{lemma}

\label{lemma-finite-fibre}
Laissez $f : X \to S$ être un morphisme de schémas. Laissez $s \in S$. Supposez que
\begin{enumerate}
\item $f$ est un emplacement de type fini, et
\item $f^{-1}(\{s\})$ est un ensemble fini.
\end{enumerate}
Ensuite $X_s$ est un espace topologique discret fini, et $f$ est quasi-fini à chaque point de $X$ situé sur $s$.
\end{lemma}

\begin{proof}
Supposons que $T$ est un schéma qui (a) est localement de type fini sur un champ $k$, et (b) a de nombreux points finis. Puis lemme \ref{lemma-ubiquity-Jacobson-schemes} montre $T$ est un schéma de Jacobson. Un espace Jacobson fini est discret, voir Topologie, lemme \ref{topology-lemma-finite-jacobson}. Appliquez cette remarque sur la fibre $X_s$ qui est localement de type fini sur $\Spec(\kappa(s))$ pour voir la première affirmation. Enfin, appliquez lemme \ref{lemma-quasi-finite-at-point-characterize} pour voir la seconde.
\end{proof}

\begin{lemma}

\label{lemma-locally-quasi-finite-fibres}

\begin{slogan}

Les morphismes de type fini avec des fibres discrètes sont quasi-finis.

\end{slogan}

Laisser $f : X \to S$ être un morphisme des schémas. $f$ est localement de type fini. Ensuite, les éléments suivants sont équivalents

\begin{enumerate}

\item  $f$ est localement quasi-finie,
\item pour chaque $s \in S$ les fibres $X_s$ est un espace topologique discret, et
\item pour chaque morphisme $\Spec(k) \to S$ où $k$ est un champ le changement de base $X_k$ a un espace topologique discret sous-jacent.

\end{enumerate}

\end{lemma}

\begin{proof}
Il est immédiat que (3) implique (2). lemme \ref{lemma-quasi-finite-at-point-characterize} montre que (2) est équivalent à (1). Supposons que (2) et laisser $\Spec(k) \to S$ être comme dans (3). Notons $s \in S$ l'image de $\Spec(k) \to S$. Alors $X_k$ est le changement de base de $X_s$ via $\Spec(k) \to \Spec(\kappa(s))$. Par conséquent, chaque point de $X_k$ est fermé par lemme \ref{lemma-base-change-closed-point-fibre-locally-finite-type}. Comme $X_k \to \Spec(k)$ est localement de type fini (par lem \ref{lemma-base-change-finite-type}), nous pouvons appliquer lemme \ref{lemma-quasi-finite-at-point-characterize} pour conclure que chaque point de $X_k$ est isolé, c.-à-d., $X_k$ a un espace topologique sous-jacent discret.
\end{proof}

\begin{lemma}

\label{lemma-quasi-finite-locally-quasi-compact}
Laissez $f : X \to S$ être un morphisme de schémas. Alors $f$ est quasi-fini si et seulement si $f$ est localement quasi-fini et quasi-compact.
\end{lemma}

\begin{proof}
Supposons que $f$ est quasi-fini. Il est quasi-compact par définition \ref{definition-finite-type}. Laissez $x \in X$. Nous voyons que $f$ est quasi-fini à $x$ par lemme \ref{lemma-quasi-finite-at-point-characterize}. Par conséquent, $f$ est quasi-compact et localement quasi-fini.

\medskip\noindent
Supposons que $f$ est quasi-compact et localement quasi-fini. Alors $f$ est de type fini. Laissez $x \in X$ être un point. Par lemme \ref{lemma-quasi-finite-at-point-characterize} nous voyons que $x$ est un point isolé de sa fibre. Le lemme est prouvé.
\end{proof}

\begin{lemma}

\label{lemma-quasi-finite}

Laisser $f : X \to S$ morphisme des régimes. Les conditions suivantes sont équivalentes:

\begin{enumerate}

\item  $f$ est quasi-fini, et
\item  $f$ est localement de type fini, quasi compact, et a des fibres finies.

\end{enumerate}

\end{lemma}

\begin{proof}

Présumer $f$ est quasi-finie. $f$ est localement de type fini et quasi-compact (puisque c'est de type fini). $s \in S$. Depuis chaque $x \in X_s$ est isolé dans $X_s$ nous voyons que
$X_s = \bigcup_{x \in X_s} \{x\}$ est un recouvrement ouvert. $f$
est quasi-compact, la fibre $X_s$ est quasi-compact. Donc nous voyons que $X_s$ C'est fini.

\medskip\noindent
Inversement, supposez que $f$ est localement de type fini, quasi-compact et a des fibres finies. Ensuite, il est localement quasi-fini par lemme \ref{lemma-finite-fibre}. Donc il est quasi-fini par lemme \ref{lemma-quasi-finite-locally-quasi-compact}.
\end{proof}

\noindent
Rappelez-vous qu'un homomorphisme d'anneaux $R \to A$ est quasi-fini s'il est de type fini et quasi-fini à {\it tous} les premiers de $A$, voir Algèbre, définition \ref{algebra-definition-quasi-finite}.

\begin{lemma}

\label{lemma-locally-quasi-finite-characterize}

Laisser $f : X \to S$ morphisme des régimes. Les suivants sont équivalents

\begin{enumerate}

\item Le morphisme $f$ est localement quasi-finie.
\item Pour chaque paire d'ouvertures affines $U \subset X$, $V \subset S$
avec $f(U) \subset V$ l'homomorphisme d'anneaux
$\mathcal{O}_S(V) \to \mathcal{O}_X(U)$ C'est presque fini.
\item Il existe une ouverture de recouvrement $S = \bigcup_{j \in J} V_j$
et couvertures ouvertes $f^{-1}(V_j) = \bigcup_{i \in I_j} U_i$ de sorte que chacun des morphismes $U_i \to V_j$, $j\in J, i\in I_j$
est localement quasi-finie.
\item Il existe un recouvrement ouvert affine $S = \bigcup_{j \in J} V_j$
et revêtements pour affines d'ouverture $f^{-1}(V_j) = \bigcup_{i \in I_j} U_i$ tel que l'homomorphisme d'anneaux $\mathcal{O}_S(V_j) \to \mathcal{O}_X(U_i)$ est quasi-fini, pour tous $j\in J, i\in I_j$.

\end{enumerate}

En outre, si $f$ est localement quasi-finie puis pour tout sous-système ouvert $U \subset X$, $V \subset S$ avec $f(U) \subset V$
la restriction $f|_U : U \to V$ est localement quasi-finie.

\end{lemma}

\begin{proof}

Pour un homomorphisme d'anneaux $R \to A$ définissons
$P(R \to A)$ de signifier «»$R \to A$ (voir remarque ci-dessus lemme). $P$ est une propriété locale des homomorphismes d'anneaux. Nous vérifions les conditions (a), (b) et (c) de définition
\ref{definition-property-local}. Dans la démonstration de lemme \ref{lemma-locally-finite-type-characterize}
nous avons vu que (a), (b) et (c) tenir pour la propriété d'être ``de type fini''. Notez que, pour un type fini homomorphisme d'anneaux
$R \to A$, la propriété $R \to A$ est quasi-fini à $\mathfrak q$
dépend seulement de l'anneau local $A_{\mathfrak q}$ comme un algèbre $R_{\mathfrak p}$ où $\mathfrak p = R \cap \mathfrak q$
En utilisant ces remarques a), b) et c) de la définition \ref{definition-property-local} suivez immédiatement. Par exemple, supposez $R \to A$ est un homomorphisme d'anneaux tel que tous les homomorphismes d'anneaux $R \to A_{a_i}$ sont quasi-finis pour $a_1, \ldots, a_n \in A$ générant l'idéal de l'unité. $R \to A$ est de type fini. Aussi, pour tout $\mathfrak q \subset A$ l'anneau local $A_{\mathfrak q}$
est isomorphe comme un $R$-l'algèbre à l'anneau local
$(A_{a_i})_{\mathfrak q_i}$ pour certains $i$ et certains
$\mathfrak q_i \subset A_{a_i}$. Nous concluons donc que
$R \to A$ est quasi-fini à $\mathfrak q$.

\medskip\noindent
Nous concluons que lemme \ref{lemma-locally-P} s'applique avec $P$ comme dans le paragraphe précédent. Par conséquent, il suffit de prouver que $f$ est localement quasi-fini est équivalent à $f$ est localement de type $P$. Puisque $P(R \to A)$ est ``$R \to A$ est quasi-fini''', ce qui signifie que $R \to A$ est quasi-fini à chaque prime de $A$, cela découle de lemme \ref{lemma-quasi-finite-at-point-characterize}.
\end{proof}

\begin{lemma}

\label{lemma-composition-quasi-finite}

La composition de deux morphismes qui sont localement quasi-finis est localement quasi-finis. Il en est de même pour les morphismes quasi-finis.

\end{lemma}

\begin{proof}

Dans la démonstration de lemme \ref{lemma-locally-quasi-finite-characterize}
On a vu ça. $P = $"quasi-fini" est une propriété locale des homomorphismes d'anneaux, et qu'un morphisme des schémas est localement quasi-fini si et seulement s'il est localement de type $P$ comme dans la définition \ref{definition-locally-P}. D'où la première déclaration du lemme suit de lemme \ref{lemma-composition-type-P} combiné avec le fait qu'être quasi-fini est une propriété des homomorphismes d'anneaux qui est stable sous la composition, voir Algèbre, lemme \ref{algebra-lemma-quasi-finite-composition}Par ce qui précède, lemme \ref{lemma-quasi-finite-locally-quasi-compact}
et le fait que les compositions des morphismes quasi-compacts sont quasi-compactes, voir Schémas, lemme \ref{schemes-lemma-composition-quasi-compact}
nous voyons que la composition des morphismes quasi-finis est quasi-finie.

\end{proof}

\noindent
Nous verrons plus tard (lemme \ref{lemma-quasi-finite-points-open}) que le jeu $U$ du lemme suivant est ouvert.

\begin{lemma}

\label{lemma-base-change-quasi-finite}

\begin{slogan}

(Localement) les morphismes quasi-finis sont stables sous changement de base.

\end{slogan}

Laisser $f : X \to S$ être un morphisme de schémas. $g : S' \to S$ être un morphisme de schémas. $f' : X' \to S'$ le changement de base de $f$ par $g$
et indiquer $g' : X' \to X$ la projection. Supposons que $X$ est localement de type fini sur $S$.

\begin{enumerate}

\item Laisser $U \subset X$ (resp.\ $U' \subset X'$) être l'ensemble des points où $f$ (resp.\ $f'$) est quasi-finie. $U' = U \times_S S' = (g')^{-1}(U)$.
\item Le changement de base d'un morphisme local quasi-fini est local quasi-fini.
\item Le changement de base d'un morphisme quasi-fini est quasi-fini.

\end{enumerate}

\end{lemma}

\begin{proof}

La première et la deuxième affirmation suivent du résultat correspondant de l'algèbre, voir Algèbre, lemme \ref{algebra-lemma-quasi-finite-base-change}
(combiné avec le fait que $f'$ est également localement de type fini par lemme \ref{lemma-base-change-finite-type}Par ce qui précède, lemme \ref{lemma-quasi-finite-locally-quasi-compact}
et le fait qu'un changement de base d'un morphisme quasi-compact est quasi-compact, voir Schémas, lemme \ref{schemes-lemma-quasi-compact-preserved-base-change}
nous voyons que le changement de base d'un morphisme quasi-fini est quasi-fini.

\end{proof}

\begin{lemma}

\label{lemma-quasi-finite-at-a-finite-number-of-points}
Laissez $f : X \to S$ être un morphisme de schémas de type fini. Laissez $s \in S$. Il y a au plus finiment de nombreux points de $X$ au-dessus de $s$ où $f$ est quasi-fini.
\end{lemma}

\begin{proof}
La fibre $X_s$ est un schéma de type fini sur un champ, donc noethérien (lemme \ref{lemma-finite-type-noetherian}). D'où la topologie sur $X_s$ est noethérien (Propriétés, lemme \ref{properties-lemma-Noetherian-topology}) et peut avoir au plus un nombre fini de points isolés (par topologie élémentaire). Ainsi notre lemme suit de lemme \ref{lemma-quasi-finite-at-point-characterize}.
\end{proof}

\begin{lemma}

\label{lemma-monomorphism-loc-finite-type-loc-quasi-finite}
Laissez $f : X \to Y$ être un morphisme de schémas. Si $f$ est localement de type fini et un monomorphisme, alors $f$ est séparé et localement quasi-fini.
\end{lemma}

\begin{proof}
Un monomorphisme est séparé par Schémas, lemme \ref{schemes-lemma-monomorphism-separated}. Un monomorphisme est injectable, donc nous obtenons $f$ est quasi-fini à chaque $x \in X$ par exemple par lemme \ref{lemma-quasi-finite-at-point-characterize}.
\end{proof}

\begin{lemma}

\label{lemma-immersion-locally-quasi-finite}

Toute immersion est localement quasi-finie.

\end{lemma}

\begin{proof}

C'est vrai parce qu'une immersion ouverte est un isomorphisme local et une immersion fermée est clairement quasi-finie.

\end{proof}

\begin{lemma}

\label{lemma-permanence-quasi-finite}
Laissez $X \to Y$ être un morphisme de schémas sur un schéma de base $S$. Laissez $x \in X$. Si $X \to S$ est quasi-fini à $x$, alors $X \to Y$ est quasi-fini à $x$. Si $X$ est localement quasi-fini sur $S$, alors $X \to Y$ est localement quasi-fini.
\end{lemma}

\begin{proof}
Via lemme \ref{lemma-locally-quasi-finite-characterize} ceci se traduit par le fait d'algèbre suivant: Compte tenu des homomorphismes d'anneaux $A \to B \to C$ de telle sorte que $A \to C$ soit quasi-fini, alors $B \to C$ est quasi-fini. Ceci suit d'Algèbre, lemme \ref{algebra-lemma-four-rings} avec $R = A$, $S = S' = C$ et $R' = B$.
\end{proof}

\begin{lemma}

\label{lemma-quasi-finite-local-source}
Laissez $f : X \to Y$ et $g : Y \to S$ être des morphismes de schémas. Si $f$ est surjectif, $g \circ f$ localement quasi-fini, et $g$ localement de type fini, alors $g : Y \to S$ est localement quasi-fini.
\end{lemma}

\begin{proof}

Laisser $x \in X$ avec des images $y \in Y$ et $s \in S$Depuis $g \circ f$ est localement quasi-fini par lemme \ref{lemma-residue-field-quasi-finite}
l'extension $\kappa(x)/\kappa(s)$ C'est fini. $\kappa(y)/\kappa(s)$ C'est fini.
$y$ est un point fermé de $Y_s$ par lemme \ref{lemma-algebraic-residue-field-extension-closed-point-fibre}Depuis $f$ est surjectif, nous voyons que chaque point de $Y$ est fermé dans sa fibre sur $S$. Ainsi par lemme \ref{lemma-quasi-finite-at-point-characterize}
nous concluons que $g$ est quasi-fini à chaque point.

\end{proof}














\section{Morphismes de présentation finie}
\label{section-finite-presentation}

\noindent
Rappelez qu'un homomorphisme d'anneaux $R \to A$ est de présentation finie si $A$ est isomorphe à $R[x_1, \ldots, x_n]/(f_1, \ldots, f_m)$ comme algèbre $R$ pour certains $n, m$ et certains polynomes $f_j$, voir Algèbre, définition \ref{algebra-definition-finite-type}.

\begin{definition}

\label{definition-finite-presentation}

Laisser $f : X \to S$ être un morphisme de schémas.

\begin{enumerate}

\item Nous disons que $f$ est de {\it Présentation finale au $x \in X$} s'il existe un quartier affiné d'ouverture $\Spec(A) = U \subset X$
de $x$ et ouvert affine $\Spec(R) = V \subset S$
avec $f(U) \subset V$ tel que l'homomorphisme d'anneaux induit
$R \to A$ est de présentation finie.
\item Nous disons que $f$ est {\it localement de présentation finie} si c'est de présentation finie à chaque point $X$.
\item Nous disons que $f$ est de {\it Présentation finie} s'il s'agit d'un lieu de présentation finie, quasi-compact et quasi-séparé.

\end{enumerate}

\end{definition}\noindent Notons qu'un morphisme de présentation finie n'est {\bf pas} simplement un morphisme quasi-compact localement de présentation finie. Plus loin, nous caractériserons les morphismes localement de présentation finie comme ceux tels que $$
\colim \Mor_S(T_i, X) = \Mor_S(\lim T_i, X)
$$
pour tout système dirigé de régimes d'affinage $T_i$ terminé $S$Voir Limites, proposition \ref{limits-proposition-characterize-locally-finite-presentation}. Dans Limites, section \ref{limits-section-descending-relative} nous montrons que, si $S = \lim_i S_i$ est une limite des régimes d'affine, tout régime $X$ de présentation finie sur $S$ descend à un régime
$X_i$ terminé $S_i$ pour certains $i$.

\begin{lemma}

\label{lemma-locally-finite-presentation-characterize}

Laisser $f : X \to S$ morphisme des régimes. Les suivants sont équivalents

\begin{enumerate}

\item Le morphisme $f$ est le lieu de présentation finie.
\item Pour tous les ouverts $U \subset X$, $V \subset S$
avec $f(U) \subset V$ l'homomorphisme d'anneaux
$\mathcal{O}_S(V) \to \mathcal{O}_X(U)$ est de présentation finie.
\item Il existe un recouvrement ouvert $S = \bigcup_{j \in J} V_j$
et couvertures ouvertes $f^{-1}(V_j) = \bigcup_{i \in I_j} U_i$ de sorte que chacun des morphismes $U_i \to V_j$, $j\in J, i\in I_j$
est le lieu de présentation finie.
\item Il existe un recouvrement ouvert affine $S = \bigcup_{j \in J} V_j$
et revêtements pour affines d'ouverture $f^{-1}(V_j) = \bigcup_{i \in I_j} U_i$ tel que l'homomorphisme d'anneaux $\mathcal{O}_S(V_j) \to \mathcal{O}_X(U_i)$ est de présentation finie, pour tous $j\in J, i\in I_j$.

\end{enumerate}

En outre, si $f$ est localement de présentation finie puis pour tout sous-système ouvert $U \subset X$, $V \subset S$ avec $f(U) \subset V$
la restriction $f|_U : U \to V$ est le lieu de présentation finie.

\end{lemma}

\begin{proof}

Cela s'explique par le fait qu'il s'agit d'une \ref{lemma-locally-P-characterize} si nous montrons que la propriété ``$R \to A$ Nous vérifions les conditions (a), (b) et (c) de la définition
\ref{definition-property-local}. Par Algèbre, lemme \ref{algebra-lemma-base-change-finiteness}
être de présentation finie est stable sous changement de base et donc nous concluons (a) tient.
\ref{algebra-lemma-compose-finite-type} être de présentation finie est stable sous composition et trivialement pour n'importe quel anneau
$R$ l'homomorphisme d'anneaux $R \to R_f$ est de présentation finie. Nous concluons (b) détient. Enfin, la propriété (c) est vraie selon Algèbre, lemme \ref{algebra-lemma-cover-upstairs}.

\end{proof}

\begin{lemma}

\label{lemma-composition-finite-presentation}

La composition de deux morphismes qui sont localement de présentation finie est localment de présentation finie. Il en est de même pour les morphismes de présentation finie.

\end{lemma}

\begin{proof}

Dans la démonstration de lemme \ref{lemma-locally-finite-presentation-characterize}
nous avons vu qu'être de présentation finie est une propriété locale des homomorphismes d'anneaux. \ref{lemma-composition-type-P} combiné avec le fait qu'être de présence finie est une propriété d'homomorphismes d'anneaux stable sous composition, voir Algèbre, lemme \ref{algebra-lemma-compose-finite-type}. Par ce qui précède et le fait que les compositions de quasi-compact, quasi-morphismes séparés sont quasi-compact et quasi-séparés, voir Schémas, lemmes \ref{schemes-lemma-composition-quasi-compact}
et \ref{schemes-lemma-separated-permanence}
nous voyons que la composition des morphismes de présentation finie est de présentation finie.

\end{proof}

\begin{lemma}

\label{lemma-base-change-finite-presentation}

Le changement de base d'un morphisme qui est localement de présentation finie est localment de présentation finie. Il en est de même pour les morphismes de présentation finie.

\end{lemma}

\begin{proof}

Dans la démonstration de lemme \ref{lemma-locally-finite-presentation-characterize}
nous avons vu qu'être de présentation finie est une propriété locale des homomorphismes d'anneaux. \ref{lemma-composition-type-P} combiné avec le fait qu'être de présentation finie est une propriété d'homomorphismes d'anneaux qui est stable sous changement de base, voir Algèbre, lemme \ref{algebra-lemma-base-change-finiteness}. Par ce qui précède et le fait qu'un changement de base d'un quasi-compact, quasi-morphisme séparé est quasi-compact et quasi-séparé, voir Schémas, lemmes \ref{schemes-lemma-quasi-compact-preserved-base-change}
et \ref{schemes-lemma-separated-permanence}
nous voyons que le changement de base d'un morphisme de présentation finie est un morphisme de présentation finie.

\end{proof}

\begin{lemma}

\label{lemma-open-immersion-locally-finite-presentation}

Toute immersion ouverte est localement de présentation finie.

\end{lemma}

\begin{proof}

C'est vrai parce qu'une immersion ouverte est un isomorphisme local.

\end{proof}

\begin{lemma}

\label{lemma-quasi-compact-open-immersion-finite-presentation}
Toute immersion ouverte est de présentation finie si et seulement si elle est quasi compacte.
\end{lemma}

\begin{proof}
Nous avons vu (lemme \ref{lemma-open-immersion-locally-finite-presentation}) qu'une immersion ouverte est localement de présentation finie. Nous avons vu (Schémas, lemme \ref{schemes-lemma-immersions-monomorphisms}) qu'une immersion est séparée et donc quasi-séparée. De cette et définition \ref{definition-finite-presentation} la lemme suit.
\end{proof}

\begin{lemma}

\label{lemma-closed-immersion-finite-presentation}

\begin{slogan}

immersions fermées de présentation finie correspondent à des faisceaux quasi-cohérents d'idéaux de type fini.

\end{slogan}

Une fermée d'immersion $i : Z \to X$ est de présence finie si et seulement si le faisceau associé quasi-cohérent des idéaux
$\mathcal{I} = \Ker(\mathcal{O}_X \to i_*\mathcal{O}_Z)$
est de type fini (en tant que $\mathcal{O}_X$- Module).

\end{lemma}

\begin{proof}
Sur n'importe quelle ouvert affine $\Spec(R) \subset X$ nous avons $i^{-1}(\Spec(R)) = \Spec(R/I)$ et $\mathcal{I} = \widetilde{I}$. En outre, $\mathcal{I}$ est de type fini si et seulement si $I$ est un fini $R$-module pour chaque ouvert affine (voir Propriétés, lemme \ref{properties-lemma-finite-type-module}). Et $R/I$ est de présentation finie sur $R$ si et seulement si $I$ est un fini $R$-module. Par conséquent, nous gagnons.
\end{proof}

\begin{lemma}

\label{lemma-finite-presentation-finite-type}
Un morphisme qui est localement de présentation finie est localement de type fini. Un morphisme de présentation finie est de type fini.
\end{lemma}

\begin{proof}
Démonstration omise.
\end{proof}

\begin{lemma}

\label{lemma-noetherian-finite-type-finite-presentation}

\begin{slogan}
Au-dessus d'une base de noéthérien local, le type fini est présentation finie.
\end{slogan}
Laissez $f : X \to S$ être un morphisme.
\begin{enumerate}
\item Si $S$ est localisation noéthérien et $f$ localisation de type fini alors $f$ est localisation de présentation finie.
\item Si $S$ est localisation noéthérien et $f$ de type fini alors $f$ est de présentation finie.
\end{enumerate}

\end{lemma}

\begin{proof}

La première affirmation découle du fait qu'un anneau de type fini sur un anneau Noétherien est de présentation finie, voir Algèbre, lemme \ref{algebra-lemma-Noetherian-finite-type-is-finite-presentation}. Supposons que $f$ est de type fini et $S$ est localement noethérien. $f$ est quasi-compact et localement de présentation finie par (1). Il suffit donc de prouver que $f$ est quasi-séparée. Ceci résulte de \ref{lemma-finite-type-Noetherian-quasi-separated}
(et lemme \ref{lemma-finite-presentation-finite-type}).

\end{proof}

\begin{lemma}

\label{lemma-finite-presentation-quasi-compact-quasi-separated}

Laisser $S$ être un système quasi compact et quasi séparé. $X$ est de présentation finie sur $S$Alors $X$ est quasi-compact et quasi-séparé.

\end{lemma}

\begin{proof}
Démonstration omise.
\end{proof}

\begin{lemma}

\label{lemma-finite-presentation-permanence}

Laisser $f : X \to Y$ être un morphisme des schémas sur $S$.

\begin{enumerate}

\item Si $X$ est localement de présentation finie over $S$ et
$Y$ est localement de type fini sur $S$Alors $f$ est le lieu de présentation finie.
\item Si $X$ est de présentation finie sur $S$ et $Y$ est quasi-séparé et localisation de type fini sur $S$Alors $f$ est de présentation finie.

\end{enumerate}

\end{lemma}

\begin{proof}
démonstration de (1). Via lemme \ref{lemma-locally-finite-presentation-characterize} cela se traduit par le fait d'algèbre suivant: Compte tenu des homomorphismes d'anneaux $A \to B \to C$ de telle sorte que $A \to C$ est de présentation finie et $A \to B$ est de type fini, alors $B \to C$ est de présentation finie. Voir Algèbre, lemme \ref{algebra-lemma-compose-finite-type}.

\medskip\noindent
La partie (2) est la suivante: (1) et Schémas, lemmes \ref{schemes-lemma-compose-after-separated} et \ref{schemes-lemma-quasi-compact-permanence}.
\end{proof}

\begin{lemma}

\label{lemma-diagonal-morphism-finite-type}
Laissez $f : X \to Y$ être un morphisme de schémas avec la diagonale $\Delta : X \to X \times_Y X$. Si $f$ est localement de type fini alors $\Delta$ est localment de présentation finie. Si $f$ est quasi-séparé et localment de type fini, alors $\Delta$ est de présentation finie.
\end{lemma}

\begin{proof}

Notez que $\Delta$ est un morphisme de schémas sur $X$ (par la deuxième projection) $X \times_Y X \to X$). Supposons que $f$ est localement de type fini. Notez que $X$ est de présentation finie sur $X$ et $X \times_Y X$ est localement de type fini sur $X$ (par la suite) \ref{lemma-base-change-finite-type}Ainsi, la première déclaration tient par lemme \ref{lemma-finite-presentation-permanence}. La deuxième affirmation suit de la première, les définitions, et le fait qu'un morphisme diagonale est un monomorphisme, donc séparé (Schémas, lemme \ref{schemes-lemma-monomorphism-separated}).

\end{proof}








\section{Ensembles constructibles}
\label{section-constructible}

\noindent
Des ensembles constructibles et des ensembles de localisation ont été discutés dans la section \ref{properties-section-constructible} de Propriétés. Dans cette section, nous montrons quelques résultats concernant des images et des images inverses des ensembles (localement) constructibles. Le résultat principal est le théorème de Chevalley qui affirme que l'image d'un ensemble de localisation constructible sous un morphisme de présentation finie est localement constructible.

\begin{lemma}

\label{lemma-inverse-image-constructible}
Laissez $f : X \to Y$ être un morphisme de schémas. Laissez $E \subset Y$ être un sous-ensemble. Si $E$ est (localement) constructible dans $Y$, alors $f^{-1}(E)$ est (localement) constructible dans $X$.
\end{lemma}

\begin{proof}

Pour montrer que l'image inverse de chaque sous-ensemble constructible est constructible, il suffit de montrer que l'image inverse de chaque rétrocompact ouvert $V$
de $Y$ est rétrocompact dans $X$, voir Topologie, me laisse \ref{topology-lemma-inverse-images-constructibles}. L'importance de $V$ être rétrocompact dans $Y$ c'est juste que l'immersion ouverte $V \to Y$ est quasi-compact. D'où le changement de base $f^{-1}(V) = X \times_Y V \to X$ est quasi compacte aussi, voir Schémas, lemme \ref{schemes-lemma-quasi-compact-preserved-base-change}. C'est pourquoi nous voyons $f^{-1}(V)$ est rétrocompact dans $X$. Supposons que $E$ est localement constructible dans $Y$. Choisissez $x \in X$. Choisissez un quartier affiné $V$ de $f(x)$ et un quartier affiné $U \subset X$ de $x$ de telle manière que $f(U) \subset V$. Ainsi nous pensons à $f|_U : U \to V$ comme un morphisme en $V$. Par Propriétés, me laisse \ref{properties-lemma-locally-constructible}
nous voyons que $E \cap V$ est constructible dans $V$. Par le cas constructible nous voyons que $(f|_U)^{-1}(E \cap V)$ est constructible dans $U$Depuis $(f|_U)^{-1}(E \cap V) = f^{-1}(E) \cap U$ Nous avons gagné.

\end{proof}

\begin{lemma}

\label{lemma-chevalley}
Laissez $f : X \to Y$ être un morphisme de schémas. Supposons que
\begin{enumerate}
\item $f$ est quasi-compact et localisation de présentation finie, et
\item $Y$ est quasi-compact et quasi-séparé.
\end{enumerate}
Ensuite, l'image de chaque sous-ensemble constructible de $X$ est constructible dans $Y$.
\end{lemma}

\begin{proof}

Par Propriétés, lemme \ref{properties-lemma-constructible-quasi-compact-quasi-separated}
il suffit de prouver ce lemme au cas où $Y$ Dans ce cas, $X$ est quasi-compact. Donc nous pouvons écrire
$X = U_1 \cup \ldots \cup U_n$ avec chacun $U_i$ ouvert affine dans $X$Si $E \subset X$ est constructible, puis chacun $E \cap U_i$ est constructible aussi, voir Topologie, lemme \ref{topology-lemma-open-immersion-constructible-inverse-image}. Par conséquent, depuis $f(E) = \bigcup f(E \cap U_i)$ et comme les unions finies d'ensembles constructibles sont constructibles, cela nous réduit au cas où
$X$ est affine. Dans ce cas, le résultat est Algèbre, théorème \ref{algebra-theorem-chevalley}.

\end{proof}

\begin{theorem}
[Théorème de Chevalley]
\label{theorem-chevalley}

\begin{reference}
\cite[IV, théorème 1.8.4]{EGA}
\end{reference}
Laissez $f : X \to Y$ être un morphisme de schémas. Supposons que $f$ est quasi-compact et localement de présentation finie. Puis l'image de chaque sous-ensemble constructible local est localement constructible.
\end{theorem}

\begin{proof}

Laisser $E \subset X$ Nous devons montrer que $f(E)$ est localement constructible aussi. Nous allons montrer que $f(E) \cap V$ est constructible pour n'importe quelle ouvert affine $V \subset Y$. Ainsi nous réduisons au cas où $Y$ Dans ce cas, $X$ est quasi-compact. Donc nous pouvons écrire
$X = U_1 \cup \ldots \cup U_n$ avec chacun $U_i$ ouvert affine dans $X$Si $E \subset X$ est localement constructible, puis chaque $E \cap U_i$
est constructible, voir Propriétés, lemme \ref{properties-lemma-locally-constructible}. Par conséquent, depuis $f(E) = \bigcup f(E \cap U_i)$ et comme les unions finies d'ensembles constructibles sont constructibles, cela nous réduit au cas où $X$
est affine. Dans ce cas, le résultat est Algèbre, théorème \ref{algebra-theorem-chevalley}.

\end{proof}

\begin{lemma}

\label{lemma-constructible-containing-open}
Laissez $X$ être un schéma. Laissez $x \in X$. Laissez $E \subset X$ être un sous-ensemble constructible localement. Si $\{x' \mid x' \leadsto x\} \subset E$, alors $E$ contient un quartier ouvert de $x$.
\end{lemma}

\begin{proof}

Présumer $\{x' \mid x' \leadsto x\} \subset E$. Nous pouvons supposer $X$ Dans ce cas, $E$ est constructible, voir Propriétés, lemme \ref{properties-lemma-locally-constructible}. En particulier, le complément $E^c$ par Algèbre, lemme \ref{algebra-lemma-constructible-is-image}
nous pouvons trouver un morphisme de schémas affines $f : Y \to X$ de telle manière que
$E^c = f(Y)$. Laisse $Z \subset X$ être l'image schématique de $f$. Par lemme \ref{lemma-reach-points-scheme-theoretic-image}
et l ' hypothèse $\{x' \mid x' \leadsto x\} \subset E$
nous voyons que $x \not \in Z$. D'où $X \setminus Z \subset E$ est un quartier ouvert de $x$ figurant dans $E$.

\end{proof}




\section{Morphismes ouverts}
\label{section-open}

\begin{definition}

\label{definition-open}
Laissez $f : X \to S$ être un morphisme.
\begin{enumerate}
\item Nous disons que $f$ est {\it ouvert} si l'application sur les espaces topologiques sous-jacents est ouverte.
\item Nous disons que $f$ est {\it universel ouvert} si pour un morphisme de schémas $S' \to S$ le changement de base $f' : X_{S'} \to S'$ est ouvert.
\end{enumerate}

\end{definition}

\noindent
Selon Topologie, les générisations lemme \ref{topology-lemma-closed-open-map-specialization} remontent le long de certains types de cartes ouvertes des espaces topologiques. En fait, les générisations remontent le long de tout morphisme ouvert des schémas (voir lemme \ref{lemma-open-generizing}). De plus, nous verrons que les générisations remontent le long des morphismes plats des schémas (lemme \ref{lemma-generalizations-lift-flat}). Cela implique parfois que le morphisme est ouvert.

\begin{lemma}

\label{lemma-locally-finite-presentation-universally-open}
Laissez $f : X \to S$ être un morphisme.
\begin{enumerate}
\item Si $f$ est une localisation de présentation finie et des générisations remontent le long de $f$, alors $f$ est ouvert.
\item Si $f$ est une localisation de présentation finie et des générisations remontent le long de chaque changement de base de $f$, alors $f$ est une universalisation ouvert.
\end{enumerate}

\end{lemma}

\begin{proof}
Il suffit de prouver la première affirmation. Cela réduit au cas où les deux $X$ et $S$ sont affines. Dans ce cas, le résultat suit d'Algèbre, lemme \ref{algebra-lemma-going-up-down-specialization} et la proposition \ref{algebra-proposition-fppf-open}.
\end{proof}

\noindent
Voir aussi lemme \ref{lemma-fppf-open} pour le cas d'un morphisme plat de présentation finie.

\begin{lemma}

\label{lemma-composition-open}

Une composition de morphismes ouverts (universellement) est (universellement) ouverte.

\end{lemma}

\begin{proof}
Démonstration omise.
\end{proof}

\begin{lemma}

\label{lemma-scheme-over-field-universally-open}
Laissez $k$ être un champ. Laissez $X$ être un schéma sur $k$. Le morphisme structure $X \to \Spec(k)$ est universel ouvert.
\end{lemma}

\begin{proof}
Laissez $S \to \Spec(k)$ être un morphisme. Nous devons montrer que le changement de base $X_S \to S$ est ouvert. La question est locale sur $S$ et $X$, donc nous pouvons supposer que $S$ et $X$ sont affines. Dans ce cas, le résultat est Algèbre, lemme \ref{algebra-lemma-map-into-tensor-algebra-open}.
\end{proof}

\begin{lemma}

\label{lemma-open-generizing}

\begin{reference}

Suit de l'implication (a) $\Rightarrow$(b) dans \cite[IV, Corollary 1.10.4]{EGA}\end{reference}

Laisser $\varphi : X \to Y$ être un morphisme de schémas. $\varphi$ est ouvert, alors $\varphi$ est génératrice (c.-à-d. que les générisations $\varphi$) Si $\varphi$ est universellement ouvert, alors $\varphi$ est universellement génératrice.

\end{lemma}

\begin{proof}

Présumer $\varphi$ est ouvert. $y' \leadsto y$ une spécialisation des points de $Y$. Laisse $x \in X$ avec $\varphi(x) = y$. Choisissez les ouverts affines $U \subset X$ et $V \subset Y$ de telle manière que
$\varphi(U) \subset V$ et $x \in U$. Puis aussi $y' \in V$. Par conséquent, nous pouvons remplacer $X$ par $U$ et $Y$ par $V$ et supposer $X$, $Y$ affine. L'affaire affine est Algèbre, lemme \ref{algebra-lemma-open-going-down}
(combiné avec Algèbre, lemme \ref{algebra-lemma-going-up-down-specialization}).

\end{proof}

\begin{lemma}

\label{lemma-descent-quasi-compact}
Laissez $f : X \to Y$ être un morphisme de schémas. Laissez $g : Y' \to Y$ être ouvert et surjectif de telle sorte que le changement de base $f' : X' \to Y'$ soit quasi-compact. Alors $f$ est quasi-compact.
\end{lemma}

\begin{proof}
Laissez $V \subset Y$ être un quasi-compact ouvert. Comme $g$ est ouvert et surjectif, nous pouvons trouver un quasi-compact ouvert $W' \subset W$ tel que $g(W') = V$. Par hypothèse $(f')^{-1}(W')$ est quasi-compact. L'image de $(f')^{-1}(W')$ dans $X$ est égale à $f^{-1}(V)$, voir lemme \ref{lemma-when-point-maps-to-pair}. Par conséquent $f^{-1}(V)$ est quasi-compact comme l'image d'un espace quasi-compact, voir Topologie, lemme \ref{topology-lemma-image-quasi-compact}. Ainsi $f$ est quasi-compact.
\end{proof}





\section{Morphismes submersifs}
\label{section-submersive}

\begin{definition}

\label{definition-submersive}
Laissez $f : X \to Y$ être un morphisme de schémas.
\begin{enumerate}
\item Nous disons que $f$ est {\it submersif}\footnote{Ceci est très différent de la notion d'une submersion de multiples différentiels.} si l'application continue des espaces topologiques sous-jacents est submersif, voir Topologie, définition \ref{topology-definition-submersive}.
\item Nous disons que $f$ est {\it universellement submersif} si pour chaque morphisme de schémas $Y' \to Y$ le changement de base $Y' \times_Y X \to Y'$ est submersif.
\end{enumerate}

\end{definition}

\noindent
Nous notons qu'un submersif morphisme est particulièrement surjectif.

\begin{lemma}

\label{lemma-base-change-universally-submersive}

Le changement de base d'un submersif universel de schémas par tout morphisme de schémas est submersif universel.

\end{lemma}

\begin{proof}
Ceci est immédiat de la définition.
\end{proof}

\begin{lemma}

\label{lemma-composition-universally-submersive}
La composition d'une paire de morphismes (universellement) submersifs de schémas est (universellement) submersive.
\end{lemma}

\begin{proof}
Démonstration omise.
\end{proof}










\section{Morphismes plats}
\label{section-flat}

\noindent
La platitude est l'un des outils techniques les plus importants dans la géométrie algébrique. Dans cette section nous introduisons cette notion. Nous limitons intentionnellement la discussion à des observations simples, en dehors de lemme \ref{lemma-fppf-open}. Une classe très importante de résultats, à savoir les critères de planéité, sont discutés dans Algèbre, Sections \ref{algebra-section-criteria-flatness}, \ref{algebra-section-flatness-artinian}, \ref{algebra-section-more-flatness-criteria}, et Compléments sur les morphismes, section \ref{more-morphisms-section-criterion-flat-fibres}. Il y a un chapitre dédié au matériel avancé sur les morphismes plats de schémas, à savoir Plus sur la planéité, section \ref{flat-section-introduction}.

\medskip\noindent
Rappelez qu'un module $M$ sur un anneau $R$ est {\it plat} si le functeur $-\otimes_R M : \text{Mod}_R \to \text{Mod}_R$ est exact. Un homomorphisme d'anneaux $R \to A$ est dit {\it plat} si $A$ est plat comme un module $R$. Voir Algèbre, définition \ref{algebra-definition-flat}.

\begin{definition}

\label{definition-flat}

Laisser $f : X \to S$ être un morphisme de schémas. $\mathcal{F}$ être un faisceau quasi-cohérent de $\mathcal{O}_X$- Des modules.

\begin{enumerate}

\item Nous disons $f$ est {\it plat à un point $x \in X$} si l'anneau local $\mathcal{O}_{X, x}$ est plat au-dessus de l'Anneau local
$\mathcal{O}_{S, f(x)}$.
\item Nous disons que $\mathcal{F}$ est {\it plat sur $S$ à un moment donné $x \in X$} si la tige $\mathcal{F}_x$ est un plat $\mathcal{O}_{S, f(x)}$- Module.
\item Nous disons $f$ est {\it plat} si $f$ est plat à chaque point de $X$.
\item Nous disons que $\mathcal{F}$ est {\it plat sur $S$} si $\mathcal{F}$ est plat sur $S$ à tous les points $x$ de $X$.

\end{enumerate}

\end{definition}

\noindent
Ainsi, nous voyons que $f$ est plat si et seulement si la matrice de structure $\mathcal{O}_X$ est plat sur $S$.

\begin{lemma}

\label{lemma-flat-module-characterize}

Laisser $f : X \to S$ être un morphisme de schémas. $\mathcal{F}$ être un faisceau quasi-cohérent de $\mathcal{O}_X$-modules. Les conditions suivantes sont équivalentes

\begin{enumerate}

\item La gerbe $\mathcal{F}$ est plat sur $S$.
\item Pour tous les ouverts $U \subset X$, $V \subset S$
avec $f(U) \subset V$ les $\mathcal{O}_S(V)$-module $\mathcal{F}(U)$ C'est plat.
\item Il existe une ouverture de recouvrement $S = \bigcup_{j \in J} V_j$
et couvertures ouvertes $f^{-1}(V_j) = \bigcup_{i \in I_j} U_i$ de telle sorte que chacun des modules $\mathcal{F}|_{U_i}$ est plat sur $V_j$, pour tous $j\in J, i\in I_j$.
\item Il existe un recouvrement ouvert affine $S = \bigcup_{j \in J} V_j$
et revêtements pour affines d'ouverture $f^{-1}(V_j) = \bigcup_{i \in I_j} U_i$ de telle manière que $\mathcal{F}(U_i)$ est un plat $\mathcal{O}_S(V_j)$- Module, pour tous
$j\in J, i\in I_j$.

\end{enumerate}

En outre, si $\mathcal{F}$ est plat sur $S$ puis pour les sous-systèmes ouverts $U \subset X$, $V \subset S$ avec $f(U) \subset V$
la restriction $\mathcal{F}|_U$ est plat sur $V$.

\end{lemma}

\begin{proof}

Laisser $R \to A$ être un homomorphisme d'anneaux. $M$ être un $A$- Module. $M$ est $R$-plat, puis pour tous les premiers
$\mathfrak q$ le module $M_{\mathfrak q}$ est plat sur $R_{\mathfrak p}$
avec $\mathfrak p$ le premier de $R$ couché sous $\mathfrak q$. Inversement, si
$M_{\mathfrak q}$ est plat sur $R_{\mathfrak p}$ pour tous les primes $\mathfrak q$
de $A$Alors $M$ est plat sur $R$Voir Algèbre, lemme \ref{algebra-lemma-flat-localization}Cette équivalence implique aisément les affirmations du lemme.

\end{proof}

\begin{lemma}

\label{lemma-flat-characterize}

Laisser $f : X \to S$ morphisme des régimes. Les suivants sont équivalents

\begin{enumerate}

\item Le morphisme $f$ C'est plat.
\item Pour tous les ouverts $U \subset X$, $V \subset S$
avec $f(U) \subset V$ l'homomorphisme d'anneaux
$\mathcal{O}_S(V) \to \mathcal{O}_X(U)$ C'est plat.
\item Il existe une ouverture de recouvrement $S = \bigcup_{j \in J} V_j$
et couvertures ouvertes $f^{-1}(V_j) = \bigcup_{i \in I_j} U_i$ de sorte que chacun des morphismes $U_i \to V_j$, $j\in J, i\in I_j$
C'est plat.
\item Il existe un recouvrement ouvert affine $S = \bigcup_{j \in J} V_j$
et revêtements pour affines d'ouverture $f^{-1}(V_j) = \bigcup_{i \in I_j} U_i$ de telle manière que $\mathcal{O}_S(V_j) \to \mathcal{O}_X(U_i)$ est plat, pour tous
$j\in J, i\in I_j$.

\end{enumerate}

En outre, si $f$ est plat puis pour les sous-systèmes ouverts $U \subset X$, $V \subset S$ avec $f(U) \subset V$
la restriction $f|_U : U \to V$ C'est plat.

\end{lemma}

\begin{proof}
C'est un cas spécial de \ref{lemma-flat-module-characterize} lemme ci-dessus.
\end{proof}

\begin{lemma}

\label{lemma-pushforward-flat-affine}
Laissez $f : X \to Y$ être un mode morphisme de schémas sur un schéma de base $S$. Laissez $\mathcal{F}$ être un module quasi-cohérent $\mathcal{O}_X$. Puis $\mathcal{F}$ est plat sur $S$ si et seulement si $f_*\mathcal{F}$ est plat sur $S$.
\end{lemma}

\begin{proof}
Par lemme \ref{lemma-flat-module-characterize} et le fait que $f$ est un morphisme affine, cela nous réduit à l'affine. Écrivons $X \to Y \to S$ correspond aux homomorphismes d'anneaux $C \leftarrow B \leftarrow A$.Soit $N$ le module $C$ correspondant à $\mathcal{F}$. Rappelez que $f_*\mathcal{F}$ correspond à $N$ vu comme un module $B$, voir Schémas, lemme \ref{schemes-lemma-widetilde-pullback}. Ainsi le résultat est clair.
\end{proof}

\begin{lemma}

\label{lemma-composition-module-flat}
Laissez $X \to Y \to Z$ être des morphismes de schémas. Laissez $\mathcal{F}$ être un module $\mathcal{O}_X$ quasi-cohérent. Laissez $x \in X$ avec l'image $y$ dans $Y$. Si $\mathcal{F}$ est plat sur $Y$ à $x$, et $Y$ est plat sur $Z$ à $y$, alors $\mathcal{F}$ est plat sur $Z$ à $x$.
\end{lemma}
X

\begin{proof}
Voir Algèbre, lemme \ref{algebra-lemma-composition-flat}.
\end{proof}

\begin{lemma}

\label{lemma-composition-flat}
La composition des plateaux morphismes est plate.
\end{lemma}

\begin{proof}
C'est un cas spécial de \ref{lemma-composition-module-flat} lemme.
\end{proof}

\begin{lemma}

\label{lemma-base-change-module-flat}

Laisser $f : X \to S$ être un morphisme de schémas. $\mathcal{F}$ être un faisceau quasi-cohérent de $\mathcal{O}_X$- Les modules. $g : S' \to S$ être un morphisme de schémas. $g' : X' = X_{S'} \to X$ la projection. $x' \in X'$ être un point avec image $x = g'(x') \in X$Si $\mathcal{F}$ est plat sur $S$ au $x$Alors
$(g')^*\mathcal{F}$ est plat sur $S'$ au $x'$. En particulier, si $\mathcal{F}$ est plat sur $S$Alors
$(g')^*\mathcal{F}$ est plat sur $S'$.

\end{lemma}

\begin{proof}
Voir Algèbre, lemme \ref{algebra-lemma-flat-base-change}.
\end{proof}

\begin{lemma}

\label{lemma-base-change-flat}
Le changement de base d'un plateau morphisme est plat.
\end{lemma}

\begin{proof}
C'est un cas spécial de \ref{lemma-base-change-module-flat} lemme.
\end{proof}

\begin{lemma}

\label{lemma-generalizations-lift-flat}
Laissez $f : X \to S$ être un plateau morphisme de schémas. Puis générisations remontent le long de $f$, voir Topologie, définition \ref{topology-definition-lift-specializations}.
\end{lemma}

\begin{proof}
Voir Algèbre, section \ref{algebra-section-going-up}.
\end{proof}

\begin{lemma}

\label{lemma-fppf-open}

Un morphisme plat localement de présentation finie est universelment ouvert.

\end{lemma}

\begin{proof}
Ceci découle de lemmes \ref{lemma-generalizations-lift-flat} et lemme \ref{lemma-locally-finite-presentation-universally-open} ci-dessus. Nous pouvons également argumenter directement comme suit.

\medskip\noindent

Laisser $f : X \to S$ être plat et localement de présentation finie. Par lemmes \ref{lemma-base-change-flat} et
\ref{lemma-base-change-finite-presentation} le changement de base de $f$ est plat et localement de présentation finie. Il suffit donc de montrer $f$ est ouvert. $f$ est ouvert il suffit de montrer que nous pouvons couvrir
$X$ par des affines ouvertes $X = \bigcup U_i$ de telle manière que $U_i \to S$
C'est ouvert. $X$ par ouverts affines $U_i \subset X$ de telle sorte que chacun
$U_i$ cartes dans un ouvert affine $V_i \subset S$ et de telle sorte que l'homomorphisme d'anneaux induit $\mathcal{O}_S(V_i) \to \mathcal{O}_X(U_i)$ est plat et de présentation finie (lemmes \ref{lemma-flat-characterize} et
\ref{lemma-locally-finite-presentation-characterize}). Puis $U_i \to V_i$ est ouvert par Algèbre, proposition \ref{algebra-proposition-fppf-open}
et la démonstration est complète.

\end{proof}

\begin{lemma}

\label{lemma-pf-flat-module-open}
Laissez $f : X \to Y$ être un morphisme de schémas. Laissez $\mathcal{F}$ être un quasi-cohérent $\mathcal{O}_X$-module. Supposons que $f$ localement présentation finie, $\mathcal{F}$ de type fini, $X = \text{Supp}(\mathcal{F})$, et $\mathcal{F}$ plat sur $Y$. Alors $f$ est universel ouvert.
\end{lemma}

\begin{proof}
Par lemmes \ref{lemma-base-change-module-flat}, \ref{lemma-base-change-finite-presentation}, et \ref{lemma-support-finite-type} les hypothèses sont conservées sous changement de base. Par lemme \ref{lemma-locally-finite-presentation-universally-open} il suffit de montrer que les générisations se lèvent le long de $f$. Ceci découle d'Algèbre, lemme \ref{algebra-lemma-going-down-flat-module}.
\end{proof}

\begin{lemma}

\label{lemma-fpqc-quotient-topology}

\begin{reference}
\cite[Expose VIII, Corollaire 4.3]{SGA1} et \cite[IV, Corollaire 2.3.12]{EGA}
\end{reference}
Laissez $f : X \to Y$ être un plat quasi-compact, surjectif, morphisme. Un sous-ensemble $T \subset Y$ est ouvert (resp.\ closed) si et seulement $f^{-1}(T)$ est ouvert (resp.\ closed). En d'autres termes, $f$ est un submersif morphisme.
\end{lemma}

\begin{proof}
La question est locale sur $Y$, donc nous pouvons supposer que $Y$ est affine. Dans ce cas, $X$ est quasi-compact comme $f$ est quasi-compact. Écrivez $X = X_1 \cup \ldots \cup X_n$ comme une union finie d'ouverts affines. Alors $f' : X' = X_1 \amalg \ldots \amalg X_n \to Y$ est un plateau morphisme surjectif de schémas affins. Notez que pour $T \subset Y$ nous avons $(f')^{-1}(T) = f^{-1}(T) \cap X_1 \amalg \ldots \amalg f^{-1}(T) \cap X_n$. Par conséquent, $f^{-1}(T)$ est ouvert si et seulement si $(f')^{-1}(T)$ est ouvert. Ainsi, nous pouvons supposer que les deux $X$ et $Y$ sont affines.

\medskip\noindent

Laisser $f : \Spec(B) \to \Spec(A)$ être un morphisme surjectif de schémas affines correspondant à un homomorphisme plat d'anneaux $A \to B$. Supposons que $f^{-1}(T)$ est fermé, disons $f^{-1}(T) = V(J)$ pour $J \subset B$ un idéal. Alors
$T = f(f^{-1}(T)) = f(V(J))$ est l'image de $\Spec(B/J) \to \Spec(A)$
(ici nous utilisons que $f$ d'autre part, les générisations le long $f$ (lemme) \ref{lemma-generalizations-lift-flat}D'où par Topologie, lemme \ref{topology-lemma-lift-specializations-images}
nous voyons que $Y \setminus T = f(X \setminus f^{-1}(T))$ est stable pargénérisation. $T$ est stable par spécialisation (Topologie, lemme \ref{topology-lemma-open-closed-specialization}). Ainsi $T$ est fermé par Algèbre, lemme \ref{algebra-lemma-image-stable-specialization-closed}.

\end{proof}

\begin{lemma}

\label{lemma-flat-permanence}

Laisser $h : X \to Y$ être un morphisme des schémas sur $S$. Laisse $\mathcal{G}$ être un faisceau quasi-cohérent sur $Y$. Laisse $x \in X$ avec $y = h(x) \in Y$Si $h$ est plat à $x$Alors
$$
\mathcal{G}\text{ flat over }S\text{ at }y
\Leftrightarrow
h^*\mathcal{G}\text{ flat over }S\text{ at }x.
$$
En particulier: $h$ est surjectif et plat, alors
$\mathcal{G}$ est plat sur $S$, si et seulement si
$h^*\mathcal{G}$ est plat sur $S$Si $h$ est surjectif et plat, et $X$ est plat sur $S$Alors $Y$ est plat sur $S$.

\end{lemma}

\begin{proof}

Vous pouvez le prouver en appliquant Algèbre, lemme \ref{algebra-lemma-flatness-descends-more-general}. Voici une preuve directe. $s \in S$ être l'image de $y$. Considérez les cartes locales de l'Anneau
$\mathcal{O}_{S, s} \to \mathcal{O}_{Y, y} \to \mathcal{O}_{X, x}$. Par hypothèse l'homomorphisme d'anneaux $\mathcal{O}_{Y, y} \to \mathcal{O}_{X, x}$
est Fidèlement Plat, voir Algèbre, lemme \ref{algebra-lemma-local-flat-ff}. Laisse $N = \mathcal{G}_y$. Notez que
$h^*\mathcal{G}_x = N \otimes_{\mathcal{O}_{Y, y}} \mathcal{O}_{X, x}$, voir Faisceaux, lemme \ref{sheaves-lemma-stalk-pullback-modules}. Laisse $M' \to M$ une injection de $\mathcal{O}_{S, s}$-modules. Par la plateur fidèle mentionnée ci-dessus, nous avons

\begin{align*}
\Ker(
M' \otimes_{\mathcal{O}_{S, s}} N \to M \otimes_{\mathcal{O}_{S, s}} N)
\otimes_{\mathcal{O}_{Y, y}} \mathcal{O}_{X, x} \\
=
\Ker(
M' \otimes_{\mathcal{O}_{S, s}} N
\otimes_{\mathcal{O}_{Y, y}} \mathcal{O}_{X, x}
\to
M \otimes_{\mathcal{O}_{S, s}} N
\otimes_{\mathcal{O}_{Y, y}} \mathcal{O}_{X, x})
\end{align*}

D'où l'équivalence du lemme résulte de la deuxième caractérisation de la planéité dans Algèbre, lemme \ref{algebra-lemma-flat}.

\end{proof}

\begin{lemma}

\label{lemma-flat-pullback-support}
Laissez $f : Y \to X$ être un morphisme de schémas. Laissez $\mathcal{F}$ être un type quasi-cohérent $\mathcal{O}_X$-module avec support théorique de schéma $Z \subset X$. Si $f$ est plat, alors $f^{-1}(Z)$ est le schéma de support de $f^*\mathcal{F}$.
\end{lemma}

\begin{proof}
En utilisant la caractérisation du support schématique sur les affines donnée dans lemme \ref{lemma-scheme-theoretic-support} nous réduisons à Algèbre, lemme \ref{algebra-lemma-annihilator-flat-base-change}.
\end{proof}

\begin{lemma}

\label{lemma-flat-morphism-scheme-theoretically-dense-open}
Laissez $f : X \to Y$ être un plateau morphisme de schémas. Laissez $V \subset Y$ être un rétrocompact ouvert qui est dense de schématiquement. Puis $f^{-1}V$ est dense de schématiquement dans $X$.
\end{lemma}

\begin{proof}

Nous utiliserons la caractérisation de lemme \ref{lemma-characterize-scheme-theoretically-dense}. Nous devons montrer que pour toute ouverture $U \subset X$ l'application
$\mathcal{O}_X(U) \to \mathcal{O}_X(U \cap f^{-1}V)$ Il suffit de le prouver quand $U$ est un ouvert affine qui cartographie dans un ouvert affine $W \subset Y$. Dites $W = \Spec(A)$ et $U = \Spec(B)$. Ensuite $V \cap W = D(f_1) \cup \ldots \cup D(f_n)$ pour certains
$f_i \in A$, voir Algèbre, lemme \ref{algebra-lemma-qc-open}. Ainsi, nous devons montrer que
$B \to B_{f_1} \times \ldots \times B_{f_n}$ est injectable. On nous donne que $A \to A_{f_1} \times \ldots \times A_{f_n}$ est injectable et que $A \to B$ est plat. Depuis $B_{f_i} = A_{f_i} \otimes_A B$ Nous avons gagné.

\end{proof}

\begin{lemma}

\label{lemma-flat-base-change-scheme-theoretic-image}

\begin{slogan}
Prendre des images théoriques de schéma commute avec le changement plat de base dans le cas quasi-compact
\end{slogan}
Laisser $f : X \to Y$ être un plateau morphiste de schémas. Laisser $g : V \to Y$ être un morphisme quasi-compact de schémas. Laisser $Z \subset Y$ être l'image théorique de schéma de $g$ et laisser $Z' \subset X$ être l'image schématique du changement de base $V \times_Y X \to X$. Puis $Z' = f^{-1}Z$.
\end{lemma}

\begin{proof}

Rappelez-vous que $Z$ est coupé par
$\mathcal{I} = \Ker(\mathcal{O}_Y \to g_*\mathcal{O}_V)$
et $Z'$ est coupé par
$\mathcal{I}' = \Ker(\mathcal{O}_X \to
(V \times_Y X \to X)_*\mathcal{O}_{V \times_Y X})$, voir lemme \ref{lemma-quasi-compact-scheme-theoretic-image}. Par conséquent, la question est locale sur $X$ et $Y$ et nous pouvons supposer $X$ et $Y$
affine. Notez que nous pouvons remplacer $V$ par $\coprod V_i$ où
$V = V_1 \cup \ldots \cup V_n$ est un recouvrement fini ouvert affine. $g$ Dans ce cas,
$(V \times_Y X \to X)_*\mathcal{O}_{V \times_Y X}$ est le retrait de $g_*\mathcal{O}_V$ par $f$Depuis $f$ est plat nous concluons que
$f^*\mathcal{I} = \mathcal{I}'$ et le lemme tient.

\end{proof}




\section{Immersions fermées plates}
\label{section-flat-closed-immersions}

\noindent
Les composants connectés des schémas ne sont pas toujours ouverts. Mais ils ont toujours une structure de schéma canonique. Nous l'expliquons dans cette section.

\begin{lemma}

\label{lemma-characterize-flat-closed-immersions}
Laissez $X$ être un schéma. La règle qui associe à un sous-système fermé de $X$ son sous-ensemble fermé sous-jacent définit un bijection $$
\left\{
\begin{matrix}
\text{closed subschemes }Z \subset X \\
\text{such that }Z \to X\text{ is flat}
\end{matrix}
\right\}
\leftrightarrow
\left\{
\begin{matrix}
\text{closed subsets }Z \subset X \\
\text{closed under generalizations}
\end{matrix}
\right\}
$$ Si $Z \subset X$ est un tel sous-système fermé, chaque morphisme de schémas $g : Y \to X$ avec $g(Y) \subset Z$ ensemble théoriquement des facteurs (système théoriquement) par $Z$.
\end{lemma}

\begin{proof}

Le boîtier affiné du bijection est Algèbre, lemme \ref{algebra-lemma-pure-open-closed-specializations}. Pour les régimes généraux $X$ le bijection suit en couvrant $X$ omis. Pour l'affirmation finale, observez que la projection $Z \times_{X, g} Y \to Y$ est un plat (lemme \ref{lemma-base-change-flat}) immersion fermée qui est bijective sur les espaces topologiques sous-jacents et doit donc être un isomorphisme par le bijection établi dans la première partie de la démonstration.

\end{proof}

\begin{lemma}

\label{lemma-flat-closed-immersions-finite-presentation}

Une immersion à plat fermée de présence finie est l'ouverture d'une immersion ouverte et fermée.

\end{lemma}

\begin{proof}

L'affaire affine est Algèbre, lemme \ref{algebra-lemma-finitely-generated-pure-ideal}. En général, le lemme suit en couvrant $X$ Les détails sont omis.

\end{proof}

\noindent
Notez qu'un composant connecté $T$ d'un schéma $X$ est une pargénérisation stable en sous-ensemble fermé.

\begin{definition}

\label{definition-scheme-structure-connected-component}
Laissez $X$ être un schéma. Laissez $T \subset X$ être un composant connecté. La structure de schéma canonique {\it sur $T$} est la structure de schéma unique sur $T$ de sorte que l'immersion fermée $T \to X$ est plat, voir lemme \ref{lemma-characterize-flat-closed-immersions}.
\end{definition}

\noindent
Il s'avère que nous pouvons déterminer quand chaque module $\mathcal{O}_X$ plat fini est libre localement en utilisant le lemme précédent.

\begin{lemma}

\label{lemma-finite-flat-is-finite-locally-free}
Laissez $X$ être un schéma. Ce qui suit sont équivalents
\begin{enumerate}
\item chaque module $\mathcal{O}_X$ plat quasi-cohérent fini est libre localement, et
\item chaque sous-ensemble fermé $Z \subset X$ qui est fermé sous générisations est ouvert.
\end{enumerate}

\end{lemma}

\begin{proof}
Dans le cas affine c'est Algèbre, lemme \ref{algebra-lemma-finite-flat-module-finitely-presented}. Le cas des schémas ne découle pas directement du cas affine, donc nous répétons simplement les arguments.

\medskip\noindent
Supposons (1). Considérez une immersion fermée $i : Z \to X$ telle que $i$ est plat. Alors $i_*\mathcal{O}_Z$ est quasi-cohérent et plat, donc fini localement libre par (1). Ainsi $Z = \text{Supp}(i_*\mathcal{O}_Z)$ est également ouvert et nous voyons que (2) tient. D'où l'implication (1) $\Rightarrow$ (2) découle de la caractérisation des fermées d'immersion plate dans leme \ref{lemma-characterize-flat-closed-immersions}.

\medskip\noindent

Pour l'inverse, supposer que $X$ (2). Laisser $\mathcal{F}$ être un plat fini quasi-cohérent $\mathcal{O}_X$- Module. Le soutien $Z = \text{Supp}(\mathcal{F})$ de $\mathcal{F}$ est fermé, voir Modules, lemme \ref{modules-lemma-support-finite-type-closed}. D'autre part, si $x \leadsto x'$ est une spécialisation, puis par Algèbre, lemme \ref{algebra-lemma-finite-flat-local}
le module $\mathcal{F}_{x'}$ est libre $\mathcal{O}_{X, x'}$,
$$
\mathcal{F}_x =
\mathcal{F}_{x'} \otimes_{\mathcal{O}_{X, x'}} \mathcal{O}_{X, x}.
$$
C'est pourquoi
$x' \in \text{Supp}(\mathcal{F}) \Rightarrow x \in \text{Supp}(\mathcal{F})$, c'est-à-dire que le soutien est fermé sous générisation. $X$ (2) nous voyons que le soutien de $\mathcal{F}$
est ouvert et fermé. Les modules $\wedge^i(\mathcal{F})$, $i = 1, 2, 3, \ldots$
sont quasi-cohérents plats finis $\mathcal{O}_X$-modules aussi, voir Modules, section \ref{modules-section-symmetric-exterior}. Notez que
$\text{Supp}(\wedge^{i + 1}(\mathcal{F})) \subset
\text{Supp}(\wedge^i(\mathcal{F}))$. Ainsi, nous voyons qu'il existe une décomposition
$$
X = U_0 \amalg U_1 \amalg U_2 \amalg \ldots
$$
par des sous-ensembles ouverts et fermés de sorte que le soutien de
$\wedge^i(\mathcal{F})$ est $U_i \cup U_{i + 1} \cup \ldots$ pour tous $i$. Laisse $x$ être un point de vue $X$, et dire $x \in U_r$. Notez que
$\wedge^i(\mathcal{F})_x \otimes \kappa(x) =
\wedge^i(\mathcal{F}_x \otimes \kappa(x))$. Par conséquent, $x \in U_r$ implique que $\mathcal{F}_x \otimes \kappa(x)$
est un espace vectoriel de dimension $r$. Par le lemme de Nakayama, voir Algèbre, lemme \ref{algebra-lemma-NAK}
nous pouvons choisir un quartier affiné ouvert $U \subset U_r \subset X$
de $x$ et sections $s_1, \ldots, s_r \in \mathcal{F}(U)$ telle que l'application induite
$$
\mathcal{O}_U^{\oplus r} \longrightarrow \mathcal{F}|_U, \quad
(f_1, \ldots, f_r) \longmapsto \sum f_i s_i
$$
c'est surjectif. Cela signifie que
$\wedge^r(\mathcal{F}|_U)$ est un plat fini quasi-cohérent
$\mathcal{O}_U$-module dont le soutien est tout de $U$. Par ce qui précède, il est généré par un seul élément, à savoir
$s_1 \wedge \ldots \wedge s_r$. D'où
$\wedge^r(\mathcal{F}|_U) \cong \mathcal{O}_U/\mathcal{I}$
pour quelques faisceau quasi-cohérents d'idéaux $\mathcal{I}$
de telle manière que $\mathcal{O}_U/\mathcal{I}$ est plat sur $\mathcal{O}_U$ et de telle sorte que $V(\mathcal{I}) = U$Il s'ensuit que $\mathcal{I} = 0$ par l'application \ref{lemma-characterize-flat-closed-immersions}. Ainsi $s_1 \wedge \ldots \wedge s_r$ est une base pour
$\wedge^r(\mathcal{F}|_U)$ et il s'ensuit que l'application affichée est aussi bien injective que surjective. Cela prouve que $\mathcal{F}$ est fini localement libre comme désiré.

\end{proof}






\section{Platitude générique}
\label{section-generic-flatness}

\noindent
Un schéma de type fini sur une base intégrale est plat sur une ouverture dense de la base. Dans Algèbre, section \ref{algebra-section-generic-flatness} nous avons prouvé une version noethérien, une version pour les morphismes de présentation finie, et une version générale. Nous n'affirmons et prouvons ici la version générale. Cependant, il s'avère que cela sera remplacé par la proposition \ref{proposition-generic-flatness-reduced} qui montre le résultat tient si nous supposons que la base est réduite.

\begin{proposition}
[platitude générale]

\label{proposition-generic-flatness}

Laisser $f : X \to S$ être un morphisme de schémas. $\mathcal{F}$ être un faisceau quasi-cohérent de $\mathcal{O}_X$- Modules.

\begin{enumerate}

\item  $S$ fait partie intégrante,
\item  $f$ est de type fini, et
\item  $\mathcal{F}$ est un type fini $\mathcal{O}_X$- Module.

\end{enumerate}

Puis il existe un sous-système dense ouvert $U \subset S$ de telle manière que
$X_U \to U$ est plat et de présentation finie et telle que
$\mathcal{F}|_{X_U}$ est plat sur $U$ et de présentation finie sur $\mathcal{O}_{X_U}$.

\end{proposition}

\begin{proof}

Comme $S$ Il est irréductible (voir Propriétés, lemme \ref{properties-lemma-characterize-integral}) et toute ouverture non vide est dense.
$S$ par une affinité d'ouverture de $S$ et supposer que $S = \Spec(A)$ C'est affine. $S$ fait partie intégrante nous voyons que $A$ est un domaine. $f$ est de type fini, il est quasi-compact, donc $X$ c'est quasi-compact. On peut donc trouver une couverture ouvert affine finie
$X = \bigcup_{i = 1, \ldots, n} X_i$. Écrire $X_i = \Spec(B_i)$. Ensuite $B_i$ est un type fini $A$-l'algèbre, voyez-moi. \ref{lemma-locally-finite-type-characterize}. De plus, il y a le type fini
$B_i$-modules $M_i$ de telle manière que $\mathcal{F}|_{X_i}$ est le faisceau quasi-cohérent associé au $B_i$-module $M_i$, voir Propriétés, laisse \ref{properties-lemma-finite-type-module}. Suivant, pour chaque paire d'indices $i, j$ choisir un idéal $I_{ij} \subset B_i$
de telle manière que $X_i \setminus X_i \cap X_j = V(I_{ij})$ à l'intérieur
$X_i = \Spec(B_i)$. Prêt $M_{ij} = B_i/I_{ij}$ et y penser comme un $B_i$- Module. $V(I_{ij}) = \text{Supp}(M_{ij})$
et $M_{ij}$ est un fini $B_i$- Module.

\medskip\noindent
A ce stade, nous appliquons Algèbre, lemme \ref{algebra-lemma-generic-flatness} les paires $(A \to B_i, M_{ij})$ et les paires $(A \to B_i, M_i)$. Ainsi, nous obtenons $f_{ij}, f_i \in A$ non zéro de telle sorte que (a) $A_{f_{ij}} \to B_{i, f_{ij}}$ est plat et de présentation finie et $M_{ij, f_{ij}}$ est plat sur $A_{f_{ij}}$ et de présentation finie sur $B_{i, f_{ij}}$, et (b) $B_{i, f_i}$ est plat et de présentation finie sur $A_f$ et $M_{i, f_i}$ est plat et de présentation finie sur $B_{i, f_i}$. Posons $f = (\prod f_i) (\prod f_{ij})$. Nous affirmons que la prise de $U = D(f)$ fonctionne.

\medskip\noindent

Pour prouver notre revendication, nous pouvons remplacer $A$ par $A_f$, c'est-à-dire effectuer le changement de base par $U = \Spec(A_f) \to S$. Après ce changement de base nous voyons que chacun de $A \to B_i$ est plat et de présentation finie et que $M_i$, $M_{ij}$ sont plats sur $A$
et de présentation finie sur $B_i$. Cela prouve déjà que $X \to S$ est quasi-compact, localment de présentation finie, plat, et que $\mathcal{F}$
est plat sur $S$ et de présentation finie sur $\mathcal{O}_X$, voir lemme \ref{lemma-locally-finite-presentation-characterize}
et Propriétés, lemme \ref{properties-lemma-finite-presentation-module}Depuis $M_{ij}$ est de présentation finie sur $B_i$ nous voyons que
$X_i \cap X_j = X_i \setminus \text{Supp}(M_{ij})$ est un quasi-compact ouvert de $X_i$, voir Algèbre, lemme \ref{algebra-lemma-support-finite-presentation-constructible}. Nous voyons donc que $X \to S$ est quasi-séparé par Schémas, lemme \ref{schemes-lemma-characterize-quasi-separated}. Cela prouve la proposition.

\end{proof}

\noindent
Il s'avère en fait qu'il y a aussi une version de planéité générique sur une base réduite arbitraire.

\begin{proposition}
[platitude générale, cas réduit]

\label{proposition-generic-flatness-reduced}

Laisser $f : X \to S$ être un morphisme de schémas. $\mathcal{F}$ être un faisceau quasi-cohérent de $\mathcal{O}_X$- Modules.

\begin{enumerate}

\item  $S$ est réduit,
\item  $f$ est de type fini, et
\item  $\mathcal{F}$ est un type fini $\mathcal{O}_X$- Module.

\end{enumerate}

Puis il existe un sous-système dense ouvert $U \subset S$ de telle manière que
$X_U \to U$ est plat et de présentation finie et telle que
$\mathcal{F}|_{X_U}$ est plat sur $U$ et de présentation finie sur $\mathcal{O}_{X_U}$.

\end{proposition}

\begin{proof}
Pour le lecteur impatient: Cette démonstration est une répétition de la démonstration de la proposition \ref{proposition-generic-flatness} en utilisant Algèbre, lemme \ref{algebra-lemma-generic-flatness-reduced} au lieu d'Algèbre, lemme \ref{algebra-lemma-generic-flatness}.

\medskip\noindent

Comme être plat et être de présentation finie est local sur la base, voir lemmes \ref{lemma-flat-module-characterize} et
\ref{lemma-locally-finite-presentation-characterize}, nous pouvons travailler localement sur $S$. Ainsi, nous pouvons supposer que
$S = \Spec(A)$où: $A$ est une bague réduite (voir Propriétés, \ref{properties-lemma-characterize-reduced}) . $f$ est de type fini, il est quasi-compact, donc $X$ c'est quasi-compact. On peut donc trouver une couverture ouvert affine finie
$X = \bigcup_{i = 1, \ldots, n} X_i$. Écrire $X_i = \Spec(B_i)$. Ensuite $B_i$ est un type fini $A$-l'algèbre, voyez-moi. \ref{lemma-locally-finite-type-characterize}. De plus, il y a le type fini
$B_i$-modules $M_i$ de telle manière que $\mathcal{F}|_{X_i}$ est le faisceau quasi-cohérent associé au $B_i$-module $M_i$, voir Propriétés, laisse \ref{properties-lemma-finite-type-module}. Suivant, pour chaque paire d'indices $i, j$ choisir un idéal $I_{ij} \subset B_i$
de telle manière que $X_i \setminus X_i \cap X_j = V(I_{ij})$ à l'intérieur
$X_i = \Spec(B_i)$. Prêt $M_{ij} = B_i/I_{ij}$ et y penser comme un $B_i$- Module. $V(I_{ij}) = \text{Supp}(M_{ij})$
et $M_{ij}$ est un fini $B_i$- Module.

\medskip\noindent
A ce stade, nous appliquons Algèbre, lemme \ref{algebra-lemma-generic-flatness-reduced} aux paires $(A \to B_i, M_{ij})$ et aux paires $(A \to B_i, M_i)$. Ainsi, nous obtenons des ouvertures denses $U(A \to B_i, M_{ij}) \subset S$ et denses $U(A \to B_i, M_i) \subset S$ avec notation comme en Algèbre, Équation (\ref{algebra-equation-good-locus}). Comme une intersection finie d'ouvertures denses est dense ouverte, nous voyons que $$
U =
\bigcap\nolimits_{i, j} U(A \to B_i, M_{ij})
\quad\cap\quad
\bigcap\nolimits_i U(A \to B_i, M_i)
$$ est ouvert et dense dans $S$. Nous affirmons que $U$ est l'ouverture souhaitée.

\medskip\noindent

Choisissez $u \in U$. Par définition des loci $U(A \to B_i, M_{ij})$
et $U(A \to B, M_i)$ il existe $f_{ij}, f_i \in A$ de telle manière que a) $u \in D(f_i)$ et $u \in D(f_{ij})$, (b) $A_{f_{ij}} \to B_{i, f_{ij}}$ est plat et de présentation finie et $M_{ij, f_{ij}}$ est plat sur $A_{f_{ij}}$ et de présentation finie sur $B_{i, f_{ij}}$, et c) $B_{i, f_i}$ est plat et de présentation finie sur $A_{f_i}$ et
$M_{i, f_i}$ est de présentation finie sur $B_{i, f_i}$ et plat sur
$A_{f_i}$. Prêt
$f = (\prod f_i) (\prod f_{ij})$. Maintenant il suffit de prouver que $X \to S$ est plat et de présentation finie sur $D(f)$ et que $\mathcal{F}$ limités à $X_{D(f)}$ est plat sur $D(f)$ et de présentation finie sur la $X_{D(f)}$.

\medskip\noindent

C'est pourquoi nous pouvons remplacer $A$ par $A_f$, c'est-à-dire effectuer le changement de base par $\Spec(A_f) \to S$. Après ce changement de base nous voyons que chacun de $A \to B_i$ est plat et de présentation finie et que $M_i$, $M_{ij}$ sont plats sur $A$
et de présentation finie sur $B_i$. Cela prouve déjà que $X \to S$ est quasi-compact, localment de présentation finie, plat, et que $\mathcal{F}$
est plat sur $S$ et de présentation finie sur $\mathcal{O}_X$, voir lemme \ref{lemma-locally-finite-presentation-characterize}
et Propriétés, lemme \ref{properties-lemma-finite-presentation-module}Depuis $M_{ij}$ est de présentation finie sur $B_i$ nous voyons que
$X_i \cap X_j = X_i \setminus \text{Supp}(M_{ij})$ est un quasi-compact ouvert de $X_i$, voir Algèbre, lemme \ref{algebra-lemma-support-finite-presentation-constructible}. Nous voyons donc que $X \to S$ est quasi-séparé par Schémas, lemme \ref{schemes-lemma-characterize-quasi-separated}. Cela prouve la proposition.

\end{proof}









\begin{remark}

\label{remark-flattening}

Les résultats ci-dessus sont une première étape vers des techniques d'aplanissement plus raffinées pour les morphismes des schémas. \cite{GruRay} par Raynaud et Gruson contient de nombreux résultats merveilleux dans cette direction.

\end{remark}







\section{Morphismes et dimensions des fibres}
\label{section-dimension-fibres}

\noindent
Laissez $X$ être un espace topologique, et $x \in X$. Rappelons que nous avons défini $\dim_x(X)$ comme le minimum des dimensions des quartiers ouverts de $x$ dans $X$. Voir Topologie, définition \ref{topology-definition-Krull}.

\begin{lemma}

\label{lemma-dimension-fibre-at-a-point}
Laisser $f : X \to S$ être un morphisme de schémas. Laisser $x \in X$ et définir $s = f(x)$. Supposons que $f$ est localement de type fini. Puis $$
\dim_x(X_s) =
\dim(\mathcal{O}_{X_s, x}) + \text{trdeg}_{\kappa(s)}(\kappa(x)).
$$
\end{lemma}

\begin{proof}
Cela réduit immédiatement au cas $S = s$, et $X$ affine. Dans ce cas, le résultat suit d'Algèbre, lemme \ref{algebra-lemma-dimension-at-a-point-finite-type-field}.
\end{proof}

\begin{lemma}

\label{lemma-dimension-fibre-at-a-point-additive}
Laisser $f : X \to Y$ et $g : Y \to S$ être des morphismes de schémas. Laisser $x \in X$ et définir $y = f(x)$, $s = g(y)$. Supposons que $f$ et $g$ localement de type fini. Puis $$
\dim_x(X_s) \leq \dim_x(X_y) + \dim_y(Y_s).
$$ En outre, l'égalité est maintenue si $\mathcal{O}_{X_s, x}$ est plat sur $\mathcal{O}_{Y_s, y}$, qui tient par exemple si $\mathcal{O}_{X, x}$ est plat sur $\mathcal{O}_{Y, y}$.
\end{lemma}

\begin{proof}
Notez que $\text{trdeg}_{\kappa(s)}(\kappa(x)) =
\text{trdeg}_{\kappa(y)}(\kappa(x)) + \text{trdeg}_{\kappa(s)}(\kappa(y))$. Ainsi par lemme \ref{lemma-dimension-fibre-at-a-point} l'énoncé est équivalent à $$
\dim(\mathcal{O}_{X_s, x})
\leq
\dim(\mathcal{O}_{X_y, x}) + \dim(\mathcal{O}_{Y_s, y}).
$$ Pour ceci voir Algèbre, lemme \ref{algebra-lemma-dimension-base-fibre-total}. Pour le cas plat voir Algèbre, lemme \ref{algebra-lemma-dimension-base-fibre-equals-total}.
\end{proof}

\begin{lemma}

\label{lemma-dimension-fibre-after-base-change}

Laisser
$$
\xymatrix{
X' \ar[r]_{g'} \ar[d]_{f'} & X \ar[d]^f \\
S' \ar[r]^g & S
}
$$
être un produit fibré diagramme des schémas. Supposons que $f$ localement de type fini. Supposons que $x' \in X'$, $x = g'(x')$, $s' = f'(x')$ et
$s = g(s') = f(x)$. Ensuite

\begin{enumerate}

\item  $\dim_x(X_s) = \dim_{x'}(X'_{s'})$,
\item si $F$ est la fibre du morphisme $X'_{s'} \to X_s$
terminé $x$Alors
$$
\dim(\mathcal{O}_{F, x'}) =
\dim(\mathcal{O}_{X'_{s'}, x'}) - \dim(\mathcal{O}_{X_s, x}) =
\text{trdeg}_{\kappa(s)}(\kappa(x)) -
\text{trdeg}_{\kappa(s')}(\kappa(x'))
$$
En particulier $\dim(\mathcal{O}_{X'_{s'}, x'}) \geq \dim(\mathcal{O}_{X_s, x})$
et $\text{trdeg}_{\kappa(s)}(\kappa(x)) \geq
\text{trdeg}_{\kappa(s')}(\kappa(x'))$.
\item données $s', s, x$ il existe un choix de $x'$ de telle manière que
$\dim(\mathcal{O}_{X'_{s'}, x'}) = \dim(\mathcal{O}_{X_s, x})$ et
$\text{trdeg}_{\kappa(s)}(\kappa(x)) = \text{trdeg}_{\kappa(s')}(\kappa(x'))$.

\end{enumerate}

\end{lemma}

\begin{proof}
La partie 1) suit immédiatement d'Algèbre, lemme \ref{algebra-lemma-dimension-at-a-point-preserved-field-extension}. Les parties 2) et 3) d'Algèbre, lemme \ref{algebra-lemma-inequalities-under-field-extension}.
\end{proof}

\noindent
Le lemme suivant suit d'un résultat algébrique non trivial. C'est-à-dire, la version algébrique du théorème principal de Zariski.

\begin{lemma}

\label{lemma-openness-bounded-dimension-fibres}

\begin{reference}
\cite[IV théorème 13.1.3]{EGA}
\end{reference}
Laissez $f : X \to S$ être un morphisme de schémas. Laissez $n \geq 0$. Supposez que $f$ est localement de type fini. Le jeu $$
U_n = \{x \in X \mid \dim_x X_{f(x)} \leq n\}
$$ est ouvert dans $X$.
\end{lemma}

\begin{proof}
Ceci est immédiat d'Algèbre, lemme \ref{algebra-lemma-dimension-fibres-bounded-open-upstairs}
\end{proof}

\begin{lemma}

\label{lemma-morphism-finite-type-bounded-dimension}
Laissez $f : X \to Y$ être un morphisme de type fini avec $Y$ quasi-compact. Puis la dimension des fibres de $f$ est limitée.
\end{lemma}

\begin{proof}

Par moi \ref{lemma-openness-bounded-dimension-fibres}
l'ensemble $U_n \subset X$ des points où la dimension de la fibre est $\leq n$ est ouvert. Depuis $f$ est de type fini, chaque point est contenu dans certains $U_n$ (parce que la dimension d'un type d'algèbre fini sur un champ est finie). $Y$ est quasi compacte et $f$ est de type fini, nous voyons que $X$ est quasi-compact. $X = U_n$ pour certains $n$.

\end{proof}

\begin{lemma}

\label{lemma-openness-bounded-dimension-fibres-finite-presentation}

Laisser $f : X \to S$ être un morphisme de schémas. $n \geq 0$Supposons que $f$ est localement de présentation finie.
$$
U_n = \{x \in X \mid \dim_x X_{f(x)} \leq n\}
$$
de l'eau de mer \ref{lemma-openness-bounded-dimension-fibres} est rétrocompact dans $X$(Voir Topologie, définition) \ref{topology-definition-quasi-compact}.)

\end{lemma}

\begin{proof}

L'espace topologique $X$ a une base pour sa topologie consistant en ouverts affines $U \subset X$ tel que le morphisme induit
$f|_U : U \to S$ facteurs à travers un ouvert affine $V \subset S$. Il suffit donc de montrer que $U \cap U_n$ est quasi-compact pour un tel $U$. Notez que $U_n \cap U$ est la même que l'ouverture
$\{x \in U \mid \dim_x U_{f(x)} \leq n\}$. Cela nous réduit au cas où $X$ et $S$ dans ce cas le lemme suit d'Algèbre, lemme \ref{algebra-lemma-dimension-fibres-bounded-quasi-compact-open-upstairs}
(et lemme \ref{lemma-locally-finite-presentation-characterize}).

\end{proof}

\begin{lemma}

\label{lemma-dimension-fibre-specialization}

Laisser $f : X \to S$ être un morphisme de schémas. $x \leadsto x'$ être une spécialisation non triviale des points dans $X$
couché sur le même point $s \in S$Supposons que $f$ est localement de type fini. Puis

\begin{enumerate}

\item  $\dim_x(X_s) \leq \dim_{x'}(X_s)$,
\item  $\dim(\mathcal{O}_{X_s, x}) < \dim(\mathcal{O}_{X_s, x'})$,
\item  $\text{trdeg}_{\kappa(s)}(\kappa(x)) >
\text{trdeg}_{\kappa(s)}(\kappa(x'))$.

\end{enumerate}

\end{lemma}

\begin{proof}
Part (1) découle du fait que toute ouverture de $X_s$ contenant $x'$ contient également $x$. Part (2) suit puisque $\mathcal{O}_{X_s, x}$ est une localisation de $\mathcal{O}_{X_s, x'}$ à un idéal premier, d'où toute chaîne d'idéaux principaux dans $\mathcal{O}_{X_s, x}$ fait partie d'une chaîne de primes strictement plus longue dans $\mathcal{O}_{X_s, x'}$. La dernière inégalité suit d'Algèbre, lemme \ref{algebra-lemma-tr-deg-specialization}.
\end{proof}






\section{Morphismes de dimension relative donnée}
\label{section-relative-dimension}

\noindent
Afin de pouvoir parler confortablement des morphismes d'une dimension donnée relative, nous introduisons la notion suivante.

\begin{definition}

\label{definition-relative-dimension-d}

Laisser $f : X \to S$ être un morphisme des schémas. $f$ est localement de type fini.

\begin{enumerate}

\item Nous disons $f$ est de {\it dimension relative $\leq d$ au $x$} si $\dim_x(X_{f(x)}) \leq d$.
\item Nous disons $f$ est de {\it dimension relative $\leq d$} si $\dim_x(X_{f(x)}) \leq d$ pour tous $x \in X$.
\item Nous disons $f$ est de {\it dimension relative $d$} si toutes les fibres non vides $X_s$ sont équidimensionnelles $d$.

\end{enumerate}

\end{definition}

\noindent
Ce n'est pas une notion particulièrement bien conduite, mais elle fonctionne bien dans un certain nombre de situations.

\begin{lemma}

\label{lemma-base-change-relative-dimension-d}
Laissez $f : X \to S$ être un morphisme de schémas qui est localement de type fini. Si $f$ a une dimension relative $d$, alors tout changement de base de $f$. Idem pour la dimension relative $\leq d$.
\end{lemma}

\begin{proof}
Ceci est immédiat de lemme \ref{lemma-dimension-fibre-after-base-change}.
\end{proof}

\begin{lemma}

\label{lemma-composition-relative-dimension-d}

Laisser $f : X \to Y$, $g : Y \to Z$ être localement de type fini. Si $f$ a dimension relative $\leq d$ et $g$ a dimension relative $\leq e$
alors $g \circ f$ a dimension relative $\leq d + e$Si

\begin{enumerate}

\item  $f$ a dimension relative $d$,
\item  $g$ a dimension relative $e$,
\item  $f$ est plat,

\end{enumerate}

alors $g \circ f$ a dimension relative $d + e$.

\end{lemma}

\begin{proof}
Ceci est immédiat de lemme \ref{lemma-dimension-fibre-at-a-point-additive}.
\end{proof}

\noindent
En général, il n'est pas possible de décomposer un morphisme en morceaux où la dimension relative est donnée. Cependant, il est possible si le morphisme a des fibres de Cohen--Macaulay et est plat de présentation finie.

\begin{lemma}

\label{lemma-flat-finite-presentation-CM-fibres-relative-dimension}

\begin{slogan}

de Cohen--Les morphismes Macaulay se décomposent en clopens de dimension pure relative

\end{slogan}

Laisser $f : X \to S$ être un morphisme de schémas. Supposons que

\begin{enumerate}

\item  $f$ est plat,
\item  $f$ est localement de présentation finie, et
\item pour tous $s \in S$ les fibres $X_s$ est de Cohen--Macaulay (Propriétés, définition) \ref{properties-definition-Cohen-Macaulay})

\end{enumerate}

Ensuite, il existe des sous-systèmes ouverts et fermés $X_d \subset X$
de telle manière que $X = \coprod_{d \geq 0} X_d$ et $f|_{X_d} : X_d \to S$
a dimension relative $d$.

\end{lemma}

\begin{proof}
Ceci est immédiat d'Algèbre, lemme \ref{algebra-lemma-relative-dimension-CM}.
\end{proof}

\begin{lemma}

\label{lemma-locally-quasi-finite-rel-dimension-0}
Laissez $f : X \to S$ être un morphisme de schémas. Supposons que $f$ est localement de type fini. Laissez $x \in X$ avec $s = f(x)$. Alors $f$ est quasi-fini à $x$ si et seulement si $\dim_x(X_s) = 0$. En particulier, $f$ est localement quasi-fini si et seulement si $f$ a une dimension relative $0$.
\end{lemma}

\begin{proof}

Première preuve. $f$ est quasi-fini à $x$ alors $\kappa(x)$ est une extension finie de
$\kappa(s)$ (par la suite) \ref{lemma-residue-field-quasi-finite}) et $x$ est isolé dans $X_s$ (par la suite) \ref{lemma-quasi-finite-at-point-characterize}) et donc $\dim_x(X_s) = 0$ par lemme \ref{lemma-dimension-fibre-at-a-point}. Inversement, si $\dim_x(X_s) = 0$ puis par moi \ref{lemma-dimension-fibre-at-a-point}
nous voyons $\kappa(s) \subset \kappa(x)$ est algébrique et il n'y a pas d'autres points de $X_s$ spécialisée dans $x$. D'où $x$ est fermée dans sa fibre par Lemme \ref{lemma-algebraic-residue-field-extension-closed-point-fibre}
et par la suite \ref{lemma-quasi-finite-at-point-characterize} 3) nous concluons que $f$ est quasi-fini à $x$.

\medskip\noindent
Deuxième épreuve. La fibre $X_s$ est un schéma localement de type fini sur un champ, donc localement noethérien (lemme \ref{lemma-finite-type-noetherian}). Le résultat suit maintenant de lemme \ref{lemma-quasi-finite-at-point-characterize} et Propriétés, lemme \ref{properties-lemma-isolated-dim-0-locally-noetherian}.
\end{proof}

\begin{lemma}

\label{lemma-rel-dimension-dimension}
Laissez $f : X \to Y$ être un morphisme des schémas de localement noethérien qui est plat, localement de type fini et de dimension relative $d$. Pour chaque point $x$ dans $X$ avec l'image $y$ dans $Y$ nous avons $\dim_x(X) = \dim_y(Y) + d$.
\end{lemma}

\begin{proof}
Après avoir réduit $X$ et $Y$ aux quartiers ouverts de $x$ et $y$, nous pouvons supposer que $\dim(X) = \dim_x(X)$ et $\dim(Y) = \dim_y(Y)$, par définition de la dimension d'un schéma à un point (Propriétés, définition \ref{properties-definition-dimension}). Le morphisme $f$ est ouvert par lemmes \ref{lemma-noetherian-finite-type-finite-presentation} et \ref{lemma-fppf-open}. Par conséquent, nous pouvons réduire $Y$ pour arranger que $f$ est surjectif. Il reste à montrer que $\dim(X) = \dim(Y) + d$.

\medskip\noindent
Laissez $a$ être un point dans $X$ avec l'image $b$ dans $Y$. Par Algèbre, lemme \ref{algebra-lemma-dimension-base-fibre-equals-total}, $$
\dim(\mathcal{O}_{X,a}) = \dim(\mathcal{O}_{Y,b}) + \dim(\mathcal{O}_{X_b, a}).
$$ Prendre le supremum sur tous les points $a$ dans $X$, il s'ensuit que $\dim(X) = \dim(Y) + d$, comme nous le voulons, voir Propriétés, lemme \ref{properties-lemma-dimension}.
\end{proof}











\section{Morphismes syntomiques}
\label{section-syntomic}

\noindent

Une algèbre $A$ sur un champ $k$ est appelé un {\it intersection globale complète sur $k$} si $A \cong k[x_1, \ldots, x_n]/(f_1, \ldots, f_c)$ et
$\dim(A) = n - c$. Une algèbre $A$ sur un champ $k$ est appelé un {\it intersection locale complète} si $\Spec(A)$
peuvent être couverts par des ouvertures standard dont chacune sont des intersections complètes mondiales sur $k$. Voir Algèbre, section
\ref{algebra-section-lci}. Rappelons qu'un homomorphisme d'anneaux $R \to A$
est {\it syntomique} s'il s'agit de présentation finie, plat avec des anneaux d'intersection complets locaux en fibres, voir Algèbre, définition \ref{algebra-definition-lci}.

\begin{definition}

\label{definition-syntomic}

Laisser $f : X \to S$ être un morphisme de schémas.

\begin{enumerate}

\item Nous disons que $f$ est {\it syntomique à $x \in X$} s'il existe un quartier affiné d'ouverture $\Spec(A) = U \subset X$
de $x$ et ouvert affine $\Spec(R) = V \subset S$
avec $f(U) \subset V$ tel que l'homomorphisme d'anneaux induit
$R \to A$ C'est syntomique.
\item Nous disons que $f$ est {\it syntomique} si elle est syntomique à chaque point $X$.
\item Si $S = \Spec(k)$ et $f$ est syntomique, alors nous disons que
$X$ est un {\it intersection locale complète sur $k$}.
\item Un morphisme des schémas affinés $f : X \to S$
est appelé {\it syntomique standard} s'il existe une intersection globale relative complète
$R \to R[x_1, \ldots, x_n]/(f_1, \ldots, f_c)$ (voir Algèbre, définition) \ref{algebra-definition-relative-global-complete-intersection}) de telle sorte que $X \to S$ est isomorphe à
$$
\Spec(R[x_1, \ldots, x_n]/(f_1, \ldots, f_c)) \to \Spec(R).
$$

\end{enumerate}

\end{definition}

\noindent
Dans la littérature, un morphisme syntomique est parfois appelé {\it morphisme local complet d'intersection}. Il s'avère qu'il s'agit d'une classe commode de morphismes. Par exemple, on peut définir une topologie syntomique en utilisant celles-ci, qui est plus fine que les topologies lisse et \'etale, mais a beaucoup des mêmes propriétés formelles.

\medskip\noindent
Une intersection globale relative complète (que nous avons utilisée pour définir les homomorphismes syntomiques standard d'anneaux) est particulièrement plate. Dans Compléments sur les morphismes, section \ref{more-morphisms-section-lci}, nous considérerons les morphismes $X \to S$ qui sont localement de la forme $$
\Spec(R[x_1, \ldots, x_n]/(f_1, \ldots, f_c)) \to \Spec(R).
$$ pour certaines séquences régulières de Koszul $f_1, \ldots, f_r$ dans $R[x_1, \ldots, x_n]$. Un tel morphisme sera appelé un {\it local complet morphisme d'intersection}. Une fois que nous aurons cette définition en place, il sera le cas qu'un morphisme est syntomique si et seulement s'il s'agit d'un morphisme plat, local complet morphisme d'intersection.

\medskip\noindent
Notez qu'il n'y a pas d'hypothèses de séparation ou de quasi-compactitude dans la définition d'un morphisme syntomique. D'où la question d'être syntomique est de nature locale sur la source. Voici le résultat précis.

\begin{lemma}

\label{lemma-syntomic-characterize}

Laisser $f : X \to S$ morphisme des régimes. Les suivants sont équivalents

\begin{enumerate}

\item Le morphisme $f$ C'est syntomique.
\item Pour tous les ouverts $U \subset X$, $V \subset S$
avec $f(U) \subset V$ l'homomorphisme d'anneaux
$\mathcal{O}_S(V) \to \mathcal{O}_X(U)$ C'est syntomique.
\item Il existe une ouverture de recouvrement $S = \bigcup_{j \in J} V_j$
et couvertures ouvertes $f^{-1}(V_j) = \bigcup_{i \in I_j} U_i$ de sorte que chacun des morphismes $U_i \to V_j$, $j\in J, i\in I_j$
C'est syntomique.
\item Il existe un recouvrement ouvert affine $S = \bigcup_{j \in J} V_j$
et revêtements pour affines d'ouverture $f^{-1}(V_j) = \bigcup_{i \in I_j} U_i$ tel que l'homomorphisme d'anneaux $\mathcal{O}_S(V_j) \to \mathcal{O}_X(U_i)$ est syntomique, pour tous $j\in J, i\in I_j$.

\end{enumerate}

En outre, si $f$ est syntomique alors pour tout sous-système ouvert $U \subset X$, $V \subset S$ avec $f(U) \subset V$
la restriction $f|_U : U \to V$ C'est syntomique.

\end{lemma}

\begin{proof}

Cela s'explique par le fait qu'il s'agit d'une \ref{lemma-locally-P} si nous montrons que la propriété ``$R \to A$ Nous vérifions les conditions (a), (b) et (c) de définition
\ref{definition-property-local}. Par Algèbre, lemme \ref{algebra-lemma-base-change-syntomic}
être syntomique est stable sous changement de base et donc nous concluons (a) tient. Par Algèbre, lemme \ref{algebra-lemma-composition-syntomic}
être syntomique est stable sous la composition et trivialement pour n'importe quel anneau
$R$ l'homomorphisme d'anneaux $R \to R_f$ est syntomique. Nous concluons b) détient. Enfin, la propriété (c) est vraie selon Algèbre, lemme \ref{algebra-lemma-local-syntomic}.

\end{proof}

\begin{lemma}

\label{lemma-composition-syntomic}

La composition de deux morphismes syntomiques est syntomique.

\end{lemma}

\begin{proof}
Dans la démonstration de lemme \ref{lemma-syntomic-characterize} nous avons vu qu'être syntomique est une propriété locale d'homomorphismes d'anneaux. D'où la première affirmation du lemme suit de lemme \ref{lemma-composition-type-P} combinée au fait qu'être syntomique est une propriété d'homomorphismes d'anneaux qui est stable sous composition, voir Algèbre, lemme \ref{algebra-lemma-composition-syntomic}.
\end{proof}

\begin{lemma}

\label{lemma-base-change-syntomic}

Le changement de base d'un morphisme syntomique est syntomique.

\end{lemma}

\begin{proof}
Dans la démonstration de lemme \ref{lemma-syntomic-characterize} nous avons vu qu'être syntomique est une propriété locale d'homomorphismes d'anneaux. C'est pourquoi le lemme suit de lemme \ref{lemma-composition-type-P} combiné au fait qu'être syntomique est une propriété d'homomorphismes d'anneaux qui est stable sous changement de base, voir Algèbre, lemme \ref{algebra-lemma-base-change-syntomic}.
\end{proof}

\begin{lemma}

\label{lemma-open-immersion-syntomic}
Toute immersion ouverte est syntomique.
\end{lemma}

\begin{proof}

C'est vrai parce qu'une immersion ouverte est un isomorphisme local.

\end{proof}

\begin{lemma}

\label{lemma-syntomic-locally-finite-presentation}

Un morphisme syntomique est localement de présentation finie.

\end{lemma}

\begin{proof}

C'est vrai parce qu'un homomorphisme syntomique d'anneaux est de présentation finie par définition.

\end{proof}

\begin{lemma}

\label{lemma-syntomic-flat}
Un morphisme syntomique est plat.
\end{lemma}

\begin{proof}

C'est vrai parce qu'un homomorphisme syntomique d'anneaux est plat par définition.

\end{proof}

\begin{lemma}

\label{lemma-syntomic-open}
Un morphisme syntomique est universel ouvert.
\end{lemma}

\begin{proof}
Combiner les lemmes \ref{lemma-syntomic-locally-finite-presentation}, \ref{lemma-syntomic-flat} et \ref{lemma-fppf-open}.
\end{proof}

\noindent
Laissez $k$ être un champ. Laissez $A$ être un $k$-algebra local essentiellement de type fini sur $k$. Rappelez-vous que $A$ est appelé une {\it intersection complète sur $k$} si nous pouvons écrire $A \cong R/(f_1, \ldots, f_c)$ où $R$ est un anneau local régulier essentiellement de type fini sur $k$, et $f_1, \ldots, f_c$ est une séquence régulière dans $R$, voir Algobre, définition \ref{algebra-definition-lci-local-ring}.

\begin{lemma}

\label{lemma-local-complete-intersection}

Laisser $k$ être un champ. $X$ être un système localement de type fini sur $k$. Les conditions suivantes sont équivalentes:

\begin{enumerate}

\item  $X$ est une intersection locale complète sur $k$,
\item pour chaque $x \in X$ il existe un ouvert affine
$U = \Spec(R) \subset X$ quartier de $x$
de telle manière que $R \cong k[x_1, \ldots, x_n]/(f_1, \ldots, f_c)$
est une intersection globale complète sur $k$,
\item pour chaque $x \in X$ l'anneau local $\mathcal{O}_{X, x}$
est une intersection complète sur $k$.

\end{enumerate}

\end{lemma}

\begin{proof}
Les résultats correspondants de l'algèbre se trouvent dans Algèbre, lemmes \ref{algebra-lemma-lci-at-prime} et \ref{algebra-lemma-lci-global}.
\end{proof}

\noindent
Le lemme suivant dit localement que tout morphisme syntomique est syntomique standard. Par conséquent, nous pouvons utiliser les morphismes syntomiques standard comme un modèle local {\it} pour un morphisme syntomique. De plus, il dit qu'un plat morphisme de présentation finie est syntomique si et seulement si les fibres sont des schémas d'intersection complète locale.

\begin{lemma}

\label{lemma-syntomic-locally-standard-syntomic}

Laisser $f : X  \to S$ être un morphisme de schémas. $x \in X$ être un point avec image $s = f(x)$. Laisse $V \subset S$ être un quartier affiné ouvert de $s$. Les conditions suivantes sont équivalentes

\begin{enumerate}

\item Le morphisme $f$ est syntomique à $x$.
\item Il existe un ouvert affine $U \subset X$ avec $x \in U$ et
$f(U) \subset V$ de telle manière que $f|_U : U \to V$ est syntomique standard.
\item Le morphisme $f$ est de présentation finie à $x$, l'application locale de l'Anneau
$\mathcal{O}_{S, s} \to \mathcal{O}_{X, x}$
est plat et $\mathcal{O}_{X, x}/\mathfrak m_s \mathcal{O}_{X, x}$
est une intersection complète sur $\kappa(s)$ (voir Algèbre, définition) \ref{algebra-definition-lci-local-ring}).

\end{enumerate}

\end{lemma}

\begin{proof}
Suit des définitions et Algèbre, lemme \ref{algebra-lemma-syntomic}.
\end{proof}

\begin{lemma}

\label{lemma-syntomic-flat-fibres}
Laissez $f : X \to S$ être un morphisme de schémas. Si $f$ est plat, localment de présentation finie, et toutes les fibres $X_s$ sont des intersections complètes locales, alors $f$ est syntomique.
\end{lemma}

\begin{proof}
Supprimer des lemmes \ref{lemma-local-complete-intersection} et \ref{lemma-syntomic-locally-standard-syntomic} et les isomorphismes des anneaux locaux $
\mathcal{O}_{X, x}/\mathfrak m_s \mathcal{O}_{X, x}
\cong
\mathcal{O}_{X_s, x}
$.
\end{proof}

\begin{lemma}

\label{lemma-set-points-where-fibres-lci}

Laisser $f : X \to S$ être un morphisme des schémas. $f$ localement de type fini. Formation de l'ensemble
$$
T = \{x \in X \mid \mathcal{O}_{X_{f(x)}, x}
\text{ is a complete intersection over }\kappa(f(x))\}
$$
commutations avec changement arbitraire de base: Pour tout morphisme $g : S' \to S$, considérer le changement de base $f' : X' \to S'$ de $f$ et la projection $g' : X' \to X$. Ensuite, l'ensemble correspondant $T'$ pour le morphisme $f'$ est égal à $T' = (g')^{-1}(T)$. En particulier, si $f$ est supposé plat, et localement de présentation finie puis le même tient pour l'ensemble ouvert de points où $f$ C'est syntomique.

\end{lemma}

\begin{proof}
Laissez $s' \in S'$ être un point, et laissez $s = g(s')$. Ensuite, nous avons $$
X'_{s'} =
\Spec(\kappa(s')) \times_{\Spec(\kappa(s))} X_s
$$ En d'autres termes, les fibres du changement de base sont les changements de base des fibres. C'est pourquoi la première partie est équivalente à Algèbre, lemme \ref{algebra-lemma-lci-field-change-local}. La deuxième partie suit de la première parce que dans ce cas $T$ est l'ensemble de points où $f$ est syntomique selon lemme \ref{lemma-syntomic-locally-standard-syntomic}.
\end{proof}

\begin{lemma}

\label{lemma-standard-syntomic-relative-dimension}
Laissez $R$ être une bague. Laissez $R \to A = R[x_1, \ldots, x_n]/(f_1, \ldots, f_c)$ être une intersection globale complète relative. Posons $S = \Spec(R)$ et $X = \Spec(A)$. Considérez le morphisme $f : X \to S$ associé à l'homomorphisme d'anneaux $R \to A$. La fonction $x \mapsto \dim_x(X_{f(x)})$ est constante avec la valeur $n - c$.
\end{lemma}

\begin{proof}
Par Algèbre, définition \ref{algebra-definition-relative-global-complete-intersection} $R \to A$ étant une intersection globale complète relative signifie que tous les anneaux de fibres non nulles ont une dimension $n - c$. Ainsi, pour un $\mathfrak p$ de $R$, l'anneau de fibre $\kappa(\mathfrak p)[x_1, \ldots, x_n]/(\overline{f}_1, \ldots, \overline{f}_c)$ est soit zéro, soit un anneau d'intersection global complet de dimension $n - c$. Par la discussion qui suit Algèbre, définition \ref{algebra-definition-lci-field} cela implique qu'il est équidimensionnel de dimension $n - c$. D'où la lemme.
\end{proof}
X

\begin{lemma}

\label{lemma-syntomic-relative-dimension}
Laissez $f : X \to S$ être un morphisme syntomique. La fonction $x \mapsto \dim_x(X_{f(x)})$ est localement constante sur $X$.
\end{lemma}

\begin{proof}
Par lemme \ref{lemma-syntomic-locally-standard-syntomic} le morphisme $f$ ressemble localement à un morphisme syntomique standard d'affines. D'où le résultat suivant de lemme \ref{lemma-standard-syntomic-relative-dimension}.
\end{proof}

\noindent
lemme \ref{lemma-syntomic-relative-dimension} dit que la définition suivante a un sens.

\begin{definition}

\label{definition-syntomic-relative-dimension}
Laissez $d \geq 0$ être un entier. Nous disons un morphisme de schémas $f : X \to S$ est {\it syntomic de dimension relative $d$} si $f$ est syntomic et la fonction $\dim_x(X_{f(x)}) = d$ pour tous $x \in X$.
\end{definition}

\noindent
En d'autres termes, $f$ est syntomique et les fibres non vides sont équidimensionnelles de dimension $d$.

\begin{lemma}

\label{lemma-syntomic-permanence}

Laisser
$$
\xymatrix{
X \ar[rr]_f \ar[rd]_p & &
Y \ar[dl]^q \\
& S
}
$$
être un diagramme commutatif des morphismes des schémas.

\begin{enumerate}

\item  $f$ est surjectif et syntomique,
\item  $p$ est syntomique, et
\item  $q$ est localement de présentation finie\footnote{En fait, si $f$ est surjectif, plat et localement de présentation finie et $p$ est syntomique, puis les deux $q$ et $f$ sont syntomiques, voir Descente, lemme \ref{descent-lemma-syntomic-permanence}.}.

\end{enumerate}

Alors $q$ C'est syntomique.

\end{lemma}

\begin{proof}

Par moi \ref{lemma-flat-permanence} nous voyons que $q$ est plat. Il suffit donc de montrer que les fibres $Y \to S$ sont des intersections locales complètes, voir lemme \ref{lemma-syntomic-flat-fibres}. Laisse $s \in S$. Considérez le morphisme $X_s \to Y_s$. C'est un changement de base du morphisme $X \to Y$ et donc surjectif, et syntomique (lemme \ref{lemma-base-change-syntomic}Pour la même raison $X_s$ est syntomique sur $\kappa(s)$. En outre, $Y_s$ est localement de type fini sur $\kappa(s)$
(lemme) \ref{lemma-base-change-finite-type}). De cette façon, nous réduisons au cas où $S$ est le spectre d'un champ.

\medskip\noindent

Présumer $S = \Spec(k)$. Laisse $y \in Y$. Choisissez un ouvert affine $\Spec(A) \subset Y$ quartier de $y$. Laisse
$\Spec(B) \subset X$ être un ouvert affine telle que
$f(\Spec(B)) \subset \Spec(A)$, contenant un point $x \in X$ de telle manière que $f(x) = y$. Choisissez une surjection
$k[x_1, \ldots, x_n] \to A$ avec noyau $I$. Choisissez une surjection $A[y_1, \ldots, y_m] \to B$, qui provoque à son tour une surjection $k[x_i, y_j] \to B$ avec noyau $J$. Laisse $\mathfrak q \subset k[x_i, y_j]$ être le premier correspondant à $y \in \Spec(B)$ et laisser $\mathfrak p \subset k[x_i]$ le premier correspondant à: $x \in \Spec(A)$Depuis $x$ cartes à $y$ Nous avons $\mathfrak p = \mathfrak q \cap k[x_i]$. Considérez le diagramme commutatif suivant des anneaux locaux :
$$
\xymatrix{
\mathcal{O}_{X, x} \ar@{=}[r] &
B_{\mathfrak q} &
k[x_1, \ldots, x_n, y_1, \ldots, y_m]_{\mathfrak q} \ar[l] \\
\mathcal{O}_{Y, y} \ar@{=}[r] \ar[u] & A_{\mathfrak p} \ar[u] &
k[x_1, \ldots, x_n]_{\mathfrak p} \ar[l] \ar[u]
}
$$
Nous prétendons que les hypothèses d'Algèbre, lemme \ref{algebra-lemma-lci-permanence-initial} sont satisfaites. Les conditions (1) et (2) sont triviales. La condition (4) suit comme suit :
$X \to Y$ est plat. L'état (3) suit comme les anneaux
$\mathcal{O}_{Y, y}$ et
$\mathcal{O}_{X_y, x} = \mathcal{O}_{X, x}/\mathfrak m_y\mathcal{O}_{X, x}$
sont des anneaux d'intersection complets par nos hypothèses que
$f$ et $p$ sont syntomiques, voir \ref{lemma-syntomic-locally-standard-syntomic}. La sortie d'Algèbre, lemme \ref{algebra-lemma-lci-permanence-initial}
C'est exactement ça $\mathcal{O}_{Y, y}$ c'est un anneau d'intersection complet ! \ref{lemma-syntomic-locally-standard-syntomic}
Encore une fois nous voyons que $Y$ est syntomique sur $k$ au $y$ comme souhaité.

\end{proof}









\section{Faisceau conormal d'une immersion}
\label{section-conormal-sheaf}

\noindent

Laisser $i : Z \to X$ être une immersion fermée.
$\mathcal{I} \subset \mathcal{O}_X$ être le faisceau quasi-cohérent correspondant des idéaux. Considérez la suite exacte court
$$
0 \to \mathcal{I}^2 \to \mathcal{I} \to \mathcal{I}/\mathcal{I}^2 \to 0
$$
de faisceaux quasi-cohérents sur $X$. Depuis la gerbe $\mathcal{I}/\mathcal{I}^2$
est anéanti par $\mathcal{I}$ il correspond à une gerbe sur $Z$ par lemme \ref{lemma-i-star-equivalence}. Cette quasi-cohérente
$\mathcal{O}_Z$-module est appelé {\it fait conormal de $Z$ en $X$} et est souvent simplement dénotée
$\mathcal{I}/\mathcal{I}^2$ par l ' abus de notation mentionné dans la section \ref{section-closed-immersions-quasi-coherent}.

\medskip\noindent
Dans le cas où $i : Z \to X$ est une immersion (localement fermée) nous définissons le faisceau conormal de $i$ comme le faisceau conormal de l'immersion fermée $i : Z \to X \setminus \partial Z$, où $\partial Z = \overline{Z} \setminus Z$. Il est souvent dénoté $\mathcal{I}/\mathcal{I}^2$ où $\mathcal{I}$ est le faisceau d'idéaux de l'immersion fermée $i : Z \to X \setminus \partial Z$.

\begin{definition}

\label{definition-conormal-sheaf}
Laissez $i : Z \to X$ être une immersion. Le {\it faisceau conormal $\mathcal{C}_{Z/X}$ de $Z$ dans $X$} ou le {\it faisceau conormal de $i$} est le quasi-cohérent $\mathcal{O}_Z$-module $\mathcal{I}/\mathcal{I}^2$ décrit ci-dessus.
\end{definition}

\noindent
Dans \cite[IV définition 16.1.2]{EGA} cette feuille est dénotée $\mathcal{N}_{Z/X}$. Nous ne suivons pas cette convention puisque nous aimerions réserver la notation $\mathcal{N}_{Z/X}$ pour la feuille normale {\it de l'immersion}. Elle est définie comme $$
\mathcal{N}_{Z/X} =
\SheafHom_{\mathcal{O}_Z}(\mathcal{C}_{Z/X}, \mathcal{O}_Z) =
\SheafHom_{\mathcal{O}_Z}(\mathcal{I}/\mathcal{I}^2, \mathcal{O}_Z)
$$ à condition que le faisceau conormal soit de présentation finie (autrement, la feuille normale ne peut même pas être quasi-cohérente). Nous reviendrons plus tard à la feuille normale (insérer ici une référence future).

\begin{lemma}

\label{lemma-affine-conormal}
Laissez $i : Z \to X$ être une immersion. Le faisceau conormal de $i$ a les propriétés suivantes:
\begin{enumerate}
\item Laissez $U \subset X$ être un sous-système ouvert tel que $i$ facteurs comme $Z \xrightarrow{i'} U \to X$ où $i'$ est une immersion fermée. Laissez $\mathcal{I} = \Ker((i')^\sharp) \subset \mathcal{O}_U$. Puis $$
\mathcal{C}_{Z/X} = (i')^*\mathcal{I}\quad\text{and}\quad
i'_*\mathcal{C}_{Z/X} = \mathcal{I}/\mathcal{I}^2
$$
\item Pour toute ouverture affine $\Spec(R) = U \subset X$ tel que $Z \cap U = \Spec(R/I)$ il y a un isomorphisme canonique $\Gamma(Z \cap U, \mathcal{C}_{Z/X}) = I/I^2$.
\end{enumerate}

\end{lemma}

\begin{proof}

La plupart du temps clair à partir des définitions. Notez qu'en raison d'une bague $R$ et un idéal $I$ de $R$ Nous avons $I/I^2 = I \otimes_R R/I$. Détails omis.

\end{proof}

\begin{lemma}

\label{lemma-conormal-functorial}
Laissez $$
\xymatrix{
Z \ar[r]_i \ar[d]_f & X \ar[d]^g \\
Z' \ar[r]^{i'} & X'
}
$$ être un diagramme commutatif dans la catégorie des schémas. Supposons que $i$, $i'$ immersions. Il y a une application canonique de $\mathcal{O}_Z$-modules $$
f^*\mathcal{C}_{Z'/X'}
\longrightarrow
\mathcal{C}_{Z/X}
$$ caractérisée par la propriété suivante: Pour chaque paire d'ouvertures affines $(\Spec(R) = U \subset X, \Spec(R') = U' \subset X')$ avec $g(U) \subset U'$ telle que $Z \cap U = \Spec(R/I)$ et $Z' \cap U' = \Spec(R'/I')$ l'application induite $$
\Gamma(Z' \cap U', \mathcal{C}_{Z'/X'}) = I'/I'^2
\longrightarrow
I/I^2 = \Gamma(Z \cap U, \mathcal{C}_{Z/X})
$$ est celle induite par l'homomorphisme d'anneaux $g^\sharp : R' \to R$ qui possède la propriété $g^\sharp(I') \subset I$.
\end{lemma}

\begin{proof}
Laissez $\partial Z' = \overline{Z'} \setminus Z'$ et $\partial Z = \overline{Z} \setminus Z$. Ce sont des sous-ensembles fermés de $X'$ et de $X$. Remplaçant $X'$ par $X' \setminus \partial Z'$ et $X$ par $X \setminus \big(g^{-1}(\partial Z') \cup \partial Z\big)$ nous voyons que nous pouvons supposer que $i$ et $i'$ sont des immersions fermées.

\medskip\noindent
Le fait que les facteurs $g \circ i$ par $i'$ impliquent que $g^*\mathcal{I}'$ cartographie dans $\mathcal{I}$ sous la application canonique $g^*\mathcal{I}' \to \mathcal{O}_X$, voir Schémas, lemmes \ref{schemes-lemma-characterize-closed-subspace} et \ref{schemes-lemma-restrict-map-to-closed}. Par conséquent, nous obtenons une application induite de faisceaux quasi-cohérents $g^*(\mathcal{I}'/(\mathcal{I}')^2) \to \mathcal{I}/\mathcal{I}^2$. Retirer par $i$ donne $i^*g^*(\mathcal{I}'/(\mathcal{I}')^2) \to i^*(\mathcal{I}/\mathcal{I}^2)$. Notez que $i^*(\mathcal{I}/\mathcal{I}^2) = \mathcal{C}_{Z/X}$. D'autre part, $i^*g^*(\mathcal{I}'/(\mathcal{I}')^2) =
f^*(i')^*(\mathcal{I}'/(\mathcal{I}')^2) = f^*\mathcal{C}_{Z'/X'}$. Ceci donne l'application désirée.

\medskip\noindent

Vérifier que l'application est décrite localement comme l'application donnée
$I'/(I')^2 \to I/I^2$ est une question de démantèlement des définitions et est omis.
$x \in i(Z)$ il existe des quartiers affines ouverts $U$, $U'$
avec $f(U) \subset U'$ et $Z \cap U$ ainsi que $U' \cap Z'$
fermés de manière à ce que $x \in U$. démonstration omise. C'est pourquoi l'exigence du lemme caractérise effectivement l'application (et aurait pu être utilisée pour la définir).

\end{proof}

\begin{lemma}

\label{lemma-conormal-functorial-flat}
Laissez $$
\xymatrix{
Z \ar[r]_i \ar[d]_f & X \ar[d]^g \\
Z' \ar[r]^{i'} & X'
}
$$ être un diagramme fibré produit dans la catégorie des schémas avec $i$, $i'$ immersions. Puis l'application canonique $f^*\mathcal{C}_{Z'/X'} \to \mathcal{C}_{Z/X}$ de lemme \ref{lemma-conormal-functorial} est surjective. Si $g$ est plat, alors c'est un isomorphisme.
\end{lemma}

\begin{proof}
Laissez $R' \to R$ être un homomorphisme d'anneaux, et $I' \subset R'$ un idéal. Posons $I = I'R$. Alors $I'/(I')^2 \otimes_{R'} R \to I/I^2$ est surjectif. Si $R' \to R$ est plat, alors $I = I' \otimes_{R'} R$ et $I^2 = (I')^2 \otimes_{R'} R$ et nous voyons que l'application est un isomorphisme.
\end{proof}

\begin{lemma}

\label{lemma-transitivity-conormal}

Laisser $Z \to Y \to X$ puis il y a une suite exacte canonique
$$
i^*\mathcal{C}_{Y/X} \to
\mathcal{C}_{Z/X} \to
\mathcal{C}_{Z/Y} \to 0
$$
d'où viennent les cartes de Lemme \ref{lemma-conormal-functorial}
et $i : Z \to Y$ est le premier morphisme.

\end{lemma}

\begin{proof}
Via lemme \ref{lemma-conormal-functorial} ceci se traduit par l'algèbre suivante. Supposons que $C \to B \to A$ sont des homomorphismes surjectifs d'anneaux. Laissez $I = \Ker(B \to A)$, $J = \Ker(C \to A)$ et $K = \Ker(C \to B)$. Puis il y a une suite exacte $$
K/K^2 \otimes_B A \to J/J^2 \to I/I^2 \to 0.
$$ Ceci suit immédiatement de l'observation que $I = J/K$.
\end{proof}










\section{Faisceau des différentielles d'un morphisme}
\label{section-sheaf-differentials}

\noindent
Nous suggérons au lecteur d'examiner la section correspondante du chapitre sur l'algèbre commutative (Algèbre, section \ref{algebra-section-differentials}) et la section correspondante du chapitre sur les faisceaux de modules (Modules, section \ref{modules-section-differentials}).

\begin{definition}

\label{definition-sheaf-differentials}
Laissez $f : X \to S$ être un morphisme de schémas. Le {\it faisceau des différences $\Omega_{X/S}$ de $X$ sur $S$} est le faisceau des différences de $f$ vu comme un morphisme d'espaces annelés (modules, définition \ref{modules-definition-differentials}) équipés de sa {\it universal $S$-derivation} $$
\text{d}_{X/S} : \mathcal{O}_X \longrightarrow \Omega_{X/S}.
$$
\end{definition}

\noindent
Il s'avère que $\Omega_{X/S}$ est un module quasi-cohérent $\mathcal{O}_X$ par exemple car il est isomorphique au faisceau conormal de la diagonale morphisme $\Delta : X \to X \times_S X$ (lemme \ref{lemma-differentials-diagonal}). Nous avons défini le module des différentiels de $X$ sur $S$ en utilisant une propriété universelle, à savoir comme le récipient de la dérivation universelle. Si vous avez une autre construction de la tige de différentiels relatifs qui satisfait cette propriété universelle, alors, par la lemme de Yoneda, il sera canoniquement isomorphique à celle définie ci-dessus. Pour plus de commodité, nous réaffirmons la propriété universelle ici.

\begin{lemma}

\label{lemma-universal-derivation-universal}
Laissez $f : X \to S$ être un morphisme de schémas. La carte $$
\Hom_{\mathcal{O}_X}(\Omega_{X/S}, \mathcal{F})
\longrightarrow
\text{Der}_S(\mathcal{O}_X, \mathcal{F}),\quad
\alpha \longmapsto \alpha \circ \text{d}_{X/S}
$$ est un isomorphisme de functeurs $\textit{Mod}(\mathcal{O}_X) \to \textit{Sets}$.
\end{lemma}

\begin{proof}

Ce n'est qu'une reformulation de la définition.

\end{proof}

\begin{lemma}

\label{lemma-differentials-restrict-open}
Laissez $f : X \to S$ être un morphisme de schémas. Laissez $U \subset X$, $V \subset S$ être des sous-systèmes ouverts tels que $f(U) \subset V$. Ensuite, il ya un isomorphisme unique $\Omega_{X/S}|_U = \Omega_{U/V}$ de $\mathcal{O}_U$-modules tels que $\text{d}_{X/S}|_U = \text{d}_{U/V}$.
\end{lemma}

\begin{proof}
Il s'agit d'un cas particulier de modules, lemme \ref{modules-lemma-localize-differentials} si nous utilisons l'identification canonique $f^{-1}\mathcal{O}_S|_U = (f|_U)^{-1}\mathcal{O}_V$.
\end{proof}

\noindent
À partir de maintenant, nous allons utiliser ces identifications canoniques et simplement écrire $\Omega_{U/S}$ ou $\Omega_{U/V}$ pour la restriction de $\Omega_{X/S}$ à $U$.

\begin{lemma}

\label{lemma-affine-case-derivation}
Laissez $R \to A$ être un homomorphisme d'anneaux. Laissez $\mathcal{F}$ être une bande de $\mathcal{O}_X$-modules sur $X = \Spec(A)$. Posons $S = \Spec(R)$. La règle qui s'associe à une $S$-derivation sur $\mathcal{F}$ son action sur les sections globales définit un bijection entre l'ensemble de $S$-derivations de $\mathcal{F}$ et l'ensemble de $R$-derivations sur $M = \Gamma(X, \mathcal{F})$.
\end{lemma}

\begin{proof}

Laisser $D : A \to M$ être un $R$Nous devons montrer qu'il existe un $S$-dérivation sur $\mathcal{F}$ qui donne lieu à
$D$ sur les sections mondiales. $U = D(f) \subset X$ être un ouvert affine standard. $\Gamma(U, \mathcal{O}_X)$ est de la forme
$a/f^n$ pour certains $a \in A$ et $n \geq 0$. Par la règle de Leibniz nous avons
$$
D(a)|_U = a/f^n D(f^n)|_U + f^n D(a/f^n)
$$
en $\Gamma(U, \mathcal{F})$Depuis $f$ agit de manière inversée sur $\Gamma(U, \mathcal{F})$ cela détermine complètement la valeur de $D(a/f^n) \in \Gamma(U, \mathcal{F})$. Cela prouve l'unicité. L'existence suit simplement en définissant
$$
D(a/f^n) := (1/f^n) D(a)|_U - a/f^{2n} D(f^n)|_U
$$
et prouver que cela a toutes les propriétés désirées (sur la base de la norme ouvre de $X$). Détails omis.

\end{proof}

\begin{lemma}

\label{lemma-differentials-affine}
Laissez $f : X \to S$ être un morphisme de schémas. Pour n'importe quelle paire d'ouvertures affines $\Spec(A) = U \subset X$, $\Spec(R) = V \subset S$ avec $f(U) \subset V$, il y a un isomorphisme unique $$
\Gamma(U, \Omega_{X/S}) = \Omega_{A/R}.
$$ compatible avec $\text{d}_{X/S}$ et $\text{d} : A \to \Omega_{A/R}$.
\end{lemma}

\begin{proof}

Par moi \ref{lemma-differentials-restrict-open} Nous pourrions remplacer
$X$ et $S$ par $U$ et $V$. Ainsi, nous pouvons supposer $X = \Spec(A)$
et $S = \Spec(R)$ et nous devons montrer le lemme avec $U = X$
et $V = S$. Considérez la $A$-module $M = \Gamma(X, \Omega_{X/S})$
avec le $R$-dérivé $\text{d}_{X/S} : A \to M$. Laisse $N$ être un autre $A$-module et indiquer $\widetilde{N}$ la quasi-cohérente
$\mathcal{O}_X$-module associé à $N$, voir Schémas, section \ref{schemes-section-quasi-coherent-affine}. Précomposition par $\text{d}_{X/S} : A \to M$ Nous avons une flèche
$$
\alpha : \Hom_A(M, N) \longrightarrow \text{Der}_R(A, N)
$$
Utilisation de lemmes \ref{lemma-universal-derivation-universal} et
\ref{lemma-affine-case-derivation} nous obtenons des identifications
$$
\Hom_{\mathcal{O}_X}(\Omega_{X/S}, \widetilde{N}) =
\text{Der}_S(\mathcal{O}_X, \widetilde{N}) =
\text{Der}_R(A, N)
$$
Prise de sections globales détermine une flèche
$\Hom_{\mathcal{O}_X}(\Omega_{X/S}, \widetilde{N}) \to \Hom_R(M, N)$. En combinant cette flèche et les identifications ci-dessus nous obtenons une flèche
$$
\beta : \text{Der}_R(A, N) \longrightarrow \Hom_R(M, N)
$$
En vérifiant ce qui se passe sur les sections mondiales, nous trouvons que $\alpha$
et $\beta$ sont inversés les uns les autres.
$\text{d}_{X/S} : A \to M$ satisfait à la même propriété universelle que $\text{d} : A \to \Omega_{A/R}$, voir Algèbre, lemme \ref{algebra-lemma-universal-omega}. Ainsi le Yoneda lemme (Catégories, lemme \ref{categories-lemma-yoneda}) implique un isomorphisme unique de $A$-modules
$M \cong \Omega_{A/R}$ compatible avec les dérivations.

\end{proof}

\begin{remark}

\label{remark-differentials-glue}
Le lemme ci-dessus donne une deuxième façon de construire le module des différentiels. A savoir, laissez $f : X \to S$ être un morphisme de schémas. Considérez la collection de tous les ouverts affines $U \subset X$ qui mapent dans un ouvert affine de $S$. Ces derniers forment une base pour la topologie sur $X$. Ainsi, il suffit de définir $\Gamma(U, \Omega_{X/S})$ pour $U$. Nous avons simplement défini $\Gamma(U, \Omega_{X/S}) = \Omega_{A/R}$ si $A$, $R$ sont comme dans le lemme \ref{lemma-differentials-affine} ci-dessus. Cela fonctionne, mais il faut un peu plus de préliminaires algébriques pour construire les cartes de restriction et pour vérifier l'état de chanvre avec cette ansatz.
\end{remark}

\noindent
Le lemme suivant donne encore une autre façon de définir la gerbe des différentiels et il montre en particulier que $\Omega_{X/S}$ est quasi-cohérent si $X$ et $S$ sont des schémas.

\begin{lemma}

\label{lemma-differentials-diagonal}
Laissez $f : X \to S$ être un morphisme de schémas. Il y a un isomorphisme canonique entre $\Omega_{X/S}$ et le faisceau conormal de la diagonale morphisme $\Delta_{X/S} : X \longrightarrow X \times_S X$.
\end{lemma}

\begin{proof}
Nous établissons d'abord l'existence de deux gerbes « globaux » et de cartes globales de gerbes, et plus bas nous décrivons les constructions sur certains ouverts affines.

\medskip\noindent

Rappelez-vous que $\Delta = \Delta_{X/S} : X \to X \times_S X$
est une immersion, voir Schémas, lemme \ref{schemes-lemma-diagonal-immersion}. Laisse $\mathcal{J}$ être le faisceau d'idéaux de l'immersion qui vit sur un sous-système ouvert $W$ de $X \times_S X$
de telle manière que $\Delta(X) \subset W$ est fermé. Prenons celui qui a été trouvé dans la démonstration de Schémas, lemme \ref{schemes-lemma-diagonal-immersion}. Notez que le faisceau d'anneaux $\mathcal{O}_W/\mathcal{J}^2$
est soutenu sur $\Delta(X)$. De plus, il est situé dans une suite exacte courtois de gerbes
$$
0 \to \mathcal{J}/\mathcal{J}^2
\to \mathcal{O}_W/\mathcal{J}^2
\to \Delta_*\mathcal{O}_X
\to 0.
$$
Utilisation $\Delta^{-1}$ Nous pouvons penser à cela comme une surjection de gerbes de $f^{-1}\mathcal{O}_S$-algèbres à noyau le faisceau conormal $\Delta$ (voir définition) \ref{definition-conormal-sheaf}
et lemme \ref{lemma-affine-conormal}).
$$
0 \to \mathcal{C}_{X/X \times_S X}
\to \Delta^{-1}(\mathcal{O}_W/\mathcal{J}^2)
\to \mathcal{O}_X
\to 0
$$
Cela nous place dans la situation de Modules, lemme \ref{modules-lemma-double-structure-gives-derivation}. Les morphismes de projection $p_i : X \times_S X \to X$, $i = 1, 2$ induire des cartes des faisceaux d'anneaux
$(p_i)^\sharp : (p_i)^{-1}\mathcal{O}_X \to \mathcal{O}_{X \times_S X}$. Nous pouvons nous limiter à $W$ et quotient par $\mathcal{J}^2$ pour obtenir
$(p_i)^{-1}\mathcal{O}_X \to \mathcal{O}_W/\mathcal{J}^2$Depuis $\Delta^{-1}p_i^{-1}\mathcal{O}_X = \mathcal{O}_X$
nous obtenons des cartes
$$
s_i : \mathcal{O}_X \to \Delta^{-1}(\mathcal{O}_W/\mathcal{J}^2).
$$
Les deux $s_1$ et $s_2$ sont des sections de l'application
$\Delta^{-1}(\mathcal{O}_W/\mathcal{J}^2) \to \mathcal{O}_X$, comme dans les modules, laisser \ref{modules-lemma-double-structure-gives-derivation}. Ainsi nous obtenons un $S$-dérivé
$\text{d} = s_2 - s_1 : \mathcal{O}_X \to \mathcal{C}_{X/X \times_S X}$. Par la propriété universelle du module des différences, nous trouvons un $\mathcal{O}_X$-carte linéaire
$$
\Omega_{X/S} \longrightarrow \mathcal{C}_{X/X \times_S X},\quad
f\text{d}g \longmapsto fs_2(g) - fs_1(g)
$$
Pour voir l'application est un isomorphisme, résolvons ce problème sur des ouverts adaptés. $X$ par ouverts affines $\Spec(A) = U \subset X$
dont l'image est contenue dans un ouvert affine $\Spec(R) = V \subset S$. Selon la démonstration de Schémas, lemme \ref{schemes-lemma-diagonal-immersion}

$U \times_V U \subset X \times_S X$ est un ouvert affine contenue dans le $W$ mentionné ci-dessus.
$U \times_V U = \Spec(A \otimes_R A)$. La gerbe $\mathcal{J}$ correspond à l'idéal
$J = \Ker(A \otimes_R A \to A)$. La suite exacte courtoisie à la suite exacte courtoisie de $A \otimes_R A$-modules
$$
0 \to J/J^2 \to (A \otimes_R A)/J^2 \to A \to 0
$$
Les sections $s_i$ correspondent aux homomorphismes d'anneaux
$$
A \longrightarrow (A \otimes_R A)/J^2,\quad
s_1 : a \mapsto a \otimes 1,\quad
s_2 : a \mapsto 1 \otimes a.
$$
Par moi \ref{lemma-affine-conormal} Nous avons
$\Gamma(U, \mathcal{C}_{X/X \times_S X}) = J/J^2$ et par la suite \ref{lemma-differentials-affine}
Nous avons $\Gamma(U, \Omega_{X/S}) = \Omega_{A/R}$. La carte ci-dessus est l'application $a \text{d}b \mapsto a \otimes b - ab \otimes 1$
qui est montré comme un isomorphisme dans Algèbre, lemme \ref{algebra-lemma-differentials-diagonal}.

\end{proof}

\begin{lemma}

\label{lemma-functoriality-differentials}
Laissez $$
\xymatrix{
X' \ar[d] \ar[r]_f & X \ar[d] \\
S' \ar[r] & S
}
$$ être un diagramme commutatif de schémas. L'application canonique $\mathcal{O}_X \to f_*\mathcal{O}_{X'}$ composée avec l'application $f_*\text{d}_{X'/S'} : f_*\mathcal{O}_{X'} \to f_*\Omega_{X'/S'}$ est une dérivée $S$. Ainsi, nous obtenons une application canonique de $\mathcal{O}_X$-modules $\Omega_{X/S} \to f_*\Omega_{X'/S'}$, et par adjonction de $f_*$ et $f^*$ un homomorphisme canonique $\mathcal{O}_{X'}$-module $$
c_f : f^*\Omega_{X/S} \longrightarrow \Omega_{X'/S'}.
$$ Il est caractérisé de façon unique par la propriété que $f^*\text{d}_{X/S}(h)$ mappe à $\text{d}_{X'/S'}(f^* h)$ pour n'importe quelle section locale $h$ de $\mathcal{O}_X$.
\end{lemma}

\begin{proof}

Il s'agit d'un cas particulier de modules, lemme \ref{modules-lemma-functoriality-differentials-ringed-spaces}. Dans le cas des régimes, nous pouvons également utiliser la func tionnalité des gerbes conormales (voir lemme \ref{lemma-conormal-functorial}) et lemme \ref{lemma-differentials-diagonal} à définir $c_f$. Ou nous pouvons utiliser la caractérisation dans la dernière ligne du lemme pour coller les cartes définies sur les taches affines (voir Algèbre, Équation (\ref{algebra-equation-functorial-omega})).

\end{proof}

\begin{lemma}

\label{lemma-triangle-differentials}
Laissez $f : X \to Y$, $g : Y \to S$ être des morphismes de schémas. Ensuite, il y a une suite exacte canonique $$
f^*\Omega_{Y/S} \to \Omega_{X/S} \to \Omega_{X/Y} \to 0
$$ où les cartes proviennent des applications de lemme \ref{lemma-functoriality-differentials}.
\end{lemma}

\begin{proof}
C'est la version sheafified d'Algèbre, lemme \ref{algebra-lemma-exact-sequence-differentials}. Alternativement, il y a une version générale pour les morphismes des espaces annelés, voir Modules, lemme \ref{modules-lemma-triangle-differentials}.
\end{proof}

\begin{lemma}

\label{lemma-base-change-differentials}
Laissez $X \to S$ être un morphisme de schémas. Laissez $g : S' \to S$ être un morphisme de schémas. Laissez $X' = X_{S'}$ être le changement de base de $X$. Notez $g' : X' \to X$ la projection. Puis l'application $$
(g')^*\Omega_{X/S} \to \Omega_{X'/S'}
$$ de lemme \ref{lemma-functoriality-differentials} est un isomorphisme.
\end{lemma}

\begin{proof}
Voici la version sheafified d'Algèbre, lemme \ref{algebra-lemma-differentials-base-change}.
\end{proof}

\begin{lemma}

\label{lemma-differential-product}
Laissez $f : X \to S$ et $g : Y \to S$ être des morphismes de schémas avec la même cible. Laissez $p : X \times_S Y \to X$ et $q : X \times_S Y \to Y$ être les morphismes de projection. Les cartes de lemme \ref{lemma-functoriality-differentials} $$
p^*\Omega_{X/S} \oplus q^*\Omega_{Y/S}
\longrightarrow
\Omega_{X \times_S Y/S}
$$ donnent un isomorphisme.
\end{lemma}

\begin{proof}
Par lemme \ref{lemma-base-change-differentials} la composition $p^*\Omega_{X/S} \to \Omega_{X \times_S Y/S} \to \Omega_{X \times_S Y/Y}$ est un isomorphisme, et de même pour $q$. De plus, le cokernel de $p^*\Omega_{X/S} \to \Omega_{X \times_S Y/S}$ est $\Omega_{X \times_S Y/X}$ par lemme \ref{lemma-triangle-differentials}. Le résultat suit.
\end{proof}

\begin{lemma}

\label{lemma-finite-type-differentials}
Laissez $f : X \to S$ être un morphisme de schémas. Si $f$ est localement de type fini, alors $\Omega_{X/S}$ est un type fini $\mathcal{O}_X$-module.
\end{lemma}

\begin{proof}
Immédiate d'Algèbre, lemme \ref{algebra-lemma-differentials-finitely-generated}, lemme \ref{lemma-differentials-affine}, lemme \ref{lemma-locally-finite-type-characterize}, et Propriétés, lemme \ref{properties-lemma-finite-type-module}.
\end{proof}

\begin{lemma}

\label{lemma-finite-presentation-differentials}
Laissez $f : X \to S$ être un morphisme de schémas. Si $f$ est localement de présentation finie, alors $\Omega_{X/S}$ est un $\mathcal{O}_X$-module de présentation finie.
\end{lemma}

\begin{proof}
Immédiate d'Algèbre, lemme \ref{algebra-lemma-differentials-finitely-presented}, lemme \ref{lemma-differentials-affine}, lemme \ref{lemma-locally-finite-presentation-characterize}, et Propriétés, lemme \ref{properties-lemma-finite-presentation-module}.
\end{proof}

\begin{lemma}

\label{lemma-immersion-differentials}

Si $X \to S$ est une immersion, ou plus généralement un monomorphisme, alors
$\Omega_{X/S}$ C'est zéro.

\end{lemma}

\begin{proof}
C'est vrai parce que $\Delta_{X/S}$ est un isomorphisme dans ce cas et a donc un faisceau banal conormal. D'où $\Omega_{X/S} = 0$ par lemme \ref{lemma-differentials-diagonal}. La version algébrique est Algèbre, lemme \ref{algebra-lemma-trivial-differential-surjective}.
\end{proof}

\begin{lemma}

\label{lemma-differentials-relative-immersion}
Laissez $i : Z \to X$ être une immersion de schémas sur $S$. Il ya une suite exacte canonique $$
\mathcal{C}_{Z/X} \to i^*\Omega_{X/S} \to \Omega_{Z/S} \to 0
$$ où la première flèche est induite par $\text{d}_{X/S}$ et la seconde flèche vient de lemme \ref{lemma-functoriality-differentials}.
\end{lemma}

\begin{proof}

Voici la version sheafified d'Algèbre, lemme \ref{algebra-lemma-differential-seq}. Cependant, nous devons nous assurer de pouvoir définir la première flèche à l'échelle mondiale. $\text{d}_{X/S}$'' ici. A savoir, nous pouvons supposer que $i$ est une immersion fermée par rétrécissement $X$. Laisse $\mathcal{I} \subset \mathcal{O}_X$
soit le faisceau d'idées correspondant à $Z \subset X$. Ensuite $\text{d}_{X/S} : \mathcal{I} \to \Omega_{X/S}$
cartographie le sous-chapeau $\mathcal{I}^2 \subset \mathcal{I}$ à
$\mathcal{I}\Omega_{X/S}$. C'est pourquoi il induit une application
$\mathcal{I}/\mathcal{I}^2 \to \Omega_{X/S}/\mathcal{I}\Omega_{X/S}$
qui est $\mathcal{O}_X/\mathcal{I}$- linéaire. Par lemme \ref{lemma-i-star-equivalence} cela correspond à une application
$\mathcal{C}_{Z/X} \to i^*\Omega_{X/S}$ comme souhaité.

\end{proof}

\begin{lemma}

\label{lemma-differentials-relative-immersion-section}
Laissez $i : Z \to X$ être une immersion de schémas sur $S$, et supposez que $i$ (localement) a un inverse de gauche. Ensuite, la séquence canonique $$
0 \to \mathcal{C}_{Z/X} \to i^*\Omega_{X/S} \to \Omega_{Z/S} \to 0
$$ de lemme \ref{lemma-differentials-relative-immersion} est (localement) divisée exacte. En particulier, si $s : S \to X$ est une section du morphisme de structure $X \to S$ alors l'application $\mathcal{C}_{S/X} \to s^*\Omega_{X/S}$ induite par $\text{d}_{X/S}$ est un isomorphisme.
\end{lemma}

\begin{proof}

Il s'agit d'Algèbre, lemme \ref{algebra-lemma-differential-seq-split}
que le résultat tient lorsque $X, Z, S$ En général, nous concluons que la séquence est localement fractionnée exacte. $g : X \to Z$ est un inverse gauche de $i$Alors $i^*c_g$ est un droit inverse de l'application
$i^*\Omega_{X/S} \to \Omega_{Z/S}$ par modules, lemme \ref{modules-lemma-check-functoriality-differentials}. Cela implique que la séquence est divisée par Homologie, lemme \ref{homology-lemma-ses-split}Enfin, si $s$ est une section, alors c'est une immersion $s : Z = S \to X$
terminé $S$ (voir Schémas, lemme \ref{schemes-lemma-section-immersion}) et dans ce cas $\Omega_{Z/S} = 0$.

\end{proof}

\begin{remark}

\label{remark-differentials-diagonal}

Laisser $X \to S$ être un morphisme des schémas. \ref{lemma-differential-product}
Nous avons
$$
\Omega_{X \times_S X/S} =
\text{pr}_1^*\Omega_{X/S} \oplus \text{pr}_2^*\Omega_{X/S}
$$
D'autre part, la diagonale morphisme $\Delta : X \to X \times_S X$
est une immersion, qui a localement un inverse gauche. \ref{lemma-differentials-relative-immersion-section}
nous obtenons une suite canonique exacte courte
$$
0 \to \mathcal{C}_{X/X \times_S X} \to \Omega_{X/S} \oplus \Omega_{X/S}
\to \Omega_{X/S} \to 0
$$
Notez que la flèche droite est $(1, 1)$ Ce qui est en effet une surjection fractionnée. \ref{lemma-differentials-diagonal}
Nous avons une identification $\Omega_{X/S} = \mathcal{C}_{X/X \times_S X}$Parce que nous avons choisi $\text{d}_{X/S}(f) = s_2(f) - s_1(f)$ dans cette identification il s'avère que la flèche de gauche est l'application
$(-1, 1)$\footnote{Nomment, la section locale $\text{d}_{X/S}(f) = 1 \otimes f - f \otimes 1$ du faisceau d'idées de $\Delta$ cartes via $\text{d}_{X \times_S X/X}$ à la section locale
$1 \otimes 1 \otimes 1 \otimes f - 1 \otimes f \otimes 1 \otimes 1
-1 \otimes 1 \otimes f \otimes 1 + f \otimes 1 \otimes 1 \otimes 1 =
\text{pr}_2^*\text{d}_{X/S}(f) - \text{pr}_1^*\text{d}_{X/S}(f)$.}.

\end{remark}

\begin{lemma}

\label{lemma-two-immersions}
Laissez $$
\xymatrix{
Z \ar[r]_i \ar[rd]_j & X \ar[d] \\
& Y
}
$$ être un diagramme commutatif des schémas où $i$ et $j$ sont des immersions. Ensuite, il y a une suite exacte canonique $$
\mathcal{C}_{Z/Y} \to
\mathcal{C}_{Z/X} \to
i^*\Omega_{X/Y} \to 0
$$ où la première flèche vient de lemme \ref{lemma-conormal-functorial} et la seconde de lemme \ref{lemma-differentials-relative-immersion}.
\end{lemma}

\begin{proof}
La version algébrique est Algèbre, lemme \ref{algebra-lemma-application-NL}.
\end{proof}










\section{Opérateurs différentiels d'ordre fini}
\label{section-differential-operators}

\noindent
Nous suggérons au lecteur d'examiner la section correspondante du chapitre sur l'algèbre commutative (Algèbre, section \ref{algebra-section-differential-operators}) et la section correspondante du chapitre sur les faisceaux de modules (Modules, section \ref{modules-section-differential-operators}).

\begin{lemma}

\label{lemma-affine-case-differential-operators}
Laissez $R \to A$ être un homomorphisme d'anneaux. Notez $f : X \to S$ le morphisme correspondant des schémas d'affine. Laissez $\mathcal{F}$ et $\mathcal{G}$ être $\mathcal{O}_X$-modules. Si $\mathcal{F}$ est quasi-cohérent alors l'application $$
\text{Diff}^k_{X/S}(\mathcal{F}, \mathcal{G}) \to
\text{Diff}^k_{A/R}(\Gamma(X, \mathcal{F}), \Gamma(X, \mathcal{G}))
$$ envoyer un opérateur différentiel à son action sur les sections globales est bijective.
\end{lemma}

\begin{proof}

Écrire $\mathcal{F} = \widetilde{M}$ pour certains $A$-module $M$. Prêt
$N = \Gamma(X, \mathcal{G})$. Laisse $D : M \to N$ être un opérateur différentiel d'ordre $k$. Nous devons montrer qu'il existe un opérateur différentiel unique $\mathcal{F} \to \mathcal{G}$ de l ' ordre $k$
qui donne lieu à $D$ sur les sections mondiales. $U = D(f) \subset X$
être un ouvert affine standard. $\mathcal{F}(U) = M_f$ est la localisation. Par Algèbre, lemme \ref{algebra-lemma-invert-system-differential-operators}
l'opérateur différentiel $D$ s'étend à un opérateur différentiel unique
$$
D_f : \mathcal{F}(U) = \widetilde{M}(U) = M_f \to N_f = \widetilde{N}(U)
$$
L'unicité montre que ces cartes $D_f$ colle pour donner une application des gerbes
$\widetilde{M} \to \widetilde{N}$ sur la base de toutes les normes ouvertes $X$. Nous obtenons donc une application unique de gerbes
$\widetilde{D} : \widetilde{M} \to \widetilde{N}$ en accord avec ces cartes par le matériel de Faisceaux, section \ref{sheaves-section-bases}Depuis $\widetilde{D}$ est donné par les opérateurs différents de l'ordre $k$ sur la norme s'ouvre, nous trouvons que $\widetilde{D}$ est un opérateur différentiel d'ordre $k$ (petit détail omis). Enfin, nous pouvons post-composer avec le canonique $\mathcal{O}_X$-module carte
$c : \widetilde{N} \to \mathcal{G}$
(Schémas, lemme) \ref{schemes-lemma-compare-constructions}) pour obtenir $c \circ \widetilde{D} : \mathcal{F} \to \mathcal{G}$
qui est un opérateur différentiel d'ordre $k$ par modules, lemme \ref{modules-lemma-composition-differential-operators}. Cela prouve l'existence. Nous omettons la démonstration de l'unicité.

\end{proof}

\begin{lemma}

\label{lemma-base-change-differential-operators}

Laisser $a : X \to S$ et $b : Y \to S$ les morphismes des schémas. $\mathcal{F}$ et $\mathcal{F}'$ être quasi-cohérent $\mathcal{O}_X$- Les modules. $D : \mathcal{F} \to \mathcal{F}'$ être un opérateur différentiel d'ordre $k$ le $X/S$. Laisse $\mathcal{G}$ être un quasi-cohérent
$\mathcal{O}_Y$-module. Puis il y a un opérateur différentiel unique
$$
D' :
\text{pr}_1^*\mathcal{F} \otimes_{\mathcal{O}_{X \times_S Y}}
\text{pr}_2^*\mathcal{G}
\longrightarrow
\text{pr}_1^*\mathcal{F}' \otimes_{\mathcal{O}_{X \times_S Y}}
\text{pr}_2^*\mathcal{G}
$$
de l ' ordre $k$ le $X \times_S Y / Y$ de telle manière que
$
D'(s \otimes t) = D(s) \otimes t
$
pour les sections locales $s$ de $\mathcal{F}$ et $t$ de $\mathcal{G}$.

\end{lemma}

\begin{proof}
Dans le cas où $X$, $Y$, et $S$ sont affines, ceci suit, via lemme \ref{lemma-affine-case-differential-operators}, du résultat algèbre correspondant, voir Algèbre, lemme \ref{algebra-lemma-base-change-differential-operators}. En général, on utilise des revêtements par affines (par exemple dans Schémas, lemme \ref{schemes-lemma-affine-covering-fibre-product}) pour construire $D'$ globalement. Détails omis.
\end{proof}

\begin{remark}

\label{remark-base-change-differential-operators}

Laisser $a : X \to S$ et $b : Y \to S$ morphismes des schémas. $p : X \times_S Y \to X$ et $q : X \times_S Y \to Y$ les projections. Dans cette remarque, $\mathcal{O}_X$-module $\mathcal{F}$
et une $\mathcal{O}_Y$-module $\mathcal{G}$ Prenons le coup.
$$
\mathcal{F} \boxtimes \mathcal{G} =
p^*\mathcal{F} \otimes_{\mathcal{O}_{X \times_S Y}} q^*\mathcal{G}
$$
Dénotation $\mathcal{A}_{X/S}$ la catégorie d'additifs dont les objets sont quasi-cohérents $\mathcal{O}_X$-modules et dont les morphismes sont opérateurs différentiels d'ordre fini sur $X/S$. De même pour $\mathcal{A}_{Y/S}$ et $\mathcal{A}_{X \times_S Y/S}$. La construction de lemme \ref{lemma-base-change-differential-operators}
détermine un functeur
$$
\boxtimes : 
\mathcal{A}_{X/S} \times \mathcal{A}_{Y/S} \longrightarrow
\mathcal{A}_{X \times_S Y/S},
\quad
(\mathcal{F}, \mathcal{G}) \longmapsto \mathcal{F} \boxtimes \mathcal{G}
$$
qui est bilinéaire sur les morphismes. $X = \Spec(A)$, $Y = \Spec(B)$,
$S = \Spec(R)$, puis par l'identification de faisceaux quasi-cohérents avec des modules ce functeur est donné par $(M, N) \mapsto M \otimes_R N$
sur les objets et envoie le morphisme $(D, D') : (M, N) \to (M', N')$ à
$D \otimes D' : M \otimes_R N \to M' \otimes_R N'$.

\end{remark}








\section{Morphismes lisses}
\label{section-smooth}

\noindent
Laissez $f : X \to Y$ être une application continue des espaces topologiques. Considérez la condition suivante : Pour chaque $x \in X$ il existe des quartiers ouverts $x \in U \subset X$ et $f(x) \in V \subset Y$, et un entier $d$ tel que $f(U) \subset V$ et tel que nous obtenons un diagramme commutatif $$
\xymatrix{
X \ar[d] &
U \ar[l] \ar[d] \ar[r]_-\pi &
V \times \mathbf{R}^d \ar[ld] \\
Y & V \ar[l]
}
$$ où $\pi$ est un homéomorphisme sur un sous-ensemble ouvert. morphismes de schémas sont l'analogue de ces cartes dans la catégorie de schémas. Voir lemme \ref{lemma-smooth-locally-standard-smooth} et lemme \ref{lemma-smooth-etale-over-affine-space}.

\medskip\noindent
Contrairement aux attentes (peut-être) la notion d'un homomorphisme lisse d'anneaux n'est pas définie uniquement en termes du module des différences. N'oublions pas que $R \to A$ est un {\it homomorphisme lisse d'anneaux} si $A$ est de présentation finie sur $R$ et si le complexe de cotangente naïf de $A$ sur $R$ est quasi-isomorphique à un module projectif placé dans le degré $0$, voir Algobre, définition \ref{algebra-definition-smooth}.

\begin{definition}

\label{definition-smooth}

Laisser $f : X \to S$ être un morphisme de schémas.

\begin{enumerate}

\item Nous disons que $f$ est {\it lisse à $x \in X$} s'il existe un quartier affiné d'ouverture $\Spec(A) = U \subset X$
de $x$ et ouvert affine $\Spec(R) = V \subset S$
avec $f(U) \subset V$ tel que l'homomorphisme d'anneaux induit
$R \to A$ C'est lisse.
\item Nous disons que $f$ est {\it lisse} si elle est lisse à chaque point de $X$.
\item Un morphisme des schémas affinés $f : X \to S$
est appelé {\it standard lisse} s'il existe un homomorphisme standard lisse d'anneaux $R \to R[x_1, \ldots, x_n]/(f_1, \ldots, f_c)$ (voir Algèbre, définition) \ref{algebra-definition-standard-smooth}) de telle sorte que $X \to S$ est isomorphe à
$$
\Spec(R[x_1, \ldots, x_n]/(f_1, \ldots, f_c)) \to \Spec(R).
$$

\end{enumerate}

\end{definition}

\noindent
Une caractéristique agréable de cette définition est que l'ensemble de points où un morphisme est lisse est automatiquement ouvert.

\medskip\noindent
Notez qu'il n'y a pas d'hypothèses de séparation ou de quasi-compactitude dans la définition. C'est pourquoi la question d'être lisse est de nature locale sur la source. Voici le résultat précis.

\begin{lemma}

\label{lemma-smooth-characterize}

Laisser $f : X \to S$ morphisme des régimes. Les suivants sont équivalents

\begin{enumerate}

\item Le morphisme $f$ C'est lisse.
\item Pour tous les ouverts $U \subset X$, $V \subset S$
avec $f(U) \subset V$ l'homomorphisme d'anneaux
$\mathcal{O}_S(V) \to \mathcal{O}_X(U)$ C'est lisse.
\item Il existe une ouverture de recouvrement $S = \bigcup_{j \in J} V_j$
et couvertures ouvertes $f^{-1}(V_j) = \bigcup_{i \in I_j} U_i$ de sorte que chacun des morphismes $U_i \to V_j$, $j\in J, i\in I_j$
C'est lisse.
\item Il existe un recouvrement ouvert affine $S = \bigcup_{j \in J} V_j$
et revêtements pour affines d'ouverture $f^{-1}(V_j) = \bigcup_{i \in I_j} U_i$ tel que l'homomorphisme d'anneaux $\mathcal{O}_S(V_j) \to \mathcal{O}_X(U_i)$ est lisse, pour tous $j\in J, i\in I_j$.

\end{enumerate}

En outre, si $f$ est lisse alors pour tous les sous-systèmes ouverts $U \subset X$, $V \subset S$ avec $f(U) \subset V$
la restriction $f|_U : U \to V$ C'est lisse.

\end{lemma}

\begin{proof}

Cela s'explique par le fait qu'il s'agit d'une \ref{lemma-locally-P} si nous montrons que la propriété ``$R \to A$ Nous vérifions les conditions (a), (b) et (c) de définition
\ref{definition-property-local}. Par Algèbre, lemme \ref{algebra-lemma-base-change-smooth}
être lisse est stable sous changement de base et donc nous concluons (a) tient. Par Algèbre, lemme \ref{algebra-lemma-compose-smooth}
être lisse est stable sous composition et pour n'importe quel anneau
$R$ l'homomorphisme d'anneaux $R \to R_f$ est (standard) lisse. Nous concluons (b) détient. Enfin, la propriété (c) est vraie selon Algèbre, lemme \ref{algebra-lemma-locally-smooth}.

\end{proof}

\noindent
Le lemme suivant caractérise un morphisme lisse comme un morphisme plat, de présentation finie avec des fibres lisses. Notez que les schémas lisses sur un champ sont discutés plus en détail dans Variétés, section \ref{varieties-section-smooth}.

\begin{lemma}

\label{lemma-smooth-flat-smooth-fibres}
Laissez $f : X \to S$ être un morphisme de schémas. Si $f$ est plat, localment de présentation finie, et toutes les fibres $X_s$ sont lisses, alors $f$ est lisse.
\end{lemma}

\begin{proof}
Suit d'Algèbre, lemme \ref{algebra-lemma-flat-fibre-smooth}.
\end{proof}

\begin{lemma}

\label{lemma-composition-smooth}

La composition de deux morphismes lisses est lisse.

\end{lemma}

\begin{proof}
Dans la démonstration de lemme \ref{lemma-smooth-characterize} nous avons vu qu'être lisse est une propriété locale d'homomorphismes d'anneaux. D'où la première affirmation de lemme suit de lemme \ref{lemma-composition-type-P} combinée au fait qu'être lisse est une propriété d'homomorphismes d'anneaux qui est stable sous composition, voir Algèbre, lemme \ref{algebra-lemma-compose-smooth}.
\end{proof}

\begin{lemma}

\label{lemma-base-change-smooth}
Le changement de base d'un morphisme lisse est lisse.
\end{lemma}

\begin{proof}
Dans la démonstration de lemme \ref{lemma-smooth-characterize} nous avons vu qu'être lisse est une propriété locale d'homomorphismes d'anneaux. C'est pourquoi le lemme suit de lemme \ref{lemma-composition-type-P} combiné au fait qu'être lisse est une propriété d'homomorphismes d'anneaux qui est stable sous changement de base, voir Algèbre, lemme \ref{algebra-lemma-base-change-smooth}.
\end{proof}

\begin{lemma}

\label{lemma-open-immersion-smooth}

Toute immersion ouverte est lisse.

\end{lemma}

\begin{proof}

C'est vrai parce qu'une immersion ouverte est un isomorphisme local.

\end{proof}

\begin{lemma}

\label{lemma-smooth-syntomic}
Une lisse morphisme est syntomique.
\end{lemma}

\begin{proof}
Voir Algèbre, lemme \ref{algebra-lemma-smooth-syntomic}.
\end{proof}

\begin{lemma}

\label{lemma-smooth-locally-finite-presentation}

A morphisme lisse is localement de présentation finie.

\end{lemma}

\begin{proof}

C'est vrai parce qu'un homomorphisme lisse d'anneaux est de présentation finie par définition.

\end{proof}

\begin{lemma}

\label{lemma-smooth-flat}
Une lisse morphisme est plate.
\end{lemma}

\begin{proof}
Combiner les lemmes \ref{lemma-syntomic-flat} et \ref{lemma-smooth-syntomic}.
\end{proof}

\begin{lemma}

\label{lemma-smooth-open}
Un morphisme lisse est universel ouvert.
\end{lemma}

\begin{proof}
Combiner les lemmes \ref{lemma-smooth-flat}, \ref{lemma-smooth-locally-finite-presentation} et \ref{lemma-fppf-open}. Ou bien, combiner les lemmes \ref{lemma-smooth-syntomic}, \ref{lemma-syntomic-open}.
\end{proof}

\noindent
Le lemme suivant dit localement que n'importe quelle lisse morphisme est lisse standard. Par conséquent, nous pouvons utiliser les lisses morphismes standard comme un modèle local {\it} pour une lisse morphisme.

\begin{lemma}

\label{lemma-smooth-locally-standard-smooth}

\begin{slogan}

morphismes lisses sont localement standard lisses.

\end{slogan}

Laisser $f : X  \to S$ être un morphisme de schémas. $x \in X$ Je ne suis pas d'accord. $V \subset S$ être un quartier affiné ouvert de $f(x)$. Les conditions suivantes sont équivalentes

\begin{enumerate}

\item Le morphisme $f$ est lisse à $x$.
\item Il existe un ouvert affine $U \subset X$Avec $x \in U$ et $f(U) \subset V$ tel que le morphisme induit $f|_U : U \to V$ est standard lisse.

\end{enumerate}

\end{lemma}

\begin{proof}
Suit des définitions et Algèbre, lemmes \ref{algebra-lemma-standard-smooth} et \ref{algebra-lemma-smooth-syntomic}.
\end{proof}

\begin{lemma}

\label{lemma-smooth-omega-finite-locally-free}
Laissez $f : X \to S$ être un morphisme de schémas. Supposons que $f$ est lisse. Puis le module des différentes $\Omega_{X/S}$ de $X$ sur $S$ est fini localement libre et $$
\text{rank}_x(\Omega_{X/S}) = \dim_x(X_{f(x)})
$$ pour chaque $x \in X$.
\end{lemma}

\begin{proof}
L'énoncé est local sur $X$ et $S$. Par lemme \ref{lemma-smooth-locally-standard-smooth} ci-dessus, nous pouvons supposer que $f$ est une lisse morphisme standard d'affines. Dans ce cas, le résultat suit d'Algèbre, lemme \ref{algebra-lemma-standard-smooth} (et la définition d'une intersection globale relative complète, voir Algobre, définition \ref{algebra-definition-relative-global-complete-intersection}).
\end{proof}

\noindent
lemme \ref{lemma-smooth-omega-finite-locally-free} dit que la définition suivante a un sens.

\begin{definition}

\label{definition-smooth-relative-dimension}
Laissez $d \geq 0$ être un entier. Nous disons un morphisme de schémas $f : X \to S$ est {\it lisse de dimension relative $d$} si $f$ est lisse et $\Omega_{X/S}$ est fini localement libre de rang constant $d$.
\end{definition}

\noindent
En d'autres termes, $f$ est lisse et les fibres non vides sont équidimensionnelles de dimension $d$. Par lemme \ref{lemma-smooth-at-point} ci-dessous ceci est également le même que ce qui nécessite: (a) $f$ est localment de présentation finie, (b) $f$ est plat, (c) toutes les fibres non vides équidimensionnelles de dimension $d$, et (d) $\Omega_{X/S}$ finite localement libre de rang $d$. Il ne suffit pas de supposer simplement que $f$ est plat, de présentation finie, et $\Omega_{X/S}$ finite localement libre de rang $d$. Un contre-exemple est donné par $\Spec(\mathbf{F}_p[t]) \to \Spec(\mathbf{F}_p[t^p])$.

\medskip\noindent
Voici un critère différentiel de douceur à un point. Il y a beaucoup de variantes de ce résultat qui peuvent être utiles à un moment donné. Nous allons juste les ajouter ici au besoin.

\begin{lemma}

\label{lemma-smooth-at-point}

Laisser $f : X \to S$ être un morphisme de schémas. $x \in X$. Prêt $s = f(x)$Supposons que $f$ est localement de présentation finie. Les conditions suivantes sont équivalentes:

\begin{enumerate}

\item Le morphisme $f$ est lisse à $x$.
\item La carte locale de l'anneau $\mathcal{O}_{S, s} \to \mathcal{O}_{X, x}$
est plat et $X_s \to \Spec(\kappa(s))$ est lisse à $x$.
\item La carte locale de l'anneau $\mathcal{O}_{S, s} \to \mathcal{O}_{X, x}$
est plat et le $\mathcal{O}_{X, x}$-module $\Omega_{X/S, x}$
peut être généré par au plus $\dim_x(X_{f(x)})$ des éléments.
\item La carte locale de l'anneau $\mathcal{O}_{S, s} \to \mathcal{O}_{X, x}$
est plat et le $\kappa(x)$- espace vectoriel
$$
\Omega_{X_s/s, x} \otimes_{\mathcal{O}_{X_s, x}} \kappa(x) =
\Omega_{X/S, x} \otimes_{\mathcal{O}_{X, x}} \kappa(x)
$$
peut être généré par au plus $\dim_x(X_{f(x)})$ des éléments.
\item Il existe des ouverts affines $U \subset X$, $V \subset S$ de telle manière que $x \in U$, $f(U) \subset V$ et le morphisme induit $f|_U : U \to V$ est standard lisse.
\item Il existe des ouverts affines $\Spec(A) = U \subset X$
et $\Spec(R) = V \subset S$ avec $x \in U$ correspondant à $\mathfrak q \subset A$, $f(U) \subset V$
telle qu'il existe une présentation
$$
A = R[x_1, \ldots, x_n]/(f_1, \ldots, f_c)
$$
avec
$$
g =
\det
\left(
\begin{matrix}
\partial f_1/\partial x_1 &
\partial f_2/\partial x_1 &
\ldots &
\partial f_c/\partial x_1 \\
\partial f_1/\partial x_2 &
\partial f_2/\partial x_2 &
\ldots &
\partial f_c/\partial x_2 \\
\ldots & \ldots & \ldots & \ldots \\
\partial f_1/\partial x_c &
\partial f_2/\partial x_c &
\ldots &
\partial f_c/\partial x_c
\end{matrix}
\right)
$$
cartographie d'un élément de $A$ pas en $\mathfrak q$.

\end{enumerate}

\end{lemma}

\begin{proof}
Notez que si $f$ est lisse à $x$, alors nous voyons de lemme \ref{lemma-smooth-locally-standard-smooth} qui (5) tient, et (6) est une version légèrement affaiblie de (5). De plus, $f$ lisse implique que l'application de l'anneau $\mathcal{O}_{S, s} \to \mathcal{O}_{X, x}$ est plate (voir lemme \ref{lemma-smooth-flat}) et que $\Omega_{X/S}$ est finie localement libre de rang égal à $\dim_x(X_s)$ (voir lemme \ref{lemma-smooth-omega-finite-locally-free}). Ainsi (1) implique (3) et (4). Par lemme \ref{lemma-base-change-smooth} nous voyons également que (1) implique (2).

\medskip\noindent
Par lemme \ref{lemma-base-change-differentials} le module des différentes $\Omega_{X_s/s}$ de la fibre $X_s$ sur $\kappa(s)$ est le retrait du module des différentes $\Omega_{X/S}$ de $X$ sur $S$. D'où l'égalité affichée dans la partie (4) de la lemme. Par lemme \ref{lemma-finite-type-differentials} ces modules sont de type fini. D'où le nombre minimal de générateurs des modules $\Omega_{X/S, x}$ et $\Omega_{X_s/s, x}$ est le même et égal à la dimension de cet espace $\kappa(x)$-vector par le lemme de Nakayama (Algèbre, lemme \ref{algebra-lemma-NAK}). Ceci montre en particulier que (3) et (4) sont équivalents.

\medskip\noindent
Algèbre, lemme \ref{algebra-lemma-flat-fibre-smooth} montre que (2) implique (1). Algèbre, lemme \ref{algebra-lemma-characterize-smooth-over-field} montre que (3) et (4) implique (2). Enfin, (6) implique (5) voir par exemple Algèbre, exemple \ref{algebra-example-make-standard-smooth} et (5) implique (1) par Algèbre, lemme \ref{algebra-lemma-standard-smooth}.
\end{proof}

\begin{lemma}

\label{lemma-set-points-where-fibres-smooth}
Laissez $$
\xymatrix{
X' \ar[r]_{g'} \ar[d]_{f'} & X \ar[d]^f \\
S' \ar[r]^g & S
}
$$ être un diagramme cartésien des schémas. Laissez $W \subset X$, resp.\ $W' \subset X'$ être le sous-système ouvert des points où $f$, resp.\ $f'$ est lisse. Ensuite, $W' = (g')^{-1}(W)$ si
\begin{enumerate}
\item $f$ est plat et localement de présentation finie, ou
\item $f$ est localment de présentation finie et $g$ est plat.
\end{enumerate}

\end{lemma}

\begin{proof}

Supposons d'abord que $f$ localement de type fini. Considérez l'ensemble
$$
T = \{x \in X \mid X_{f(x)}\text{ is smooth over }\kappa(f(x))\text{ at }x\}
$$
et l'ensemble correspondant $T' \subset X'$ pour $f'$. Ensuite, nous réclamons
$T' = (g')^{-1}(T)$. A savoir, laissez $s' \in S'$ être un point, et laisser
$s = g(s')$. Alors nous avons
$$
X'_{s'} =
\Spec(\kappa(s')) \times_{\Spec(\kappa(s))} X_s
$$
En d'autres termes, les fibres du changement de base sont les modifications de base des fibres. \ref{algebra-lemma-smooth-field-change-local}.

\medskip\noindent
Ainsi le cas (1) suit parce que dans le cas (1) $T$ est l'ensemble (ouvert) de points où $f$ est lisse par lemme \ref{lemma-smooth-at-point}.

\medskip\noindent
Dans le cas (2) laissez $x' \in W'$. Puis $g'$ est plat à $x'$ (lemme \ref{lemma-base-change-module-flat}) et $g \circ f'$ est plat à $x'$ (lemme \ref{lemma-composition-module-flat}). Il en résulte que $f$ est plat à $x = g'(x')$ par lemme \ref{lemma-flat-permanence}. D'autre part, depuis $x' \in T'$ (lemme \ref{lemma-base-change-smooth}) nous voyons que $x \in T$. Par conséquent $f$ est lisse à $x$ par lemme \ref{lemma-smooth-at-point}.
\end{proof}
X

\noindent
Voici un lemme qui utilise en fait la disparition de $H^{-1}$ du complexe de cotange naïf pour un homomorphisme d'anneaux lisse.

\begin{lemma}

\label{lemma-triangle-differentials-smooth}
Laissez $f : X \to Y$, $g : Y \to S$ être des morphismes de schémas. Supposons que $f$ est lisse. Alors $$
0 \to f^*\Omega_{Y/S} \to \Omega_{X/S} \to \Omega_{X/Y} \to 0
$$ (voir lemme \ref{lemma-triangle-differentials}) est courte exacte.
\end{lemma}

\begin{proof}
La version algébrique de cette lemme est la suivante : Compte tenu des homomorphismes d'anneaux $A \to B \to C$ avec $B \to C$ lisse, puis la séquence $$
0 \to C \otimes_B \Omega_{B/A} \to \Omega_{C/A} \to \Omega_{C/B} \to 0
$$ d'Algèbre, lemme \ref{algebra-lemma-exact-sequence-differentials} est exact. Ceci est Algèbre, lemme \ref{algebra-lemma-triangle-differentials-smooth}.
\end{proof}

\begin{lemma}

\label{lemma-differentials-relative-immersion-smooth}
Laissez $i : Z \to X$ être une immersion de schémas sur $S$. Supposons que $Z$ est lisse sur $S$. Puis la suite exacte canonique $$
0 \to \mathcal{C}_{Z/X} \to i^*\Omega_{X/S} \to \Omega_{Z/S} \to 0
$$ de lemme \ref{lemma-differentials-relative-immersion} est exacte.
\end{lemma}

\begin{proof}
La version algébrique de cette lemme est la suivante : Compte tenu des homomorphismes d'anneaux $A \to B \to C$ avec $A \to C$ lisse et $B \to C$ surjectif avec noyau $J$, puis la séquence $$
0 \to J/J^2 \to C \otimes_B \Omega_{B/A} \to \Omega_{C/A} \to 0
$$ d'Algèbre, lemme \ref{algebra-lemma-differential-seq} est exacte. Ceci est Algobre, lemme \ref{algebra-lemma-differential-seq-smooth}.
\end{proof}

\begin{lemma}

\label{lemma-two-immersions-smooth}
Laissez $$
\xymatrix{
Z \ar[r]_i \ar[rd]_j & X \ar[d] \\
& Y
}
$$ être un diagramme commutatif des schémas où $i$ et $j$ sont des immersions et $X \to Y$ est lisse. Puis la suite exacte canonique $$
0 \to  \mathcal{C}_{Z/Y} \to \mathcal{C}_{Z/X} \to i^*\Omega_{X/Y} \to 0
$$ de lemme \ref{lemma-two-immersions} est exacte.
\end{lemma}

\begin{proof}
La version algébrique de cette lemme est la suivante : Compte tenu des homomorphismes d'anneaux $A \to B \to C$ avec $A \to C$ surjectif et $A \to B$ lisse, puis la séquence $$
0 \to I/I^2 \to J/J^2 \to C \otimes_B \Omega_{B/A} \to 0
$$ d'Algèbre, lemme \ref{algebra-lemma-application-NL} est exacte. Ceci est Algèbre, lemme \ref{algebra-lemma-application-NL-smooth}.
\end{proof}

\begin{lemma}

\label{lemma-smooth-permanence}

Laisser
$$
\xymatrix{
X \ar[rr]_f \ar[rd]_p & &
Y \ar[dl]^q \\
& S
}
$$
être un diagramme commutatif des morphismes des schémas.

\begin{enumerate}

\item  $f$ est surjectif, et lisse,
\item  $p$ est lisse, et
\item  $q$ est localement de présentation finie\footnote{En fait cela est implicite par (1) et (2), voir Descente, lemme \ref{descent-lemma-flat-finitely-presented-permanence}. En outre, il suffit de supposer $f$ est surjectif, plat et localement de présentation finie, voir Descente, lemme \ref{descent-lemma-smooth-permanence}.}.

\end{enumerate}

Alors $q$ C'est lisse.

\end{lemma}

\begin{proof}

Par moi \ref{lemma-flat-permanence} nous voyons que $q$ C'est plat. Choisissez un point $y \in Y$. Choisissez un point $x \in X$ cartographie vers $y$. Supposons que $f$ a dimension relative $a$ au $x$ et $p$ a dimension relative $b$ au $x$. Par lemme \ref{lemma-smooth-omega-finite-locally-free}
Cela signifie que $\Omega_{X/S, x}$ est exempt de grade $b$ et $\Omega_{X/Y, x}$
est exempt de grade $a$. Par la suite exacte courtoisie de lemme \ref{lemma-triangle-differentials-smooth}
Cela signifie que $(f^*\Omega_{Y/S})_x$ est exempt de grade $b - a$. Par le laisse de Nakayama cela implique que
$\Omega_{Y/S, y}$ peut être générée par $b - a$ éléments. Aussi, par lemme \ref{lemma-dimension-fibre-at-a-point-additive} nous voyons que
$\dim_y(Y_s) = b - a$. Nous concluons donc que
$Y \to S$ est lisse à $y$ par lemme \ref{lemma-smooth-at-point} Partie 2).

\end{proof}

\noindent
Dans la situation de la lemme suivante, l'image de $\sigma$ est localement sur $X$ découpée par une séquence régulière, voir Diviseurs, lemme \ref{divisors-lemma-section-smooth-regular-immersion}.

\begin{lemma}

\label{lemma-section-smooth-morphism}

Laisser $f : X \to S$ être un morphisme de schémas. $\sigma : S \to X$ être une section de $f$. Laisse $s \in S$ être un point tel que $f$ est lisse à $x = \sigma(s)$. Ensuite, il existe des quartiers affines ouverts
$\Spec(A) = U \subset S$ de $s$ et $\Spec(B) = V \subset X$
de $x$ de telle manière que

\begin{enumerate}

\item  $f(V) \subset U$ et $\sigma(U) \subset V$,
\item avec $I = \Ker(\sigma^\# : B \to A)$ le module $I/I^2$
est un libre $A$-module, et
\item  $B^\wedge \cong A[[x_1, \ldots, x_d]]$ comme $A$-Algèbres où
$B^\wedge$ désigne l'achèvement de $B$ en ce qui concerne $I$.

\end{enumerate}

\end{lemma}

\begin{proof}

Choisissez un ouvert affine $U \subset S$ contenant $s$
Choisissez un ouvert affine $V \subset f^{-1}(U)$ contenant $x$. Choisissez un ouvert affine $U' \subset \sigma^{-1}(V)$ contenant $s$. Notez que $V' = f^{-1}(U') \cap V$ est affiné car il est égal au produit fibré $V' = U' \times_U V$. Ensuite $U'$ et $V'$ 1) Écrire $U' = \Spec(A')$ et $V' = \Spec(B')$. Par Algèbre, lemme \ref{algebra-lemma-section-smooth}
le module $I'/(I')^2$ est fini localement libre comme un $A'$- Module. Donc après le remplacement $U'$ par un ouvert affine plus petite $U'' \subset U'$
et $V'$ par $V'' = V' \cap f^{-1}(U'')$ nous obtenons la situation où
$I''/(I'')^2$ est libre, c'est-à-dire, (2) détient. Dans ce cas (3) détient également par Algèbre, lemme \ref{algebra-lemma-section-smooth}.

\end{proof}

\noindent
La dimension d'un schéma $X$ à un point $x$ (Propriétés, définition \ref{properties-definition-dimension}) est juste la dimension de $X$ à $x$ comme un espace topologique, voir Topologie, définition \ref{topology-definition-Krull}. Ce n'est pas la dimension de l'anneau local $\mathcal{O}_{X,x}$, en général.

\begin{lemma}

\label{lemma-smoothness-dimension}
Laissez $f : X \to Y$ être un morphisme de schémas de localisation noéthérien. Pour chaque point $x$ dans $X$ avec l'image $y$ dans $Y$, $$
\dim_x(X) = \dim_y(Y) + \dim_x(X_y),
$$ où $X_y$ désigne la fibre sur $y$.
\end{lemma}

\begin{proof}
Après avoir remplacé $X$ par un quartier ouvert de $x$, il y a un numéro naturel $d$ tel que toutes les fibres de $X \to Y$ ont la dimension $d$ à chaque point, voir lemme \ref{lemma-smooth-omega-finite-locally-free}. Ensuite, $f$ est plat (lemme \ref{lemma-smooth-flat}), localement de type fini (lemme \ref{lemma-smooth-locally-finite-presentation}), et de dimension relative $d$. D'où le résultat suit de lemme \ref{lemma-rel-dimension-dimension}.
\end{proof}












\section{Morphismes non ramifiés}
\label{section-unramified}

\noindent
Nous discutons brièvement des morphismes non ramifiés avant la classe (peut-être) plus intéressante des morphismes \'etale. Rappelons qu'un homomorphisme d'anneaux $R \to A$ est {\it non ramifié} si c'est de type fini et $\Omega_{A/R} = 0$ (c'est la définition de \cite{Henselian}). Un homomorphisme d'anneaux $R \to A$ est appelé {\it G-unramifié} si c'est de présentation finie et $\Omega_{A/R} = 0$ (c'est la définition de \cite{EGA}). Voir Algèbre, définition \ref{algebra-definition-unramified}.

\begin{definition}

\label{definition-unramified}

Laisser $f : X \to S$ être un morphisme de schémas.

\begin{enumerate}

\item Nous disons que $f$ est {\it non ramifiés à $x \in X$} s'il existe un quartier affiné d'ouverture $\Spec(A) = U \subset X$
de $x$ et ouvert affine $\Spec(R) = V \subset S$
avec $f(U) \subset V$ tel que l'homomorphisme d'anneaux induit
$R \to A$ n'est pas ramifié.
\item Nous disons que $f$ est {\it G-non ramifié à $x \in X$} s'il existe un quartier affiné d'ouverture $\Spec(A) = U \subset X$
de $x$ et ouvert affine $\Spec(R) = V \subset S$
avec $f(U) \subset V$ tel que l'homomorphisme d'anneaux induit
$R \to A$ est G-non ramifié.
\item Nous disons que $f$ est {\it non ramifié} s'il n'est pas ramifié à chaque point de $X$.
\item Nous disons que $f$ est {\it G-non-ramifié} s'il est G-non-ramifié à chaque point de $X$.

\end{enumerate}

\end{definition}

\noindent
Notez qu'un G-morphisme non ramifié n'est pas ramifié. D'où tout résultat pour les morphismes non ramifiés implique le résultat correspondant pour les morphismes G-non ramifiés. De plus, si $S$ est localement noethérien alors il n'y a pas de différence entre G-unramifié et les morphismes non ramifiés, voir lemme \ref{lemma-noetherian-unramified}. Une caractéristique agréable de cette définition est que l'ensemble des points où un morphisme est non ramifié (resp.\ G-unramifié) est automatiquement ouvert.

\begin{lemma}

\label{lemma-unramified-omega-zero}
Laissez $f : X \to S$ être un morphisme de schémas. Ensuite,
\begin{enumerate}
\item $f$ n'est pas ramifié si et seulement si $f$ est localement de type fini et $\Omega_{X/S} = 0$, et
\item $f$ est G-non ramifié si et seulement si $f$ est localment de présentation finie et $\Omega_{X/S} = 0$.
\end{enumerate}

\end{lemma}

\begin{proof}
Par définition, un homomorphisme d'anneaux $R \to A$ n'est pas ramifié (resp.\ G-non ramifié) si et seulement s'il est de type fini (resp.\ finite présentation) et $\Omega_{A/R} = 0$. Par conséquent, la lemme suit directement des définitions et lemme \ref{lemma-differentials-affine}.
\end{proof}

\noindent
Notez qu'il n'y a pas d'hypothèses de séparation ou de quasi-compactitude dans la définition. D'où la question d'être non ramifié est de nature locale sur la source. Voici le résultat précis.

\begin{lemma}

\label{lemma-unramified-characterize}

Laisser $f : X \to S$ morphisme des régimes. Les suivants sont équivalents

\begin{enumerate}

\item Le morphisme $f$ n'est pas ramifié (resp.\ G-non ramifié).
\item Pour chaque ouvert affine $U \subset X$, $V \subset S$
avec $f(U) \subset V$ l'homomorphisme d'anneaux
$\mathcal{O}_S(V) \to \mathcal{O}_X(U)$ n'est pas ramifié (resp.\ G-non ramifié).
\item Il existe une ouverture de recouvrement $S = \bigcup_{j \in J} V_j$
et couvertures ouvertes $f^{-1}(V_j) = \bigcup_{i \in I_j} U_i$ de sorte que chacun des morphismes $U_i \to V_j$, $j\in J, i\in I_j$
n'est pas ramifié (resp.\ G-non ramifié).
\item Il existe un recouvrement ouvert affine $S = \bigcup_{j \in J} V_j$
et revêtements pour affines d'ouverture $f^{-1}(V_j) = \bigcup_{i \in I_j} U_i$ tel que l'homomorphisme d'anneaux $\mathcal{O}_S(V_j) \to \mathcal{O}_X(U_i)$ n'est pas ramifié (resp.\ G-non ramifié), pour tous $j\in J, i\in I_j$.

\end{enumerate}

En outre, si $f$ n'est pas ramifié (resp.\ G-non ramifié) puis pour tous les sous-systèmes ouverts $U \subset X$, $V \subset S$ avec $f(U) \subset V$
la restriction $f|_U : U \to V$ n'est pas ramifié (resp.\ G-non ramifié).

\end{lemma}

\begin{proof}
Ceci découle de lemme \ref{lemma-locally-P} si nous montrons que la propriété ``$R \to A$ n'est pas ramifiée'' est locale. Nous vérifions les conditions (a), (b) et (c) de définition \ref{definition-property-local}. Ces propriétés sont prouvées dans Algèbre, lemme \ref{algebra-lemma-unramified}.
\end{proof}

\begin{lemma}

\label{lemma-composition-unramified}

La composition de deux morphismes non ramifiés n'est pas ramifiée. La même chose vaut pour G-morphismes non ramifiés.

\end{lemma}

\begin{proof}
La preuve de lemme \ref{lemma-unramified-characterize} montre qu'être non ramifié (resp.\ G-non ramifié) est une propriété locale d'homomorphismes d'anneaux. D'où la première affirmation de la lemme suit de lemme \ref{lemma-composition-type-P} combinée avec le fait d'être non ramifié (resp.\ G-non ramifié) est une propriété d'homomorphismes d'anneaux qui est stable sous composition, voir Algèbre, lemme \ref{algebra-lemma-unramified}.
\end{proof}

\begin{lemma}

\label{lemma-base-change-unramified}

Le changement de base d'un morphisme non ramifié n'est pas ramifié. Il en va de même pour les G-morphismes non ramifiés.

\end{lemma}

\begin{proof}
La preuve de lemme \ref{lemma-unramified-characterize} montre qu'être non ramifié (resp.\ G-non ramifié) est une propriété locale d'homomorphismes d'anneaux. C'est pourquoi le lemme suit de lemme \ref{lemma-base-change-type-P} combiné avec le fait d'être non ramifié (resp.\ G-non ramifié) est une propriété d'homomorphismes d'anneaux qui est stable sous changement de base, voir Algobre, lemme \ref{algebra-lemma-unramified}.
\end{proof}

\begin{lemma}

\label{lemma-noetherian-unramified}
Laissez $f : X \to S$ être un morphisme de schémas. Supposons que $S$ est localement noethérien. Alors $f$ est non ramifié si et seulement si $f$ est G-non ramifié.
\end{lemma}

\begin{proof}
Suit des définitions et laisse \ref{lemma-noetherian-finite-type-finite-presentation}.
\end{proof}

\begin{lemma}

\label{lemma-open-immersion-unramified}
Toute ouverture d'immersion n'est pas G.
\end{lemma}

\begin{proof}

C'est vrai parce qu'une immersion ouverte est un isomorphisme local.

\end{proof}

\begin{lemma}

\label{lemma-closed-immersion-unramified}

Une fermée d'immersion $i : Z \to X$ Il est G-non-ramifié si et seulement si le faisceau associé quasi-cohérent des idéaux $\mathcal{I} = \Ker(\mathcal{O}_X \to i_*\mathcal{O}_Z)$
est de type fini (en tant que $\mathcal{O}_X$- Module).

\end{lemma}

\begin{proof}
Suit de \ref{lemma-closed-immersion-finite-presentation} et d'Algèbre, de \ref{algebra-lemma-unramified}.
\end{proof}

\begin{lemma}

\label{lemma-unramified-locally-finite-type}
Un morphisme non ramifie est localement de type fini. Un G-morphisme non ramifie est localment de présentation finie.
\end{lemma}

\begin{proof}

Un homomorphisme d'anneaux non ramifié est de type fini par définition. Un homomorphisme d'anneaux non ramifié est de présentation finie par définition.

\end{proof}

\begin{lemma}

\label{lemma-unramified-quasi-finite}
Laissez $f : X \to S$ être un morphisme de schémas. Si $f$ n'est pas ramifié à $x$ alors $f$ est quasi-fini à $x$. En particulier, un morphisme non ramifié est localement quasi-fini.
\end{lemma}

\begin{proof}
Voir Algèbre, lemme \ref{algebra-lemma-unramified-quasi-finite}.
\end{proof}

\begin{lemma}

\label{lemma-unramified-over-field}

Fibres de morphismes non ramifiés.

\begin{enumerate}

\item Laisser $X$ être un régime sur un terrain $k$. Le morphisme structurel $X \to \Spec(k)$ n'est pas ramifié si et seulement si $X$ est une union disjointe des spectres des extensions de champ séparables finies de $k$.
\item Si $f : X \to S$ est un morphisme non ramifié alors pour chaque $s \in S$
les fibres $X_s$ est une union disjointe des spectres des extensions de champ séparables finies de $\kappa(s)$.

\end{enumerate}

\end{lemma}

\begin{proof}

La partie 2 suit la partie 1 et laisse \ref{lemma-base-change-unramified}. Prouvons la partie (1). Nous utilisons d'abord Algèbre, lemme \ref{algebra-lemma-characterize-unramified}. Ce lemme implique que si $X$ est une union disjointe des spectres des extensions de champ séparables finies de $k$ alors $X \to \Spec(k)$ est non ramifié. Inversement, supposez que $X \to \Spec(k)$ par Algèbre, lemme \ref{algebra-lemma-unramified-at-prime} pour chaque $x \in X$
l'extension du corps résiduel $\kappa(x)/k$ est séparable fini. $X \to \Spec(k)$ est localement quasi-finie (lemme \ref{lemma-unramified-quasi-finite}) nous voyons que tous les points de $X$ sont des points fermés isolés, voir lemme \ref{lemma-quasi-finite-at-point-characterize}. Ainsi $X$ est un espace discret, en particulier l'union disjointe des spectres de ses anneaux locaux. \ref{algebra-lemma-unramified-at-prime} Encore une fois, ces anneaux locaux sont des champs, et nous gagnons.

\end{proof}

\noindent
Le lemme suivant caractérise un morphisme non ramifiés comme des morphismes locaux de type fini avec des fibres non ramifiées.

\begin{lemma}

\label{lemma-unramified-etale-fibres}

Laisser $f : X \to S$ être un morphisme de schémas.

\begin{enumerate}

\item Si $f$ est non ramifié alors pour n'importe quel $x \in X$ l'extension du champ
$\kappa(x)/\kappa(f(x))$ est séparable fini.
\item Si $f$ est localement de type fini, et pour chaque
$s \in S$ les fibres $X_s$ est une union disjointe des spectres des extensions de champ séparables finies de $\kappa(s)$ alors $f$ n'est pas ramifié.
\item Si $f$ est localement de présentation finie, et pour chaque
$s \in S$ les fibres $X_s$ est une union disjointe des spectres des extensions de champ séparables finies de $\kappa(s)$ alors $f$ est G-non ramifié.

\end{enumerate}

\end{lemma}

\begin{proof}
Suit d'Algèbre, lemmes \ref{algebra-lemma-unramified-at-prime} et \ref{algebra-lemma-characterize-unramified}.
\end{proof}

\noindent
Voici une caractérisation des morphismes non ramifiés en termes de diagonale du morphisme.

\begin{lemma}

\label{lemma-diagonal-unramified-morphism}
Laissez $f : X \to S$ être un morphisme.
\begin{enumerate}
\item Si $f$ n'est pas ramifié, alors la diagonale morphisme $\Delta : X \to X \times_S X$ est une immersion ouverte.
\item Si $f$ est localement de type fini et $\Delta$ est une immersion ouverte, alors $f$ est non ramifié.
\item Si $f$ est localment de présentation finie et $\Delta$ est une immersion ouverte, alors $f$ est G-unramifié.
\end{enumerate}

\end{lemma}

\begin{proof}

La première déclaration découle d'Algèbre, lemme \ref{algebra-lemma-diagonal-unramified}. La deuxième déclaration du fait que $\Omega_{X/S}$
est le faisceau conormal de la diagonale morphisme (lemme \ref{lemma-differentials-diagonal}) et donc clairement zéro si $\Delta$ est une immersion ouverte.

\end{proof}

\begin{lemma}

\label{lemma-unramified-at-point}

Laisser $f : X \to S$ être un morphisme de schémas. $x \in X$. Prêt $s = f(x)$Supposons que $f$ est localement de type fini (resp.\ localement de présentation finie). Les conditions suivantes sont équivalentes:

\begin{enumerate}

\item Le morphisme $f$ n'est pas ramifié (resp.\ G-non ramifié) à $x$.
\item La fibre $X_s$ n'est pas ramifié $\kappa(s)$ au $x$.
\item Les $\mathcal{O}_{X, x}$-module $\Omega_{X/S, x}$ C'est zéro.
\item Les $\mathcal{O}_{X_s, x}$-module $\Omega_{X_s/s, x}$ C'est zéro.
\item Les $\kappa(x)$- espace vectoriel
$$
\Omega_{X_s/s, x} \otimes_{\mathcal{O}_{X_s, x}} \kappa(x) =
\Omega_{X/S, x} \otimes_{\mathcal{O}_{X, x}} \kappa(x)
$$
C'est zéro.
\item Nous avons $\mathfrak m_s\mathcal{O}_{X, x} = \mathfrak m_x$
et l ' extension sur le terrain $\kappa(x)/\kappa(s)$ est séparable fini.

\end{enumerate}

\end{lemma}

\begin{proof}

Notez que si $f$ est non ramifié à $x$, alors nous voyons que $\Omega_{X/S} = 0$ dans un quartier de $x$
par les définitions et les résultats sur les modules de différences dans la section \ref{section-sheaf-differentials}. Par conséquent (1) implique (3) et la disparition de l'espace vecteur de la main droite dans (5). Il implique aussi (2) parce que par lemme \ref{lemma-base-change-differentials}
le module des différences $\Omega_{X_s/s}$ des fibres $X_s$
terminé $\kappa(s)$ est le retrait du module des différences
$\Omega_{X/S}$ de $X$ terminé $S$. Ce fait sur les modules de différences implique également l'égalité affichée des espaces vectoriels en partie (4). \ref{lemma-finite-type-differentials}
les modules $\Omega_{X/S, x}$ et $\Omega_{X_s/s, x}$ sont de type fini. D'où les modules $\Omega_{X/S, x}$ et $\Omega_{X_s/s, x}$ sont zéro si et seulement si le correspondant $\kappa(x)$-l'espace vectoriel en (4) est zéro par le lemme de Nakayama (Algèbre, lemme \ref{algebra-lemma-NAK}), ce qui montre en particulier que (3), (4) et (5) sont équivalents. $\Omega_{X/S}$ est fermé en $X$, voir Modules, laisser \ref{modules-lemma-support-finite-type-closed}. L'hypothèse (3) implique que $x$ n'est pas dans le soutien. $\Omega_{X/S}$ est zéro dans un quartier de $x$, ce qui implique (1). L'équivalence de (1) et (3) appliquée à $X_s \to s$
Nous avons vu à ce stade que (1) -- (5) sont équivalents.

\medskip\noindent
Vous pouvez également utiliser Algèbre, lemme \ref{algebra-lemma-unramified} pour voir l'équivalence de (1) -- (5) plus directement.

\medskip\noindent

L'équivalence des points 1) et 6) découle de la
\ref{lemma-unramified-etale-fibres}. Il suit aussi plus directement d'Algèbre, lemmes \ref{algebra-lemma-unramified-at-prime} et
\ref{algebra-lemma-characterize-unramified}.

\end{proof}

\begin{lemma}

\label{lemma-set-points-where-fibres-unramified}

Laisser $f : X \to S$ être un morphisme des schémas. $f$ localement de type fini. Formation de l'ensemble ouvert

\begin{align*}
T
& =
\{x \in X \mid X_{f(x)}\text{ is unramified over }\kappa(f(x))\text{ at }x\} \\
& =
\{x \in X \mid X\text{ is unramified over }S\text{ at }x\}
\end{align*}

commutations avec changement arbitraire de base: Pour tout morphisme $g : S' \to S$, considérer le changement de base $f' : X' \to S'$ de $f$ et la projection $g' : X' \to X$. Ensuite, l'ensemble correspondant $T'$ pour le morphisme $f'$ est égal à $T' = (g')^{-1}(T)$Si $f$ est supposé localement de présentation finie puis le même tient pour l'ensemble ouvert de points où $f$ est G-non ramifié.

\end{lemma}

\begin{proof}
Laissez $s' \in S'$ être un point, et laissez $s = g(s')$. Ensuite, nous avons $$
X'_{s'} =
\Spec(\kappa(s')) \times_{\Spec(\kappa(s))} X_s
$$ En d'autres termes, les fibres du changement de base sont les changements de base des fibres. En particulier $$
\Omega_{X_s/s, x} \otimes_{\mathcal{O}_{X_s, x}} \kappa(x')
=
\Omega_{X'_{s'}/s', x'} \otimes_{\mathcal{O}_{X'_{s'}, x'}} \kappa(x')
$$ voir lemme \ref{lemma-base-change-differentials}. D'où $x' \in T'$ si et seulement si $x \in T$ par lemme \ref{lemma-unramified-at-point}. La deuxième partie suit de la première parce que dans ce cas $T$ est l'ensemble (ouvert) de points où $f$ est G-unramifié selon lemme \ref{lemma-unramified-at-point}.
\end{proof}

\begin{lemma}

\label{lemma-unramified-permanence}
Laissez $f : X \to Y$ être un morphisme de schémas sur $S$.
\begin{enumerate}
\item Si $X$ n'est pas ramifié sur $S$, alors $f$ n'est pas ramifié.
\item Si $X$ n'est pas ramifié sur $S$ et $Y$ est localement de type fini sur $S$, alors $f$ est ramifié.
\end{enumerate}

\end{lemma}

\begin{proof}
Supposons que $X$ n'est pas ramifié sur $S$. Par lemme \ref{lemma-permanence-finite-type} nous voyons que $f$ est localisation de type fini. Par l'hypothèse que nous avons $\Omega_{X/S} = 0$. D'où $\Omega_{X/Y} = 0$ par lemme \ref{lemma-triangle-differentials}. Ainsi $f$ n'est pas ramifié. Si $X$ est G-unramifié sur $S$ et $Y$ est localement de type fini sur $S$, puis par lemme \ref{lemma-finite-presentation-permanence} nous voyons que $f$ est localisation de présentation finie et nous concluons que $f$ est ramifié.
\end{proof}

\begin{lemma}

\label{lemma-value-at-one-point}

Laisser $S$ Être un plan. $X$, $Y$ les régimes sont plus $S$. Laisse $f, g : X \to Y$ les morphismes sur $S$. Laisse $x \in X$. Supposons que

\begin{enumerate}

\item le morphisme structurel $Y \to S$ n'est pas ramifié,
\item  $f(x) = g(x)$ en $Y$, dis $y = f(x) = g(x)$,
\item les cartes induites $f^\sharp, g^\sharp : \kappa(y) \to \kappa(x)$
Ils sont égaux.

\end{enumerate}

Ensuite, il existe un quartier ouvert de $x$ en $X$ sur lequel
$f$ et $g$ Ils sont égaux.

\end{lemma}

\begin{proof}
Considérez le morphisme $(f, g) : X \to Y \times_S Y$. Par hypothèse (1) et lemme \ref{lemma-diagonal-unramified-morphism} l'image inverse de $\Delta_{Y/S}(Y)$ est ouverte dans $X$. Et les hypothèses (2) et (3) impliquent que $x$ est dans ce sous-ensemble ouvert.
\end{proof}










\section{Morphismes étales}
\label{section-etale}

\noindent
La topologie Zariski d'un schéma est une topologie très grossière. Ceci est particulièrement clair lorsque l'on regarde les variétés sur $\mathbf{C}$. Il s'avère que déclarer un morphisme \'etale comme l'analogue d'un isomorphisme local en topologie introduit une topologie beaucoup plus fine. Sur les variétés sur $\mathbf{C}$ cette topologie donne lieu au nombre «correct» de Betti lors de l'informatique de la cohomologie avec des coefficients finis. Un autre observable est que si $f : X \to Y$ est un \'etale morphisme des variétés sur $\mathbf{C}$, et si $x$ est un point fermé de $X$, alors $f$ induit un isomorphisme $\mathcal{O}^{\wedge}_{Y, f(x)} \to \mathcal{O}^{\wedge}_{X, x}$ de cycles locaux complets.

\medskip\noindent
Dans cette section, nous commençons notre étude de ces questions. En fait, nous restreignons délibérément notre discussion au minimum puisque nous discuterons des résultats plus intéressants ailleurs. Rappelons qu'un homomorphisme d'anneaux $R \to A$ est dit {\it \'etale} si elle est lisse et $\Omega_{A/R} = 0$, voir Algèbre, définition \ref{algebra-definition-etale}.

\begin{definition}

\label{definition-etale}

Laisser $f : X \to S$ être un morphisme de schémas.

\begin{enumerate}

\item Nous disons que $f$ est {\it  \'de l ' état $x \in X$} s'il existe un quartier affiné d'ouverture $\Spec(A) = U \subset X$
de $x$ et ouvert affine $\Spec(R) = V \subset S$
avec $f(U) \subset V$ tel que l'homomorphisme d'anneaux induit
$R \to A$ est \'C'est une histoire.
\item Nous disons que $f$est {\it \'etale} s'il est étale en tout point de $X$.
\item Un morphisme des schémas affinés $f : X \to S$est appelé {\it standard \'etale} si $X \to S$ est isomorphe à
$$
\Spec(R[x]_h/(g)) \to \Spec(R)
$$
où $R \to R[x]_h/(g)$ est une norme \'homomorphisme étale d'anneaux, voir Algèbre, définition \ref{algebra-definition-standard-etale}, i.e., $g$ est monique et $g'$ inversible dans $R[x]_h/(g)$.

\end{enumerate}

\end{definition}

\noindent
Un morphisme est \'etal si et seulement s'il est lisse de dimension relative $0$ (voir définition \ref{definition-smooth-relative-dimension}). Une caractéristique agréable de la définition est que l'ensemble de points où un morphisme est \'etal est automatiquement ouvert.

\medskip\noindent
Notez qu'il n'y a pas d'hypothèse de séparation ou de quasi-compactitude dans la définition. D'où la question d'être \'etale est de nature locale sur la source. Voici le résultat précis.

\begin{lemma}

\label{lemma-etale-characterize}

Laisser $f : X \to S$ morphisme des régimes. Les suivants sont équivalents

\begin{enumerate}

\item Le morphisme $f$ est \'C'est une histoire.
\item Pour tous les ouverts $U \subset X$, $V \subset S$
avec $f(U) \subset V$ l'homomorphisme d'anneaux
$\mathcal{O}_S(V) \to \mathcal{O}_X(U)$ est \'C'est une histoire.
\item Il existe une ouverture de recouvrement $S = \bigcup_{j \in J} V_j$
et couvertures ouvertes $f^{-1}(V_j) = \bigcup_{i \in I_j} U_i$ de sorte que chacun des morphismes $U_i \to V_j$, $j\in J, i\in I_j$
est \'C'est une histoire.
\item Il existe un recouvrement ouvert affine $S = \bigcup_{j \in J} V_j$
et revêtements pour affines d'ouverture $f^{-1}(V_j) = \bigcup_{i \in I_j} U_i$ tel que l'homomorphisme d'anneaux $\mathcal{O}_S(V_j) \to \mathcal{O}_X(U_i)$ est
\'et, pour tous $j\in J, i\in I_j$.

\end{enumerate}

En outre, si $f$ est \'etale ensuite pour les sous-systèmes ouverts $U \subset X$, $V \subset S$ avec $f(U) \subset V$
la restriction $f|_U : U \to V$ est \'C'est une histoire.

\end{lemma}

\begin{proof}
Ceci découle de lemme \ref{lemma-locally-P} si nous montrons que la propriété ``$R \to A$ est \'etale''' est locale. Nous vérifions les conditions (a), (b) et (c) de définition \ref{definition-property-local}. Ces conditions suivent toutes Algèbre, lemme \ref{algebra-lemma-etale}.
\end{proof}

\begin{lemma}

\label{lemma-composition-etale}
La composition de deux morphismes qui sont \'etale est \'etale.
\end{lemma}

\begin{proof}
Dans la démonstration de lemme \ref{lemma-etale-characterize} nous avons vu qu'être \'etale est une propriété locale des homomorphismes d'anneaux. D'où la première affirmation de la lemme suit de lemme \ref{lemma-composition-type-P} combinée avec le fait qu'être \'etale est une propriété des homomorphismes d'anneaux qui est stable sous composition, voir Algèbre, lemme \ref{algebra-lemma-etale}.
\end{proof}

\begin{lemma}

\label{lemma-base-change-etale}
Le changement de base d'un morphisme qui est \'etal est \'etal.
\end{lemma}

\begin{proof}
Dans la démonstration de lemme \ref{lemma-etale-characterize} nous avons vu qu'être \'etale est une propriété locale des homomorphismes d'anneaux. C'est pourquoi le lemme suit de lemme \ref{lemma-composition-type-P} combiné avec le fait qu'être \'etale est une propriété des homomorphismes d'anneaux qui est stable sous changement de base, voir Algobre, lemme \ref{algebra-lemma-etale}.
\end{proof}

\begin{lemma}

\label{lemma-etale-smooth-unramified}
Laissez $f : X \to S$ être un morphisme de schémas. Laissez $x \in X$. Alors $f$ est \'etale à $x$ si et seulement si $f$ est lisse et non ramifié à $x$.
\end{lemma}

\begin{proof}
Ceci suit immédiatement des définitions.
\end{proof}

\begin{lemma}

\label{lemma-etale-locally-quasi-finite}

Une \'morphisme étale est localement quasi-fini.

\end{lemma}

\begin{proof}

Par moi \ref{lemma-etale-smooth-unramified}
et \'morphisme stale n'est pas ramifié. \ref{lemma-unramified-quasi-finite}
un morphisme non ramifié est localement quasi-fini.

\end{proof}

\begin{lemma}

\label{lemma-etale-over-field}

\begin{slogan}

Description de la \'sur les champs et les fibres \'morphismes stales.

\end{slogan}

Fibres de \'morphismes stales.

\begin{enumerate}

\item Laisser $X$ être un régime sur un terrain $k$. Le morphisme structurel $X \to \Spec(k)$ est \'si et seulement si $X$ est une union disjointe des spectres des extensions de champ séparables finies de $k$.
\item Si $f : X \to S$ est un \'morphisme stale, puis pour chaque $s \in S$ les fibres $X_s$ est une union disjointe des spectres des extensions de champ séparables finies de $\kappa(s)$.

\end{enumerate}

\end{lemma}

\begin{proof}
Vous pouvez déduire ceci de lemme \ref{lemma-unramified-over-field} via lemme \ref{lemma-etale-smooth-unramified} ci-dessus. Voici une preuve directe.

\medskip\noindent

Nous utiliserons Algèbre, lemme \ref{algebra-lemma-etale-over-field}. Il est donc clair que si $X$ est une union disjointe des spectres des extensions de champ séparables finies de $k$ alors $X \to \Spec(k)$ est \'à l'inverse, supposer que $X \to \Spec(k)$ est \'Et puis pour n'importe quelle ouvert affine $U \subset X$ nous voyons que $U$ est une union disjointe finie des spectres des extensions de champ séparables finies de $k$. D'où tous les points de $X$
les points fermés (voir lemme \ref{lemma-algebraic-residue-field-extension-closed-point-fibre}
Par exemple). Ainsi $X$ c'est un espace discret et nous gagnons.

\end{proof}

\noindent
Le lemme suivant caractérise un morphisme \'etale comme un morphisme plat, de présentation finie avec des fibres \'etale.

\begin{lemma}

\label{lemma-etale-flat-etale-fibres}
Laissez $f : X \to S$ être un morphisme de schémas. Si $f$ est plat, localement de présentation finie, et pour chaque $s \in S$ la fibre $X_s$ est une union disjointe des spectres des extensions de champ séparables finies de $\kappa(s)$, alors $f$ est \'etal.
\end{lemma}

\begin{proof}
Vous pouvez déduire ceci d'Algèbre, lemme \ref{algebra-lemma-characterize-etale}. Voici une autre preuve.

\medskip\noindent
Par lemme \ref{lemma-etale-over-field} une fibre $X_s$ est \'etale et donc lisse sur $s$. Par lemme \ref{lemma-smooth-flat-smooth-fibres} nous voyons que $X \to S$ est lisse. Par lemme \ref{lemma-unramified-etale-fibres} nous voyons que $f$ est non ramifié. Nous concluons par lemme \ref{lemma-etale-smooth-unramified}.
\end{proof}

\begin{lemma}

\label{lemma-open-immersion-etale}
Toute immersion ouverte est \'etale.
\end{lemma}

\begin{proof}

C'est vrai parce qu'une immersion ouverte est un isomorphisme local.

\end{proof}

\begin{lemma}

\label{lemma-etale-syntomic}

Une \'morphisme stale est syntomique.

\end{lemma}

\begin{proof}
Voir Algèbre, lemme \ref{algebra-lemma-smooth-syntomic} et utiliser qu'un morphisme \'etale est le même qu'une lisse morphisme de dimension relative $0$.
\end{proof}

\begin{lemma}

\label{lemma-etale-locally-finite-presentation}

Une \'morphisme stale is localement de présentation finie.

\end{lemma}

\begin{proof}

C'est vrai parce que \'homomorphisme étale d'anneaux est de présentation finie par définition.

\end{proof}

\begin{lemma}

\label{lemma-etale-flat}
Un morphisme \'etale est plat.
\end{lemma}

\begin{proof}
Combiner les lemmes \ref{lemma-syntomic-flat} et \ref{lemma-etale-syntomic}.
\end{proof}

\begin{lemma}

\label{lemma-etale-open}
Un morphisme \'etale est ouvert.
\end{lemma}

\begin{proof}
Combiner les lemmes \ref{lemma-etale-flat}, \ref{lemma-etale-locally-finite-presentation} et \ref{lemma-fppf-open}.
\end{proof}

\noindent
Le lemme suivant dit localement que tout morphisme \'etale est standard \'etale. Il s'agit en fait d'une sorte de résultat délicat à prouver en généralité complète. Les parties délicates sont cachées dans le chapitre sur l'algèbre commutative. Par conséquent, un morphisme standard \'etale est un modèle local {\it} pour un morphisme général \'etale.

\begin{lemma}

\label{lemma-etale-locally-standard-etale}
Laissez $f : X  \to S$ être un morphisme de schémas.Soit $x \in X$ un point.Soit $V \subset S$ un quartier affiné d'ouverture de $f(x)$. Les conditions suivantes sont équivalentes
\begin{enumerate}
\item Le morphisme $f$ est \'etale à $x$.
\item Il existe un morphisme affiné d'ouverture $U \subset X$ avec $x \in U$ et $f(U) \subset V$ de telle sorte que le morphisme induit $f|_U : U \to V$ est standard \'etale (voir définition \ref{definition-etale}).
\end{enumerate}

\end{lemma}

\begin{proof}
Suit des définitions et Algèbre, proposition \ref{algebra-proposition-etale-locally-standard}.
\end{proof}

\noindent
Voici un critère différentiel de \'etaleness à un point. Il y a beaucoup de variantes de ce résultat qui peuvent être utiles à un moment donné. Nous allons juste les ajouter ici au besoin.

\begin{lemma}

\label{lemma-etale-at-point}

Laisser $f : X \to S$ être un morphisme de schémas. $x \in X$. Prêt $s = f(x)$Supposons que $f$ est localement de présentation finie. Les conditions suivantes sont équivalentes:

\begin{enumerate}

\item Le morphisme $f$ est \'de l ' état $x$.
\item La carte locale de l'anneau $\mathcal{O}_{S, s} \to \mathcal{O}_{X, x}$
est plat et $X_s \to \Spec(\kappa(s))$ est \'de l ' état $x$.
\item La carte locale de l'anneau $\mathcal{O}_{S, s} \to \mathcal{O}_{X, x}$
est plat et $X_s \to \Spec(\kappa(s))$ est non ramifié à $x$.
\item La carte locale de l'anneau $\mathcal{O}_{S, s} \to \mathcal{O}_{X, x}$
est plat et le $\mathcal{O}_{X, x}$-module $\Omega_{X/S, x}$
C'est zéro.
\item La carte locale de l'anneau $\mathcal{O}_{S, s} \to \mathcal{O}_{X, x}$
est plat et le $\kappa(x)$- espace vectoriel
$$
\Omega_{X_s/s, x} \otimes_{\mathcal{O}_{X_s, x}} \kappa(x) =
\Omega_{X/S, x} \otimes_{\mathcal{O}_{X, x}} \kappa(x)
$$
C'est zéro.
\item La carte locale de l'anneau $\mathcal{O}_{S, s} \to \mathcal{O}_{X, x}$
est plat, nous avons $\mathfrak m_s\mathcal{O}_{X, x} = \mathfrak m_x$ et l ' extension sur le terrain $\kappa(x)/\kappa(s)$ est séparable fini.
\item Il existe des ouverts affines $U \subset X$, $V \subset S$ de telle manière que $x \in U$, $f(U) \subset V$ et le morphisme induit $f|_U : U \to V$ est standard lisse de dimension relative $0$.
\item Il existe des ouverts affines $\Spec(A) = U \subset X$
et $\Spec(R) = V \subset S$ avec $x \in U$ correspondant à $\mathfrak q \subset A$, $f(U) \subset V$
telle qu'il existe une présentation
$$
A = R[x_1, \ldots, x_n]/(f_1, \ldots, f_n)
$$
avec
$$
g =
\det
\left(
\begin{matrix}
\partial f_1/\partial x_1 &
\partial f_2/\partial x_1 &
\ldots &
\partial f_n/\partial x_1 \\
\partial f_1/\partial x_2 &
\partial f_2/\partial x_2 &
\ldots &
\partial f_n/\partial x_2 \\
\ldots & \ldots & \ldots & \ldots \\
\partial f_1/\partial x_n &
\partial f_2/\partial x_n &
\ldots &
\partial f_n/\partial x_n
\end{matrix}
\right)
$$
cartographie d'un élément de $A$ pas en $\mathfrak q$.
\item Il existe des ouverts affines $U \subset X$, $V \subset S$ de telle manière que $x \in U$, $f(U) \subset V$ et le morphisme induit $f|_U : U \to V$ est standard \'C'est une histoire.
\item Il existe des ouverts affines $\Spec(A) = U \subset X$
et $\Spec(R) = V \subset S$ avec $x \in U$ correspondant à $\mathfrak q \subset A$, $f(U) \subset V$
telle qu'il existe une présentation
$$
A = R[x]_Q/(P) = R[x, 1/Q]/(P)
$$
avec $P, Q \in R[x]$, $P$ et de $P' = \text{d}P/\text{d}x$ cartographie d'un élément de $A$ pas en $\mathfrak q$.

\end{enumerate}

\end{lemma}

\begin{proof}
Utilisez lemme \ref{lemma-etale-locally-standard-etale} et les définitions pour voir que (1) implique toutes les autres conditions. Pour chacune des conditions (2) -- (10) combinez lemmes \ref{lemma-smooth-at-point} et \ref{lemma-unramified-at-point} pour voir que (1) tient en montrant $f$ est à la fois lisse et non ramifié à $x$ et appliquant lemme \ref{lemma-etale-smooth-unramified}. Quelques détails omis.
\end{proof}

\begin{lemma}

\label{lemma-flat-unramified-etale}

Un morphisme est \'est à un point si et seulement s'il est plat et G-non ramifié à ce point. \'si et seulement s'il est plat et G-non ramifié.

\end{lemma}

\begin{proof}
Ceci est clair de lemmes \ref{lemma-etale-at-point} et \ref{lemma-unramified-at-point}.
\end{proof}

\begin{lemma}

\label{lemma-set-points-where-fibres-etale}
Laissez $$
\xymatrix{
X' \ar[r]_{g'} \ar[d]_{f'} & X \ar[d]^f \\
S' \ar[r]^g & S
}
$$ être un diagramme cartésien des schémas. Laissez $W \subset X$, rép.\ $W' \subset X'$ être le sous-système ouvert des points où $f$, rép.\ $f'$ est \'etale. Ensuite, $W' = (g')^{-1}(W)$ si
\begin{enumerate}
\item $f$ est plat et localement de présentation finie, ou
\item $f$ est localement de présentation finie et $g$ZX est plat.
\end{enumerate}

\end{lemma}

\begin{proof}
Supposons que d'abord que $f$ localement de type fini. Considérez l'ensemble $$
T = \{x \in X \mid f\text{ is unramified at }x\}
$$ et l'ensemble correspondant $T' \subset X'$ pour $f'$. Puis $T' = (g')^{-1}(T)$ par lemme \ref{lemma-set-points-where-fibres-unramified}.

\medskip\noindent
Ainsi le cas (1) suit parce que dans le cas (1) $T$ est l'ensemble (ouvert) de points où $f$ est \'etale par lemme \ref{lemma-flat-unramified-etale}.

\medskip\noindent
Dans le cas (2) laissez $x' \in W'$. Ensuite, $g'$ est plat à $x'$ (lemme \ref{lemma-base-change-module-flat}) et $g \circ f'$ est plat à $x'$ (lemme \ref{lemma-composition-module-flat}). Il en résulte que $f$ est plat à $x = g'(x')$ par lemme \ref{lemma-flat-permanence}. D'autre part, depuis $x' \in T'$ (lemme \ref{lemma-base-change-smooth}) nous voyons que $x \in T$. D'où $f$X est \'X à $x$ par lemme \ref{lemma-etale-at-point}.
\end{proof}
X

\begin{lemma}

\label{lemma-etale-permanence-better}

\begin{reference}
Voir par exemple \cite[remarque 18.30 (2).]{GWII}.
\end{reference}
Laissez $f : X \to Y$ être un morphisme de schémas sur $S$. Si $X$X est \'etale sur $S$ et $Y$ est non ramifié sur $S$, alors $f$ est \'etale.
\end{lemma}
X

\begin{proof}

Considérez la factorisation $X \to X \times_S Y \to Y$, où la première flèche est donnée $\text{id}_X$ et $f$ et la deuxième flèche est la projection. Nous affirmons que les deux flèches sont \'et donc $f$ est \'Par la suite, la Commission a adopté une proposition de directive sur la protection des données à caractère personnel. \ref{lemma-composition-etale}. A savoir, la projection est \'en tant que changement de base de $X \to S$, voir lemme \ref{lemma-base-change-etale}. La première flèche est le changement de base de la diagonale morphisme
$Y \to Y \times_S Y$ parce que le carré
$$
\xymatrix{
X \ar[d] \ar[r] & X \times_S Y \ar[d] \\
Y \ar[r] & Y \times_S Y
}
$$
est cartésien. La diagonale $Y \to Y \times_S Y$ est une immersion ouverte par lemme \ref{lemma-diagonal-unramified-morphism}. Le changement de base d'une ouverture d'immersion est une ouverture d'immersion (Schémas, lemme \ref{schemes-lemma-base-change-immersion}) et une immersion ouverte est \'(lemme) \ref{lemma-open-immersion-etale}).

\end{proof}

\begin{lemma}

\label{lemma-etale-permanence}

\begin{slogan}
Loi d'annulation pour \'etale morphismes
\end{slogan}
Laissez $f : X \to Y$ être un morphisme de schémas sur $S$. Si $X$ et $Y$X sont \'etale sur $S$, alors $f$X est \'etale.
\end{lemma}

\begin{proof}
Puisque $Y \to S$ n'est pas ramifié par lemme \ref{lemma-etale-smooth-unramified}, nous pouvons appliquer lemme \ref{lemma-etale-permanence-better}. Voir aussi Algèbre, lemme \ref{algebra-lemma-map-between-etale}.
\end{proof}

\begin{lemma}

\label{lemma-etale-permanence-two}

Laisser
$$
\xymatrix{
X \ar[rr]_f \ar[rd]_p & &
Y \ar[dl]^q \\
& S
}
$$
être un diagramme commutatif des morphismes des schémas.

\begin{enumerate}

\item  $f$ est surjectif, et \'d'une durée de trois ans,
\item  $p$ est \'d'une part, et
\item  $q$ est localement de présentation finie\footnote{En fait cela est implicite par (1) et (2), voir Descente, lemme \ref{descent-lemma-flat-finitely-presented-permanence}. En outre, il suffit de supposer que $f$ est surjectif, plat et localement de présentation finie, voir Descente, lemme \ref{descent-lemma-smooth-permanence}.}.

\end{enumerate}

Alors $q$ est \'C'est une histoire.

\end{lemma}

\begin{proof}
Par lemme \ref{lemma-smooth-permanence} nous voyons que $q$ est lisse. Ainsi, nous avons seulement besoin de voir que $q$ a une dimension relative $0$. Ceci découle de lemme \ref{lemma-dimension-fibre-at-a-point-additive} et du fait que $f$ et $p$ ont une dimension relative $0$.
\end{proof}

\noindent
Une caractérisation finale des morphismes lisses est qu'un morphisme lisse $f : X \to S$ est localement la composition d'un morphisme \'etale par une projection $\mathbf{A}_S^d \to S$.

\begin{lemma}

\label{lemma-smooth-etale-over-affine-space}

\begin{slogan}
Les schémas lisses sont \'etale-localement comme des espaces affinés.
\end{slogan}
Laissez $\varphi : X \to Y$ être un morphisme de schémas. Laissez $x \in X$. Laissez $V \subset Y$ être un quartier affiné ouvert de $\varphi(x)$. Si $\varphi$ est lisse à $x$, alors il existe un entier $d \geq 0$ et un objet affiné $U \subset X$ avec $x \in U$ et $\varphi(U) \subset V$ de sorte qu'il existe un diagramme commutatif $$
\xymatrix{
X \ar[d] & U \ar[l] \ar[d] \ar[r]_-\pi & \mathbf{A}^d_V \ar[ld] \\
Y & V \ar[l]
}
$$ où $\pi$ est \'etale.
\end{lemma}

\begin{proof}

Par moi \ref{lemma-smooth-locally-standard-smooth}
nous pouvons trouver un ouvert affine $U$ comme dans le lemme telle que
$\varphi|_U : U \to V$ est standard lisse.
$U = \Spec(A)$ et $V = \Spec(R)$ pour que nous puissions écrire
$$
A = R[x_1, \ldots, x_n]/(f_1, \ldots, f_c)
$$
avec
$$
g =
\det
\left(
\begin{matrix}
\partial f_1/\partial x_1 &
\partial f_2/\partial x_1 &
\ldots &
\partial f_c/\partial x_1 \\
\partial f_1/\partial x_2 &
\partial f_2/\partial x_2 &
\ldots &
\partial f_c/\partial x_2 \\
\ldots & \ldots & \ldots & \ldots \\
\partial f_1/\partial x_c &
\partial f_2/\partial x_c &
\ldots &
\partial f_c/\partial x_c
\end{matrix}
\right)
$$
cartographie vers un élément inversible de $A$. Alors il est clair que
$R[x_{c + 1}, \ldots, x_n] \to A$ est standard lisse de dimension relative $0$. Par conséquent, il est lisse de dimension relative $0$. En d'autres termes, l'homomorphisme d'anneaux $R[x_{c + 1}, \ldots, x_n] \to A$
est \'en tant que $\mathbf{A}^{n - c}_V = \Spec(R[x_{c + 1}, \ldots, x_n])$
le lemme avec $d = n - c$.

\end{proof}
















\section{Faisceaux relativement amples}
\label{section-relatively-ample}

\noindent
Laissez $X$ être un schéma et $\mathcal{L}$ un filet inversible sur $X$. Alors $\mathcal{L}$ est ample sur $X$ si $X$ est quasi-compact et chaque point de $X$ est contenu dans une affine ouvert de la forme $X_s$, où $s \in \Gamma(X, \mathcal{L}^{\otimes n})$ et $n \geq 1$, voir Propriétés, définition \ref{properties-definition-ample}. Nous transformons ceci en une notion relative comme suit.

\begin{definition}

\label{definition-relatively-ample}

\begin{reference}

\cite[II définition 4.6.1]{EGA}

\end{reference}

Laisser $f : X \to S$ être un morphisme de schémas. $\mathcal{L}$ être un inversible $\mathcal{O}_X$- Module. $\mathcal{L}$ est {\it relative ample}, ou {\it  $f$-relatif ample}, ou {\it ample sur $X/S$}, ou {\it  $f$-ample} si $f : X \to S$
est quasi-compact, et si pour chaque ouvert affine $V \subset S$
la restriction de $\mathcal{L}$ au sous-système ouvert
$f^{-1}(V)$ de $X$ C'est ample.

\end{definition}

\noindent
Nous notons que l'existence d'une gausse relative sur $X$ ne force pas le morphisme $X \to S$ à être de type fini.

\begin{lemma}

\label{lemma-ample-power-ample}
Laissez $X \to S$ être un morphisme de schémas. Laissez $\mathcal{L}$ être un module $\mathcal{O}_X$ inversible. Laissez $n \geq 1$. Puis $\mathcal{L}$ est $f$-ample si et seulement si $\mathcal{L}^{\otimes n}$ est $f$-ample.
\end{lemma}

\begin{proof}
Ceci découle des Propriétés, deme \ref{properties-lemma-ample-power-ample}.
\end{proof}

\begin{lemma}

\label{lemma-relatively-ample-separated}
Laissez $f : X \to S$ être un morphisme de schémas. S'il existe un $f$-faisceau inversible ample, alors $f$ est séparé.
\end{lemma}

\begin{proof}
Être séparé est local sur la base (voir Schémas, lemme \ref{schemes-lemma-characterize-separated} par exemple; il en découle aussi facilement de la définition). Par conséquent, nous pouvons supposer que $S$ est affine et $X$ a un faisceau inversible ample. Dans ce cas, le résultat suit de Propriétés, lemme \ref{properties-lemma-ample-separated}.
\end{proof}

\noindent
Il existe de nombreuses façons de caractériser le rapport amples gerbes inversibles, analogues aux conditions équivalentes dans Propriétés, proposition \ref{properties-proposition-characterize-ample}. Nous les ajouterons ici au besoin.

\begin{lemma}

\label{lemma-characterize-relatively-ample}

\begin{reference}\cite[II, Proposition 4.6.3]{EGA}\end{reference}

Laisser $f : X \to S$ être un morphisme quasi-compact des schémas. $\mathcal{L}$ être une gerbe inversible sur $X$. Les conditions suivantes sont équivalentes:

\begin{enumerate}

\item La gerbe inversible $\mathcal{L}$ est $f$- L'échantillon.
\item Il existe une ouverture de recouvrement $S = \bigcup V_i$
de telle sorte que chacun $\mathcal{L}|_{f^{-1}(V_i)}$ est ample par rapport à $f^{-1}(V_i) \to V_i$.
\item Il existe un recouvrement ouvert affine $S = \bigcup V_i$
de telle sorte que chacun $\mathcal{L}|_{f^{-1}(V_i)}$ C'est ample.
\item Il existe une classification quasi-cohérente $\mathcal{O}_S$-Algèbre
$\mathcal{A}$ et une application de classement $\mathcal{O}_X$-Algèbres
$\psi : f^*\mathcal{A} \to \bigoplus_{d \geq 0} \mathcal{L}^{\otimes d}$
de telle manière que $U(\psi) = X$ et
$$
r_{\mathcal{L}, \psi} :
X
\longrightarrow
\underline{\text{Proj}}_S(\mathcal{A})
$$
est une immersion ouverte (voir Constructions, lemme
\ref{constructions-lemma-invertible-map-into-relative-proj} pour la notation).
\item Le morphisme $f$ est quasi-séparée et la partie (4) ci-dessus tient avec
$\mathcal{A} = f_*(\bigoplus_{d \geq 0} \mathcal{L}^{\otimes d})$
et $\psi$ la cartographie d'adjonction.
\item Même que (4) mais exigeant seulement $r_{\mathcal{L}, \psi}$
être une immersion.

\end{enumerate}

\end{lemma}

\begin{proof}
Il est immédiat de la définition que (1) implique (2) et (2) implique (3). Il est clair que (5) implique (4).

\medskip\noindent

Supposons que (3) tient pour le recouvrement ouvert affine $S = \bigcup V_i$. Nous allons montrer (5) tient. Puisque chacun $f^{-1}(V_i)$ a un faisceau inversible ample nous voyons que $f^{-1}(V_i)$ est séparé (Propriétés, lemme \ref{properties-lemma-ample-separated}). D'où $f$ est séparé. Par Schémas, me laisse \ref{schemes-lemma-push-forward-quasi-coherent}
nous voyons que $\mathcal{A} = f_*(\bigoplus_{d \geq 0} \mathcal{L}^{\otimes d})$
est un grade quasi-cohérent $\mathcal{O}_S$-l'algèbre. $\psi : f^*\mathcal{A} \to \bigoplus_{d \geq 0} \mathcal{L}^{\otimes d}$
la cartographie d'adjonction. La description de l'open $U(\psi)$ dans Constructions, section
\ref{constructions-section-invertible-relative-proj}
et de la définition de l'amplitude $\mathcal{L}|_{f^{-1}(V_i)}$ montrer que $U(\psi) = X$. En outre, Constructions, laisse \ref{constructions-lemma-invertible-map-into-relative-proj} la partie 3) montre que la restriction de $r_{\mathcal{L}, \psi}$ à
$f^{-1}(V_i)$ est le même que le morphisme de Propriétés, lemme \ref{properties-lemma-map-into-proj}
qui est une immersion ouverte selon Propriétés, lemme \ref{properties-lemma-ample-immersion-into-proj}. C'est pourquoi (5) tient.

\medskip\noindent

Montreons que (4) implique (1). Supposons que (4). Denoter $\pi : \underline{\text{Proj}}_S(\mathcal{A}) \to S$
le morphisme structurel. Choisissez $V \subset S$ ouvert affine. Par Constructions, définition \ref{constructions-definition-relative-proj}
nous voyons que $\pi^{-1}(V) \subset \underline{\text{Proj}}_S(\mathcal{A})$
est égal à $\text{Proj}(A)$ où $A = \mathcal{A}(V)$
comme une bague notée. $r_{\mathcal{L}, \psi}$ cartes
$f^{-1}(V)$ isomorphiquement sur un quasi-compact ouvert de $\text{Proj}(A)$. En outre, $\mathcal{L}^{\otimes d}$ est isomorphe à la traction de
$\mathcal{O}_{\text{Proj}(A)}(d)$ pour certains $d \geq 1$. (Voir partie 3) de Constructions, \ref{constructions-lemma-invertible-map-into-relative-proj}
et l ' état final des travaux de construction, \ref{constructions-lemma-invertible-map-into-proj}.) Cela implique que $\mathcal{L}|_{f^{-1}(V)}$ est ample par Propriétés, lemmes \ref{properties-lemma-open-in-proj-ample}
et \ref{properties-lemma-ample-power-ample}.

\medskip\noindent

Par l'équivalence de (1) - (5) ci-dessus, nous voyons que la propriété d'être relatif amplement sur $X/S$ est local sur $S$. Nous pouvons donc supposer que $S$ est affine, et nous devons montrer que
$\mathcal{L}$ est ample sur $X$. Dans ce cas le morphisme
$r_{\mathcal{L}, \psi}$ est identifié avec le morphisme, également dénoté
$r_{\mathcal{L}, \psi} : X \to \text{Proj}(A)$ associé à l'application
$\psi : A = \mathcal{A}(V) \to \Gamma_*(X, \mathcal{L})$. (Voir les références ci-dessus.) Comme ci-dessus, nous voyons aussi que
$\mathcal{L}^{\otimes d}$ est le retrait de la gerbe
$\mathcal{O}_{\text{Proj}(A)}(d)$ pour certains $d \geq 1$. En outre, depuis $X$ est quasi-compact nous voyons que $X$ est identifié avec un sous-système fermé d'un sous-système ouvert quasi-compact $Y \subset \text{Proj}(A)$. Par Constructions, lemme
\ref{constructions-lemma-ample-on-proj}
(voir aussi Propriétés,
\ref{properties-lemma-open-in-proj-ample}) nous voyons que $\mathcal{O}_Y(d')$ est un faisceau inversible ample sur
$Y$ pour certains $d' \geq 1$. Comme la restriction d'une ample gerbe à un sous-système fermé est ample, voir Propriétés, me
\ref{properties-lemma-ample-on-closed}
nous concluons que le retrait de
$\mathcal{O}_Y(d')$ combiner ces résultats avec Propriétés, lemme
\ref{properties-lemma-ample-power-ample}
nous concluons que $\mathcal{L}$ est ample comme désiré.

\end{proof}

\begin{lemma}

\label{lemma-ample-over-affine}

\begin{reference}
\cite[II corollaire 4.6.6]{EGA}
\end{reference}
Laisser $f : X \to S$ être un morphisme de schémas. Laisser $\mathcal{L}$ être un inversible $\mathcal{O}_X$-module. Supposons que $S$ affine. Puis $\mathcal{L}$ est $f$-relativement ample si et seulement si $\mathcal{L}$ est ample sur $X$.
\end{lemma}

\begin{proof}
Immédiate de lemme \ref{lemma-characterize-relatively-ample} et les définitions.
\end{proof}

\begin{lemma}

\label{lemma-quasi-affine-O-ample}

\begin{reference}
\cite[II proposition 5.1.6]{EGA}
\end{reference}
Laissez $f : X \to S$ être un morphisme de schémas. Alors $f$ est quasi-affiner si et seulement si $\mathcal{O}_X$ est $f$-relativement ample.
\end{lemma}

\begin{proof}
Suit de Propriétés, lemme \ref{properties-lemma-quasi-affine-O-ample} et les définitions.
\end{proof}

\begin{lemma}

\label{lemma-pullback-ample-tensor-relatively-ample}
Laissez $f : X \to Y$ être un morphisme de schémas, $\mathcal{M}$ un module $\mathcal{O}_Y$ inversible, et $\mathcal{L}$ un module $\mathcal{O}_X$ inversible.
\begin{enumerate}
\item Si $\mathcal{L}$ est $f$-ample et $\mathcal{M}$ est ample, alors $\mathcal{L} \otimes f^*\mathcal{M}^{\otimes a}$ est ample pour $a \gg 0$.
\item Si $\mathcal{M}$ est ample et $f$ quasi-affinée, alors $f^*\mathcal{M}$ est ample.
\end{enumerate}

\end{lemma}

\begin{proof}

Présumer $\mathcal{L}$ est $f$-ample et $\mathcal{M}$ ample. Par hypothèse $Y$ et $f$ sont quasi compactes (voir définition \ref{definition-relatively-ample} et Propriétés, définition \ref{properties-definition-ample}). D'où $X$ est quasi compact. Par Propriétés, lemme
\ref{properties-lemma-ample-separated} le régime $Y$ est séparé et par lemme \ref{lemma-relatively-ample-separated}
le morphisme $f$ est séparé. $X$ est séparé par Schémas, lemme \ref{schemes-lemma-separated-permanence}. Choisissez $x \in X$. Nous pouvons choisir $m \geq 1$
et $t \in \Gamma(Y, \mathcal{M}^{\otimes m})$ de telle manière que $Y_t$
est affine et $f(x) \in Y_t$Depuis $\mathcal{L}$ se limite à un faisceau inversible ample sur $f^{-1}(Y_t) = X_{f^*t}$
nous pouvons choisir $n \geq 1$ et $s \in \Gamma(X_{f^*t}, \mathcal{L}^{\otimes n})$
avec $x \in (X_{f^*t})_s$ avec $(X_{f^*t})_s$ affine. Par Propriétés, lemme \ref{properties-lemma-invert-s-sections} partie (2) dont les hypothèses sont satisfaites par ce qui précède, il existe un entier $e \geq 1$ et une section
$s' \in \Gamma(X, \mathcal{L}^{\otimes n} \otimes f^*\mathcal{M}^{\otimes em})$
qui se limite $s(f^*t)^e$ le $X_{f^*t}$. Pour tout $b > 0$
examiner la section $s'' = s'(f^*t)^b$ de
$\mathcal{L}^{\otimes n} \otimes f^*\mathcal{M}^{\otimes (e + b)m}$. Ensuite $X_{s''} = (X_{f^*t})_s$ est un ouvert affine de $X$ contenant $x$. Choix $b$ de telle manière que $n$ divisions $e + b$ nous voyons
$\mathcal{L}^{\otimes n} \otimes f^*\mathcal{M}^{\otimes (e + b)m}$
est la $n$la puissance de $\mathcal{L} \otimes f^*\mathcal{M}^{\otimes a}$
pour certains $a$ et on peut avoir n'importe quel $a$ divisible par $m$ et assez grand. Depuis $X$ est quasi-compact un nombre fini de ces ouverts affines couverture $X$. Nous concluons que pour certains $a$ suffisamment divisible et assez grande la gerbe inversible
$\mathcal{L} \otimes f^*\mathcal{M}^{\otimes a}$ est ample sur $X$. D'autre part, nous savons que $\mathcal{M}^{\otimes c}$
(et donc son retrait vers $X$) est générée à l'échelle mondiale pour tous $c \gg 0$
par Propriétés, proposition \ref{properties-proposition-characterize-ample}. Ainsi $\mathcal{L} \otimes f^*\mathcal{M}^{\otimes a + c}$ est ample (Propriétés, \ref{properties-lemma-ample-tensor-globally-generated}) pour $c \gg 0$ et (1) est prouvé.

\medskip\noindent

La partie 2 découle de l'article \ref{lemma-quasi-affine-O-ample}, Propriétés, lemme \ref{properties-lemma-ample-power-ample}, et partie (1).

\end{proof}

\begin{lemma}

\label{lemma-ample-composition}

Laisser $g : Y \to S$ et $f : X \to Y$ les morphismes des schémas. $\mathcal{M}$ être un inversible $\mathcal{O}_Y$- Module. $\mathcal{L}$ être un inversible $\mathcal{O}_X$- Module. $S$ est quasi-compact, $\mathcal{M}$ est $g$-l'échantillon, et
$\mathcal{L}$ est $f$-Ample, alors
$\mathcal{L} \otimes f^*\mathcal{M}^{\otimes a}$
est $g \circ f$- échantillon pour $a \gg 0$.

\end{lemma}

\begin{proof}
Laissez $S = \bigcup_{i = 1, \ldots, n} V_i$ être une finition de récupération ouvert affine. Par lemme \ref{lemma-characterize-relatively-ample} il suffit de prouver que $\mathcal{L} \otimes f^*\mathcal{M}^{\otimes a}$ est ample sur $(g \circ f)^{-1}(V_i)$ pour $i = 1, \ldots, n$. Ainsi, le lemme suit de lemme \ref{lemma-pullback-ample-tensor-relatively-ample}.
\end{proof}

\begin{lemma}

\label{lemma-ample-base-change}
Laissez $f : X \to S$ être un morphisme de schémas. Laissez $\mathcal{L}$ être un module $\mathcal{O}_X$ inversible. Laissez $S' \to S$ être un morphisme de schémas. Laissez $f' : X' \to S'$ être le changement de base de $f$ et dénotent $\mathcal{L}'$ le retrait de $\mathcal{L}$ vers $X'$. Si $\mathcal{L}$ est $f$-ample, alors $\mathcal{L}'$ est $f'$-ample.
\end{lemma}

\begin{proof}
Par lemme \ref{lemma-characterize-relatively-ample} il suffit de trouver un recouvrement ouvert affine $S' = \bigcup U'_i$ tel que $\mathcal{L}'$ se limite à un large filet inversible sur $(f')^{-1}(U_i')$ pour tous $i$. Nous pouvons choisir la cartographie $U'_i$ dans une affine ouvert $U_i \subset S$. Dans ce cas, le morphisme $(f')^{-1}(U'_i) \to f^{-1}(U_i)$ est un changement de base du morphisme affine $U'_i \to U_i$ (lemme \ref{lemma-base-change-affine}). Ainsi $\mathcal{L}'|_{(f')^{-1}(U'_i)}$ est ample par lemme \ref{lemma-pullback-ample-tensor-relatively-ample}.
\end{proof}

\begin{lemma}

\label{lemma-ample-permanence}
Laissez $g : Y \to S$ et $f : X \to Y$ être des morphismes de schémas. Laissez $\mathcal{L}$ être un inversible $\mathcal{O}_X$-module. Si $\mathcal{L}$ est $g \circ f$-ample et $f$ est quasi-compact\footnote{Ceci suit si $g$ est quasi-séparé par Schémas, lemme \ref{schemes-lemma-quasi-compact-permanence}.} alors $\mathcal{L}$ est $f$-ample.
\end{lemma}

\begin{proof}

Présumer $f$ est quasi compacte et $\mathcal{L}$ est $g \circ f$- A l'ampoule. $U \subset S$ être une affine ouvert et laisser $V \subset Y$ être une affinité d'ouverture avec $g(V) \subset U$. Ensuite $\mathcal{L}|_{(g \circ f)^{-1}(U)}$ est ample sur
$(g \circ f)^{-1}(U)$ par hypothèse. Depuis $f^{-1}(V) \subset (g \circ f)^{-1}(U)$ nous voyons que
$\mathcal{L}|_{f^{-1}(V)}$ est ample sur $f^{-1}(V)$ par Propriétés, lemme \ref{properties-lemma-ample-on-locally-closed}. A savoir, $f^{-1}(V) \to (g \circ f)^{-1}(U)$ est un ouverte d'immersion quasi-compact de Schémas, lemme \ref{schemes-lemma-quasi-compact-permanence}
comme $(g \circ f)^{-1}(U)$ est séparé (Propriétés, lemme \ref{properties-lemma-ample-separated}) et $f^{-1}(V)$ est quasi-compact (comme $f$ est quasi-compact). Ainsi, nous concluons que $\mathcal{L}$ est $f$-ample par lemme \ref{lemma-characterize-relatively-ample}.

\end{proof}







\section{Faisceaux très amples}
\label{section-very-ample}

\noindent
Rappelez-vous que compte tenu d'un faisceau quasi-cohérent $\mathcal{E}$ sur un schéma $S$ le {\it projective bundle} associé à $\mathcal{E}$ est le morphisme $\mathbf{P}(\mathcal{E}) \to S$, où $\mathbf{P}(\mathcal{E}) = \underline{\text{Proj}}_S(\text{Sym}(\mathcal{E}))$, voir Constructions, définition \ref{constructions-definition-projective-bundle}.

\begin{definition}

\label{definition-very-ample}

Laisser $f : X \to S$ être un morphisme de schémas. $\mathcal{L}$ être un inversible $\mathcal{O}_X$- Module. $\mathcal{L}$ est {\it relativement très ample} ou plus précisément {\it  $f$- relativement très ample}, ou {\it très ample sur $X/S$}, ou {\it  $f$- très ample} s'il y a un quasi-cohérent $\mathcal{O}_S$-module
$\mathcal{E}$ et une immersion $i : X \to \mathbf{P}(\mathcal{E})$
terminé $S$ de telle manière que
$\mathcal{L} \cong i^*\mathcal{O}_{\mathbf{P}(\mathcal{E})}(1)$.

\end{definition}



\noindent
Comme il n'y a pas d'hypothèse de quasi-compactivité dans cette définition, il n'est pas vrai en général qu'un ample faisceau relativement très inversible est un ample faisceau relativement inversible.

\begin{lemma}

\label{lemma-ample-very-ample}

\begin{reference}
\cite[II, proposition 4.6.2]{EGA}
\end{reference}
Laisser $f : X \to S$ être un morphisme de schémas. Laisser $\mathcal{L}$ être un inversible $\mathcal{O}_X$-module. Si $f$ est quasi-compact et $\mathcal{L}$ est un relativement très faisceau inversible ample, alors $\mathcal{L}$ est un relativement faisceau inversible ample.
\end{lemma}

\begin{proof}

Par définition, il existe une quasi-cohérence $\mathcal{O}_S$-module
$\mathcal{E}$ et une immersion $i : X \to \mathbf{P}(\mathcal{E})$
terminé $S$ de telle manière que
$\mathcal{L} \cong i^*\mathcal{O}_{\mathbf{P}(\mathcal{E})}(1)$. Prêt $\mathcal{A} = \text{Sym}(\mathcal{E})$Donc...
$\mathbf{P}(\mathcal{E}) = \underline{\text{Proj}}_S(\mathcal{A})$
par définition. $\mathcal{O}_S$-Algèbre $\mathcal{A}$
est équipé d'une application
$$
\psi :
\mathcal{A} \to
\bigoplus\nolimits_{n \geq 0}
\pi_*\mathcal{O}_{\mathbf{P}(\mathcal{E})}(n) \to
\bigoplus\nolimits_{n \geq 0}
f_*\mathcal{L}^{\otimes n}
$$
où la deuxième flèche utilise l'identification
$\mathcal{L} \cong i^*\mathcal{O}_{\mathbf{P}(\mathcal{E})}(1)$. Par adjonction de $f_*$ et $f^*$ nous obtenons un morphisme
$\psi : f^*\mathcal{A} \to \bigoplus_{n \geq 0}\mathcal{L}^{\otimes n}$. Nous omettons la vérification que le morphisme $r_{\mathcal{L}, \psi}$
associé à cette application est exactement l'immersion $i$. D'où le résultat de la partie (6) de lemme \ref{lemma-characterize-relatively-ample}.

\end{proof}

\noindent

Pour arriver à l'inverse correct de ce lemme nous demandons si donné un faisceau relativement inversable ample $\mathcal{L}$ il existe un entier $n \geq 1$ de telle manière que $\mathcal{L}^{\otimes n}$ est-ce relativement très ample? En général, c'est faux. Il y a plusieurs choses qui empêchent que cela soit vrai:

\begin{enumerate}

\item Même si $S$ est affine, il peut arriver qu'aucun entier fini
$n$ fonctionne parce que $X \to S$ n'est pas de type fini, voir exemple \ref{example-not-finite-type-proj}.
\item La base n'étant pas quasi-compact signifie que le résultat peut être empêché d'être vrai même avec $f$ type fini. A savoir, donné un champ $k$ il existe un régime $X_d$ de type fini sur $k$ avec un faisceau inversible ample $\mathcal{O}_{X_d}(1)$ de sorte que la plus petite puissance de tension de $\mathcal{O}_{X_d}(1)$ qui est très ample est le $d$. Voir exemple \ref{example-not-bounded}. Prendre $f$ être l'union disjointe des régimes $X_d$ cartographie à l'union disjointe des copies de $\Spec(k)$ donne un exemple.

\end{enumerate}

Pour voir notre version de l'inverse jetez un oeil à lemme \ref{lemma-finite-type-ample-very-ample} Nous ferons quelques travaux préliminaires avant de le prouver.

\begin{example}

\label{example-very-ample}

Laisser $S$ Être un plan. $\mathcal{A}$ être quasi-cohérents
$\mathcal{O}_S$-algèbre généré par $\mathcal{A}_1$ terminé $\mathcal{A}_0$. Prêt $X = \underline{\text{Proj}}_S(\mathcal{A})$. Dans ce cas
$\mathcal{O}_X(1)$ est inversible et très ample sur $X/S$. A savoir, le morphisme associé au $\mathcal{O}_S$- carte de l'algèbre
$$
\text{Sym}_{\mathcal{O}_X}^*(\mathcal{A}_1)
\longrightarrow
\mathcal{A}
$$
est une immersion fermée $X \to \mathbf{P}(\mathcal{A}_1)$ qui se retient
$\mathcal{O}_{\mathbf{P}(\mathcal{A}_1)}(1)$ à $\mathcal{O}_X(1)$, voir Constructions, lemme
\ref{constructions-lemma-surjective-generated-degree-1-map-relative-proj}.

\end{example}

\begin{example}

\label{example-not-finite-type-proj}

Laisser $k$ être un champ. Considérez la note $k$-Algèbre
$$
A = k[U, V, Z_1, Z_2, Z_3, \ldots]/I
\quad
\text{with}
\quad
I = (U^2 - Z_1^2, U^4 - Z_2^2, U^6 - Z_3^2, \ldots)
$$
avec classement donné par $\deg(U) = \deg(V) = \deg(Z_1) = 1$
et $\deg(Z_d) = d$. Notez que $X = \text{Proj}(A)$ est couvert par: $D_{+}(U)$ et
$D_{+}(V)$. D'où les gerbes $\mathcal{O}_X(n)$ sont tous inversibles et isomorphiques à $\mathcal{O}_X(1)^{\otimes n}$. En particulier $\mathcal{O}_X(1)$ est ample et $f$-ample pour le morphisme $f : X \to \Spec(k)$. Nous prétendons qu'aucune puissance de $\mathcal{O}_X(1)$ est $f$- relativement très ample. C'est-à-dire, il est facile de voir que $\Gamma(X, \mathcal{O}_X(n))$
est le degré $n$ la somme de l'algèbre $A$. Par conséquent, si $\mathcal{O}_X(n)$
étaient très amples, alors $X$ serait un sous-système fermé d'un projet d'espace sur $k$ et donc de type fini sur $k$. D'un autre côté
$D_{+}(V)$ est le spectre de
$k[t, t_1, t_2, \ldots]/(t^2 - t_1^2, t^4 - t_2^2, t^6 - t_3^2, \ldots)$
qui n'est pas de type fini sur $k$.

\end{example}

\begin{example}

\label{example-not-bounded}

Laisser $k$ être un champ infini. $\lambda_1, \lambda_2, \lambda_3, \ldots$
être en couple des éléments distincts de $k^*$. (Ceci n'est pas strictement nécessaire, et en fait l'exemple fonctionne parfaitement même si tous $\lambda_i$
sont égales à $1$.) Considérer la classe $k$-Algèbre
$$
A_d = k[U, V, Z]/I_d
\quad
\text{with}
\quad
I_d = (Z^2 - \prod\nolimits_{i = 1}^{2d} (U - \lambda_i V)).
$$
avec classement donné par $\deg(U) = \deg(V) = 1$ et $\deg(Z) = d$. Ensuite $X_d = \text{Proj}(A_d)$ a faisceau inversible ample
$\mathcal{O}_{X_d}(1)$. Nous prétendons que si $\mathcal{O}_{X_d}(n)$
est très ample, alors $n \geq d$. La raison de cela est que $Z$
a un diplôme $d$, et donc $\Gamma(X_d, \mathcal{O}_{X_d}(n)) =
k[U, V]_n$ pour $n < d$. Détails omis.

\end{example}

\begin{lemma}

\label{lemma-relatively-very-ample-separated}
Laissez $f : X \to S$ être un morphisme de schémas. Laissez $\mathcal{L}$ être une feuille inversible sur $X$. Si $\mathcal{L}$ est relativement très ample sur $X/S$ alors $f$ est séparé.
\end{lemma}

\begin{proof}
Se séparer est local sur la base (voir Schémas, section \ref{schemes-section-separation-axioms}). Une immersion est séparée (voir Schémas, lemme \ref{schemes-lemma-immersions-monomorphisms}). Ainsi la lemme suit depuis localement $X$ a une immersion dans le spectre homogène d'un anneau gradé qui est séparé, voir Constructions, lemme \ref{constructions-lemma-proj-separated}.
\end{proof}

\begin{lemma}

\label{lemma-relatively-very-ample}

Laisser $f : X \to S$ être un morphisme de schémas. $\mathcal{L}$ être une gerbe inversible sur $X$Supposons que $f$ est quasi-compact. Les conditions suivantes sont équivalentes

\begin{enumerate}

\item  $\mathcal{L}$ est relativement très ample sur $X/S$,
\item il existe un recouvrement ouvert $S = \bigcup V_j$ de telle manière que $\mathcal{L}|_{f^{-1}(V_j)}$ est relativement très ample sur $f^{-1}(V_j)/V_j$ pour tous $j$,
\item il existe un faisceau quasi-cohérent de graded
$\mathcal{O}_S$-Algèbres $\mathcal{A}$ générés en degré
$1$ terminé $\mathcal{O}_S$ et une application de classement $\mathcal{O}_X$-Algèbres
$\psi : f^*\mathcal{A} \to \bigoplus_{n \geq 0} \mathcal{L}^{\otimes n}$
de telle manière que $f^*\mathcal{A}_1 \to \mathcal{L}$ est surjectif et le morphisme associé
$r_{\mathcal{L}, \psi} : X \to \underline{\text{Proj}}_S(\mathcal{A})$
est une immersion, et
\item  $f$ est quasi-séparé, l'application canonique
$\psi : f^*f_*\mathcal{L} \to \mathcal{L}$ est surjectif, et l'application associée $r_{\mathcal{L}, \psi} : X \to \mathbf{P}(f_*\mathcal{L})$
c'est une immersion.

\end{enumerate}

\end{lemma}

\begin{proof}
Il est clair que (1) implique (2). Il est également clair que (4) implique (1); l'hypothèse de quasi-séparation dans (4) est utilisée pour garantir que $f_*\mathcal{L}$ est quasi-cohérent via Schémas, lemme \ref{schemes-lemma-push-forward-quasi-coherent}.

\medskip\noindent

Supposons (2). Nous prouverons (4). $S = \bigcup V_j$ être un recouvrement ouvert comme dans (2). $X_j = f^{-1}(V_j)$ et $f_j : X_j \to V_j$ la restriction de $f$. Nous voyons que $f$ est séparé par lemme \ref{lemma-relatively-very-ample-separated} (comme être séparé est local sur la base). $\mathcal{O}_{V_j}$-module $\mathcal{E}_j$ et une immersion
$i_j : X_j \to \mathbf{P}(\mathcal{E}_j)$ avec
$\mathcal{L}|_{X_j} \cong i_j^*\mathcal{O}_{\mathbf{P}(\mathcal{E}_j)}(1)$Le morphisme $i_j$ correspond à une surjection
$f_j^*\mathcal{E}_j \to \mathcal{L}|_{X_j}$, voir Constructions, section \ref{constructions-section-projective-bundle}. Cette application est associée à une application $\mathcal{E}_j \to f_*\mathcal{L}|_{V_j}$
telle que la composition
$$
f_j^*\mathcal{E}_j \to (f^*f_*\mathcal{L})|_{X_j} \to \mathcal{L}|_{X_j}
$$
Nous concluons que
$\psi : f^*f_*\mathcal{L} \to \mathcal{L}$ c'est surjectif.
$r_{\mathcal{L}, \psi} : X \to \mathbf{P}(f_*\mathcal{L})$ morphisme associé. Nous devons encore montrer que $r_{\mathcal{L}, \psi}$
est une immersion; nous exhortons le lecteur à le prouver par eux-mêmes. $\mathcal{O}_{V_j}$-module carte $\mathcal{E}_j \to f_*\mathcal{L}|_{V_j}$
détermine un homomorphisme sur les algèbres symétriques, qui à son tour définit un morphisme
$$
\mathbf{P}(f_*\mathcal{L}|_{V_j}) \supset U_j \longrightarrow
\mathbf{P}(\mathcal{E}_j)
$$
où $U_j$ est le sous-système ouvert de Constructions, lemme \ref{constructions-lemma-morphism-relative-proj}. La compatibilité des $\psi$ avec $\mathcal{E}_j \to f_*\mathcal{L}|_{V_j}$
montre que $r_{\mathcal{L}, \psi}(X_j) \subset U_j$ et qu'il y a une factorisation
$$
\xymatrix{
X_j \ar[r]^-{r_{\mathcal{L}, \psi}} & U_j \ar[r] &
\mathbf{P}(\mathcal{E}_j)
}
$$
Nous omettons la vérification. $r_{\mathcal{L}, \psi}$ c'est une immersion.

\medskip\noindent

A ce stade, nous voyons que (1), (2) et (4) sont équivalents. De toute évidence (4) implique (3). Supposons que (3). Nous prouverons (1). $\mathcal{A}$ être un faisceau quasi-cohérent de graded $\mathcal{O}_S$-algèbres générées en degré $1$ terminé $\mathcal{O}_S$. Considérez l'application des grades $\mathcal{O}_S$-Algèbres $\text{Sym}(\mathcal{A}_1) \to \mathcal{A}$. Ceci est surjectif par hypothèse et provoque donc une immersion fermée
$$
\underline{\text{Proj}}_S(\mathcal{A})
\longrightarrow
\mathbf{P}(\mathcal{A}_1)
$$
qui se retient $\mathcal{O}(1)$ à $\mathcal{O}(1)$, voir Constructions, lemme
\ref{constructions-lemma-surjective-generated-degree-1-map-relative-proj}. Il est donc clair que (3) implique (1).

\end{proof}

\begin{lemma}

\label{lemma-very-ample-base-change}
Laissez $f : X \to S$ être un morphisme de schémas.Soit $\mathcal{L}$ un module $\mathcal{O}_X$ inversible.Soit $S' \to S$ un morphisme de schémas.Soit $f' : X' \to S'$ le changement de base de $f$ et dénote $\mathcal{L}'$ le retrait de $\mathcal{L}$ vers $X'$. Si $\mathcal{L}$ est $f$-tres ample, alors $\mathcal{L}'$ est $f'$-tres ample.
\end{lemma}

\begin{proof}

Par définition \ref{definition-very-ample} il existe une quasi-cohérence $\mathcal{O}_S$-module $\mathcal{E}$ et une immersion
$i : X \to \mathbf{P}(\mathcal{E})$ terminé $S$ de telle manière que
$\mathcal{L} \cong i^*\mathcal{O}_{\mathbf{P}(\mathcal{E})}(1)$. Le changement de base de $\mathbf{P}(\mathcal{E})$ à $S'$ est le faisceau projectif associé au retrait $\mathcal{E}'$
de $\mathcal{E}$ et le retrait de
$\mathcal{O}_{\mathbf{P}(\mathcal{E})}(1)$
est
$\mathcal{O}_{\mathbf{P}(\mathcal{E}')}(1)$, voir Constructions, lemme \ref{constructions-lemma-relative-proj-base-change}. Enfin, le changement de base d'une immersion est une immersion (Schémas, lemme \ref{schemes-lemma-base-change-immersion}).

\end{proof}







\section{Faisceaux amples et très amples relativement aux morphismes de type fini}
\label{section-ample-finite-type}

\noindent
En fait, la plupart du matériel de cette section concerne la notion d'un (quasi-)morphisme projectif que nous n'avons pas encore défini.

\begin{lemma}

\label{lemma-very-ample-finite-type-over-affine}

Laisser $f : X \to S$ être un morphisme de schémas. $\mathcal{L}$ être une gerbe inversible sur $X$. Supposons que

\begin{enumerate}

\item la gerbe inversible $\mathcal{L}$ est très ample sur $X/S$,
\item le morphisme $X \to S$ est de type fini, et
\item  $S$ C'est affiné.

\end{enumerate}

Alors il existe un $n \geq 0$ et une immersion
$i : X \to \mathbf{P}^n_S$ terminé $S$ de telle manière que
$\mathcal{L} \cong i^*\mathcal{O}_{\mathbf{P}^n_S}(1)$.

\end{lemma}

\begin{proof}

Supposons (1), (2) et (3). $S = \Spec(R)$ pour un peu de bague $R$. Condition (1) signifie par définition qu'il existe une quasi-cohérence $\mathcal{O}_S$-module
$\mathcal{E}$ et une immersion $\alpha : X \to \mathbf{P}(\mathcal{E})$
de telle manière que $\mathcal{L} = \alpha^*\mathcal{O}_{\mathbf{P}(\mathcal{E})}(1)$. Écrire $\mathcal{E} = \widetilde{M}$ pour certains $R$-module $M$. Ainsi nous avons
$$
\mathbf{P}(\mathcal{E}) = \text{Proj}(\text{Sym}_R(M)).
$$
Depuis $\alpha$ est une immersion, et depuis la topologie de
$\text{Proj}(\text{Sym}_R(M))$ est généré par la norme s'ouvre $D_{+}(f)$, $f \in \text{Sym}_R^d(M)$, $d \geq 1$, nous pouvons trouver pour chacun $x \in X$ et
$f \in \text{Sym}_R^d(M)$, $d \geq 1$Avec $\alpha(x) \in D_{+}(f)$
de telle manière que
$$
\alpha|_{\alpha^{-1}(D_{+}(f))} : \alpha^{-1}(D_{+}(f)) \to D_{+}(f)
$$
est une immersion fermée. $X$ est quasi-compact. Ainsi, nous pouvons trouver une collection finie d'éléments
$f_j \in \text{Sym}_R^{d_j}(M)$, $d_j \geq 1$
tel que pour chacun $f = f_j$ l'application affichée ci-dessus est une immersion fermée et telle que $\alpha(X) \subset \bigcup D_{+}(f_j)$. Écrire $U_j = \alpha^{-1}(D_{+}(f_j))$. Notez que $U_j$ est affine en tant que sous-système fermé du régime affine $D_{+}(f_j)$. Écrire $U_j = \Spec(A_j)$. La condition (2) implique également que
$A_j$ est de type fini sur $R$, voir lemme \ref{lemma-locally-finite-type-characterize}. Choisissez finiment beaucoup $x_{j, k} \in A_j$ qui génèrent $A_j$ en tant que $R$- de l'algèbre. $\alpha|_{U_j}$ est une immersion fermée nous voyons que $x_{j, k}$ est l'image d'un élément
$$
f_{j, k}/f_j^{e_{j, k}} \in \text{Sym}_R(M)_{(f_j)}
=
\Gamma(D_{+}(f_j), \mathcal{O}_{\text{Proj}(\text{Sym}_R(M))}).
$$
Enfin, choisissez $n \geq 1$ et éléments $y_0, \ldots, y_n \in M$ de telle sorte que chacun des polynômes $f_j, f_{j, k} \in \text{Sym}_R(M)$ est un polynôme dans les éléments $y_t$ avec des coefficients en $R$. Considérez l'homomorphisme d'anneaux
$$
\psi : R[Y_0, \ldots, Y_n] \longrightarrow \text{Sym}_R(M),
\quad Y_i \longmapsto y_i.
$$
Dénotation $F_j$, $F_{j, k}$ les éléments de $R[Y_0, \ldots, Y_n]$ de telle manière que $\psi(F_j) = f_j$ et $\psi(F_{j, k}) = f_{j, k}$. Par Constructions, lemme \ref{constructions-lemma-morphism-proj}
nous obtenons un sous-système ouvert
$$
U(\psi) \subset \text{Proj}(\text{Sym}_R(M))
$$
et un morphisme
$r_\psi : U(\psi) \to \mathbf{P}^n_R$. Ce morphisme satisfait $r_\psi^{-1}(D_{+}(F_j)) = D_{+}(f_j)$, et donc nous voyons que $\alpha(X) \subset U(\psi)$. De plus, il est clair que
$$
i = r_\psi \circ \alpha : X \longrightarrow \mathbf{P}^n_R
$$
est encore une immersion depuis
$i^\sharp(F_{j, k}/F_j^{e_{j, k}}) = x_{j, k} \in
A_j = \Gamma(U_j, \mathcal{O}_X)$
par la construction. De plus, le morphisme $r_\psi$ est équipé d'une application
$\theta : r_\psi^*\mathcal{O}_{\mathbf{P}^n_R}(1)
\to \mathcal{O}_{\text{Proj}(\text{Sym}_R(M))}(1)|_{U(\psi)}$
qui est un isomorphisme dans ce cas (pour la construction $\theta$
Voir lemme ci-dessus; quelques détails omis). Depuis l'application originale $\alpha$ a été supposé avoir la propriété qui
$\mathcal{L} = \alpha^*\mathcal{O}_{\text{Proj}(\text{Sym}_R(M))}(1)$
Nous avons gagné.

\end{proof}

\begin{lemma}

\label{lemma-quasi-affine-finite-type-over-S}
Laissez $\pi : X \to S$ être un morphisme de schémas. Supposons que $X$ est quasi-affiné et que $\pi$ est localement de type fini. Ensuite il existe $n \geq 0$ et une immersion $i : X \to \mathbf{A}^n_S$ sur $S$.
\end{lemma}

\begin{proof}

Laisser $A = \Gamma(X, \mathcal{O}_X)$. Par hypothèse $X$ est quasi-compact et est identifié avec un sous-système ouvert de $\Spec(A)$, voir Propriétés, laisse \ref{properties-lemma-quasi-affine}. En outre, l'ensemble des ouvertures $X_f$, pour ceux $f \in A$ de telle manière que $X_f$ est affine, forme une base pour la topologie de $X$, voir la démonstration de Propriétés, lemme \ref{properties-lemma-quasi-affine}. Nous pouvons donc trouver un nombre fini de $f_j \in A$, $j = 1, \ldots, m$ de telle manière que
$X = \bigcup X_{f_j}$, et de telle sorte que $\pi(X_{f_j}) \subset V_j$ pour un ouvert affine $V_j \subset S$. Par lemme \ref{lemma-locally-finite-type-characterize}
les homomorphismes d'anneaux $\mathcal{O}(V_j) \to \mathcal{O}(X_{f_j}) = A_{f_j}$
sont de type fini. Ainsi, nous pouvons choisir $a_1, \ldots, a_N \in A$ de sorte que les éléments $a_1, \ldots, a_N, 1/f_j$ générer
$A_{f_j}$ terminé $\mathcal{O}(V_j)$ pour chaque $j$. Prenez $n = m + N$ et laisser
$$
i : X \longrightarrow \mathbf{A}^n_S
$$
être le morphisme donné par les sections mondiales
$f_1, \ldots, f_m, a_1, \ldots, a_N$ de la structure de l'écharpe $X$. Laisse $D(x_j) \subset \mathbf{A}^n_S$ être le sous-système ouvert où le
$j$la fonction de coordonnées n'est pas zéro. $1 \leq j \leq m$ Nous avons $i^{-1}(D(x_j)) = X_{f_j}$ et le morphisme induit $X_{f_j} \to D(x_j)$ facteurs à travers l'affine de l'ouverture $\Spec(\mathcal{O}(V_j)[x_1, \ldots, x_n, 1/x_j])$
de $D(x_j)$. Depuis l'homomorphisme d'anneaux
$\mathcal{O}(V_j)[x_1, \ldots, x_n, 1/x_j] \to A_{f_j}$ est surjectif par la construction nous concluons que $i^{-1}(D(x_j)) \to D(x_j)$
est une immersion comme désiré.

\end{proof}

\begin{lemma}

\label{lemma-quasi-projective-finite-type-over-S}

Laisser $f : X \to S$ être un morphisme de schémas. $\mathcal{L}$ être une gerbe inversible sur $X$. Supposons que

\begin{enumerate}

\item la gerbe inversible $\mathcal{L}$ est ample sur $X$,
\item le morphisme $X \to S$ est localement de type fini.

\end{enumerate}

Puis il existe un $d_0 \geq 1$ tel que pour chaque $d \geq d_0$
il existe un $n \geq 0$ et une immersion
$i : X \to \mathbf{P}^n_S$ terminé $S$ de telle manière que
$\mathcal{L}^{\otimes d} \cong i^*\mathcal{O}_{\mathbf{P}^n_S}(1)$.

\end{lemma}

\begin{proof}

Laisser
$A = \Gamma_*(X, \mathcal{L}) =
\bigoplus_{d \geq 0} \Gamma(X, \mathcal{L}^{\otimes d})$. Par Propriétés, proposition \ref{properties-proposition-characterize-ample}
l'ensemble des ouverts affinés $X_a$ avec $a \in A_{+}$ homogène forme une base pour la topologie de $X$. Par conséquent, nous pouvons trouver de nombreux éléments finis $a_0, \ldots, a_n \in A_{+}$ de telle manière que

\begin{enumerate}

\item Nous avons $X = \bigcup_{i = 0, \ldots, n} X_{a_i}$,
\item chaque $X_{a_i}$ est affine, et
\item chaque $X_{a_i}$ cartes dans un ouvert affine $V_i \subset S$.

\end{enumerate}

Par moi \ref{lemma-locally-finite-type-characterize}
nous voyons que les homomorphismes d'anneaux
$\mathcal{O}_S(V_i) \to \mathcal{O}_X(X_{a_i})$ sont de type fini. Ainsi, nous pouvons trouver de nombreux éléments finis $f_{ij} \in \mathcal{O}_X(X_{a_i})$, $j = 1, \ldots, n_i$
qui génèrent $\mathcal{O}_X(X_{a_i})$ comme une $\mathcal{O}_S(V_i)$- Algèbre. Par Propriétés, lemme \ref{properties-lemma-invert-s-sections}
nous pouvons écrire chacun
$f_{ij}$ comme $a_{ij}/a_i^{e_{ij}}$ pour certains
$a_{ij} \in A_{+}$ homogène. Laisser $N$ être un entier positif qui est un multiple commun de tous les degrés des éléments
$a_i$, $a_{ij}$. Examiner les éléments
$$
a_i^{N/\deg(a_i)}, \ a_{ij}a_i^{(N/\deg(a_i)) - e_{ij}} \in A_N.
$$
Par la construction, ils génèrent la gerbe inversible
$\mathcal{L}^{\otimes N}$ terminé $X$. C'est pourquoi ils donnent naissance à un morphisme
$$
j : X \longrightarrow
\mathbf{P}_S^{m}
\quad
\text{with } m = n + \sum n_i
$$
terminé $S$, voir Constructions, lemme \ref{constructions-lemma-projective-space}
et définition \ref{constructions-definition-projective-space}. En outre, $j^*\mathcal{O}_{\mathbf{P}_S}(1) = \mathcal{L}^{\otimes N}$. Nous nominons les coordonnées homogènes $T_0, \ldots, T_n, T_{ij}$
au lieu de $T_0, \ldots, T_m$. Pour $i = 0, \ldots, n$ Nous avons $j^{-1}(D_{+}(T_i)) = X_{a_i}$. En outre, retirer l'élément $T_{ij}/T_i$ par $j^\sharp$ nous obtenons l'élément $f_{ij} \in \mathcal{O}_X(X_{a_i})$D'où le morphisme $j$ limités à $X_{a_i}$
donne une immersion fermée $X_{a_i}$ dans l'affinité de l'ouverture
$D_{+}(T_i) \cap \mathbf{P}^m_{V_i}$ de $\mathbf{P}^N_S$. Par conséquent, nous concluons que le morphisme $j$ est une immersion. Cela implique que le lemme tient pour certains $d$ et $n$ ce qui est suffisant dans pratiquement toutes les applications.

\medskip\noindent

Cela prouve que pour un seul $d_2 \geq 1$
(à savoir: $d_2 = N$ ci-dessus), certains $m \geq 0$ il existe une certaine immersion $j : X \to \mathbf{P}^m_S$ données par les sections mondiales
$s'_0, \ldots, s'_m \in \Gamma(X, \mathcal{L}^{\otimes d_2})$. Par Propriétés, proposition \ref{properties-proposition-characterize-ample}
Nous savons qu'il existe un entier
$d_1$ de telle manière que $\mathcal{L}^{\otimes d}$ est générée globalement pour tous $d \geq d_1$. Prêt $d_0 = d_1 + d_2$. Nous prétendons que le lemme tient avec cette valeur de $d_0$. A savoir, donné un entier $d \geq d_0$ nous pouvons choisir $s''_1, \ldots, s''_t
\in \Gamma(X, \mathcal{L}^{\otimes d - d_2})$ qui génèrent
$\mathcal{L}^{\otimes d - d_2}$ terminé $X$. Prêt $k = (m + 1)t$ et indiquer $s_0, \ldots, s_k$ la collecte des sections
$s'_\alpha s''_\beta$, $\alpha = 0, \ldots, m$,
$\beta = 1, \ldots, t$. Ils génèrent $\mathcal{L}^{\otimes d}$
terminé $X$ et donc définir un morphisme
$$
i : X \longrightarrow \mathbf{P}^{k - 1}_S
$$
de telle manière que $i^*\mathcal{O}_{\mathbf{P}^n_S}(1) \cong \mathcal{L}^{\otimes d}$. Pour voir que $i$ est une immersion, observez que $i$ est la composition
$$
X \longrightarrow \mathbf{P}^m_S \times_S \mathbf{P}^{t - 1}_S
\longrightarrow \mathbf{P}^{k - 1}_S
$$
où se trouve le premier morphisme $(j, j')$ avec $j'$ données par
$s''_1, \ldots, s''_t$ et le second morphisme est l'intégration Segre (Constructions, lemme \ref{constructions-lemma-segre-embedding}Depuis $j$ est une immersion, ainsi est $(j, j')$
(Appliquer \ref{lemma-immersion-permanence}
à $X \to \mathbf{P}^m_S \times_S \mathbf{P}^{t - 1}_S
\to \mathbf{P}^m_S$). Ainsi $i$ est une composition d'immersions et donc une immersion (Schémas, lemme \ref{schemes-lemma-composition-immersion}).

\end{proof}

\begin{lemma}

\label{lemma-finite-type-over-affine-ample-very-ample}

Laisser $f : X \to S$ être un morphisme de schémas. $\mathcal{L}$ être un inversible $\mathcal{O}_X$- Module. $S$ affine et $f$ de type fini. Les valeurs suivantes sont équivalentes

\begin{enumerate}

\item  $\mathcal{L}$ est ample sur $X$,
\item  $\mathcal{L}$ est $f$- l'échantillon,
\item  $\mathcal{L}^{\otimes d}$ est $f$- très ample pour certains $d \geq 1$,
\item  $\mathcal{L}^{\otimes d}$ est $f$- très ample pour tous $d \gg 1$,
\item pour certains $d \geq 1$ il existe $n \geq 1$ et une immersion
$i : X \to \mathbf{P}^n_S$ de telle manière que
$\mathcal{L}^{\otimes d} \cong i^*\mathcal{O}_{\mathbf{P}^n_S}(1)$,
\item pour tous $d \gg 1$ il existe $n \geq 1$ et une immersion
$i : X \to \mathbf{P}^n_S$ de telle manière que
$\mathcal{L}^{\otimes d} \cong i^*\mathcal{O}_{\mathbf{P}^n_S}(1)$.

\end{enumerate}

\end{lemma}

\begin{proof}
L'équivalence de (1) et (2) est lemme \ref{lemma-ample-over-affine}. L'implication (2) $\Rightarrow$ (6) est lemme \ref{lemma-quasi-projective-finite-type-over-S}. Trivialement (6) implique (5). Comme $\mathbf{P}^n_S$ est un bundle projectif sur $S$ (voir Constructions, lemme \ref{constructions-lemma-projective-space-bundle}) nous voyons que (5) implique (3) et (6) implique (4) de la définition d'un gabarit relativement très large. Trivialement (4) implique (3). Pour terminer, nous devons montrer que (3) implique (2) qui découle de lemme \ref{lemma-ample-very-ample} et lemme \ref{lemma-ample-power-ample}.
\end{proof}

\begin{lemma}

\label{lemma-finite-type-ample-very-ample}

Laisser $f : X \to S$ être un morphisme de schémas. $\mathcal{L}$ être un inversible $\mathcal{O}_X$- Module. $S$ quasi-compact et $f$ de type fini. Les valeurs suivantes sont équivalentes

\begin{enumerate}

\item  $\mathcal{L}$ est $f$- l'échantillon,
\item  $\mathcal{L}^{\otimes d}$ est $f$- très ample pour certains $d \geq 1$,
\item  $\mathcal{L}^{\otimes d}$ est $f$- très ample pour tous $d \gg 1$.

\end{enumerate}

\end{lemma}

\begin{proof}
Trivialement (3) implique (2). lemme \ref{lemma-ample-very-ample} garantit que (2) implique (1) puisque un morphisme de type fini est quasi-compact par définition. Supposons que $\mathcal{L}$ est $f$-ample. Choisissez une finie affine récupérer ouvert $S = V_1 \cup \ldots \cup V_m$. Écrivez $X_i = f^{-1}(V_i)$. Par lemme \ref{lemma-finite-type-over-affine-ample-very-ample} ci-dessus nous voyons qu'il existe un $d_0$ de telle sorte que $\mathcal{L}^{\otimes d}$ est relativement très large sur $X_i/V_i$ pour tous $d \geq d_0$. Par conséquent nous concluons (1) implique (3) par lemme \ref{lemma-relatively-very-ample}.
\end{proof}

\noindent
Les deux lemmes suivantes fournissent les caractérisations les plus utilisées et les plus utiles de gerbes relativement très amples et relativement amples inversibles lorsque le morphisme est de type fini.

\begin{lemma}

\label{lemma-characterize-very-ample-on-finite-type}

Laisser $f : X \to S$ être un morphisme de schémas. $\mathcal{L}$ être une gerbe inversible sur $X$Supposons que $f$ est de type fini. Les conditions suivantes sont équivalentes:

\begin{enumerate}

\item  $\mathcal{L}$ est $f$- relativement très ample, et
\item il existe un recouvrement ouvert $S = \bigcup V_j$, pour chaque $j$ un nombre entier $n_j$, et les immersions
$$
i_j :
X_j = f^{-1}(V_j) = V_j \times_S X
\longrightarrow
\mathbf{P}^{n_j}_{V_j}
$$
terminé $V_j$ de telle manière que
$\mathcal{L}|_{X_j} \cong i_j^*\mathcal{O}_{\mathbf{P}^{n_j}_{V_j}}(1)$.

\end{enumerate}

\end{lemma}

\begin{proof}
Nous voyons que (1) implique (2) en prenant une affine de récupération de $S$ et en appliquant lemme \ref{lemma-very-ample-finite-type-over-affine} à chacune des restrictions de $f$ et $\mathcal{L}$. Nous voyons que (2) implique (1) par lemme \ref{lemma-relatively-very-ample}.
\end{proof}

\begin{lemma}

\label{lemma-characterize-ample-on-finite-type}
Laissez $f : X \to S$ être un morphisme de schémas. Laissez $\mathcal{L}$ être un filet inversible sur $X$. Supposons que $f$ est de type fini. Les conditions suivantes sont équivalentes:
\begin{enumerate}
\item $\mathcal{L}$ est $f$-relativement ample, et
\item il existe un recouvrement ouvert $S = \bigcup V_j$, pour chaque $j$ un entier $d_j \geq 1$, $n_j \geq 0$, et immersions $$
i_j :
X_j = f^{-1}(V_j) = V_j \times_S X
\longrightarrow
\mathbf{P}^{n_j}_{V_j}
$$ sur $V_j$ tel que $\mathcal{L}^{\otimes d_j}|_{X_j} \cong
i_j^*\mathcal{O}_{\mathbf{P}^{n_j}_{V_j}}(1)$.
\end{enumerate}

\end{lemma}

\begin{proof}
Nous voyons que (1) implique (2) en prenant une affine de récupération de $S$ et en appliquant lemme \ref{lemma-finite-type-over-affine-ample-very-ample} à chacune des restrictions de $f$ et $\mathcal{L}$. Nous voyons que (2) implique (1) par lemme \ref{lemma-characterize-relatively-ample}.
\end{proof}

\begin{lemma}

\label{lemma-invertible-add-enough-ample-very-ample}
Laisser $f : X \to S$ être un morphisme de schémas. Laisser $\mathcal{N}$, $\mathcal{L}$ être inversible $\mathcal{O}_X$-modules. Supposons que $S$ est quasi-compact, $f$ est de type fini, et $\mathcal{L}$ est $f$-ample. Puis $\mathcal{N} \otimes_{\mathcal{O}_X} \mathcal{L}^{\otimes d}$ est $f$-tres large pour tous les $d \gg 1$.
\end{lemma}

\begin{proof}

Par moi \ref{lemma-characterize-very-ample-on-finite-type}
nous réduisons à l'affaire $S$ est affine. \ref{lemma-finite-type-over-affine-ample-very-ample} et Propriétés, proposition \ref{properties-proposition-characterize-ample}
On peut trouver un entier $d_0$ de telle manière que
$\mathcal{N} \otimes \mathcal{L}^{\otimes d_0}$
est générée globalement. Choisissez des sections globales
$s_0, \ldots, s_n$ de $\mathcal{N} \otimes \mathcal{L}^{\otimes d_0}$
qui le génèrent. Ceci détermine un morphisme
$j : X \to \mathbf{P}^n_S$ terminé $S$. Par lemme \ref{lemma-finite-type-over-affine-ample-very-ample}
nous pouvons aussi choisir un entier $d_1$ tel que pour tous $d \geq d_1$
il existe des sections $t_{d, 0}, \ldots, t_{d, n(d)}$
de $\mathcal{L}^{\otimes d}$ qui le génèrent et définissent une immersion
$$
j_d = \varphi_{\mathcal{L}^{\otimes d}, t_{d, 0}, \ldots, t_{d, n(d)}} :
X
\longrightarrow
\mathbf{P}^{n(d)}_S
$$
terminé $S$. Puis pour $d \geq d_0 + d_1$ nous pouvons considérer le morphisme
$$
\varphi_{\mathcal{N} \otimes \mathcal{L}^{\otimes d},
s_j \otimes t_{d - d_0, i}} :
X
\longrightarrow
\mathbf{P}^{(n + 1)(n(d - d_0) + 1) - 1}_S
$$
Ce morphisme est une immersion comme il est la composition
$$
X \to \mathbf{P}^n_S \times_S \mathbf{P}^{n(d - d_0)}_S
\to \mathbf{P}^{(n + 1)(n(d - d_0) + 1) - 1}_S
$$
où se trouve le premier morphisme $(j, j_{d - d_0})$ et le second est le Segre incorporant (Constructions, lemme \ref{constructions-lemma-segre-embedding}Depuis $j$ est une immersion, ainsi est $(j, j_{d - d_0})$
(Appliquer \ref{lemma-immersion-permanence}). Nous avons une composition d'immersions et donc une immersion (Schémas, lemme \ref{schemes-lemma-composition-immersion}).

\end{proof}










\section{Morphismes quasi-projectifs}
\label{section-quasi-projective}

\noindent
La discussion dans la section précédente suggère les définitions suivantes. Nous prenons notre définition de quasi-projecteur de \cite{EGA}. La version avec la lettre ``H'' est la définition dans \cite{H}.

\begin{definition}

\label{definition-quasi-projective}

\begin{reference}

\cite[II, définition 5.3.1]{EGA} et \cite[page 103]{H}

\end{reference}

Laisser $f : X \to S$ être un morphisme de schémas.

\begin{enumerate}

\item Nous disons $f$ est {\it quasi-projectif} si $f$ est de type fini et il existe un $f$-relatif amplement inversible $\mathcal{O}_X$- Module.
\item Nous disons $f$ est {\it H-quasi-projective} s'il existe une immersion quasi-compacte $X \to \mathbf{P}^n_S$ terminé $S$ pour certains
$n$.\footnote{Ce n'est pas exactement le même que la définition dans Hartshorne. À savoir, la définition dans Hartshorne (8e impression corrigée, 1997) est que
$f$ devrait être la composition d'un ouverte d'immersion suivi d'un projectif H-morphisme (voir définition \ref{definition-projective}), qui n'implique pas
$f$ C'est presque compact. \ref{lemma-H-quasi-projective-open-H-projective} pour l'implication dans l'autre sens.}
\item Nous disons $f$ est {\it localement quasi-projective} s'il existe une ouverture de recouvrement $S = \bigcup V_j$ de telle sorte que chacun $f^{-1}(V_j) \to V_j$
est quasi-projectif.

\end{enumerate}

\end{definition}

\noindent
Comme cette définition suggère que la propriété d'être quasi-projectif n'est pas locale sur $S$. À un stade ultérieur, nous pourrons en dire plus sur la catégorie des schémas quasi-projectifs, voir Compléments sur les morphismes, section \ref{more-morphisms-section-quasi-projective}.

\begin{lemma}

\label{lemma-base-change-quasi-projective}

Un changement de base d'un morphisme quasi-projectif est quasi-projectif.

\end{lemma}

\begin{proof}
Ceci découle de lemmes \ref{lemma-base-change-finite-type} et \ref{lemma-ample-base-change}.
\end{proof}

\begin{lemma}

\label{lemma-composition-quasi-projective}
Laissez $f : X \to Y$ et $g : Y \to S$ être des morphismes de schémas. Si $S$ est quasi-compact et $f$ et $g$ sont quasi-projectifs, alors $g \circ f$ est quasi-projectif.
\end{lemma}

\begin{proof}
Ceci découle de lemmes \ref{lemma-composition-finite-type} et \ref{lemma-ample-composition}.
\end{proof}

\begin{lemma}

\label{lemma-quasi-projective-properties}
Laissez $f : X \to S$ être un morphisme de schémas. Si $f$ est quasi-projectif, ou H-quasi-projectif ou localement quasi-projectif, alors $f$ est séparé de type fini.
\end{lemma}

\begin{proof}
Démonstration omise.
\end{proof}

\begin{lemma}

\label{lemma-H-quasi-projective-quasi-projective}

Un H-morphisme quasi-projectif est quasi-projectif.

\end{lemma}

\begin{proof}
Démonstration omise.
\end{proof}

\begin{lemma}

\label{lemma-characterize-locally-quasi-projective}
Laissez $f : X \to S$ être un morphisme de schémas. Les conditions suivantes sont équivalentes:
\begin{enumerate}
\item Le morphisme $f$ est localement quasi-projectif.
\item Il existe un recouvrement ouvert $S = \bigcup V_j$ de telle sorte que chaque $f^{-1}(V_j) \to V_j$ est H-quasi-projectif.
\end{enumerate}

\end{lemma}

\begin{proof}

Par moi \ref{lemma-H-quasi-projective-quasi-projective}
nous voyons que (2) implique (1). Supposons que (1). $S$ et donc nous pouvons supposer $S$ est affine,
$X$ de type fini sur $S$ et
$\mathcal{L}$ est un faisceau relativement inversible ample sur $X/S$. Par lemme \ref{lemma-finite-type-over-affine-ample-very-ample}
Nous pouvons supposer $\mathcal{L}$ est ample sur $X$. Par lemme \ref{lemma-quasi-projective-finite-type-over-S} nous voyons qu'il existe une immersion de $X$ dans un projet d'espace $S$, i.e., $X$ est H-quasi-projective sur $S$
comme souhaité.

\end{proof}

\begin{lemma}

\label{lemma-quasi-affine-finite-type-quasi-projective}

\begin{reference}

\cite[II, proposition 5.3.4 i)]{EGA}

\end{reference}

Un morphisme de type fini quasi-affiné est quasi-projectif.

\end{lemma}

\begin{proof}
Ceci découle de \ref{lemma-quasi-affine-O-ample}.
\end{proof}

\begin{lemma}

\label{lemma-quasi-projective-permanence}
Laissez $g : Y \to S$ et $f : X \to Y$ être des morphismes de schémas. Si $g \circ f$ est quasi-projectif et $f$ est quasi-compact\footnote{Cela suit si $g$ est quasi-séparé par Schémas, lemme \ref{schemes-lemma-quasi-compact-permanence}.}, alors $f$ est quasi-projectif.
\end{lemma}

\begin{proof}
Observez que $f$ est de type fini par lemme \ref{lemma-permanence-finite-type}. Ainsi, le lemme suit de lemme \ref{lemma-ample-permanence} et les définitions.
\end{proof}





\section{Morphismes propres}
\label{section-proper}

\noindent
La notion de morphisme propre joue un rôle important dans la géométrie algébrique. Un exemple important de morphisme propre sera le morphisme structurel $\mathbf{P}^n_S \to S$ de l'espace projectif $n$, et c'est en fait l'exemple motivant menant à la définition.

\begin{definition}

\label{definition-proper}
Laissez $f : X \to S$ être un morphisme de schémas. Nous disons que $f$ est {\it proprement dit} si $f$ est séparé, tapez fini, et universellement fermé.
\end{definition}

\noindent
Le morphisme de la ligne affine avec zéro doublé à la ligne affine est de type fini et universellement fermé, donc la condition de séparation est nécessaire dans la définition ci-dessus. Dans le reste de cette section nous prouvons certaines des propriétés fondamentales des morphismes propres et des universellement fermé morphismes.

\begin{lemma}

\label{lemma-universally-closed-local-on-the-base}
Laissez $f : X \to S$ être un morphisme de schémas. Les conditions suivantes sont équivalentes:
\begin{enumerate}
\item Le morphisme $f$ est un universellement fermé.
\item Il existe un recouvrement ouvert $S = \bigcup V_j$ de telle sorte que $f^{-1}(V_j) \to V_j$ est un universellement fermé pour tous les indices $j$.
\end{enumerate}

\end{lemma}

\begin{proof}

Cela ressort clairement de la définition.

\end{proof}

\begin{lemma}

\label{lemma-proper-local-on-the-base}
Laissez $f : X \to S$ être un morphisme de schémas. Les conditions suivantes sont équivalentes:
\begin{enumerate}
\item Le morphisme $f$ est approprié.
\item Il existe un recouvrement ouvert $S = \bigcup V_j$ de telle sorte que $f^{-1}(V_j) \to V_j$ est approprié pour tous les indices $j$.
\end{enumerate}

\end{lemma}

\begin{proof}
Démonstration omise.
\end{proof}

\begin{lemma}

\label{lemma-composition-proper}

La composition des morphismes propres est appropriée. Il en est de même pour les morphismes universels et férmés.

\end{lemma}

\begin{proof}
Une composition de morphismes fermés est fermée. Si $X \to Y \to Z$ sont des morphismes universels fermés et $Z' \to Z$ est n'importe quel morphisme, alors nous voyons que $Z' \times_Z X = (Z' \times_Z Y) \times_Y X  \to Z' \times_Z Y$ est fermé et $Z' \times_Z Y \to Z'$ est fermé. D'où le résultat pour les morphismes universels fermés. Nous avons vu que les ``séparés'' et ``type fini'' sont conservés sous des compositions (Schémas, lemme \ref{schemes-lemma-separated-permanence} et lemme \ref{lemma-composition-finite-type}). D'où le résultat pour les morphismes propres.
\end{proof}

\begin{lemma}

\label{lemma-base-change-proper}

Le changement de base d'un morphisme propre est approprié. Il en est de même pour les morphismes universels fermés.

\end{lemma}

\begin{proof}

C'est vrai par définition pour l'universellement fémé morphismes. C'est vrai pour les morphismes séparés (Schémas, lemme \ref{schemes-lemma-separated-permanence}C'est vrai pour les morphismes de type fini (lemme \ref{lemma-base-change-finite-type}C'est donc vrai pour les morphismes propres.

\end{proof}

\begin{lemma}

\label{lemma-closed-immersion-proper}

Une fermée d'immersion est appropriée, donc a fortiori universellement fermé.

\end{lemma}

\begin{proof}

Le changement de base d'une immersion fermée est une immersion fermée (Schémas, lemme \ref{schemes-lemma-base-change-immersion}). Il s'agit donc d'universellement fémé. Une fermée d'immersion est séparée (Schémas, lemme \ref{schemes-lemma-immersions-monomorphisms}Une fermée d'immersion est de type fini (lemme \ref{lemma-immersion-locally-finite-type}C'est pourquoi une immersion fermée est appropriée.

\end{proof}

\begin{lemma}

\label{lemma-image-proper-scheme-closed}

Supposons donné un diagramme commutatif des schémas
$$
\xymatrix{
X \ar[rr] \ar[rd] & &
Y \ar[ld] \\
& S &
}
$$
avec $Y$ séparé sur $S$.

\begin{enumerate}

\item Si $X \to S$ est l'universellement fémé, puis le morphisme
$X \to Y$ c'est l'universellement fémé.
\item Si $X$ C'est bien $S$, puis le morphisme $X \to Y$ C'est bien.

\end{enumerate}

En particulier, dans les deux cas, l'image de $X$ en $Y$ C'est fermé.

\end{lemma}

\begin{proof}

Supposons que $X \to S$ est universellement fémé (resp.\ morphisme comme étant $X \to X \times_S Y \to Y$. Le premier morphisme est une immersion fermée, voir Schémas, lemme \ref{schemes-lemma-semi-diagonal}. Par conséquent, le premier morphisme est approprié (lemme \ref{lemma-closed-immersion-proper}). La projection $X \times_S Y \to Y$ est le changement de base d'un universellement fermé (resp.\ morphisme et donc universellement fémé (resp.\ ), voir lemme \ref{lemma-base-change-proper}. Ainsi $X \to Y$ est universellement fémé (resp.\ propre) comme la composition de l'universellement fémé (resp.\ morphismes (lemme) \ref{lemma-composition-proper}).

\end{proof}

\noindent
La preuve de la lemme suivante est due à Bjorn Poonen, voir \href{https://mathoverflow.net/questions/23337/is-a-universally-closed-morphism-of-schemes-quasi-compact/23528#23528}{ce lieu}.

\begin{lemma}

\label{lemma-universally-closed-quasi-compact}

\begin{reference}
En raison de Bjorn Poonen.
\end{reference}
Un universellement fermé morphisme des schémas est quasi-compact.
\end{lemma}

\begin{proof}
Laissez $f : X \to S$ être un morphisme. Supposons que $f$ n'est pas quasi-compact. Notre objectif est de montrer que $f$ n'est pas quasi-compact. Par Schémas, lemme \ref{schemes-lemma-quasi-compact-affine} existe un ouvert affine $V \subset S$ de telle sorte que $f^{-1}(V)$ n'est pas quasi-compact. Pour atteindre notre objectif, il suffit de montrer que $f^{-1}(V) \to V$ n'est pas universallement ferment, donc nous pouvons supposer que $S = \Spec(A)$ pour un anneau $A$.

\medskip\noindent
Écrire $X = \bigcup_{i \in I} X_i$ où les $X_i$ sont des sous-systèmes affinés ouverts de $X$. Laissez $T = \Spec(A[y_i ; i \in I])$. Laissez $T_i = D(y_i) \subset T$. Laissez $Z$ être le jeu fermé $(X \times_S T) - \bigcup_{i \in I} (X_i \times_S T_i)$. Il suffit de prouver que l'image $f_T(Z)$ de $Z$ sous $f_T : X \times_S T \to T$ n'est pas fermée.

\medskip\noindent
Il existe un point $s \in S$ tel qu'il n'y a pas de quartier $U$ de $s$ dans $S$ tel que $X_U$ est quasi-compact. Sinon, nous pourrions couvrir $S$ avec un nombre fini de $U$ et Schémas, lemme \ref{schemes-lemma-quasi-compact-affine} impliquerait $f$ quasi-compact. Réparez un tel $s \in S$.

\medskip\noindent
Nous vérifions d'abord que $f_T(Z_s) \ne T_s$. Laissez $t \in T$ être le point sur $s$ avec $\kappa(t) = \kappa(s)$ tel que $y_i = 1$ dans $\kappa(t)$ pour tous $i$. Ensuite, $t \in T_i$ pour tous $i$, et la fibre de $Z_s \to T_s$ au-dessus de $t$ est isomorphique à $(X - \bigcup_{i \in I} X_i)_s$, qui est vide. Ainsi $t \in T_s - f_T(Z_s)$.

\medskip\noindent

Présumer $f_T(Z)$ est fermé en $T$. Puis il existe un élément
$g \in A[y_i; i \in I]$ avec $f_T(Z) \subset V(g)$ mais $t \not \in V(g)$D'où l'image de $g$ en $\kappa(t)$ est non zéro. En particulier, un certain coefficient de $g$ a une image non nulle dans $\kappa(s)$. Par conséquent, ce coefficient est inversible sur certains quartiers $U$ de $s$. Laisse $J$ soit l'ensemble fini de
$j \in I$ de telle manière que $y_j$ apparaît dans $g$Depuis $X_U$ n'est pas quasi-compact, nous pouvons choisir un point $x \in X - \bigcup_{j \in J} X_j$ au-dessus de certains
$u \in U$Depuis $g$ a un coefficient qui est inversible sur $U$, nous pouvons trouver un point $t' \in T$ au-dessus $u$ de telle manière que $t' \not \in V(g)$ et
$t' \in V(y_i)$ pour tous $i \notin J$. C'est vrai parce que
$V(y_i; i \in I, i \not\in J) = \Spec(A[y_j; j\in J])$
et de l'ensemble des points de ce régime $u$ est bijectif avec $\Spec(\kappa(u)[y_j; j \in J])$. Autrement dit $t' \notin T_i$
pour chaque $i \notin J$Par Schémas, lemme \ref{schemes-lemma-points-fibre-product}
Nous pouvons trouver un point $z$ de $X \times_S T$ cartographie vers $x \in X$ et à
$t' \in T$Depuis $x \not \in X_j$ pour $j \in J$ et $t' \not \in T_i$
pour $i \in I \setminus J$ nous voyons que $z \in Z$. D'un autre côté
$f_T(z) = t' \not \in V(g)$ qui contredit $f_T(Z) \subset V(g)$. Ainsi, l ' hypothèse " .$f_T(Z)$ fermé'' est faux et nous concluons en effet que $f_T$ n'est pas fermé, comme souhaité.

\end{proof}

\noindent
Le lemme suivant dit que l'image d'un schéma approprié (dans un schéma séparé de type fini sur la base) est appropriée.

\begin{lemma}

\label{lemma-image-proper-is-proper}
Laisser $S$ être un schéma. Laisser $f : X \to Y$ être un morphisme de schémas sur $S$. Si $X$ est universellement fermé sur $S$ et $f$ est surjectif alors $Y$ est universellement fermé sur $S$. En particulier, si également $Y$ est séparé et localement de type fini sur $S$, alors $Y$ est approprié sur $S$.
\end{lemma}

\begin{proof}

Présumer $X$ est universellement fémé et $f$ surjet. Dénoter $p : X \to S$, $q : Y \to S$ les morphismes de la structure. $S' \to S$ le changement de base.
$f' : X_{S'} \to Y_{S'}$ est surjectif (lemme \ref{lemma-base-change-surjective}), et le changement de base $p' : X_{S'} \to S'$ est fermé. Si $T \subset Y_{S'}$ est fermé, alors $(f')^{-1}(T) \subset X_{S'}$
est fermé, donc $p'((f')^{-1}(T)) = q'(T)$ est fermé. $q'$ C'est la démonstration de la première déclaration. $Y \to S$ est quasi-compact par lemme \ref{lemma-universally-closed-quasi-compact}
et donc $Y \to S$ est approprié par définition si en plus $Y \to S$ est localement de type fini et séparé.

\end{proof}

\begin{lemma}

\label{lemma-scheme-theoretic-image-is-proper}
Supposons qu'un diagramme commutatif de schémas $$
\xymatrix{
X \ar[rr]_h \ar[rd]_f & & Y \ar[ld]^g \\
& S
}
$$ Supposons que
\begin{enumerate}
\item $X \to S$ est un morphisme universel (par exemple propre) et
\item $Y \to S$ est séparé et localement de type fini.
\end{enumerate}
Ensuite, l'image schématique $Z \subset Y$ de $h$ est correcte sur $S$ et $X \to Z$ est surjective.
\end{lemma}

\begin{proof}

L'image schématique de $h$ est construit en section
\ref{section-scheme-theoretic-image}Depuis $f$ est quasi compacte (lemme \ref{lemma-universally-closed-quasi-compact}) nous trouvons que $h$ est quasi compacte (Schémas, lemme \ref{schemes-lemma-quasi-compact-permanence}). D'où $h(X) \subset Z$
est dense (lem \ref{lemma-quasi-compact-scheme-theoretic-image}). D'un autre côté $h(X)$ est fermé en $Y$
(lemme) \ref{lemma-image-proper-scheme-closed}) par conséquent $X \to Z$ C'est surprenant. $Z \to S$ est un vrai (lemme \ref{lemma-image-proper-is-proper}).

\end{proof}

\noindent
La cible d'un schéma séparé sous un universellement surjectif morphisme férmé est séparée.

\begin{lemma}

\label{lemma-image-universally-closed-separated}

Laisser $S$ Être un plan. $f : X \to Y$ être un universellement surjectif férmé morphisme des schémas sur $S$.

\begin{enumerate}

\item Si $X$ est quasi-séparé, alors $Y$ est quasi-séparée.
\item Si $X$ est séparé, puis $Y$ est séparé.
\item Si $X$ est quasi-séparée $S$Alors $Y$ est quasi-séparée $S$.
\item Si $X$ est séparé sur $S$Alors $Y$ est séparé sur $S$.

\end{enumerate}

\end{lemma}

\begin{proof}
Les parties (1) et (2) sont une conséquence de (3) et (4) pour $S = \Spec(\mathbf{Z})$ (voir Schémas, définition \ref{schemes-definition-separated}). Considérez le diagramme commutatif $$
\xymatrix{
X \ar[d] \ar[rr]_{\Delta_{X/S}} & & X \times_S X \ar[d] \\
Y \ar[rr]^{\Delta_{Y/S}} & & Y \times_S Y
}
$$ La flèche verticale de gauche est surjective (c.-à-d. universellement surjective). La flèche verticale de droite est universellement fémé en tant que composition des morphismes de l'universalisation fermé $X \times_S X \to X \times_S Y \to Y \times_S Y$. C'est pourquoi elle est aussi quasi compacte, voir lemme \ref{lemma-universally-closed-quasi-compact}.

\medskip\noindent
Supposons que $X$ est quasi-séparé par rapport à $S$, c'est-à-dire que $\Delta_{X/S}$ est quasi-compact. Si $V \subset Y \times_S Y$ est un ouvert quasi-compact, alors $V \times_{Y \times_S Y} X \to \Delta_{Y/S}^{-1}(V)$ est surjectif et $V \times_{Y \times_S Y} X$ est quasi-compact par nos remarques ci-dessus. Nous concluons que $\Delta_{Y/S}$ est quasi-compact, c'est-à-dire que $Y$ est quasi-séparé par rapport à $S$.

\medskip\noindent
Supposons que $X$ est séparé sur $S$, c'est-à-dire $\Delta_{X/S}$ est une immersion fermée. Ensuite $X \to Y \times_S Y$ est fermé comme une composition de morphismes fermés. Puisque $X \to Y$ est surjectif, il s'ensuit que $\Delta_{Y/S}(Y)$ est fermé dans $Y \times_S Y$. Par conséquent, $Y$ est séparé sur $S$ par la discussion suivant Schémas, définition \ref{schemes-definition-separated}.
\end{proof}








\section{Critères valuatifs}
\label{section-valuative-criteria}

\noindent
Nous avons déjà discuté du critère d'évaluation pour la fermeture universelle et la séparation dans Schémas, Sections \ref{schemes-section-valuative-criterion-universal-closedness} et \ref{schemes-section-valuative-separatedness}. Dans cette section, nous allons discuter de certaines conséquences et variantes. Dans Limits, section \ref{limits-section-Noetherian-valuative-criterion}, nous montrerons qu'il suffit de considérer des cycles d'évaluation discrets lorsqu'on travaille avec des schémas de localisation noethérien et des morphismes de type fini.

\begin{lemma}
[critère d'évaluation de la pertinence]

\label{lemma-characterize-proper}

\begin{reference}\cite[II Theorem 7.3.8]{EGA}\end{reference}

Laisser $S$ Être un plan. $f : X \to Y$ être un morphisme des schémas sur $S$Supposons que $f$ est de type fini et quasi-séparé. Ensuite, les éléments suivants sont équivalents

\begin{enumerate}

\item  $f$ est approprié,
\item  $f$ satisfait au critère de valeur (Schémas, définition) \ref{schemes-definition-valuative-criterion}),
\item compte tenu de tout diagramme solide commutatif
$$
\xymatrix{
\Spec(K) \ar[r] \ar[d] & X \ar[d] \\
\Spec(A) \ar[r] \ar@{-->}[ru] & Y
}
$$
où $A$ est un anneau de valorisation avec champ de fractions $K$, il existe une flèche pointillée unique faisant le déplacement du diagramme.

\end{enumerate}

\end{lemma}

\begin{proof}
Partie (3) est une reformulation de (2). Ainsi, le lemme est une conséquence formelle de Schémas, la proposition \ref{schemes-proposition-characterize-universally-closed} et lemme \ref{schemes-lemma-separated-implies-valuative}, lemme \ref{schemes-lemma-valuative-criterion-separatedness}, et les définitions.
\end{proof}

\noindent
On n'a généralement pas à considérer tous les diagrammes possibles lors de l'essai du critère de valorisation. Nous appelons un critère de valorisation comme dans la prochaine lemme un ``critère de valorisation raffiné''.

\begin{lemma}

\label{lemma-refined-valuative-criterion-universally-closed}
Laissez $f : X \to S$ et $h : U \to X$ être des morphismes de schémas. Supposons que $f$ et $h$ soient quasi-compacts et que $h(U)$ soit dense en $X$. Si vous avez un diagramme solide commutatif $$
\xymatrix{
\Spec(K) \ar[r] \ar[d] & U \ar[r]^h & X \ar[d]^f \\
\Spec(A) \ar[rr] \ar@{-->}[rru] & & S
}
$$ où $A$ est un anneau d'évaluation avec champ de fractions $K$, il existe une flèche pointillée unique faisant le déplacement du diagramme, alors $f$ est universellement fermé. Si d'ailleurs $f$ est quasi-séparé, alors $f$ est séparé.
\end{lemma}

\begin{proof}

Pour prouver $f$ est universellement fémé nous vérifierons l'existence partie du critère de valeur pour $f$ ce qui est suffisant par Schémas, proposition \ref{schemes-proposition-characterize-universally-closed}. Pour ce faire, envisager un diagramme commutatif
$$
\xymatrix{
\Spec(K) \ar[r] \ar[d] & X \ar[d] \\
\Spec(A) \ar[r] & S
}
$$
où $A$ est un anneau de valorisation et $K$ est le champ de fraction de $A$. Notez que puisque les anneaux d'évaluation et les champs sont réduits, nous pouvons remplacer $U$, $X$, $S$ par leurs réductions respectives par Schémas, \ref{schemes-lemma-map-into-reduction}. Dans ce cas, l'hypothèse que $h(U)$ est dense signifie que l'image schématique de $h : U \to X$ est $X$, voir lemme \ref{lemma-scheme-theoretic-image-reduced}. Nous pouvons également remplacer $S$ par un ouvert affine à travers laquelle le morphisme $\Spec(A) \to S$ Nous pouvons donc supposer que $S = \Spec(R)$.

\medskip\noindent

Laisser $\Spec(B) \subset X$ être une affinité d'ouverture à travers laquelle le morphisme $\Spec(K) \to X$ facteurs. Choisissez une algèbre polynôme $P$ terminé $B$ et a $B$- surjection d'algèbre $P \to K$. Ensuite $\Spec(P) \to X$ est plat. D'où l'image schématique du morphisme $U \times_X \Spec(P) \to \Spec(P)$ est $\Spec(P)$ par lemme \ref{lemma-flat-base-change-scheme-theoretic-image}. Par lemme \ref{lemma-reach-points-scheme-theoretic-image}
nous pouvons trouver un diagramme commutatif
$$
\xymatrix{
\Spec(K') \ar[r] \ar[d] & U \times_X \Spec(P) \ar[d] \\
\Spec(A') \ar[r] & \Spec(P)
}
$$
où $A'$ est un anneau de valorisation et $K'$ est le champ de fraction de $A'$
tel que le point fermé de $\Spec(A')$ cartes à $\Spec(K) \subset \Spec(P)$. En d'autres termes, il y a un $B$- carte de l'algèbre
$\varphi : K \to A'/\mathfrak m_{A'}$. Choisissez un anneau de valorisation
$A'' \subset A'/\mathfrak m_{A'}$ dominant $\varphi(A)$ avec champ de fractions $K'' = A'/\mathfrak m_{A'}$
(Algèbre, lemme) \ref{algebra-lemma-dominate}). Nous avons réglé
$$
C = \{\lambda \in A' \mid \lambda \bmod \mathfrak m_{A'} \in A''\}.
$$
qui est un anneau de valorisation par Algèbre, lemme \ref{algebra-lemma-stack-valuation-rings}. Comme $C$ est un $R$-algèbre avec champ de fractions $K'$, nous obtenons un diagramme commutatif
$$
\xymatrix{
\Spec(K') \ar[r] \ar[d] & U \ar[r] & X \ar[d] \\
\Spec(C) \ar[rr] \ar@{-->}[rru] & & S
}
$$
comme dans l'énoncé du lemme. Ainsi, une flèche pointillée s'insère dans le diagramme comme indiqué. Par l'hypothèse de l'unicité du lemme la composition $\Spec(A') \to \Spec(C) \to X$ convient avec le morphisme donné $\Spec(A') \to \Spec(P) \to \Spec(B) \subset X$. D'où la restriction du morphisme au spectre de
$C/\mathfrak m_{A'} = A''$ induit le morphisme donné
$\Spec(K'') = \Spec(A'/\mathfrak m_{A'}) \to \Spec(K) \to X$. Laisse $x \in X$ être l'image du point fermé de $\Spec(A'') \to X$. L'image de l'homomorphisme d'anneaux induit $\mathcal{O}_{X, x} \to A''$
est un subring local qui est contenu dans $K \subset K''$Depuis $A$ est maximal pour la relation de domination en $K$
et depuis $A \subset A''$, nous avons $A = K \cap A''$. Nous concluons que $\mathcal{O}_{X, x} \to A''$ des facteurs $A \subset A''$. De cette façon, nous obtenons notre flèche désirée $\Spec(A) \to X$.

\medskip\noindent

Enfin, présumons $f$ est quasi-séparée. $\Delta : X \to X \times_S X$
est quasi-compact.
$$
\xymatrix{
\Spec(K) \ar[r] \ar[d] & U \ar[r]^h & X \ar[d]^\Delta \\
\Spec(A) \ar[rr] \ar@{-->}[rru] & & X \times_S X
}
$$
où $A$ est un anneau de valorisation avec champ de fractions $K$, il existe une flèche pointillée unique faisant le déplacement du diagramme. À savoir, la flèche horizontale inférieure est la même chose qu'une paire de morphismes
$\Spec(A) \to X$ qui peut servir comme la flèche pointillée dans le diagramme du lemme. Ainsi, l'unicité requise montre que la flèche horizontale inférieure facteurs à travers $\Delta$. Nous pouvons donc appliquer le résultat que nous venons de prouver
$\Delta : X \to X \times_S X$ et $h : U \to X$ et concluent que
$\Delta$ c'est l'universellement fémé. $f$
est séparé.

\end{proof}

\begin{remark}

\label{remark-check-val-on-open}

L'hypothèse de l'unicité des flèches pointillées en lemme \ref{lemma-refined-valuative-criterion-universally-closed}
est nécessaire (détails omis). Bien sûr, l'unicité est garantie si
$f$ est séparé (Schémas, lemme \ref{schemes-lemma-separated-implies-valuative}).

\end{remark}

\begin{lemma}

\label{lemma-morphism-defined-local-ring}

Laisser $S$ Être un plan. $X$, $Y$ les régimes sont plus $S$. Laisse $s \in S$ et $x \in X$, $y \in Y$ points sur $s$.

\begin{enumerate}

\item Laisser $f, g : X \to Y$ les morphismes sur $S$ de telle manière que
$f(x) = g(x) = y$ et
$f^\sharp_x = g^\sharp_x : \mathcal{O}_{Y, y} \to \mathcal{O}_{X, x}$. Ensuite, il y a un quartier ouvert $U \subset X$ avec
$f|_U = g|_U$ dans les cas suivants

\begin{enumerate}

\item  $Y$ est localement de type fini sur $S$,
\item  $X$ fait partie intégrante,
\item  $X$ est localement noethérien, ou
\item  $X$ est réduit avec de nombreux composantes irréductibles finis.

\end{enumerate}

\item Laisser $\varphi : \mathcal{O}_{Y, y} \to \mathcal{O}_{X, x}$
être un local $\mathcal{O}_{S, s}$- carte de l'algèbre. Puis il existe un quartier ouvert $U \subset X$ de $x$ et un morphisme $f : U \to Y$
cartographie $x$ à $y$ avec $f^\sharp_x = \varphi$ dans les cas suivants

\begin{enumerate}

\item  $Y$ est localement de présentation finie over $S$,
\item  $Y$ est localement de type fini et $X$ fait partie intégrante,
\item  $Y$ est localement de type fini et $X$ est localement noethérien, ou
\item  $Y$ est localement de type fini et $X$ est réduit avec de nombreux composantes irréductibles finis.

\end{enumerate}

\end{enumerate}

\end{lemma}

\begin{proof}

(1). Nous pouvons remplacer $X$, $Y$, $S$ par des quartiers affinés d'ouverture appropriés de $x$, $y$, $s$ et réduire au problème de l'algèbre suivant: donné un anneau $R$Deux $R$- cartes de l'algèbre $\varphi, \psi : B \to A$
de telle manière que

\begin{enumerate}

\item  $R \to B$ est de type fini, ou $A$ est un domaine, ou $A$
est Noétherien, ou $A$ est réduit et a finiment beaucoup de primes minimales,
\item les deux cartes $B \to A_\mathfrak p$ sont les mêmes pour certains premiers $\mathfrak p \subset A$,

\end{enumerate}

montrer que $\varphi, \psi$ définir la même carte $B \to A_g$ pour un $g \in A$, $g \not \in \mathfrak p$Si $R \to B$ est de type fini, let $t_1, \ldots, t_m \in B$ les générateurs de $B$
comme une $R$- de l'algèbre. $j$ Nous pouvons trouver
$g_j \in A$, $g_j \not \in \mathfrak p$
de telle manière que $\varphi(t_j)$ et $\psi(t_j)$ avoir la même image dans
$A_{g_j}$. Ensuite, nous avons mis $g = \prod g_j$Dans les autres cas (si $A$ est un domaine, Noétherien, ou réduit avec finitement beaucoup de primes minimales), nous pouvons trouver un $g \in A$,
$g \not \in \mathfrak p$ de telle manière que $A_g \subset A_\mathfrak p$Voir Algèbre, lemme \ref{algebra-lemma-subring-of-local-ring}. Ainsi les cartes $B \to A_g$ sont égaux comme désiré.

\medskip\noindent

la démonstration de (2). Pour ce faire, nous pouvons remplacer $X$, $Y$, $S$ par des ouverts adaptés affines. $X = \Spec(A)$, $Y = \Spec(B)$, $S = \Spec(R)$. Laisse $\mathfrak p \subset A$ être l'idéal premier correspondant à $x$. Laisse $\mathfrak q \subset B$ être le premier correspondant à $y$. Ensuite $\varphi$ est un local $R$- carte de l'algèbre
$\varphi : B_\mathfrak q \to A_\mathfrak p$Si $R \to B$ est un homomorphisme d'anneaux de présentation finie, puis il existe un
$g \in A \setminus \mathfrak p$ et une $R$- carte de l'algèbre $B \to A_g$ de telle manière que
$$
\xymatrix{
B_\mathfrak q \ar[r]_\varphi & A_\mathfrak p \\
B \ar[u] \ar[r] & A_g \ar[u]
}
$$
déplacements, voir Algèbre, lemmes \ref{algebra-lemma-characterize-finite-presentation} et
\ref{algebra-lemma-localization-colimit}. Le morphisme induit $\Spec(A_g) \to \Spec(B)$ fonctionne. Si $B$ est de type fini sur $R$, laissez $t_1, \ldots, t_m \in B$ les générateurs de $B$ comme une $R$- de l'algèbre.
$g_j \in A$, $g_j \not \in \mathfrak p$
de telle manière que $\varphi(t_j) \in \Im(A_{g_j} \to A_\mathfrak p)$. Ainsi après avoir remplacé $A$ par $A[1/\prod g_j]$ nous pouvons supposer que $B$ des cartes dans l'image de $A \to A_\mathfrak p$. Si nous pouvons trouver un $g \in A$, $g \not \in \mathfrak p$
de telle manière que $A_g \to A_\mathfrak p$ est injectable, alors nous obtiendrons le désir $R$- carte de l'algèbre $B \to A_g$. Ainsi ce qui achève la démonstration par une autre application de See Algèbre, lemme \ref{algebra-lemma-subring-of-local-ring}.

\end{proof}

\begin{lemma}

\label{lemma-extend-across}

Laisser $S$ Être un plan. $X$, $Y$ les régimes sont plus $S$. Laisse $x \in X$. Laisse $U \subset X$ être un ouvert et laisser $f : U \to Y$ être un morphisme sur $S$Supposons que

\begin{enumerate}

\item  $x$ est à la fermeture de $U$,
\item  $X$ est réduit avec finiment de nombreux composantes irréductibles ou
$X$ est Noétherien,
\item  $\mathcal{O}_{X, x}$ est un anneau de valorisation,
\item  $Y \to S$ est approprié

\end{enumerate}

Puis il y a une ouverture $U \subset U' \subset X$ contenant
$x$ et une $S$- morphisme $f' : U' \to Y$ prolongement $f$.

\end{lemma}

\begin{proof}

Il est inoffensif de remplacer $X$ par un quartier ouvert de $x$ en $X$
(petit détail omis). Par Propriétés, laisse
\ref{properties-lemma-ring-affine-open-injective-local-ring}
Nous pouvons supposer $X$ est affine avec
$\Gamma(X, \mathcal{O}_X) \subset \mathcal{O}_{X, x}$. En particulier $X$ fait partie intégrante d'un point général unique $\xi$
dont le corps résiduel est le champ de fractions $K$ de l'anneau de valorisation $\mathcal{O}_{X, x}$Depuis $x$ est à la fermeture de $U$ nous voyons que $U$ n'est pas vide, donc $U$ contient $\xi$. Ainsi par le critère de valorisation de la justesse (lemme \ref{lemma-characterize-proper}) il y a un morphisme $t : \Spec(\mathcal{O}_{X, x}) \to Y$
s'insérer dans un diagramme commutatif
$$
\xymatrix{
\Spec(K) \ar[d]_\xi \ar[r] & \Spec(\mathcal{O}_{X, x}) \ar[d]_t \\
U \ar[r]^f & Y
}
$$
des morphismes des régimes $S$. Appliquer lemme \ref{lemma-morphism-defined-local-ring}
avec $y = t(x)$ et $\varphi = t^\sharp_x$ nous obtenons un quartier ouvert $V \subset X$ de $x$ et un morphisme $g : V \to Y$
terminé $S$ qui envoie $x$ à $y$ et de telle sorte que $g^\sharp_x = t^\sharp_x$. Comme $Y \to S$ est séparé, l'égaliseur $E$ de $f|_{U \cap V}$
et $g|_{U \cap V}$ est un sous-système fermé de $U \cap V$, voir Schémas, lemme \ref{schemes-lemma-where-are-they-equal}Depuis $f$ et $g$ déterminer le même morphisme $\Spec(K) \to Y$
par la construction nous voyons que $E$ contient le point général du régime intégral $U \cap V$. D'où $E = U \cap V$ et nous concluons que $f$ et $g$ colle à un morphisme $U' = U \cup V \to Y$
comme souhaité.

\end{proof}





\section{Morphismes projectifs}
\label{section-projective}

\noindent
Nous utiliserons la définition d'un projectif morphisme à partir de \cite{EGA}. La version de la définition avec le ``H'' est celle de \cite{H}. Les définitions résultantes sont différentes. Les deux sont utiles.

\begin{definition}

\label{definition-projective}

Laisser $f : X \to S$ être un morphisme de schémas.

\begin{enumerate}

\item Nous disons $f$ est {\it projectif} si $X$ est isomorphe comme un $S$-régime d'un sous-régime fermé d'un groupe projectif $\mathbf{P}(\mathcal{E})$
pour certains quasi-cohérents, type fini $\mathcal{O}_S$-module $\mathcal{E}$.
\item Nous disons $f$ est {\it H-projective} s'il y a un entier $n$ et une immersion fermée $X \to \mathbf{P}^n_S$ terminé $S$.
\item Nous disons $f$ est {\it localement projective} s'il existe une ouverture de recouvrement $S = \bigcup U_i$ de telle sorte que chacun $f^{-1}(U_i) \to U_i$ est projectif.

\end{enumerate}

\end{definition}

\noindent
Comme prévu, un morphisme projectif est quasi-projectif, voir lemme \ref{lemma-projective-quasi-projective}. Inversement, les morphismes quasi-projectifs sont souvent des compositions d'immersions ouvertes et morphismes projectifs, voir lemme \ref{lemma-quasi-projective-open-projective}. Pour une vue d'ensemble des propriétés des morphismes projectifs sur une base quasi-projective, voir Compléments sur les morphismes, section \ref{more-morphisms-section-projective}.

\begin{example}

\label{example-projective}

Laisser $S$ Être un plan. $\mathcal{A}$ être quasi-cohérents
$\mathcal{O}_S$-algèbre généré par $\mathcal{A}_1$ terminé $\mathcal{A}_0$. Supposons en outre que $\mathcal{A}_1$ est de type fini sur
$\mathcal{O}_S$. Prêt $X = \underline{\text{Proj}}_S(\mathcal{A})$. Dans ce cas $X \to S$ est projetif. C'est-à-dire, le morphisme associé au $\mathcal{O}_S$- carte de l'algèbre
$$
\text{Sym}_{\mathcal{O}_X}^*(\mathcal{A}_1)
\longrightarrow
\mathcal{A}
$$
est une immersion fermée, voir Constructions, lemme
\ref{constructions-lemma-surjective-generated-degree-1-map-relative-proj}.

\end{example}
 
\begin{lemma}

\label{lemma-H-projective}

Un projectif H-morphisme est projectif H-quasi. Un projectif H-morphisme est projectif.

\end{lemma}

\begin{proof}
La première déclaration est immédiate des définitions. La seconde tient comme $\mathbf{P}^n_S$ est un bundle projectif sur $S$, voir Constructions, lemme \ref{constructions-lemma-projective-space-bundle}.
\end{proof}

\begin{lemma}

\label{lemma-characterize-locally-projective}
Laissez $f : X \to S$ être un morphisme de schémas. Les conditions suivantes sont équivalentes:
\begin{enumerate}
\item Le morphisme $f$ est localement projectif.
\item Il existe un recouvrement ouvert $S = \bigcup U_i$ de telle sorte que chaque $f^{-1}(U_i) \to U_i$ est H-projectif.
\end{enumerate}

\end{lemma}

\begin{proof}

Par moi \ref{lemma-H-projective} nous voyons que (2) implique (1). Supposons que (1). Pour chaque point $s \in S$ Nous pouvons trouver $\Spec(R) = U \subset S$
un quartier affiné d'ouverture de $s$ de telle manière que $X_U$ est isomorphe à un sous-système fermé de $\mathbf{P}(\mathcal{E})$ pour certains types fini, faisceau quasi-cohérent de $\mathcal{O}_U$-modules $\mathcal{E}$. Écrire $\mathcal{E} = \widetilde{M}$ pour certains types de fini
$R$-module $M$ (voir Propriétés, lemme \ref{properties-lemma-finite-type-module}). Choisissez des générateurs $x_0, \ldots, x_n \in M$ de $M$ comme une $R$Considérez la surjection classée $R$- carte de l'algèbre
$$
R[X_0, \ldots, X_n] \longrightarrow \text{Sym}_R(M).
$$
Selon Constructions, \ref{constructions-lemma-surjective-graded-rings-map-proj}
le morphisme correspondant
$$
\mathbf{P}(\mathcal{E}) \to \mathbf{P}^n_R
$$
est une immersion fermée. $f^{-1}(U)$ est isomorphe à un sous-système fermé de $\mathbf{P}^n_U$ (en tant que régime $U$En d'autres termes : 2) tient.

\end{proof}

\begin{lemma}

\label{lemma-locally-projective-proper}
Un projet de morphisme local est approprié.
\end{lemma}

\begin{proof}

Laisser $f : X \to S$ être localement projective. Pour montrer que $f$ est approprié nous pouvons travailler localement sur la base, voir lemme \ref{lemma-proper-local-on-the-base}. Par conséquent, par \ref{lemma-characterize-locally-projective}
ci-dessus nous pouvons supposer qu'il existe une immersion fermée $X \to \mathbf{P}^n_S$. Par lemmes \ref{lemma-composition-proper}
et \ref{lemma-closed-immersion-proper} il suffit de prouver que
$\mathbf{P}^n_S \to S$ C'est bien.
$\mathbf{P}^n_S \to S$ est le changement de base de
$\mathbf{P}^n_{\mathbf{Z}} \to \Spec(\mathbf{Z})$ il suffit de montrer que $\mathbf{P}^n_{\mathbf{Z}} \to \Spec(\mathbf{Z})$
C'est bien, voyez-moi. \ref{lemma-base-change-proper}. Par Constructions, lemme \ref{constructions-lemma-proj-separated} le régime
$\mathbf{P}^n_{\mathbf{Z}}$ est séparé. Par Constructions, laisse \ref{constructions-lemma-proj-quasi-compact} le régime
$\mathbf{P}^n_{\mathbf{Z}}$ est quasi compact. Il est clair que $\mathbf{P}^n_{\mathbf{Z}} \to \Spec(\mathbf{Z})$
est localement de type fini depuis $\mathbf{P}^n_{\mathbf{Z}}$ est couvert par les affines des ouverts $D_{+}(X_i)$ dont chacun est le spectre du type fini $\mathbf{Z}$-Algèbre
$$
\mathbf{Z}[X_0/X_i, \ldots, X_n/X_i].
$$
Enfin, nous devons montrer que
$\mathbf{P}^n_{\mathbf{Z}} \to \Spec(\mathbf{Z})$
est universellement fémé. Ceci découle de Constructions, lemme \ref{constructions-lemma-proj-valuative-criterion}
et le critère de valorisation (voir Schémas, proposition \ref{schemes-proposition-characterize-universally-closed}).

\end{proof}

\begin{lemma}

\label{lemma-proper-ample-locally-projective}
Laissez $f : X \to S$ être un morphisme propre de schémas. S'il existe un $f$-faisceau inversible ample sur $X$, alors $f$ est localement projectif.
\end{lemma}

\begin{proof}
S'il existe un $f$-faisceau inversible ample, alors nous pouvons localement sur $S$ trouver une immersion $i : X \to \mathbf{P}^n_S$, voir lemme \ref{lemma-finite-type-over-affine-ample-very-ample}. Puisque $X \to S$ est approprié le morphisme $i$ est une immersion fermée, voir lemme \ref{lemma-image-proper-scheme-closed}.
\end{proof}

\begin{lemma}

\label{lemma-H-projective-composition}

Une composition de H-morphismes projectifs est H-projective.

\end{lemma}

\begin{proof}
Supposons que $X \to Y$ et $Y \to Z$ soient projectifs H. Ensuite, il existe des immersions fermées $X \to \mathbf{P}^n_Y$ sur $Y$, et $Y \to \mathbf{P}^m_Z$ sur $Z$. Considérez le schéma suivant $$
\xymatrix{
X \ar[r] \ar[d] &
\mathbf{P}^n_Y \ar[r] \ar[dl] &
\mathbf{P}^n_{\mathbf{P}^m_Z} \ar[dl] \ar@{=}[r] &
\mathbf{P}^n_Z \times_Z \mathbf{P}^m_Z \ar[r] &
\mathbf{P}^{nm + n + m}_Z \ar[ddllll] \\
Y \ar[r] \ar[d] & \mathbf{P}^m_Z \ar[dl] & \\
Z & &
}
$$ Ici, la flèche horizontale la plus à droite est l'intégration de Segre, voir Constructions, lemme \ref{constructions-lemma-segre-embedding}. Le diagramme identifie $X$ comme un sous-système fermé de $\mathbf{P}^{nm + n + m}_Z$ selon le souhait.
\end{proof}

\begin{lemma}

\label{lemma-H-projective-base-change}

Un changement de base d'un projectif H-morphisme est H-projectif.

\end{lemma}

\begin{proof}
C'est vrai parce que le changement de base d'espace projetif sur un schéma est espace projetif, et le fait que le changement de base d'une immersion fermée est une immersion fermée, voir Schémas, lemme \ref{schemes-lemma-base-change-immersion}.
\end{proof}

\begin{lemma}

\label{lemma-base-change-projective}

Un changement de base d'un projet morphisme (localement) est (localement) projectif.

\end{lemma}

\begin{proof}
C'est vrai parce que le changement de base d'un paquet projectif sur un schéma est un paquet projectif, le retrait d'un type fini $\mathcal{O}$-module est de type fini (Modules, lemme \ref{modules-lemma-pullback-finite-type}) et le fait que le changement de base d'une immersion fermée est une immersion fermée, voir Schémas, lemme \ref{schemes-lemma-base-change-immersion}. Quelques détails omis.
\end{proof}

\begin{lemma}

\label{lemma-projective-quasi-projective}

Un morphisme projectif est quasi-projectif.

\end{lemma}

\begin{proof}

Laisser $f : X \to S$ être un projet morphisme. Choisissez une immersion fermée
$i : X \to \mathbf{P}(\mathcal{E})$ où $\mathcal{E}$ est un quasi-cohérent, type fini $\mathcal{O}_S$- Module.
$\mathcal{L} = i^*\mathcal{O}_{\mathbf{P}(\mathcal{E})}(1)$ est $f$- Très ample. $f$ est approprié (lemme \ref{lemma-locally-projective-proper}) c'est quasi-compact. \ref{lemma-ample-very-ample} implique que $\mathcal{L}$ est $f$- Depuis. $f$ C'est vrai qu'il est de type fini. Ainsi, nous avons vérifié toutes les propriétés de définition des cales quasi-projectives et nous gagnons.

\end{proof}

\begin{lemma}

\label{lemma-H-quasi-projective-open-H-projective}
Laissez $f : X \to S$ être un quasi-projectif H-morphisme. Puis $f$ facteurs comme $X \to X' \to S$ où $X \to X'$ est un ouverte d'immersion et $X' \to S$ est H-projectif.
\end{lemma}

\begin{proof}
Par définition, nous pouvons factoriser $f$ comme une immersion quasi-compacte $i : X \to \mathbf{P}^n_S$ suivie de la projection $\mathbf{P}^n_S \to S$. Par lemme \ref{lemma-quasi-compact-immersion} il existe un sous-système fermé $X' \subset \mathbf{P}^n_S$ de telle sorte que $i$ facteurs par une immersion ouverte $X \to X'$. La lemme suit.
\end{proof}

\begin{lemma}

\label{lemma-quasi-projective-open-projective}
Laissez $f : X \to S$ être un quasi-projectif morphisme avec $S$ quasi-compact et quasi-séparé. Puis $f$ facteurs comme $X \to X' \to S$ où $X \to X'$ est une ouverture d'immersion et $X' \to S$ est projectif.
\end{lemma}

\begin{proof}

Laisser $\mathcal{L}$ être $f$- Depuis. $f$ est de type fini et $S$ est quasi compacte $\mathcal{L}^{\otimes n}$ est $f$- très ample pour certains $n > 0$, voir lemme \ref{lemma-finite-type-ample-very-ample}. Remplacer $\mathcal{L}$ par $\mathcal{L}^{\otimes n}$. Écrire $\mathcal{F} = f_*\mathcal{L}$. C'est une quasi-cohérence
$\mathcal{O}_S$-module de Schémas, me laisse \ref{schemes-lemma-push-forward-quasi-coherent}
(les morphismes quasi-projectifs sont quasi-compacts et séparés, voir lemme \ref{lemma-quasi-projective-properties}) Par Propriétés, lemme \ref{properties-lemma-directed-colimit-finite-presentation}
nous pouvons trouver un ensemble dirigé $I$ et un système de type fini quasi-cohérent $\mathcal{O}_S$-modules $\mathcal{E}_i$
terminé $I$ de telle manière que $\mathcal{F} = \colim \mathcal{E}_i$. Considérer les compositions
$\psi_i : f^*\mathcal{E}_i \to f^*\mathcal{F} \to \mathcal{L}$. Choisissez une affine finie de récupération ouvert $S = \bigcup_{j = 1, \ldots, m} V_j$. Pour chaque $j$ nous pouvons choisir des sections
$$
s_{j, 0}, \ldots, s_{j, n_j} \in
\Gamma(f^{-1}(V_j), \mathcal{L}) = f_*\mathcal{L}(V_j) = \mathcal{F}(V_j)
$$
qui génèrent $\mathcal{L}$ terminé $f^{-1}V_j$ et définir une immersion
$$
f^{-1}V_j \longrightarrow \mathbf{P}^{n_j}_{V_j},
$$
voir lemme \ref{lemma-very-ample-finite-type-over-affine}. Choisissez $i$ de telle sorte qu'il existe des sections $e_{j, t} \in \mathcal{E}_i(V_j)$
cartographie vers $s_{j, t}$ en $\mathcal{F}$ pour tous $j = 1, \ldots, m$ et
$t = 1, \ldots, n_j$. Puis l'application $\psi_i$ est surjectif comme les sections $f^*e_{j, t}$ avoir la même image que les sections $s_{j, t}$
qui génèrent $\mathcal{L}|_{f^{-1}V_j}$. D'où nous obtenons un morphisme
$$
r_{\mathcal{L}, \psi_i} : X \longrightarrow \mathbf{P}(\mathcal{E}_i)
$$
terminé $S$ de telle sorte que $V_j$ nous avons une factorisation
$$
f^{-1}V_j \to \mathbf{P}(\mathcal{E}_i)|_{V_j} \to \mathbf{P}^{n_j}_{V_j}
$$
de l'immersion ci-dessus. Il s'ensuit que $r_{\mathcal{L}, \psi_i}|_{V_j}$
est une immersion, voir lemme \ref{lemma-immersion-permanence}Depuis $S = \bigcup V_j$ nous concluons que $r_{\mathcal{L}, \psi_i}$
est une immersion. Notez que $r_{\mathcal{L}, \psi_i}$ est quasi-compact comme
$X \to S$ est quasi compacte et $\mathbf{P}(\mathcal{E}_i) \to S$ est séparé (voir Schémas, lemme \ref{schemes-lemma-quasi-compact-permanence}). Par lemme \ref{lemma-quasi-compact-immersion} il existe un sous-système fermé $X' \subset \mathbf{P}(\mathcal{E}_i)$ de telle manière que $i$ facteurs par une immersion ouverte $X \to X'$. Ensuite $X' \to S$ est projetive par définition et nous gagnons.

\end{proof}

\begin{lemma}

\label{lemma-projective-is-quasi-projective-proper}

Laisser $S$ être un système quasi compact et quasi séparé. $f : X \to S$ être un morphisme de plans.

\begin{enumerate}

\item  $f$ est projectif si et seulement si $f$ est quasi-projectif et approprié, et
\item  $f$ est H-projectif si et seulement si $f$ est H-quasi-projective et approprié.

\end{enumerate}

\end{lemma}

\begin{proof}

Si $f$ est projetive, alors $f$ est quasi-projectif par lemme \ref{lemma-projective-quasi-projective} et propre par lemme \ref{lemma-locally-projective-proper}. Inversement, si
$X \to S$ est quasi-projectif et approprié, alors nous pouvons choisir une immersion ouverte $X \to X'$ avec $X' \to S$ projetive par lemme \ref{lemma-quasi-projective-open-projective}Depuis $X \to S$ est correct, nous voyons que $X$ est fermé en $X'$
(lemme) \ref{lemma-image-proper-scheme-closed}), i.e.,
$X \to X'$ est une (ouverte et) immersion fermée. $X'$ est isomorphe à un sous-système fermé d'un faisceau projectif $S$ (définition \ref{definition-projective}) nous voyons que la même chose est vraie pour $X$, i.e.,
$X \to S$ est un projectif morphisme. Cela prouve (1). La preuve de (2) est la même, sauf qu'il utilise lemmes \ref{lemma-H-projective} et
\ref{lemma-H-quasi-projective-open-H-projective}.

\end{proof}

\begin{lemma}

\label{lemma-composition-projective}
Laissez $f : X \to Y$ et $g : Y \to S$ être des morphismes de schémas. Si $S$ est quasi-compact et quasi-séparé et $f$ et $g$ sont projectifs, alors $g \circ f$ est projectif.
\end{lemma}

\begin{proof}
Par lemmes \ref{lemma-projective-quasi-projective} et \ref{lemma-locally-projective-proper} nous voyons que $f$ et $g$ sont quasi-projectifs et appropriés. Par lemmes \ref{lemma-composition-proper} et \ref{lemma-composition-quasi-projective} nous voyons que $g \circ f$ est approprié et quasi-projectif. Ainsi, $g \circ f$ est projectif par lemme \ref{lemma-projective-is-quasi-projective-proper}.
\end{proof}

\begin{lemma}

\label{lemma-projective-permanence}
Laissez $g : Y \to S$ et $f : X \to Y$ être des morphismes de schémas. Si $g \circ f$ est projectif et $g$ est séparé, alors $f$ est projectif.
\end{lemma}

\begin{proof}

Choisissez une fermée immersion $X \to \mathbf{P}(\mathcal{E})$ où $\mathcal{E}$
est un quasi-cohérent, type fini $\mathcal{O}_S$Alors on aura un morphisme. $X \to \mathbf{P}(\mathcal{E}) \times_S Y$. Ce morphisme est une immersion fermée parce que c'est la composition
$$
X \to X \times_S Y \to \mathbf{P}(\mathcal{E}) \times_S Y
$$
où le premier morphisme est une immersion fermée de Schémas, lemme \ref{schemes-lemma-semi-diagonal}
(et le fait que $g$ est séparé) et le second comme le changement de base d'une immersion fermée. Enfin, le produit fibré $\mathbf{P}(\mathcal{E}) \times_S Y$
est isomorphe à $\mathbf{P}(g^*\mathcal{E})$ et le retrait préserve quasi-cohérent, les modules fini de type.

\end{proof}

\begin{lemma}

\label{lemma-projective-over-quasi-projective-is-H-projective}
Laissez $S$ être un schéma qui admet un faisceau inversible ample. Puis
\begin{enumerate}
\item tout morphisme projetif $X \to S$ est H-projectif, et
\item tout morphisme quasi-projectif $X \to S$ est H-quasi-projectif.
\end{enumerate}

\end{lemma}

\begin{proof}

Les hypothèses sur $S$ implique que $S$ est quasi-compact et séparé, voir Propriétés, définition \ref{properties-definition-ample} et lemme \ref{properties-lemma-ample-immersion-into-proj}
et des travaux de construction \ref{constructions-lemma-proj-separated}. Par conséquent, lemme \ref{lemma-quasi-projective-open-projective}
applique et nous voyons que (1) implique (2). $\mathcal{E}$ être un type fini quasi-cohérent $\mathcal{O}_S$-module. Par notre définition des morphismes projetés, il suffit de montrer que
$\mathbf{P}(\mathcal{E}) \to S$ est H-projective. Si $\mathcal{E}$ est générée par de nombreuses sections globales finies, puis la surjection correspondante $\mathcal{O}_S^{\oplus n} \to \mathcal{E}$
induit une immersion fermée
$$
\mathbf{P}(\mathcal{E}) \longrightarrow
\mathbf{P}(\mathcal{O}_S^{\oplus n}) = \mathbf{P}^{n - 1}_S
$$
comme souhaité. En général, laisser $\mathcal{L}$ être un faisceau inversible ample sur $S$. Par Propriétés, proposition \ref{properties-proposition-characterize-ample}
il existe un entier $n$ de telle manière que
$\mathcal{E} \otimes_{\mathcal{O}_S} \mathcal{L}^{\otimes n}$
est générée globalement par de nombreuses sections finies.
$\mathbf{P}(\mathcal{E}) =
\mathbf{P}(\mathcal{E} \otimes_{\mathcal{O}_S} \mathcal{L}^{\otimes n})$ par Constructions, lemme \ref{constructions-lemma-twisting-and-proj}
ça finit la démonstration.

\end{proof}

\begin{lemma}

\label{lemma-proper-ample-is-proj}
Laissez $f : X \to S$ être un morphisme fermé universel. Laissez $\mathcal{L}$ être un module inversible $f$-ample $\mathcal{O}_X$. Puis le morphisme canonique $$
r : X
\longrightarrow
\underline{\text{Proj}}_S
\left(
\bigoplus\nolimits_{d \geq 0} f_*\mathcal{L}^{\otimes d}
\right)
$$ de lemme \ref{lemma-characterize-relatively-ample} est un isomorphisme.
\end{lemma}

\begin{proof}

Observer que $f$ est quasi-compact parce que l'existence d'un $f$-module amples forces inversibles $f$ par le lemme cité le morphisme $r$ est une immersion ouverte. D'autre part, l'image de $r$ est fermé par lemme \ref{lemma-image-proper-scheme-closed}
(l'objectif de $r$ est séparé sur $S$ par Constructions, lemme \ref{constructions-lemma-relative-proj-separated}Enfin, l'image de $r$ est dense par Propriétés, lemme \ref{properties-lemma-ample-immersion-into-proj}
(nous utilisons également ici qu'il a été montré dans la démonstration de \ref{lemma-characterize-relatively-ample}
que le morphisme $r$ sur les ouverts affines de $S$
est donné par le morphisme canonique de Propriétés, lemme \ref{properties-lemma-map-into-proj}Ainsi, nous concluons que $r$ est une immersion surjective ouverte, c'est-à-dire un isomorphisme.

\end{proof}

\begin{lemma}

\label{lemma-proper-ample-delete-affine}
Laissez $f : X \to S$ être un morphisme férmé universel. Laissez $\mathcal{L}$ être un module inversible $f$-ample $\mathcal{O}_X$. Laissez $s \in \Gamma(X, \mathcal{L})$. Puis $X_s \to S$ est un morphisme affine.
\end{lemma}

\begin{proof}

La question est locale sur $S$ (lemme) \ref{lemma-characterize-affine}) donc nous pouvons supposer $S$ C'est affiné. \ref{lemma-proper-ample-is-proj}
nous pouvons écrire $X = \text{Proj}(A)$ où $A$ est une bague notée et $s$ correspond à $f \in A_1$ et
$X_s = D_+(f)$ (Propriétés, lemme) \ref{properties-lemma-map-into-proj}) qui prouve le lemme par la construction de $\text{Proj}(A)$, voir Constructions, section \ref{constructions-section-proj}.

\end{proof}











\section{Morphismes entiers et finis}
\label{section-integral}

\noindent
Rappelez qu'un homomorphisme d'anneaux $R \to A$ est dit {\it intégrale} si chaque élément de $A$ satisfait à une équation monique avec des coefficients dans $R$. Rappelez-vous qu'un homomorphisme d'anneaux $R \to A$ est dit {\it fini} si $A$ est fini comme un module $R$. Voir Algèbre, définition \ref{algebra-definition-integral-ring-map}.

\begin{definition}

\label{definition-integral}

Laisser $f : X \to S$ être un morphisme de schémas.

\begin{enumerate}

\item Nous disons que $f$ est {\it intégrale} si $f$ est affine et si pour chaque affine ouvert $\Spec(R) = V \subset S$
avec image inverse $\Spec(A) = f^{-1}(V) \subset X$
l'homomorphisme d'anneaux associé $R \to A$ C'est une partie intégrante.
\item Nous disons que $f$ est {\it fini} si $f$ est affine et si pour chaque affine ouvert $\Spec(R) = V \subset S$
avec image inverse $\Spec(A) = f^{-1}(V) \subset X$
l'homomorphisme d'anneaux associé $R \to A$ C'est fini.

\end{enumerate}

\end{definition}

\noindent
Il est clair que les morphismes finis sont séparés et quasi-compacts. Il est également clair qu'un morphisme fini est un morphisme de type fini. La plupart des lemmes de cette section sont complètement standard. Mais notez le fun lemme \ref{lemma-integral-universally-closed} à la fin de la section.

\begin{lemma}

\label{lemma-integral-local}

Laisser $f : X \to S$ morphisme des régimes. Les conditions suivantes sont équivalentes:

\begin{enumerate}

\item Le morphisme $f$ C'est une partie intégrante.
\item Il existe un recouvrement ouvert affine $S = \bigcup U_i$ de telle sorte que chacun $f^{-1}(U_i)$ est affine et
$\mathcal{O}_S(U_i) \to \mathcal{O}_X(f^{-1}(U_i))$ C'est une partie intégrante.
\item Il existe une ouverture de recouvrement $S = \bigcup U_i$
de telle sorte que chacun $f^{-1}(U_i) \to U_i$ C'est une partie intégrante.

\end{enumerate}

En outre, si $f$ fait partie intégrante de chaque sous-système ouvert
$U \subset S$ le morphisme $f : f^{-1}(U) \to U$ C'est une partie intégrante.

\end{lemma}

\begin{proof}
Voir Algèbre, lemme \ref{algebra-lemma-integral-local}. Quelques détails omis.
\end{proof}

\begin{lemma}

\label{lemma-finite-local}

Laisser $f : X \to S$ morphisme des régimes. Les conditions suivantes sont équivalentes:

\begin{enumerate}

\item Le morphisme $f$ C'est fini.
\item Il existe un recouvrement ouvert affine $S = \bigcup U_i$ de telle sorte que chacun $f^{-1}(U_i)$ est affine et
$\mathcal{O}_S(U_i) \to \mathcal{O}_X(f^{-1}(U_i))$ C'est fini.
\item Il existe une ouverture de recouvrement $S = \bigcup U_i$
de telle sorte que chacun $f^{-1}(U_i) \to U_i$ C'est fini.

\end{enumerate}

En outre, si $f$ est fini alors pour chaque sous-système ouvert
$U \subset S$ le morphisme $f : f^{-1}(U) \to U$ C'est fini.

\end{lemma}

\begin{proof}
Voir Algèbre, lemme \ref{algebra-lemma-integral-local}. Quelques détails omis.
\end{proof}

\begin{lemma}

\label{lemma-finite-integral}

Un morphisme fini fait partie intégrante. Un morphisme entier qui est localement de type fini est fini.

\end{lemma}

\begin{proof}
Voir Algèbre, lemme \ref{algebra-lemma-finite-is-integral} et lemme \ref{algebra-lemma-characterize-finite-in-terms-of-integral}.
\end{proof}

\begin{lemma}

\label{lemma-composition-finite}

Une composition de morphismes finis est finie. Il en est de même pour les morphismes entrants.

\end{lemma}

\begin{proof}
Voir Algèbre, lemmes \ref{algebra-lemma-finite-transitive} et \ref{algebra-lemma-integral-transitive}.
\end{proof}

\begin{lemma}

\label{lemma-base-change-finite}

Un changement de base d'un morphisme fini est fini. Il en est de même pour les morphismes entrants.

\end{lemma}

\begin{proof}
Voir Algèbre, lemme \ref{algebra-lemma-base-change-integral}.
\end{proof}

\begin{lemma}

\label{lemma-integral-universally-closed}

\begin{slogan}
intégrale $=$ affine $+$ universallement fermé
\end{slogan}
Laissez $f : X \to S$ être un morphisme de schémas. Les conditions suivantes sont équivalentes
\begin{enumerate}
\item $f$ est intégrale, et
\item $f$ est affiné et universalement fermé.
\end{enumerate}

\end{lemma}

\begin{proof}
Supposons (1). Un morphisme entier est affin par définition. Un changement de base d'un morphisme entier est donc intégral afin de prouver (2) il suffit de montrer qu'un morphisme entier est fermé. Ceci découle d'Algèbre, lemmes \ref{algebra-lemma-integral-going-up} et \ref{algebra-lemma-going-up-closed}.

\medskip\noindent

Supposons (2). Nous pouvons supposer $f$ est le morphisme
$f : \Spec(A) \to \Spec(R)$ provenant d'un homomorphisme d'anneaux
$R \to A$. Laisse $a$ être un élément de $A$. Nous devons montrer que $a$ fait partie intégrante $R$, c'est-à-dire que dans le noyau
$I$ de l'application $R[x] \to A$ envoi $x$ à $a$ il y a un polynôme monique.
$B = A[x]/(ax -1)$ et laisser $J$ être le noyau de la composition $R[x]\to A[x] \to B$Si $f\in J$ il existe
$q\in A[x]$ de telle manière que $f = (ax-1)q$ en $A[x]$ ainsi si
$f = \sum_i f_ix^i$ et $q = \sum_iq_ix^i$, pour tous $i \geq 0$
Nous avons $f_i = aq_{i-1} - q_i$. Pour $n \geq \deg q + 1$ le polynôme
$$
\sum\nolimits_{i \geq 0} f_i x^{n - i} =
\sum\nolimits_{i \geq 0} (a q_{i - 1} - q_i) x^{n - i} =
(a - x) \sum\nolimits_{i \geq 0} q_i x^{n - i - 1}
$$
est clairement dans $I$; si $f_0 = 1$ ce polynôme est aussi monique, donc nous sommes réduits pour prouver que $J$ contient un polynôme à terme constant $1$. Nous le faisons en prouvant $\Spec(R[x]/(J + (x))$ est vide.

\medskip\noindent

Depuis $f$ est universellement férmé le changement de base $\Spec(A[x]) \to \Spec(R[x])$
est fermé. D'où l'image du sous-ensemble fermé $\Spec(B) \subset \Spec(A[x])$
est le sous-ensemble fermé $\Spec(R[x]/J) \subset \Spec(R[x])$, voir exemple \ref{example-scheme-theoretic-image} et lemme \ref{lemma-quasi-compact-scheme-theoretic-image}. En particulier $\Spec(B) \to \Spec(R[x]/J)$ est surjectif. Considérez le diagramme suivant où chaque carré est un retrait:
$$
\xymatrix{
\Spec(B) \ar@{->>}[r]^g &
\Spec(R[x]/J)  \ar[r] &
\Spec(R[x])\\
\emptyset \ar[u] \ar[r] &
\Spec(R[x]/(J + (x)))\ar[u] \ar[r] &
\Spec(R) \ar[u]^0
}
$$
Le coin inférieur gauche est vide parce que c'est le spectre de
$R\otimes_{R[x]} B$ où l'application $R[x]\to B$ envoyer $x$ à un élément inversible et $R[x]\to R$ envoyer $x$ à $0$Depuis $g$
est surjectif cela implique $\Spec(R[x]/(J + (x)))$ est vide, comme nous voulions le montrer.

\end{proof}

\begin{lemma}

\label{lemma-integral-fibres}
Laissez $f : X \to S$ être un morphisme entier. Puis chaque point de $X$ est fermé dans sa fibre.
\end{lemma}

\begin{proof}
Voir Algèbre, lemme \ref{algebra-lemma-integral-no-inclusion}.
\end{proof}

\begin{lemma}

\label{lemma-integral-dimension}
Laissez $f : X \to Y$ être un morphisme entier. Puis $\dim(X) \leq \dim(Y)$. Si $f$ est surjectif puis $\dim(X) = \dim(Y)$.
\end{lemma}

\begin{proof}
Puisque la dimension de $X$ et $Y$ est le supremo des dimensions des membres d'un recouvrement ouvert affine, nous pouvons supposer que $Y$ et $X$ sont affines. L'inégalité suit d'Algèbre, lemme \ref{algebra-lemma-integral-dim-up}. L'égalité suit ensuite d'Algèbre, lemmes \ref{algebra-lemma-dimension-going-up} et \ref{algebra-lemma-integral-going-up}.
\end{proof}

\begin{lemma}

\label{lemma-finite-quasi-finite}

Un morphisme fini est quasi-fini.

\end{lemma}

\begin{proof}
Ceci est implicite par Algèbre, lemme \ref{algebra-lemma-quasi-finite} et lemme \ref{lemma-quasi-finite-locally-quasi-compact}. Alternativement, tous les points dans les fibres sont des points fermés par lemme \ref{lemma-integral-fibres} (et le fait qu'un morphisme fini est intégral) et utiliser lemme \ref{lemma-quasi-finite-at-point-characterize} (3) pour voir que $f$ est quasi-fini à $x$ pour tous les $x \in X$.
\end{proof}

\begin{lemma}

\label{lemma-finite-proper}
Laissez $f : X \to S$ être un morphisme de schémas. Les suivants sont équivalents
\begin{enumerate}
\item $f$ est fini, et
\item $f$ est affine et propre.
\end{enumerate}

\end{lemma}

\begin{proof}

Cela découle formellement de l'expérience acquise dans le domaine de l'éducation et de la formation. \ref{lemma-integral-universally-closed}, le fait qu'un morphisme fini est intégral et séparé, le fait qu'un morphisme propre est la même chose qu'un type fini, séparé, universellement fermé morphisme, et le fait qu'un morphisme de type fini intégral est fini (lemme \ref{lemma-finite-integral}).

\end{proof}

\begin{lemma}

\label{lemma-closed-immersion-finite}

Une fermée d'immersion est finie (et une intégrale fortiori).

\end{lemma}

\begin{proof}

C'est vrai parce qu'une immersion fermée est affine (lemme \ref{lemma-closed-immersion-affine}) et un homomorphisme surjectif d'anneaux est fini et intégral.

\end{proof}

\begin{lemma}

\label{lemma-finite-union-finite}
Laissez $X_i \to Y$, $i = 1, \ldots, n$ être des morphismes finis des schémas. Alors $X_1 \amalg \ldots \amalg X_n \to Y$ est fini aussi.
\end{lemma}

\begin{proof}
Suit du fait de l'algèbre que si $R \to A_i$, $i = 1, \ldots, n$ sont des homomorphismes finis d'anneaux, alors $R \to A_1 \times \ldots \times A_n$ est fini aussi.
\end{proof}

\begin{lemma}

\label{lemma-finite-permanence}
Laissez $f : X \to Y$ et $g : Y \to Z$ être des morphismes.
\begin{enumerate}
\item Si $g \circ f$ est fini et $g$ séparé alors $f$ est fini.
\item Si $g \circ f$ est intégral et $g$ séparé alors $f$ est intégré.
\end{enumerate}

\end{lemma}

\begin{proof}

Présumer $g \circ f$ est fini (resp.\ partie intégrante) et $g$ Le changement de base $X \times_Z Y \to Y$ est fini (resp.\ ), par lemme \ref{lemma-base-change-finite}Le morphisme $X \to X \times_Z Y$ est une immersion fermée $Y \to Z$ est séparé, voir Schémas, lemme \ref{schemes-lemma-section-immersion}. Une fermée d'immersion est finie (resp.\ ), voir l'annexe II, partie II, du règlement (CEE) n° \ref{lemma-closed-immersion-finite}. La composition du fini (resp.\ intégrale) morphismes est fini (resp.\ ), voir l'annexe II, partie II, du règlement (CEE) n° \ref{lemma-composition-finite}. Ainsi nous gagnons.

\end{proof}

\begin{lemma}

\label{lemma-finite-monomorphism-closed}
Laissez $f : X \to Y$ être un morphisme de schémas. Si $f$ est fini et un monomorphisme, alors $f$ est une immersion fermée.
\end{lemma}

\begin{proof}
Cela réduit à Algèbre, lemme \ref{algebra-lemma-finite-epimorphism-surjective}.
\end{proof}

\begin{lemma}

\label{lemma-finite-projective}

Un morphisme fini est projectif.

\end{lemma}

\begin{proof}

Laisser $f : X \to S$ être un morphisme fini. $f_*\mathcal{O}_X$ est un quasi-cohérent $\mathcal{O}_S$-module (lemme) \ref{lemma-affine-equivalence-algebras}) de type fini (par notre définition de morphismes finis et Propriétés, lemme \ref{properties-lemma-finite-type-module}). Nous prétendons qu'il y a une immersion fermée
$$
\sigma :
X
\longrightarrow
\mathbf{P}(f_*\mathcal{O}_X) =
\underline{\text{Proj}}_S(\text{Sym}^*_{\mathcal{O}_S}(f_*\mathcal{O}_X))
$$
terminé $S$, qui termine la démonstration. A savoir, nous laissons $\sigma$ soit le morphisme qui correspond (via Constructions, lemme \ref{constructions-lemma-apply-relative}) à la surjection
$$
f^*f_*\mathcal{O}_X \longrightarrow \mathcal{O}_X
$$
provenant de l'application d'adjonction $f^*f_* \to \text{id}$. Ensuite $\sigma$
est une immersion fermée. $S$ nous pouvons écrire
$X = \Spec(A)$ et $S = \Spec(R)$Depuis $X$ est fini sur $S$
nous pouvons choisir $a_1, \ldots, a_n \in A$ génératrice $A$ comme une $R$- Module. $R$- carte de l'algèbre $R[T_0, T_1, \ldots, T_n] \to \text{Sym}_R^*(A)$
envoi $T_0$ à $1$ et $T_i$ à $a_i$ pour $1 \leq i \leq n$ c'est surjectif. D'où $\mathbf{P}(f_*\mathcal{O}_X)$ est un sous-système fermé de
$\mathbf{P}^n_S$ par Constructions, lemme
\ref{constructions-lemma-surjective-graded-rings-map-proj}. Il suffit donc de prouver que le morphisme induit
$X \to \mathbf{P}^n_S$ est une immersion fermée (voir par exemple, \ref{lemma-closed-immersion-ideals}Le lecteur vérifie que $X \to \mathbf{P}^n_S$ a l'image contenue dans le $D_+(T_0) \cong \mathbf{A}^n_S$ et que $X \to \mathbf{A}^n_S$
correspond au surjet $R$- carte de l'algèbre
$R[x_1, \ldots, x_n] \to \text{Sym}_R^*(A)$ envoi $x_i$ à $a_i$. D'où $X \to \mathbf{P}^n_S$ est une immersion (comme une composition d'une immersion fermée et une ouverture d'immersion).
$X$ est fini sur $S$ l'image de cette immersion est fermée (ceci utilise lemme \ref{lemma-finite-proper}, lemme \ref{lemma-image-proper-scheme-closed}, et Constructions, \ref{constructions-lemma-projective-space-separated}). Nous concluons par Schémas, \ref{schemes-lemma-immersion-when-closed}.

\end{proof}







\section{Homéomorphismes universels}
\label{section-universal-homeomorphisms}

\noindent
La définition suivante est vraiment superflue puisque l'homéomorphisme universel n'est qu'un surjectif intégral, universel et morphiste, voir lemme \ref{lemma-universal-homeomorphism}.

\begin{definition}

\label{definition-universal-homeomorphism}
Un morphisme $f : X \to Y$ de schémas est appelé un {\it homéomorphisme universel} si le changement de base $f' : Y' \times_Y X \to Y'$ est un homéomorphisme pour chaque morphisme $Y' \to Y$.
\end{definition}

\noindent
Tout d'abord nous disons les lemmes obligatoires.

\begin{lemma}

\label{lemma-base-change-universal-homeomorphism}

Le changement de base d'un homéomorphisme universel de schémas par n'importe quel morphisme de schémas est un homéomorphisme universel.

\end{lemma}

\begin{proof}
Ceci est immédiat de la définition.
\end{proof}

\begin{lemma}

\label{lemma-composition-universal-homeomorphism}

La composition d'une paire d'homéomorphismes universels de schémas est un homéomorphisme universel.

\end{lemma}

\begin{proof}
Démonstration omise.
\end{proof}

\noindent
La simple lemme suivante est la clé pour caractériser les homéomorphismes universels.

\begin{lemma}

\label{lemma-homeomorphism-affine}
Laissez $f : X \to Y$ être un morphisme de schémas. Si $f$ est un homéomorphisme sur un sous-ensemble fermé de $Y$, alors $f$ est affiné.
\end{lemma}

\begin{proof}

Laisser $y \in Y$ être un point. Si $y \not \in f(X)$, puis il existe un quartier affiné de $y$ qui est dissociée de $f(X)$Si $y \in f(X)$, laissez $x \in X$ être le point unique de $X$ cartographie vers $y$. Laisse $y \in V$ être un quartier affiné ouvert. $U \subset X$ être un quartier affiné ouvert de $x$ dans laquelle les cartes $V$Depuis $f(U) \subset V \cap f(X)$ est ouvert dans la topologie induite par notre hypothèse sur $f$ nous pouvons choisir un
$h \in \Gamma(V, \mathcal{O}_Y)$ de telle manière que $y \in D(h)$
et $D(h) \cap f(X) \subset f(U)$. Dénotation $h' \in \Gamma(U, \mathcal{O}_X)$
la restriction de $f^\sharp(h)$ à $U$. Alors nous voyons que
$D(h') \subset U$ est égal à $f^{-1}(D(h))$. En d'autres termes, chaque point de $Y$ a un quartier ouvert dont l'image inverse est affinée. $f$ est affine, voir lemme \ref{lemma-characterize-affine}.

\end{proof}

\begin{lemma}

\label{lemma-universal-homeomorphism}

Laisser $f : X \to Y$ morphisme des régimes. Les conditions suivantes sont équivalentes:

\begin{enumerate}

\item  $f$ est un homéomorphisme universel, et
\item  $f$ est intégrale, universellement injectable et surjective.

\end{enumerate}

\end{lemma}

\begin{proof}
Supposons que $f$ est un homéomorphisme universel. Par lemme \ref{lemma-homeomorphism-affine} nous voyons que $f$ est affine. Puisque $f$ est clairement universellement fermé, nous voyons que $f$ est intégrale par lemme \ref{lemma-integral-universally-closed}. Il est également clair que $f$ est universellement injectable et surjectif.

\medskip\noindent

Présumer $f$ est intégrale, universellement injectable et surjective. \ref{lemma-integral-universally-closed}

$f$ est universellement fémé. Comme elle est aussi universellement bijective (voir lemme \ref{lemma-base-change-surjective}) nous voyons qu'il s'agit d'un homéomorphisme universel.

\end{proof}

\begin{lemma}

\label{lemma-reduction-universal-homeomorphism}
Laissez $X$ être un schéma. L'immersion canonique fermentée $X_{red} \to X$ (voir Schémas, définition \ref{schemes-definition-reduced-induced-scheme}) est un homéomorphisme universel.
\end{lemma}

\begin{proof}
Démonstration omise.
\end{proof}

\begin{lemma}

\label{lemma-check-closed-infinitesimally}
Laissez $f : X \to S$ et $S' \to S$ être des morphismes de schémas. Supposons que
\begin{enumerate}
\item $S' \to S$ est une immersion fermée,
\item $S' \to S$ est bijectif sur des points,
\item $X \times_S S' \to S'$ est une immersion fermée, et
\item $X \to S$ est de type fini ou $S' \to S$ est de présentation finie.
\end{enumerate}
Alors $f : X \to S$ est une immersion fermée.
\end{lemma}

\begin{proof}

Les hypothèses 1) et 2) impliquent que $S' \to S$ est un homéomorphisme universel (par exemple parce que $S_{red} = S'_{red}$ et en utilisant lemme \ref{lemma-reduction-universal-homeomorphism}) . Par conséquent, (3) implique que $X \to S$ est l'homéomorphisme sur un sous-ensemble fermé de $S$. Ensuite $X \to S$ est affiné par Lemme \ref{lemma-homeomorphism-affine}. Laisse $U \subset S$ être un ouvert affine, disons $U = \Spec(A)$. Ensuite $S' = \Spec(A/I)$
par (1) pour un idéal localement nilpotent $I$ par (2). $f$ est affine nous voyons que
$f^{-1}(U) = \Spec(B)$. L'hypothèse (4) nous dit $B$ est un type fini $A$-algèbre (lemme) \ref{lemma-locally-finite-type-characterize}) ou que $I$ est finiment généré (lemme \ref{lemma-closed-immersion-finite-presentation}). L'hypothèse (3) est que : $A/I \to B/IB$ de l'Algèbre, lemme \ref{algebra-lemma-surjective-mod-locally-nilpotent}
si $A \to B$ est de type fini ou Algèbre, lemme \ref{algebra-lemma-NAK} si $I$ est finiement généré et donc nilpotent nous déduisons que $A \to B$ c'est surjectif. Cela signifie que $f$ est une immersion fermée, voir lemme \ref{lemma-closed-immersion}.

\end{proof}

\begin{lemma}

\label{lemma-permanence-universal-homeomorphism}

Laisser $f : X \to Z$ soit la composition de deux morphismes 
$g : X \to Y$ et $h : Y \to Z$Si deux des morphismes $\{f, g, h\}$ sont des homéomorphismes universels, tout comme le troisième morphisme.

\end{lemma}

\begin{proof}
Si les deux $g$ et $h$ sont des homéomorphismes universels, $f$ par lemme \ref{lemma-composition-universal-homeomorphism} en est ainsi.

\medskip\noindent
Supposons que les deux $f$ et $g$ sont des homéomorphismes universels. Nous voulons montrer que $h$ est aussi. Maintenant changement de base du diagramme le long d'un morphisme arbitraire $\alpha : Z' \to Z$ des schémas, nous obtenons le diagramme suivant avec tous les carrés Cartesian: $$
\xymatrix{
X' \ar[r]^{g'} \ar[d] & Y' \ar[r]^{h'} \ar[d] & Z' \ar[d] \\
X \ar[r]^{g} & Y \ar[r]^{h} & Z.
}
$$ Notre hypothèse implique que la composition $f'= h' \circ g' : X' \to Z'$ et $g' : X' \to Y'$ sont des homéomorphismes, donc $h'$. Ceci termine la démonstration de $h$ étant un homéomorphisme universel.

\medskip\noindent

Enfin, présumons $f$ et $h$ Nous voulons montrer que $g$ est un homéomorphisme universel. $\beta : Y' \to Y$ Nous obtenons le diagramme suivant avec tous les carrés Cartesian:
$$
\xymatrix{
X' \ar[r]^{g'} \ar[d] & Y' \ar[d]^{\gamma} & \\
X'' \ar[r]^{g''} \ar[d] & Y'' \ar[r]^{h''} \ar[d] & Y' \ar[d]^{h \circ \beta} \\
X \ar[r]^{g} & Y \ar[r]^{h} & Z.
}
$$
Ici le morphisme $\gamma : Y' \to Y''$ est définie par la propriété universelle des produits fibreux et les deux morphismes
$id_{Y'} : Y' \to Y'$ et $\beta : Y' \to Y$. Nous prouverons que $g'$ puisque la propriété d'être un homéomorphisme a une propriété 2-de-3, nous voyons que $g''$ c'est un homéomorphisme. En regardant la place supérieure, il suffit de prouver que $\gamma$ est un homéomorphisme universel. $h''$ c'est un homéomorphisme, nous voyons qu'il est un morphisme affin par lemme \ref{lemma-homeomorphism-affine}
et a fortiori séparés (lemme \ref{lemma-affine-separated}Depuis $h'' \circ \gamma$ est l'identité, nous voyons que $\gamma$
est une immersion fermée de Schémas, lemme \ref{schemes-lemma-section-immersion}Depuis $h''$ est bijectif, il s'ensuit que $\gamma$ est une immersion fermée bijective et donc un homéomorphisme universel (par exemple par la caractérisation en lemme \ref{lemma-universal-homeomorphism}) comme souhaité.

\end{proof}




\section{Homéomorphismes universels de schémas affines}
\label{section-universal-homeomorphisms-affine}

\noindent
Dans cette section, nous caractérisons les homéomorphismes universels des schémas affines.

\begin{lemma}

\label{lemma-subalgebra-inherits}
Laissez $A \to B$ être un homomorphisme d'anneaux tel que le morphisme induit des schémas $f : \Spec(B) \to \Spec(A)$ est un homéomorphisme universel, rép.\ a homéomorphisme universel induisant des isomorphismes sur les corps résiduels, rép.\ universellement fermé, rép.\ universellement fermé et universellement injectable. Puis pour n'importe quel $A$-subalgèbre $B' \subset B$ la même chose est vraie pour $f' : \Spec(B') \to \Spec(A)$.
\end{lemma}

\begin{proof}
Si $f$ est universellement fermé, alors $B$ est intégrale sur $A$ par lemme \ref{lemma-integral-universally-closed}. Par conséquent, $B'$ est intégrale sur $A$ et $f'$ est universellement fermé (par le même lemme). Cela prouve le cas où $f$ est universellement fermé.

\medskip\noindent
En continuant, nous constatons que $B$ fait partie intégrante de $B'$ (Algèbre, lemme \ref{algebra-lemma-integral-permanence}) ce qui implique que $\Spec(B) \to \Spec(B')$ est surjectif (Algèbre, lemme \ref{algebra-lemma-integral-overring-surjective}). Ainsi, si $A \to B$ induit des extensions purement inséparables de corps résiduels, alors il en va de même pour $A \to B'$. Cela prouve le cas où $f$ est universelment férmé et universellement injectable, voir lemme \ref{lemma-universally-injective}.

\medskip\noindent
Le cas où $f$ est un homéomorphisme universel suit des remarques ci-dessus, lemme \ref{lemma-universal-homeomorphism}, et l'observation évidente que si $f$ est surjectif, alors $f'$ est aussi.

\medskip\noindent

Si $A \to B$ induit des isomorphismes sur les champs résiduaires, alors $A \to B'$ (voir argument au deuxième paragraphe). De cette façon, nous voyons que le lemme tient dans le cas restant.

\end{proof}

\begin{lemma}

\label{lemma-colimit-inherits}
Laissez $A$ être un anneau. Laissez $B = \colim B_\lambda$ être un filtrant colimite de $A$-algèbres. Si chaque $f_\lambda : \Spec(B_\lambda) \to \Spec(A)$ est un homéomorphisme universel, rép.\ a homéomorphisme universel induisant des isomorphismes sur les corps résiduels, rép.\ universellement fermé, rép.\ universellement fermé et universellement injectable, alors la même chose est vraie pour $f : \Spec(B) \to \Spec(A)$.
\end{lemma}

\begin{proof}
Si $f_\lambda$ est universellement fermé, alors $B_\lambda$ est intégrale sur $A$ par lemme \ref{lemma-integral-universally-closed}. Par conséquent, $B$ est intégrale sur $A$ et $f$ est universellement fermé (par le même lemme). Cela prouve le cas où chaque $f_\lambda$ est universellement fermé.

\medskip\noindent
Pour une première $\mathfrak q \subset B$ située sur $\mathfrak p \subset A$ dénotent $\mathfrak q_\lambda \subset B_\lambda$ l'image inverse. Puis $\kappa(\mathfrak q) = \colim \kappa(\mathfrak q_\lambda)$. Ainsi, si $A \to B_\lambda$ induit des extensions purement inséparables de corps résiduels, alors il en va de même pour $A \to B$. Cela prouve le cas où $f_\lambda$ est universelment férmé et universellement injectable, voir lemme \ref{lemma-universally-injective}.

\medskip\noindent
Le cas où $f$ est un homéomorphisme universel suit des remarques ci-dessus et lemme \ref{lemma-universal-homeomorphism} combiné au fait que les idéaux premiers dans $B$ sont la même chose que les séquences compatibles des idéaux premiers dans tous les $B_\lambda$.

\medskip\noindent

Si $A \to B_\lambda$ induit des isomorphismes sur les champs résiduaires, alors $A \to B$ (voir argument au deuxième paragraphe). De cette façon, nous voyons que le lemme tient dans le cas restant.

\end{proof}

\begin{lemma}

\label{lemma-special-elements-and-localization}
Laissez $A \subset B$ être une extension d'anneau. Laissez $S \subset A$ être un sous-ensemble multiplicatif. Laissez $n \geq 1$ et $b_i \in B$ pour $1 \leq i \leq n$. Tout $x \in S^{-1}B$ tel que $$
x \not \in S^{-1}A\text{ and } b_i x^i \in S^{-1}A\text{ for }i = 1, \ldots, n
$$ est égal à $s^{-1}y$ avec $s \in S$ et $y \in B$ tel que $$
y \not \in A\text{ and } b_i y^i \in A\text{ for }i = 1, \ldots, n
$$
\end{lemma}

\begin{proof}
Démonstration omise. Indication : dénominateurs clairs.
\end{proof}

\begin{lemma}

\label{lemma-nth-and-nplusone-implies-square-and-cube}
Laissez $A \subset B$ être une extension d'anneau. S'il existe $b \in B$, $b \not \in A$ et un entier $n \geq 2$ avec $b^n \in A$ et $b^{n + 1} \in A$, alors il existe un $b' \in B$, $b' \not \in A$ avec $(b')^2 \in A$ et $(b')^3 \in A$.
\end{lemma}

\begin{proof}
Laissez $b$ et $n$ être comme dans le lemme. Alors toutes les puissances suffisamment grandes de $b$ sont dans $A$. Autrement dit, $(b^n)^k(b^{n + 1})^i = b^{(k + i)n + i}$ qui implique toute puissance $b^m$ avec $m \geq n^2$ est dans $A$. Donc si $i \geq 1$ est le plus grand nombre entier tel que $b^i \not \in A$, puis $(b^i)^2 \in A$ et $(b^i)^3 \in A$.
\end{proof}

\begin{lemma}

\label{lemma-square-and-cube}
Laissez $A \subset B$ être une extension d'anneau telle que $\Spec(B) \to \Spec(A)$ est un homéomorphisme universel induisant des isomorphismes sur les corps résiduels. Si $A \not = B$, alors il existe un $b \in B$, $b \not \in A$ avec $b^2 \in A$ et $b^3 \in A$.
\end{lemma}

\begin{proof}

Rappelez-vous que $A \subset B$ fait partie intégrante (lemme \ref{lemma-integral-universally-closed}). Par lemme \ref{lemma-subalgebra-inherits}
nous pouvons supposer que $B$ est généré par un seul élément sur $A$. D'où $B$ est fini sur $A$
(Algèbre, lemme) \ref{algebra-lemma-characterize-finite-in-terms-of-integral}D'où l'appui de $B/A$ comme une $A$-module fermé et non vide (Algèbre, lemmes
\ref{algebra-lemma-support-closed} et \ref{algebra-lemma-support-zero}). Laisser $\mathfrak p \subset A$ être un premier minimum du support. Après avoir remplacé $A \subset B$ par
$A_\mathfrak p \subset B_\mathfrak p$ (permis par le lemme \ref{lemma-special-elements-and-localization}) nous pouvons supposer que $(A, \mathfrak m)$ est un anneau local, qui $B$ est fini sur $A$, et que $B/A$ bénéficie d'un soutien $\{\mathfrak m\}$
comme une $A$- Module. $B/A$ est un module fini, nous voyons que $I = \text{Ann}_A(B/A)$ satisfait $\mathfrak m = \sqrt{I}$
(Algèbre, lemme) \ref{algebra-lemma-support-closed}). Laisser $\mathfrak m' \subset B$ être l'idéal premier unique couché sur
$\mathfrak m$Parce que $\Spec(B) \to \Spec(A)$ est un homéomorphisme, nous trouvons que $\mathfrak m' = \sqrt{IB}$. Pour $f \in \mathfrak m'$ choisir $n \geq 1$ de telle manière que
$f^n \in IB$. Puis aussi $f^{n + 1} \in IB$Depuis $IB \subset A$ par notre choix de $I$ nous concluons que
$f^n, f^{n + 1} \in A$. Utilisation de lemme \ref{lemma-nth-and-nplusone-implies-square-and-cube}
nous concluons que notre lemme est vrai si $\mathfrak m' \not \subset A$. Toutefois, si $\mathfrak m' \subset A$Alors $\mathfrak m' = \mathfrak m$
et nous concluons que $A = B$ comme les corps résiduels sont aussi isomorphes par hypothèse. Cette contradiction termine la démonstration.

\end{proof}

\begin{lemma}

\label{lemma-pth-power-and-multiple}
Laissez $A \subset B$ être une extension d'anneau telle que $\Spec(B) \to \Spec(A)$ est un homéomorphisme universel. Si $A \not = B$, alors il existe un $b \in B$, $b \not \in A$ avec $b^2 \in A$ et $b^3 \in A$ ou il existe un numéro principal $p$ et un $b \in B$, $b \not \in A$ avec $pb \in A$ et $b^p \in A$.
\end{lemma}

\begin{proof}
L'argument est presque exactement le même que dans la démonstration de lemme \ref{lemma-square-and-cube} mais nous écrivons tout pour nous assurer qu'il fonctionne.

\medskip\noindent

Rappelez-vous que $A \subset B$ fait partie intégrante (lemme \ref{lemma-integral-universally-closed}). Par lemme \ref{lemma-subalgebra-inherits}
nous pouvons supposer que $B$ est généré par un seul élément sur $A$. D'où $B$ est fini sur $A$
(Algèbre, lemme) \ref{algebra-lemma-characterize-finite-in-terms-of-integral}D'où l'appui de $B/A$ comme une $A$-module fermé et non vide (Algèbre, lemmes
\ref{algebra-lemma-support-closed} et \ref{algebra-lemma-support-zero}). Laisser $\mathfrak p \subset A$ être un premier minimum du support. Après avoir remplacé $A \subset B$ par
$A_\mathfrak p \subset B_\mathfrak p$ (permis par le lemme \ref{lemma-special-elements-and-localization}) nous pouvons supposer que $(A, \mathfrak m)$ est un anneau local, qui $B$ est fini sur $A$, et que $B/A$ bénéficie d'un soutien $\{\mathfrak m\}$
comme une $A$- Module. $B/A$ est un module fini, nous voyons que $I = \text{Ann}_A(B/A)$ satisfait $\mathfrak m = \sqrt{I}$
(Algèbre, lemme) \ref{algebra-lemma-support-closed}). Laisser $\mathfrak m' \subset B$ être l'idéal premier unique couché sur
$\mathfrak m$Parce que $\Spec(B) \to \Spec(A)$ est un homéomorphisme, nous trouvons que $\mathfrak m' = \sqrt{IB}$. Pour $f \in \mathfrak m'$ choisir $n \geq 1$ de telle manière que
$f^n \in IB$. Puis aussi $f^{n + 1} \in IB$Depuis $IB \subset A$ par notre choix de $I$ nous concluons que
$f^n, f^{n + 1} \in A$. Utilisation de lemme \ref{lemma-nth-and-nplusone-implies-square-and-cube}
nous concluons que notre lemme est vrai si $\mathfrak m' \not \subset A$Si $\mathfrak m' \subset A$Alors $\mathfrak m' = \mathfrak m$Depuis $A \not = B$ nous concluons l'application
$\kappa = A/\mathfrak m \to B/\mathfrak m' = \kappa'$
des corps résiduels ne peuvent pas être un isomorphisme. \ref{lemma-universally-injective} nous concluons que la caractéristique de $\kappa$ est un nombre premier $p$
et que l'extension $\kappa'/\kappa$ est purement inséparable. $b \in B$ dont l ' image est $\kappa'$ est un élément non contenu dans $\kappa$ mais dont $p$la puissance est en $\kappa$. Ensuite $b \not \in A$, $b^p \in A$, $pb \in A$
(parce que $pb \in \mathfrak m' = \mathfrak m \subset A$) comme souhaité.

\end{proof}

\begin{proposition}

\label{proposition-universal-homeomorphism-equal-residue-fields}
Laissez $A \subset B$ être une extension d'anneau. Voici l'équivalent
\begin{enumerate}
\item $\Spec(B) \to \Spec(A)$ est un homéomorphisme universel induisant des isomorphismes sur les corps résiduels, et
\item chaque sous-ensemble fini $E \subset B$ est contenu dans une extension $$
A[b_1, \ldots, b_n] \subset B
$$ telle que $b_i^2, b_i^3 \in A[b_1, \ldots, b_{i - 1}]$ pour $i = 1, \ldots, n$.
\end{enumerate}

\end{proposition}

\begin{proof}

Supposons que (1). À l'aide de la récursion transfinite nous construisons pour chaque ordinal $\alpha$
et $A$- sous-algèbre $B_\alpha \subset B$ comme suit. Définir $B_0 = A$Si $\alpha$
est une limite ordinaire, alors nous configurons $B_\alpha = \colim_{\beta < \alpha} B_\beta$Si $\alpha = \beta + 1$, alors soit
$B_\beta = B$ Dans ce cas, nous avons mis $B_\alpha = B_\beta$ ou
$B_\beta \not = B$, auquel cas nous appliquons \ref{lemma-square-and-cube}
pour choisir un $b_\alpha \in B$, $b_\alpha \not \in B_\beta$ avec
$b_\alpha^2, b_\alpha^3 \in B_\beta$
et nous nous sommes mis $B_\alpha = B_\beta[b_\alpha] \subset B$. De toute évidence, $B = \colim B_\alpha$ (en fait $B = B_\alpha$ pour un peu ordinaire $\alpha$ comme on le voit en regardant les cardinalités. Nous prouverons, par induction transfinite, que (2) tient pour
$A \to B_\alpha$ pour chaque ordinal $\alpha$. Il est clair pour
$\alpha = 0$. Supposons que l'énoncé tient pour chaque $\beta < \alpha$
et laisser $E \subset B_\alpha$ être un sous-ensemble fini. Si $\alpha$ est une limite ordinaire, alors
$B_\alpha = \bigcup_{\beta < \alpha} B_\beta$ et nous voyons que $E \subset B_\beta$ pour certains $\beta < \alpha$ qui prouve le résultat dans ce cas. $\alpha = \beta + 1$Alors $B_\alpha = B_\beta[b_\alpha]$. Ainsi toute $e \in E$
peut être écrit comme un polynôme $e = \sum d_{e, i}b_\alpha^i$
avec $d_{e, i} \in B_\beta$. Laisse $D \subset B_\beta$
être l'ensemble $D = \{d_{e, i}\} \cup \{b_\alpha^2, b_\alpha^3\}$. Par hypothèse d'induction, il existe un $A$- sous-algèbre $A[b_1, \ldots, b_n] \subset B_\beta$
comme dans l'énoncé du lemme contenant $D$. Ensuite $A[b_1, \ldots, b_n, b_\alpha] \subset B_\alpha$
est un $A$-subalgèbre de $B_\alpha$ comme dans l'énoncé du lemme contenant $E$.

\medskip\noindent

Supposons que (2). Écrire $B = \colim B_\lambda$ comme le colimite de son fini $A$- des sous-algèbres. \ref{lemma-colimit-inherits} il suffit de montrer que $\Spec(B_\lambda) \to \Spec(A)$ est un homéomorphisme universel induisant des isomorphismes sur les corps résiduels. Les compositions d'universellement fermé morphismes sont universellement fermé et la même chose pour les morphismes qui induisent des isomorphismes sur les corps résiduels. $A \subset B$ et $B$ est générée par un seul élément $b$ avec $b^2, b^3 \in A$Une telle extension fait partie intégrante et, par conséquent, $\Spec(B) \to \Spec(A)$ est universellement fémé et surjectif (lemme \ref{lemma-integral-universally-closed} et d ' Algèbre, lemme \ref{algebra-lemma-integral-overring-surjective}). Notez que $(b^2)^3 = (b^3)^2$ en $A$. Pour tout homomorphisme d'anneaux $\varphi : A \to K$ vers un champ $K$ nous voyons qu'il existe un $\lambda \in K$ avec $\varphi(b^2) = \lambda^2$ et
$\varphi(b^3) = \lambda^3$. A savoir, $\lambda = 0$ si $\varphi(b^2) = 0$
et $\lambda = \varphi(b^3)/\varphi(b^2)$ si non. Ainsi
$B \otimes_A K$ est un quotient de $K[x]/(x^2 - \lambda^2, x^3 - \lambda^3)$. Cette bague a exactement un premier avec corps résiduel $K$. Cela implique que $\Spec(B) \to \Spec(A)$ est bijectif et induit des isomorphismes sur les corps résiduels. \ref{lemma-universal-homeomorphism}
et \ref{lemma-universally-injective}.

\end{proof}

\begin{proposition}

\label{proposition-universal-homeomorphism}
Laissez $A \subset B$ être une extension d'anneau. Les conditions suivantes sont équivalentes
\begin{enumerate}
\item $\Spec(B) \to \Spec(A)$ est un homéomorphisme universel, et
\item chaque sous-ensemble fini $E \subset B$ est contenu dans une extension $$
A[b_1, \ldots, b_n] \subset B
$$ telle que pour $i = 1, \ldots, n$ nous avons
\begin{enumerate}
\item $b_i^2, b_i^3 \in A[b_1, \ldots, b_{i - 1}]$, ou
\item il existe un numéro premier $p$ avec $pb_i, b_i^p \in A[b_1, \ldots, b_{i - 1}]$.
\end{enumerate}

\end{enumerate}

\end{proposition}

\begin{proof}

La démonstration est exactement la même que la démonstration de la proposition \ref{proposition-universal-homeomorphism-equal-residue-fields}
sauf pour les modifications suivantes:

\begin{enumerate}

\item Utiliser lemme \ref{lemma-pth-power-and-multiple} au lieu de lemme \ref{lemma-square-and-cube} ce qui signifie que pour chaque successeur ordinal $\alpha = \beta + 1$ Nous avons soit
$b_\alpha^2, b_\alpha^3 \in B_\beta$ ou nous avons un premier $p$ et
$pb_\alpha, b_\alpha^p \in B_\beta$.
\item Si $\alpha$ est un successeur ordinaire, puis prendre
$D = \{d_{e, i}\} \cup \{b_\alpha^2, b_\alpha^3\}$ ou prendre $D = \{d_{e, i}\} \cup \{pb_\alpha, b_\alpha^p\}$ selon le cas $\alpha$ Elle tombe dans la tête.
\item Dans la démonstration de (2) $\Rightarrow$ (1) nous devons également examiner le cas où $B$ est généré sur $A$ par un seul élément $b$
avec $pb, b^p \in B$ pour un certain nombre principal $p$. Ici $A \subset B$
induit un homéomorphisme universel sur les spectres par exemple par Algèbre, lemme \ref{algebra-lemma-p-ring-map}.

\end{enumerate}

Ceci achève la démonstration.

\end{proof}

\begin{lemma}

\label{lemma-universal-homeo-iso-if-invert-p}

Laisser $p$ être un numéro principal. $A \to B$ être un homomorphisme d'anneaux qui induit un isomorphisme $A[1/p] \to B[1/p]$
(par exemple si: $p$ est nilpotentiel dans $A$). Les conditions suivantes sont équivalentes

\begin{enumerate}

\item  $\Spec(B) \to \Spec(A)$ est un homéomorphisme universel, et
\item le noyau de $A \to B$ est un idéal localement nilpotent et pour chaque $b \in B$ il existe un $p$- puissance $q$ avec $qb$ et $b^q$
à l'image de $A \to B$.

\end{enumerate}

\end{lemma}

\begin{proof}

Si (2) tient, alors (1) tient par Algèbre, lemme \ref{algebra-lemma-p-ring-map}. Supposons que (1). Puis le noyau de $A \to B$ se compose d'éléments nilpotents par Algèbre, lemme \ref{algebra-lemma-image-dense-generic-points}. Ainsi, nous pouvons remplacer $A$ par l'image de $A \to B$ et supposer que
$A \subset B$. Par Algèbre, lemme \ref{algebra-lemma-help-with-powers}
l'ensemble
$$
B' = \{b \in B \mid p^nb, b^{p^n} \in A\text{ for some }n \geq 0\}
$$
est un $A$-subalgèbre de $B$ (être fermé sous les produits est trivial). $B' = B$Si ce n'est pas le cas, alors d'après moi \ref{lemma-pth-power-and-multiple}
il existe un $b \in B$, $b \not \in B'$ avec l'un ou l'autre
$b^2, b^3 \in B'$ ou il existe un numéro principal $\ell$
avec $\ell b, b^\ell \in B'$.- Nous montrerons que les deux cas conduisent à une contradiction, prouvant ainsi le lemme.

\medskip\noindent
Depuis $A[1/p] = B[1/p]$, nous pouvons choisir un $p$ de puissance $q$ tel que $qb \in A$.

\medskip\noindent
Si $b^2, b^3 \in B'$ alors aussi $b^q \in B'$. Par définition de $B'$ nous trouvons que $(b^q)^{q'} \in A$ pour certains $p$ puissance $q'$. Puis $qq'b, b^{qq'} \in A$ d'où $b \in B'$ ce qui est une contradiction.

\medskip\noindent
Supposons qu'il existe maintenant un numéro principal $\ell$ avec $\ell b, b^\ell \in B'$. Si $\ell \not = p$ alors $\ell b \in B'$ et $qb \in A \subset B'$ impliquent $b \in B'$ une contradiction. Ainsi $\ell = p$ et $b^p \in B'$ et nous obtenons une contradiction exactement comme avant.
\end{proof}

\begin{lemma}

\label{lemma-make-universal-homeo}
Laissez $A$ être un anneau. Laissez $x, y \in A$.
\begin{enumerate}
\item Si $x^3 = y^2$ dans $A$, alors $A \to B = A[t]/(t^2 - x, t^3 - y)$ induit des bijections sur les corps résiduels et un homéomorphisme universel sur les spectres.
\item S'il y a un numéro initial $p$ tel que $p^px = y^p$ dans $A$, alors $A \to B = A[t]/(t^p - x, pt - y)$ induit un homeomorphisme universel sur les spectres.
\end{enumerate}

\end{lemma}

\begin{proof}
Nous allons utiliser le critère de lemme \ref{lemma-universal-homeomorphism} pour vérifier cela. Dans les deux cas, l'homomorphisme d'anneaux est intégral. Ainsi, il suffit de montrer que, compte tenu d'un champ $k$ et d'un homomorphisme d'anneaux $\varphi : A \to k$, le $k$-algebra $B \otimes_A k$ a un idéal premier unique dont le corps résiduel est égal à $k$ dans le cas (1) et purement inséparable sur $k$ dans le cas (2). Voir lemme \ref{lemma-universally-injective}.

\medskip\noindent
Dans le cas (1) définir $\lambda = 0$ si $\varphi(x) = 0$ et définir $\lambda = \varphi(y)/\varphi(x)$ si non. Puis $B = k[t]/(t^2 - \lambda^2, t^3 - \lambda^2)$. Ainsi le résultat est clair.

\medskip\noindent
Dans le cas (2) si la caractéristique de $k$ est $p$, alors nous obtenons $\varphi(y) = 0$ et $B = k[t]/(t^p - \varphi(x))$ qui est une algèbre artiniane locale $k$ dont le corps résiduel est soit $k$ ou un degré $p$ extension purement inséparable de $k$. Si la caractéristique de $k$ n'est pas $p$, alors la configuration $\lambda = \varphi(y)/p$ nous voyons $B = k[t]/(t - \lambda) = k$ et nous concluons aussi.
\end{proof}

\begin{lemma}

\label{lemma-universal-homeo-limit}

Laisser $A \to B$ être un homomorphisme d'anneaux.

\begin{enumerate}

\item Si $A \to B$ induit un homéomorphisme universel sur les spectres, puis $B = \colim B_i$ est une colimite filtrante de présentation finie $A$-Algèbres
$B_i$ de telle manière que $A \to B_i$ induit un homéomorphisme universel sur les spectres.
\item Si $A \to B$ induit des isomorphismes sur les corps résiduels et un homéomorphisme universel sur les spectres, puis $B = \colim B_i$ est une colimite filtrante de présentation finie $A$-Algèbres $B_i$ de telle manière que
$A \to B_i$ induit des isomorphismes sur les corps résiduels et un homéomorphisme universel sur les spectres.

\end{enumerate}

\end{lemma}

\begin{proof}

la démonstration de (1). Nous utiliserons le critère d'Algèbre, lemme \ref{algebra-lemma-when-colimit}. Laisse $A \to C$ être de présentation finie et laisser $\varphi : C \to B$
être un $A$- Carte de l'algèbre. $B' = \varphi(C) \subset B$ être l'image. Alors $A \to B'$ induit un homéomorphisme universel sur les spectres par lemme \ref{lemma-subalgebra-inherits}. Par Algèbre, lemme \ref{algebra-lemma-ring-colimit-fp}
nous pouvons écrire $B' = \colim_{i \in I} B_i$
avec $A \to B_i$ de présentation finie et cartes de transition surjective. Par Algèbre, lemme \ref{algebra-lemma-characterize-finite-presentation}
nous pouvons choisir un index $0 \in I$
et une factorisation $C \to B_0 \to B'$ de l'application $C \to B'$. Nous prétendons que $\Spec(B_i) \to \Spec(A)$ est un homéomorphisme universel pour $i$ La revendication termine la démonstration de (1).

\medskip\noindent

la démonstration de la revendication. \ref{lemma-reduction-universal-homeomorphism}
l'homomorphisme d'anneaux $A_{red} \to B'_{red}$ induit un homéomorphisme universel sur les spectres. $A_{red} \subset B'_{red}$ par Algèbre, lemme \ref{algebra-lemma-image-dense-generic-points}. Réglage $A' = \Im(A \to B')$ Nous avons des surjections $A \to A' \to A_{red}$
induction de bijections $\Spec(A_{red}) = \Spec(A') = \Spec(A)$. Ainsi $A' \subset B'$ induit un homéomorphisme universel sur les spectres. \ref{proposition-universal-homeomorphism}
et le fait que $B'$ est le type fini sur $A'$ Nous pouvons trouver $n$ et $b'_1, \ldots, b'_n \in B'$ de telle manière que
$B' = A'[b'_1, \ldots, b'_n]$
et de telle sorte que pour $j = 1, \ldots, n$ Nous avons

\begin{enumerate}

\item  $(b'_j)^2, (b'_j)^3 \in A'[b'_1, \ldots, b'_{j - 1}]$, ou
\item il existe un nombre premier $p$ avec
$pb'_j, (b'_j)^p \in A'[b'_1, \ldots, b'_{j - 1}]$.

\end{enumerate}

Choisissez $b_1, \ldots, b_n \in B_0$ levage $b'_1, \ldots, b'_n$. Pour $i \geq 0$ indiquer $b_{j, i}$ l'image de $b_j$ en $B_i$. Pour assez grand $i$ nous aurons pour $j = 1, \ldots, n$

\begin{enumerate}

\item  $b_{j, i}^2, b_{j, i}^3 \in A_i[b_{1, i}, \ldots, b_{j - 1, i}]$, ou
\item il existe un nombre premier $p$ avec
$pb_{j, i}, b_{j, i}^p \in A_i[b_{1, i}, \ldots, b_{j - 1, i}]$.

\end{enumerate}

Ici. $A_i \subset B_i$ est l'image de $A \to B_i$. Observer que $A \to A_i$
est un homomorphisme surjectif d'anneaux dont le noyau est un idéal localement nilpotent. $i$ plus si nécessaire, nous pouvons supposer $B_i$
est générée par $b_1, \ldots, b_n$ terminé $A_i$, en d'autres termes
$B_i = A_i[b_1, \ldots, b_n]$. Par Algèbre, lemmes \ref{algebra-lemma-p-ring-map} et
\ref{algebra-lemma-2-3-ring-map} nous concluons que
$A \to A_i \to A_i[b_1] \to \ldots \to A_i[b_1, \ldots, b_n] = B_i$
induire des homéomorphismes universels sur les spectres. Ceci termine la démonstration de la revendication.

\medskip\noindent

La preuve de (2) est exactement la même.

\end{proof}







\section{Normalisation faible absolue et semi-normalisation}
\label{section-seminormalization}

\noindent
Motivée par les résultats prouvés dans la section précédente, nous donnons la définition suivante.

\begin{definition}

\label{definition-seminormal-ring}

Laisser $A$ Sois une bague.

\begin{enumerate}

\item Nous disons $A$ est {\it semi-normal} si pour tous $x, y \in A$
avec $x^3 = y^2$ il y a un unique $a \in A$ avec
$x = a^2$ et $y = a^3$.
\item Nous disons $A$ est {\it Absolument normale} si a) $A$ est semi-normale et (b) pour tout nombre principal $p$ et $x, y \in A$ avec $p^px = y^p$
il y a un unique $a \in A$ avec $x = a^p$ et $y = pa$.

\end{enumerate}

\end{definition}

\noindent

Une observation amusante, voir \cite{Costa}, est que dans la définition des anneaux semi-normaux il suffit\footnote{Laisser
$A$ être une bague telle que pour tous $x, y \in A$ avec $x^3 = y^2$ il y a un $a \in A$ avec $x = a^2$ et $y = a^3$. Ensuite $A$ est réduite: si
$x^2 = 0$Alors $x^2 = x^3$ et donc il existe un $a$ de telle manière que
$x = a^3$ et $x = a^2$. Ensuite $x = a^3 = ax = a^4 = x^2 = 0$Enfin, si $a_1^2 = a_2^2$ et $a_1^3 = a_2^3$ pour $a_1, a_2$ dans un anneau réduit, puis $(a_1 - a_2)^3 = a_1^3 - 3a_1^2a_2 +3a_1a_2^2 - a_2^3 =
(1 - 3 + 3 - 1)a_1^3 = 0$ et donc $a_1 = a_2$.} à supposer l'existence de $a$. Des régimes absolument peu normaux ont été définis dans
\cite[Annexe B]{rydh_descent}.

\begin{lemma}

\label{lemma-seminormal-local-property}
Être semi-normal ou être absolument normal est une propriété locale des anneaux, voir Propriétés, définition \ref{properties-definition-property-local}.
\end{lemma}

\begin{proof}

Supposons que $A$ est semi-normale et $f \in A$. Laisse $x', y' \in A_f$
avec $(x')^3 = (y')^2$. Écrire $x' = x/f^{2n}$ et $y' = y/f^{3n}$
pour certains $n \geq 0$ et $x, y \in A$. Après remplacement $x, y$
par $f^{2m}x, f^{3m}y$ et $n$ par $n + m$, nous voyons que
$x^3 = y^2$ en $A$. Ensuite, nous trouvons un unique $a \in A$ avec
$x = a^2$ et $y = a^3$. Réglage $a' = a/f^n$ Nous avons
$x' = (a')^2$ et $y' = (a')^3$ comme souhaité. Uniqueté $a'$ résulte de l'unicité de $a$. De la même manière que le lecteur montre que si $A$ est absolument faiblement normal, alors $A_f$ C'est tout à fait normal.

\medskip\noindent

Présumer $A$ est une bague et $f_1, \ldots, f_n \in A$ générer l'unité idéale. Supposons que $A_{f_i}$ est semi-normale pour chaque $i$. Laisse $x, y \in A$ avec $x^3 = y^2$. Pour chaque $i$ nous trouvons un unique
$a_i \in A_{f_i}$ avec $x = a_i^2$ et $y = a_i^3$ en $A_{f_i}$. Par l'unicité et le résultat du premier paragraphe (qui nous dit que $A_{f_if_j}$ est semi-normale) nous voyons que
$a_i$ et $a_j$ cartographie au même élément de $A_{f_if_j}$. Par Algèbre, lemme \ref{algebra-lemma-standard-covering}
nous trouvons un unique $a \in A$
cartographie vers $a_i$ en $A_{f_i}$ pour tous $i$. Ensuite $x = a^2$ et $y = a^3$ de la même façon. $a$ c'est unique. $A$ C'est semi-normal. $A_{f_i}$ est absolument faiblement normal, alors le même argument montre que $A$ C'est tout à fait normal.

\end{proof}

\noindent
Ensuite nous définissons des schémas semi-normales et des schémas absolument faibles.

\begin{definition}

\label{definition-seminormal}
Laissez $X$ être un schéma.
\begin{enumerate}
\item Nous disons que $X$ est {\it semi-normal} si chaque $x \in X$ a un quartier affiné à l'ouverture $\Spec(R) = U \subset X$ de telle sorte que l'anneau $R$ est semi-normal.
\item Nous disons que $X$ est {\it absolument normal} si chaque $x \in X$ a un quartier affiné à l'ouverture $\Spec(R) = U \subset X$ de telle sorte que l'anneau $R$ est absolument normal.
\end{enumerate}

\end{definition}

\noindent
Voici le lemme obligatoire.

\begin{lemma}

\label{lemma-locally-seminormal}

Laisser $X$ Les conditions suivantes sont équivalentes:

\begin{enumerate}

\item Le régime $X$ est semi-normale.
\item Pour chaque ouvert affine $U \subset X$ la bague $\mathcal{O}_X(U)$
est semi-normale.
\item Il existe un recouvrement ouvert affine $X = \bigcup U_i$ de telle sorte que chacun $\mathcal{O}_X(U_i)$ est semi-normale.
\item Il existe une ouverture de recouvrement $X = \bigcup X_j$
de telle sorte que chaque sous-système ouvert $X_j$ est semi-normale.

\end{enumerate}

En outre, si $X$ est semi-normale puis chaque sous-système ouvert est semi-normal. Les mêmes affirmations sont vraies avec ``semi-normal'' remplacé par ``absolument faiblement normal''.

\end{lemma}

\begin{proof}

Combiner les propriétés, lemme \ref{properties-lemma-locally-P} et lemme \ref{lemma-seminormal-local-property}.

\end{proof}

\begin{lemma}

\label{lemma-seminormal-reduced}

Un schéma ou un anneau semi-normal est réduit. A fortiori le même est vrai pour les schémas ou les anneaux absolument faiblement normaux.

\end{lemma}

\begin{proof}
Laissez $A$ être un anneau. Si $a \in A$ est non zéro mais $a^2 = 0$, alors $a^2 = 0^2$ et $a^3 = 0^3$ et donc $A$ n'est pas semi-normal.
\end{proof}

\begin{lemma}

\label{lemma-seminormalization-ring}

Laisser $A$ Sois une bague.

\begin{enumerate}

\item La catégorie des homomorphismes d'anneaux $A \to B$ induire un homéomorphisme universel sur les spectres a un objet final $A \to A^{awn}$.
\item Données $A \to B$ dans la catégorie de 1) l'application résultante
$B \to A^{awn}$ est un isomorphisme si et seulement si $B$ C'est tout à fait normal.
\item La catégorie des homomorphismes d'anneaux $A \to B$ induire des isomorphismes sur les corps résiduels et un homéomorphisme universel sur les spectres a un objet final $A \to A^{sn}$.
\item Données $A \to B$ dans la catégorie de (3) l'application résultante
$B \to A^{sn}$ est un isomorphisme si et seulement si $B$ est semi-normale.

\end{enumerate}

Pour tout homomorphisme d'anneaux $\varphi : A \to A'$ il y a des cartes uniques
$\varphi^{awn} : A^{awn} \to (A')^{awn}$ et
$\varphi^{sn} : A^{sn} \to (A')^{sn}$ compatible avec $\varphi$.

\end{lemma}

\begin{proof}

Nous prouvons (1) et (2) et nous omettons la démonstration de (3) et (4) et l'énoncé finale. $A$-algèbres de la forme
$$
B = A[x_1, \ldots, x_n]/J
$$
où $J$ est un idéal finiment généré tel que $A \to B$ définit un homéomorphisme universel sur les spectres. Nous affirmons que cette catégorie est dirigée (Catégories, définition \ref{categories-definition-directed}). A savoir, donné
$$
B = A[x_1, \ldots, x_n]/J
\quad\text{and}\quad
B' = A[x_1, \ldots, x_{n'}]/J'
$$
alors nous pouvons envisager
$$
B'' = A[x_1, \ldots, x_{n + n'}]/J''
$$
où $J''$ est générée par les éléments de $J$ et les éléments
$f(x_{n + 1}, \ldots, x_{n + n'})$ où $f \in J'$. Alors nous avons $A$- homomorphismes de l'algèbre $B \to B''$ et $B' \to B''$
qui induisent un isomorphisme $B \otimes_A B' \to B''$. Il s'ensuit de lemmes \ref{lemma-base-change-universal-homeomorphism} et
\ref{lemma-composition-universal-homeomorphism}
que $\Spec(B'') \to \Spec(A)$ est un homéomorphisme universel et donc $A \to B''$ est dans notre catégorie. $\varphi, \varphi' : B \to B'$
dans notre catégorie avec $B$ comme indiqué ci-dessus, alors nous considérons le quotient $B''$ de $B'$
par l'idéal généré par $\varphi(x_i) - \varphi'(x_i)$, $i = 1, \ldots, n$Depuis $\Spec(B') = \Spec(B)$ nous voyons que $\Spec(B'') \to \Spec(B')$
est une immersion bijective fermée d'où un homéomorphisme universel. $B''$ est dans notre catégorie et $\varphi, \varphi'$ sont égalisés par $B' \to B''$. Cela complète la démonstration de notre réclamation.
$$
A^{awn} = \colim B
$$
où le colimite est supérieur à la catégorie décrite ci-dessus.
$A \to A^{awn}$ induit un homéomorphisme universel sur les spectres par lemme \ref{lemma-colimit-inherits} (c'est là que nous utilisons la catégorie est dirigée).

\medskip\noindent
Avec un homomorphisme d'anneaux $A \to B$ de présentation finie induisant un homéomorphisme universel sur les spectres, nous obtenons une application canonique $B \to A^{awn}$ par la construction même de $A^{awn}$. Puisque chaque $A \to B$ comme dans (1) est une filtrante de colimite de $A \to B$ comme dans (1) de présentation finie (lemme \ref{lemma-universal-homeo-limit}), nous voyons que $A \to A^{awn}$ est final dans la catégorie (1).

\medskip\noindent

Laisser $x, y \in A^{awn}$ être des éléments tels que $x^3 = y^2$. Ensuite $A^{awn} \to A^{awn}[t]/(t^2 - x, t^3 - y)$ induit un homéomorphisme universel sur les spectres par lemme \ref{lemma-make-universal-homeo}. Ainsi $A \to A^{awn}[t]/(t^2 - x, t^3 - y)$ est dans la catégorie (1) et nous obtenons un unique $A$- carte de l'algèbre
$A^{awn}[t]/(t^2 - x, t^3 - y) \to A^{awn}$. L'image $a \in A^{awn}$ de $t$ est donc l'élément unique tel que $a^2 = x$ et $a^3 = y$ en $A^{awn}$. De la même manière, donné un premier $p$ et $x, y \in A^{awn}$
avec $p^px = y^p$ nous trouvons un unique $a \in A^{awn}$ avec
$a^p = x$ et $pq = y$. Ainsi $A^{awn}$ est par définition tout à fait normale.

\medskip\noindent

Enfin, laissons $A \to B$ être dans la catégorie (1) avec $B$ Tout à fait normale. $A^{awn} \to B^{awn}$ induit un homéomorphisme universel sur les spectres et depuis $A^{awn}$ est réduit (lemme \ref{lemma-seminormal-reduced}) nous trouvons
$A^{awn} \subset B^{awn}$
(voir Algèbre, lemme) \ref{algebra-lemma-image-dense-generic-points}Si cette inclusion n'est pas une égalité, alors lemme \ref{lemma-pth-power-and-multiple}
implique qu'il y a un élément $b \in B^{awn}$, $b \not \in A^{awn}$
tel que l'un ou l'autre $b^2, b^3 \in A^{awn}$ ou $pb, b^p \in A^{awn}$
pour un certain nombre principal $p$. Cependant, par l'existence et l'unicité dans la définition \ref{definition-seminormal-ring} cette force $b \in A^{awn}$
et donc nous obtenons la contradiction qui finit la démonstration.

\end{proof}

\begin{lemma}

\label{lemma-seminormalization}

Laisser $X$ être un régime.

\begin{enumerate}

\item La catégorie des homéomorphismes universels $Y \to X$ a un objet initial $X^{awn} \to X$.
\item Données $Y \to X$ dans la catégorie de (1) le morphisme résultant
$X^{awn} \to Y$ est un isomorphisme si et seulement si $Y$ C'est tout à fait normal.
\item La catégorie des homéomorphismes universels $Y \to X$ qui induisent des ismomorphismes sur les corps résiduels a un objet initial
$X^{sn} \to X$.
\item Données $Y \to X$ dans la catégorie de (3) le morphisme résultant
$X^{sn} \to Y$ est un isomorphisme si et seulement si $Y$ est semi-normale.

\end{enumerate}

Pour tout morphisme $h : X' \to X$ des schémas il y a des morphismes uniques
$h^{awn} : (X')^{awn} \to X^{awn}$ et $h^{sn} : (X')^{sn} \to X^{sn}$
compatible avec $h$.

\end{lemma}

\begin{proof}
Nous allons prouver (1) et (2) et omettre la démonstration de (3) et (4).Soit $h : X' \to X$ un morphisme de schémas. Si (1) tient pour $X$ et $X'$, alors $X' \times_X X^{awn} \to X'$ est un homéomorphisme universel et donc nous obtenons un morphisme unique $(X')^{awn} \to X' \times_X X^{awn}$ sur $X'$ par la propriété universelle de $(X')^{awn} \to X'$. Composé avec la projection $X' \times_X X^{awn} \to X^{awn}$ nous obtenons $h^{awn}$. Si en plus (2) tient pour $X$ et $X'$ et $h$ est une immersion ouverte, alors $X' \times_X X^{awn}$ est absolument faiblement normal (lemme \ref{lemma-locally-seminormal}) et nous déduisons que $(X')^{awn} \to X' \times_X X^{awn}$ est un isomorphisme.

\medskip\noindent

Rappelez-vous que n'importe quel homéomorphisme universel est affiné, voyez lemme \ref{lemma-homeomorphism-affine}. Ainsi si $X$ est affine puis (1) et (2) suivre immédiatement de lemme \ref{lemma-seminormalization-ring}. Laisse $X$ être un plan et laisser $\mathcal{B}$ soit l'ensemble des ouverts affines de $X$. Pour chaque $U \in \mathcal{B}$ nous obtenons $U^{awn} \to U$
et pour $V \subset U$, $V, U \in \mathcal{B}$ nous obtenons un isomorphisme canonique $\rho_{V, U} : V^{awn} \to V \times_U U^{awn}$ Ainsi, par collage relatif (Constructions, \ref{constructions-lemma-relative-glueing}) nous obtenons un morphisme $X^{awn} \to X$ qui se limite $U^{awn}$ terminé $U$ Compatible avec le $\rho_{V, U}$. Suivant, laissez $Y \to X$ être un homéomorphisme universel. $U \times_X Y \to U$ est un homéomorphisme universel pour $U \in \mathcal{B}$
et nous obtenons un morphisme unique $g_U : U^{awn} \to U \times_X Y$ terminé $U$. Ces $g_U$ sont compatibles avec les morphismes $\rho_{V, U}$; détails omis. Il y a donc un morphisme unique $g : X^{awn} \to Y$
terminé $X$ d ' accord avec $g_U$ terminé $U$, voir Constructions, remarque \ref{constructions-remark-relative-glueing-functorial}. Cela prouve (1) pour $X$. La partie 2 suit parce qu'elle est affinée localement.

\end{proof}

\begin{definition}

\label{definition-seminormalization}
Laissez $X$ être un schéma.
\begin{enumerate}
\item Le morphisme $X^{sn} \to X$ construit en lemme \ref{lemma-seminormalization} est la semi-normalisation {\it} de $X$.
\item Le morphisme $X^{awn} \to X$ construit en lemme \ref{lemma-seminormalization} est la {\it normalisationabsolue} de $X$.
\end{enumerate}

\end{definition}

\noindent
Pour être sûr, la semi-normalisation $X^{sn}$ de $X$ est un schéma semi-normal et la normalisation facile absolu $X^{awn}$ est un schéma absolument faiblement normal. De plus, pour tout morphisme $h : Y \to X$ de schémas nous obtenons un diagramme commutatif canonique $$
\xymatrix{
Y^{awn} \ar[d]^{h^{awn}} \ar[r] &
Y^{sn} \ar[d]^{h^{sn}} \ar[r] &
Y \ar[d]^h \\
X^{awn} \ar[r] &
X^{sn} \ar[r] &
X
}
$$ de schémas; les flèches $h^{sn}$ et $h^{awn}$ sont les uniques compatibles avec $h$.

\begin{lemma}

\label{lemma-characterize-seminormal}
Laissez $X$ être un schéma. Ce qui suit sont équivalents
\begin{enumerate}
\item $X$ est semi-normal,
\item $X$ est égal à sa propre semi-normalisation, c'est-à-dire que le morphisme $X^{sn} \to X$ est un isomorphisme,
\item si $\pi : Y \to X$ est un homéomorphisme universel induisant des isomorphismes sur les corps résiduels avec $Y$ réduit, alors $\pi$ est un isomorphisme.
\end{enumerate}

\end{lemma}

\begin{proof}
L'équivalence de (1) et (2) est claire de lemme \ref{lemma-seminormalization}. Si (3) tient, alors $X^{sn} \to X$ est un isomorphisme et nous voyons que (2) tient.

\medskip\noindent

Supposons que (2) tient et laissent $\pi : Y \to X$ être un homéomorphisme universel induisant des isomorphismes sur les corps résiduels avec $Y$ réduit. Ensuite, il existe une factorisation $X \to Y \to X$ de $\text{id}_X$ par lemme \ref{lemma-seminormalization}. Ensuite $X \to Y$ est une immersion fermée (par Schémas, lemme \ref{schemes-lemma-section-immersion}
et le fait que $\pi$ est séparé par exemple par lemme \ref{lemma-universally-injective-separated}Depuis $X \to Y$
est également un bijection sur les points, la réduction de $Y$ montre qu'il doit s'agir d'un isomorphisme. Cela finit par la démonstration.

\end{proof}

\begin{lemma}

\label{lemma-characterize-absolutely-weakly-normal}
Laissez $X$ être un schéma. Les conditions suivantes sont équivalentes à
\begin{enumerate}
\item $X$ est absolument faiblement normal,
\item $X$ est égal à sa propre normalisation absolue facile, c'est-à-dire que le morphisme $X^{awn} \to X$ est un isomorphisme,
\item si $\pi : Y \to X$ est un homéomorphisme universel avec $Y$ réduit, alors $\pi$ est un isomorphisme.
\end{enumerate}

\end{lemma}

\begin{proof}
Cela est prouvé exactement de la même manière que lemme \ref{lemma-characterize-seminormal}.
\end{proof}





\section{Morphismes finis localement libres}
\label{section-finite-locally-free}

\noindent
Dans de nombreux articles, les auteurs utilisent des morphismes finis plats quand ils signifient vraiment morphismes finis localement libres. La raison est que si la base est localement Noétherian alors c'est la même chose. Mais en général ce n'est pas, voir Exercices, Exercice \ref{exercises-exercise-flat-not-projective}.

\begin{definition}

\label{definition-finite-locally-free}
Laissez $f : X \to S$ être un morphisme de schémas. Nous disons que $f$ est {\it finit localement libre} si $f$ est affine et $f_*\mathcal{O}_X$ est un $\mathcal{O}_S$ fini localement libre-module. Dans ce cas, nous disons que $f$ est {\it rang} ou {\it degré} $d$ si la feuille $f_*\mathcal{O}_X$ est finie localement libre de degré $d$.
\end{definition}

\noindent
Notez que si $f : X \to S$ est fini localement libre alors $S$ est l'union disjointe des sous-systèmes ouverts et fermés $S_d$ de telle sorte que $f^{-1}(S_d) \to S_d$ est fini localement libre de degré $d$.

\begin{lemma}

\label{lemma-finite-flat}
Laissez $f : X \to S$ être un morphisme de schémas. Les conditions suivantes sont équivalentes:
\begin{enumerate}
\item $f$ est fini localement libre,
\item $f$ est fini, plat, et localement de présentation finie.
\end{enumerate}
Si $S$ est localment noethérien, ils sont également équivalents à
\begin{enumerate}
\item[(3)] $f$ est fini et plat.
\end{enumerate}

\end{lemma}

\begin{proof}

Laisser $V \subset S$ Dans les trois cas, le morphisme est donc affiné $f^{-1}(V)$ Ainsi, nous pouvons écrire $V = \Spec(R)$ et $f^{-1}(V) = \Spec(A)$
pour certains $R$-Algèbre $A$. Supposons que (1). Cela signifie que nous pouvons couvrir $S$ par ouverts affines
$V = \Spec(R)$ de telle manière que $A$ est fini libre comme un $R$- Module. $R \to A$ est de présentation finie par Algèbre, lemme \ref{algebra-lemma-finite-finite-type}. Ainsi (2) tient. Inversement, supposer (2). Pour chaque ouvert affine $V = \Spec(R)$ de $S$ l'homomorphisme d'anneaux
$R \to A$ est fini et de présentation finie et $A$ est plat comme un $R$Par Algèbre, lemme \ref{algebra-lemma-finite-finitely-presented-extension}
nous voyons que $A$ est de présentation finie comme $R$-module. Ainsi Algèbre, lemme \ref{algebra-lemma-finite-projective}
implique $A$ Ainsi (1) tient. Le boîtier Noétherien suit comme un module fini sur un anneau Noétherien est un module de présentation finie, voir Algèbre, lemme \ref{algebra-lemma-Noetherian-finite-type-is-finite-presentation}.

\end{proof}

\begin{lemma}

\label{lemma-composition-finite-locally-free}
Une composition de morphismes finis localement libres est finie localement libre.
\end{lemma}

\begin{proof}
Démonstration omise.
\end{proof}

\begin{lemma}

\label{lemma-base-change-finite-locally-free}
Un changement de base d'un morphisme fini localisation libre est fini localement libre.
\end{lemma}

\begin{proof}
Démonstration omise.
\end{proof}

\begin{lemma}

\label{lemma-finite-locally-free}
Laissez $f : X \to S$ être un morphisme fini localisation libre de schémas. Il existe une décomposition d'union disjointe $S = \coprod_{d \geq 0} S_d$ par des sous-systèmes ouverts et fermés tels que le réglage de $X_d = f^{-1}(S_d)$ les restrictions $f|_{X_d}$ sont des morphismes finis localisation libre $X_d \to S_d$ de degré $d$.
\end{lemma}

\begin{proof}

C'est vrai parce qu'une gerbe localement libre de rang fini a un rang bien défini.

\end{proof}

\begin{lemma}

\label{lemma-massage-finite}
Laissez $f : Y \to X$ être un morphisme fini avec $X$ affine. Il existe un diagramme $$
\xymatrix{
Z' \ar[rd] &
Y' \ar[l]^i \ar[d] \ar[r] &
Y \ar[d] \\
 & X' \ar[r] & X
}
$$ où
\begin{enumerate}
\item $Y' \to Y$ et $X' \to X$ sont surjectifs finis localement libres,
\item $Y' = X' \times_X Y$,
\item $i : Y' \to Z'$ est une immersion fermée,
\item $Z' \to X'$ est fini localement libre, et
\item $Z' = \bigcup_{j = 1, \ldots, m} Z'_j$ est une union finie (théorique) de sous-systèmes fermés, dont chacun des cartes isomorphiquement à $X'$.
\end{enumerate}

\end{lemma}

\begin{proof}

Écrire $X = \Spec(A)$ et $Y = \Spec(B)$. Voir aussi Compléments sur l'algèbre, section \ref{more-algebra-section-descent-flatness-integral}. Laisse $x_1, \ldots, x_n \in B$ les générateurs de $B$ terminé $A$. Pour chaque $i$ nous pouvons choisir un polynôme monique $P_i(T) \in A[T]$
de telle manière que $P(x_i) = 0$ en $B$. Par Algèbre, lemme \ref{algebra-lemma-adjoin-roots}
(appliquée $n$ temps) il existe une extension de bague localement libre de rang fini $A \subset A'$ de telle sorte que chacun $P_i$ se divise complètement:
$$
P_i(T) = \prod\nolimits_{k = 1, \ldots, d_i} (T - \alpha_{ik})
$$
pour certains $\alpha_{ik} \in A'$. Prêt
$$
C = A'[T_1, \ldots, T_n]/(P_1(T_1), \ldots, P_n(T_n))
$$
et $B' = A' \otimes_A B$La carte $C \to B'$, $T_i \mapsto 1 \otimes x_i$
est un $A'$- Surjection d'algèbre. $X' = \Spec(A')$,
$Y' = \Spec(B')$ et $Z' = \Spec(C)$ nous voyons que (1) -- (4) tenir. La partie (5) tient parce que l'ensemble théoriquement
$\Spec(C)$ est l'union des sous-systèmes fermés coupés par les idéaux
$$
(T_1 - \alpha_{1k_1}, T_2 - \alpha_{2k_2}, \ldots, T_n - \alpha_{nk_n})
$$
pour tout $1 \leq k_i \leq d_i$.

\end{proof}

\noindent
Le lemme suivant est indiqué dans la bonne généralité dans le lemme \ref{lemma-image-nowhere-dense-quasi-finite} ci-dessous.

\begin{lemma}

\label{lemma-image-nowhere-dense-finite}
Laissez $f : Y \to X$ être un morphisme fini de schémas. Laissez $T \subset Y$ être un sous-ensemble fermé nulle part dense de $Y$. Alors $f(T) \subset X$ est un sous-ensemble fermé nulle part dense de $X$.
\end{lemma}

\begin{proof}
Par lemme \ref{lemma-finite-proper} nous savons que $f(T) \subset X$ est fermé. Laissez $X = \bigcup X_i$ être un revêtement d'affine. Puisque $T$ n'est nulle part dense dans $Y$, nous voyons que $T \cap f^{-1}(X_i)$ n'est nulle part dense dans $f^{-1}(X_i)$. Donc si nous pouvons prouver le théorème dans le cas d'affine, alors nous voyons que $f(T) \cap X_i$ n'est nulle part dense. Cela implique alors que $T$ n'est nulle part dense dans $X$ par Topologie, lemme \ref{topology-lemma-nowhere-dense-local}.

\medskip\noindent

Présumer $X$ est affine. Choisissez un diagramme
$$
\xymatrix{
Z' \ar[rd] &
Y' \ar[l]^i \ar[d]^{f'} \ar[r]_a &
Y \ar[d]^f \\
 & X' \ar[r]^b & X
}
$$
comme dans le lemme \ref{lemma-massage-finite}Les morphismes $a$, $b$ sont ouverts puisqu'ils sont finis localement libres (lemmes \ref{lemma-finite-flat} et \ref{lemma-fppf-open}). D'où $T' = a^{-1}(T)$ n'est nulle part dense, voir Topologie, laisse \ref{topology-lemma-open-inverse-image-closed-nowhere-dense}Le morphisme $b$ est surjectif et ouvert. Donc, si nous pouvons prouver $f'(T') = b^{-1}(f(T))$ n'est nulle part dense, alors $f(T)$ n'est nulle part dense, voir Topologie, laisse \ref{topology-lemma-open-inverse-image-closed-nowhere-dense}. Comme $i$ est une immersion fermée, par Topologie, lemme \ref{topology-lemma-closed-image-nowhere-dense}
nous voyons que $i(T') \subset Z'$ C'est ainsi que nous avons réduit le problème à l'affaire dont il est question dans le paragraphe suivant.

\medskip\noindent

Supposons que $Y = \bigcup_{i = 1, \ldots, n} Y_i$ est une union finie de sous-ensembles fermés, chaque cartographie isomorphiquement à $X$. Considérez
$T_i = Y_i \cap T$. Si chacun des $T_i$ est nulle part dense en $Y_i$, puis chacun $f(T_i)$ est nulle part dense en $X$ comme $Y_i \to X$ est un isomorphisme. $f(T) = f(T_i)$ est une union finie de nulle part denses sous-ensembles fermés de $X$ et on gagne, voir Topologie, lemme \ref{topology-lemma-nowhere-dense}. Supposons que non, dis $T_1$ contient un ouvert non vide $V \subset Y_1$. Nous allons montrer que cela conduit à une contradiction. $Y_2 \cap V \subset V$. Il s'agit soit d'un sous-ensemble bien fermé, soit d'un sous-ensemble égal à $V$. Dans le premier cas nous remplaçons
$V$ par $V \setminus V \cap Y_2$Donc... $V \subset T_1$ est ouvert en $Y_1$ et ne se réunit pas $Y_2$. Dans le deuxième cas, nous avons
$V \subset Y_1 \cap Y_2$ est ouvert dans les deux $Y_1$ et $Y_2$. Répéter successivement avec $i = 3, \ldots, n$. Le résultat est une décomposition d'union disjointe
$$
\{1, \ldots, n\} = I_1 \amalg I_2, \quad 1 \in I_1
$$
et une ouverture $V$ de $Y_1$ figurant dans $T_1$ de telle manière que $V \subset Y_i$
pour $i \in I_1$ et $V \cap Y_i = \emptyset$ pour $i \in I_2$. Prêt
$U = f(V)$. C'est une ouverture de $X$ depuis $f|_{Y_1} : Y_1 \to X$ c'est un isomorphisme.
$$
f^{-1}(U) = V\ \amalg\ \bigcup\nolimits_{i \in I_2} (Y_i \cap f^{-1}(U))
$$
Comme $\bigcup_{i \in I_2} Y_i$ est fermé, ce qui implique que
$V \subset f^{-1}(U)$ est ouvert, donc $V \subset Y$ est ouvert. Cela contredit l'hypothèse que $T$ est nulle part dense en $Y$, comme souhaité.

\end{proof}







\section{Applications rationnelles}
\label{section-rational-maps}

\noindent
Laissez $X$ être un schéma. Notez que si $U$, $V$ sont denses ouverts dans $X$, alors $U \cap V$ est ainsi.

\begin{definition}

\label{definition-rational-map}

Laisser $X$, $Y$ les régimes.

\begin{enumerate}

\item Laisser $f : U \to Y$, $g : V \to Y$ morphismes des schémas définis sur des sous-ensembles ouverts denses $U$, $V$ de $X$. Nous disons que $f$ est {\it équivalent} à $g$ si $f|_W = g|_W$ pour certains $W \subset U \cap V$
dense ouvert dans $X$.
\item A {\it application rationnelle de $X$ à $Y$} est une classe d'équivalence pour la relation d'équivalence définie en (1).
\item Si $X$, $Y$ sont des régimes sur un régime de base $S$ nous disons qu'une application rationnelle de $X$ à $Y$ est un {\it  $S$- application rationnelle de $X$
à $Y$} s'il existe un représentant $f : U \to Y$ de la classe d'équivalence qui est une $S$- le morphisme.

\end{enumerate}

\end{definition}

\noindent
Nous disons que deux morphismes $f$, $g$ comme dans (1) de la définition définissent la même application rationnelle au lieu de dire qu'ils sont équivalents. Dans certains cas, les applications rationnelles sont déterminées par des cartes sur des anneaux locaux à des points génériques.

\begin{lemma}

\label{lemma-rational-map-finite-presentation}

Laisser $S$ Être un plan. $X$ et $Y$ les régimes sont plus $S$Supposons que $X$ a finiment beaucoup de composantes irréductibles avec des points génériques
$x_1, \ldots, x_n$. Laisse $s_i \in S$ être l'image de $x_i$. Considérez l'application
$$
\left\{
\begin{matrix}
S\text{-rational maps} \\
\text{from }X\text{ to }Y
\end{matrix}
\right\}
\longrightarrow
\left\{
\begin{matrix}
(y_1, \varphi_1, \ldots, y_n, \varphi_n)\text{ where }
y_i \in Y\text{ lies over }s_i\text{ and}\\
\varphi_i : \mathcal{O}_{Y, y_i} \to \mathcal{O}_{X, x_i}
\text{ is a local }\mathcal{O}_{S, s_i}\text{-algebra map}
\end{matrix}
\right\}
$$
qui envoie $f : U \to Y$ aux $2n$- double avec
$y_i = f(x_i)$ et $\varphi_i = f^\sharp_{x_i}$. Ensuite

\begin{enumerate}

\item Si $Y \to S$ est localement de type fini, puis l'application est injectable.
\item Si $Y \to S$ est localement de présentation finie, puis l'application est bijective.
\item Si $Y \to S$ est localement de type fini et $X$ réduit, puis l'application est bijective.

\end{enumerate}

\end{lemma}

\begin{proof}
Observez que toute ouverture dense de $X$ contient les points $x_i$ de sorte que la construction a du sens. Pour prouver (1) ou (2) nous pouvons remplacer $X$ par toute ouverture dense. Ainsi si $Z_1, \ldots, Z_n$ sont les composantes irréductibles de $X$, alors nous pouvons remplacer $X$ par $X \setminus \bigcup_{i \not = j} Z_i \cap Z_j$. Après avoir fait ce $X$ est l'union disjointe de ses composants irréductibles (visés comme des sous-systèmes ouverts et fermés). Puis tant le côté droit que le côté gauche de la flèche sont des produits sur les composantes irréductibles et nous réduisons au cas où $X$ est irréductible.

\medskip\noindent

Présumer $X$ est irréductible avec point général $x$ couché sur $s \in S$. La partie 1 découle de la partie 1 de la lettre \ref{lemma-morphism-defined-local-ring}. Les parties (2) et (3) suivent la partie (2) du même lemme.

\end{proof}

\begin{definition}

\label{definition-rational-function}
Laissez $X$ être un schéma. Une {\it fonction rationnelle sur $X$} est une application rationnelle de $X$ à $\mathbf{A}^1_{\mathbf{Z}}$.
\end{definition}

\noindent
Voir Constructions, définition \ref{constructions-definition-affine-n-space} pour la définition de la ligne affine $\mathbf{A}^1$. Laissez $X$ être un schéma sur $S$. Pour tout $U \subset X$ ouvert un morphisme $U \to \mathbf{A}^1_{\mathbf{Z}}$ est le même qu'un morphisme $U \to \mathbf{A}^1_S$ sur $S$. Par conséquent, une fonction rationnelle est également le même qu'une application rationnelle $S$ de $X$ à $\mathbf{A}^1_S$.

\medskip\noindent

Rappelez-vous que nous avons l'identification canonique
$\Mor(T, \mathbf{A}^1_{\mathbf{Z}}) = \Gamma(T, \mathcal{O}_T)$
pour tout régime $T$, voir Schémas, exemple \ref{schemes-example-global-sections}. D'où $\mathbf{A}^1_{\mathbf{Z}}$ est un objet-anneau dans la catégorie des schémas. Plus précisément, les morphismes

\begin{eqnarray*}
+ : \mathbf{A}^1_{\mathbf{Z}} \times \mathbf{A}^1_{\mathbf{Z}}
& \longrightarrow &
\mathbf{A}^1_{\mathbf{Z}} \\
(f, g) & \longmapsto & f + g \\
* : \mathbf{A}^1_{\mathbf{Z}} \times \mathbf{A}^1_{\mathbf{Z}}
& \longrightarrow &
\mathbf{A}^1_{\mathbf{Z}} \\
(f, g) & \longmapsto & fg
\end{eqnarray*}

satisfaire tous les axiomes de l'addition et de la multiplication dans un anneau (commutatif avec $1$ Ainsi aussi l'ensemble d'applications rationnelles dans $\mathbf{A}^1_{\mathbf{Z}}$ a une structure annelée naturelle.

\begin{definition}

\label{definition-ring-of-rational-functions}

Laisser $X$ soit un schéma. L'{\it anneau des fonctions rationnelles sur $X$} est l'anneau $R(X)$ dont les éléments sont des fonctions rationnelles avec addition et multiplication comme on vient de décrire.

\end{definition}

\noindent
Pour les schémas avec de nombreux composantes irréductibles finis, nous pouvons calculer cela.

\begin{lemma}

\label{lemma-integral-scheme-rational-functions}
Laissez $X$ être un schéma avec des composantes irréductibles finis $X_1, \ldots, X_n$. Si $\eta_i \in X_i$ est le point général, alors $$
R(X) = \mathcal{O}_{X, \eta_1} \times \ldots \times \mathcal{O}_{X, \eta_n}
$$ Si $X$ est réduit, cela est égal à $\prod \kappa(\eta_i)$. Si $X$ est intégral, alors $R(X) = \mathcal{O}_{X, \eta} = \kappa(\eta)$ est un champ.
\end{lemma}

\begin{proof}
Laissez $U \subset X$ être un sous-ensemble dense ouvert. Alors $U_i = (U \cap X_i) \setminus (\bigcup_{j \not = i} X_j)$ est ouvert non vide car il contenait $\eta_i$, contenu dans $X_i$, et $\bigcup U_i \subset U \subset X$ est dense. Ainsi, l'identification dans le lemme vient de la chaîne d'égalités
\begin{align*}
R(X)
& =
\colim_{U \subset X\text{ open dense}} \Mor(U, \mathbf{A}^1_\mathbf{Z}) \\
& =
\colim_{U \subset X\text{ open dense}} \mathcal{O}_X(U) \\
& =
\colim_{\eta_i \in U_i \subset X\text{ open}} \prod \mathcal{O}_X(U_i) \\
& =
\prod \colim_{\eta_i \in U_i \subset X\text{ open}} \mathcal{O}_X(U_i) \\
& =
\prod \mathcal{O}_{X, \eta_i}
\end{align*}
où la seconde égalité est Schémas, exemple \ref{schemes-example-global-sections}. L'énoncé final suit d'Algèbre, lemme \ref{algebra-lemma-minimal-prime-reduced-ring}.
\end{proof}

\begin{definition}

\label{definition-function-field}
Laissez $X$ être un schéma intégral. Le champ de fonction {\it}, ou le champ {\it des fonctions rationnelles} de $X$ est le champ $R(X)$.
\end{definition}

\noindent
Nous pouvons occasionnellement indiquer ce champ $k(X)$ au lieu de $R(X)$. Nous pouvons utiliser la notion de champ de fonction pour élucider la condition de séparation sur un schéma intégral. Notez que par lemme \ref{lemma-integral-scheme-rational-functions} sur un schéma intégral chaque anneau local $\mathcal{O}_{X, x}$ peut être vu comme un sous-anneau local de $R(X)$.

\begin{lemma}

\label{lemma-distinct-local-rings}
Laissez $X$ être un schéma séparé intégral. Laissez $Z_1$, $Z_2$ être des sous-ensembles fermés irréductibles distincts de $X$. Laissez $\eta_i$ être le point général de $Z_i$. Si $Z_1 \not\subset Z_2$, alors $\mathcal{O}_{X, \eta_1} \not \subset \mathcal{O}_{X, \eta_2}$ comme des sous-anneaux de $R(X)$. En particulier, si $Z_1 = \{x\}$ se compose d'un point fermé $x$, il existe une fonction régulière dans un quartier de $x$ qui n'est pas dans $\mathcal{O}_{X, \eta_{2}}$.
\end{lemma}

\begin{proof}

Tout d'abord, observez que dans l'hypothèse de $X$ étant séparés, il existe une application unique des régimes
$\Spec(\mathcal{O}_{X, \eta_2}) \to X$ terminé $X$
telle que la composition
$$
\Spec(R(X)) \longrightarrow
\Spec(\mathcal{O}_{X, \eta_2}) \longrightarrow X
$$
est l'application canonique $\Spec(R(X)) \to X$. A savoir, il y a l'application canonique
$can : \Spec(\mathcal{O}_{X, \eta_2}) \to X$, voir Schémas, équation (\ref{schemes-equation-canonical-morphism}). Compte tenu d'un second morphisme $a$ à $X$, nous avons ça $a$ convient avec $can$
sur le point général
$\Spec(\mathcal{O}_{X, \eta_2})$ par hypothèse. Maintenant $X$ être séparé garantit que le sous-ensemble dans
$\Spec(\mathcal{O}_{X, \eta_2})$ où ces deux cartes sont d'accord sont fermées, voir Schémas, lemme \ref{schemes-lemma-where-are-they-equal}. D'où $a = can$ sur tous les $\Spec(\mathcal{O}_{X, \eta_2})$.

\medskip\noindent
Supposons que $Z_1 \not \subset Z_2$ et supposer au contraire que $\mathcal{O}_{X, \eta_{1}} \subset \mathcal{O}_{X, \eta_{2}}$ comme des sous-anneaux de $R(X)$. Ensuite, nous obtenons un deuxième morphisme $$
\Spec(\mathcal{O}_{X, \eta_{2}}) \longrightarrow
\Spec(\mathcal{O}_{X, \eta_{1}}) \longrightarrow
X.
$$ Par ce qui précède, cette composition devrait être égale à $can$. Cela implique que $\eta_2$ se spécialise dans $\eta_1$ (voir Schémas, lemme \ref{schemes-lemma-specialize-points}). Mais cela contredit notre hypothèse $Z_1 \not \subset Z_2$.
\end{proof}

\begin{definition}

\label{definition-domain-of-definition}
Laissez $\varphi$ être une application rationnelle entre deux schémas $X$ et $Y$. Nous disons que $\varphi$ est {\it défini dans un point $x \in X$} s'il existe un $(U, f)$ représentatif de $\varphi$ avec $x \in U$. Le domaine {\it de définition} de $\varphi$ est l'ensemble de tous les points où $\varphi$ est défini.
\end{definition}

\noindent
Avec cette définition, il n'est pas vrai en général que $\varphi$ a un représentant qui est défini sur tout le domaine de définition.

\begin{lemma}

\label{lemma-rational-map-from-reduced-to-separated}
Laissez $X$ et $Y$ être des schémas. Supposons que $X$ réduit et $Y$ séparé. Laissez $\varphi$ être une application rationnelle de $X$ à $Y$ avec un domaine de définition $U \subset X$. Ensuite, il existe un morphisme unique $f : U \to Y$ représentant $\varphi$. Si $X$ et $Y$ sont des schémas sur un schéma séparé $S$ et si $\varphi$ est un rationelle d'application $S$, alors $f$ est un morphisme sur $S$.
\end{lemma}

\begin{proof}

Laisser $(V, g)$ et $(V', g')$ Être des représentants de $\varphi$. Ensuite
$g, g'$ s'accordent sur un sous-système ouvert dense $W \subset V \cap V'$. D'autre part, l'égalisateur $E$ de $g|_{V \cap V'}$ et $g'|_{V \cap V'}$
est un sous-système fermé de $V \cap V'$ (Schémas, lemme)
\ref{schemes-lemma-where-are-they-equal}). Maintenant $W \subset E$
implique que $E = V \cap V'$ de façon théorique. $V \cap V'$
est réduit nous concluons $E = V \cap V'$ en théorie, c'est-à-dire,
$g|_{V \cap V'} = g'|_{V \cap V'}$. Il s'ensuit que nous pouvons coller les représentants $g : V \to Y$ de $\varphi$ à un morphisme $f : U \to Y$, voir Schémas, lemme \ref{schemes-lemma-glue}. Nous omettons la démonstration de l'énoncé finale.

\end{proof}

\noindent
En général, il n'est pas logique de composer des applications rationnelles. La raison est que l'image d'un représentant de la première application rationnelle peut avoir une intersection vide avec le domaine de définition de la seconde. Cependant, si nous supposons que nos schémas sont irréductibles et que nous regardons les applications dominantes rationnelles, nous pouvons composer des applications rationnelles.

\begin{definition}

\label{definition-dominant-rational}
Laissez $X$ et $Y$ être des schémas irréductibles. Une application rationnelle de $X$ à $Y$ est appelée {\it dominante} si un représentant $f : U \to Y$ est un morphisme dominant des schémas.
\end{definition}

\noindent

Par moi \ref{lemma-dominant-finite-number-irreducible-components}
il est équivalent à exiger que le point général $\eta \in X$
les cartes au point général $\xi$ de $Y$, i.e., $f(\eta) = \xi$ pour tout représentant $f : U \to Y$. Nous pouvons composer une application dominante rationnelle
$\varphi$ entre régimes irréductibles $X$ et $Y$ avec une application arbitraire rationnelle $\psi$ à partir de $Y$ à $Z$. A savoir, choisir des représentants $f : U \to Y$ avec $U \subset X$ ouvert dense et $g : V \to Z$ avec $V \subset Y$ ouvert dense. Puis
$W = f^{-1}(V) \subset X$ est ouvert non vide (car il contient le point général de $X$) et nous laissons $\psi \circ \varphi$ être la classe d'équivalence de $g \circ f|_W : W \to Z$. Nous omettons de vérifier que cela est bien défini.

\medskip\noindent
Nous obtenons ainsi une catégorie dont les objets sont des schémas irréductibles et dont les morphismes sont des applications dominantes rationelles. Compte tenu d'un schéma de base $S$ nous pouvons également définir une catégorie dont les objets sont des schémas irréductibles sur $S$ et dont les morphismes sont des applications dominantes $S$- rationelles.

\begin{definition}

\label{definition-birational-schemes}

Laisser $X$ et $Y$ être irréductibles.

\begin{enumerate}

\item Nous disons $X$ et $Y$ sont {\it birationnel} si $X$ et $Y$ sont isomorphes dans la catégorie des schémas irréductibles et des applications dominantes rationnelles.
\item Présumer $X$ et $Y$ sont des régimes sur un régime de base $S$Nous disons : $X$ et $Y$ sont {\it  $S$-Birationnel} si $X$ et $Y$ sont isomorphes dans la catégorie des régimes irréductibles sur $S$ et dominant $S$- applications rationnelles.

\end{enumerate}

\end{definition}

\noindent
Si $X$ et $Y$ sont des schémas irréductibles birationnels, alors l'ensemble de applications rationnelles de $X$ à $Z$ est bijectif avec l'ensemble d'application rationnelle de $Y$ à $Z$ pour tous les schémas $Z$ (fonctionnellement dans $Z$). Pour les schémas irréductibles "générales" ce n'est qu'une définition possible. Un autre serait d'exiger $X$ et $Y$ ont des anneaux isomorphes de fonctions rationnelles. Pour les variétés, ces conditions sont équivalentes, voir leme \ref{lemma-criterion-birational-finite-presentation}.

\begin{lemma}

\label{lemma-birational-integral}

Laisser $X$ et $Y$ être irréductibles.

\begin{enumerate}

\item Les régimes $X$ et $Y$ sont birationnels si et seulement s'ils ont des ouvertures isomorphes non vides.
\item Présumer $X$ et $Y$ sont des régimes sur un régime de base $S$. Ensuite
$X$ et $Y$ sont $S$-birationnel s'il y a des ouvertures non vides $U \subset X$ et $V \subset Y$ qui sont $S$- isomorphe.

\end{enumerate}

\end{lemma}

\begin{proof}

Présumer $X$ et $Y$ Ils sont birationnels. $f : U \to Y$ et $g : V \to X$
définir les applications dominantes inverses rationnelles de $X$ à $Y$ et à partir $Y$ à $X$. Nous pouvons supposer $V$ affine. Nous pouvons remplacer $U$ par une affinité d'ouverture de $f^{-1}(V)$. Comme $g \circ f$ est l'identité comme une application dominante rationnelle, nous voyons que la composition $U \to V \to X$ est l'identité sur une ouverture dense de $U$. Ainsi après avoir remplacé $U$ par un plus petit ouvert affine nous pouvons supposer que
$U \to V \to X$ est l'inclusion de $U$ dans $X$Il s'ensuit que
$U \to V$ est une immersion (Appliquer Schémas, \ref{schemes-lemma-section-immersion}
à $U \to g^{-1}(U) \to U$). Cependant, changer les rôles de $U$ et $V$ et refait l'argument ci-dessus, nous voyons qu'il existe un non vide ouvert affine $V' \subset V$
de sorte que les facteurs d'inclusion comme $V' \to U \to V$. Ensuite $V' \to U$ est nécessairement une immersion ouverte. $V' \to f^{-1}(V') \to V'$ sont des monomorphismes (Schémas, \ref{schemes-lemma-immersions-monomorphisms}) composer à l'identité, donc isomorphismes. $V'$ est isomorphe à une ouverture des deux $X$ et $Y$. Dans $S$- cas d'applications rationnelles, le même argument fonctionne exactement.

\end{proof}

\begin{remark}

\label{remark-generalize-category}

Il s'agit ici d'une gé nération de la catégorie des régimes irréductibles et des demandes dominantes rationelles. $X$ indiquer $X^0$ la série de points $x \in X$ avec $\dim(\mathcal{O}_{X, x}) = 0$, en d'autres termes,
$X^0$ est l'ensemble de points génériques de composantes irréductibles de $X$. Ensuite, nous pouvons considérer la catégorie avec

\begin{enumerate}

\item les objets sont des schémas $X$ de telle sorte que chaque quasi-compact ouvert ait finiment de nombreux composantes irréductibles, et
\item morphismes de $X$ à $Y$ sont des applications rationnelles $f : U \to Y$
à partir de $X$ à $Y$ de telle manière que $f(U^0) = Y^0$.

\end{enumerate}

Si $U \subset X$ est une ouverture dense d'un schéma, alors
$U^0 \subset X^0$ ne doit pas être une égalité, mais si $X$ est un objet de notre catégorie, alors c'est le cas. Ainsi, étant donné deux morphismes dans notre catégorie, la composition est bien définie et un morphisme dans notre catégorie.

\end{remark}

\begin{remark}

\label{remark-pseudo-morphisms}
Il y a une variante de définition \ref{definition-rational-map} où nous considérons uniquement les morphismes $U \to Y$ définis sur les sous-systèmes ouverts denses de schématiquement $U \subset X$. Nous utilisons lemme \ref{lemma-intersection-scheme-theoretically-dense} pour voir si nous obtenons une relation d'équivalence. Une classe d'équivalence de ceux-ci est appelée un pseudo-morphisme {\it de $X$ à $Y$}. Si $X$ est réduit les deux notions coïncident.
\end{remark}
X






\section{Morphismes birationnels}
\label{section-birational}

\noindent
Vous pouvez être habitué à la notion d'une application birationnelle des variétés ayant la propriété qu'il s'agit d'un isomorphisme sur un sous-ensemble ouvert de la cible. Cependant, en général, un morphisme birationnel peut ne pas être un isomorphisme sur un ouvert non vide, voir exemple \ref{example-birational-not-iso-over-open}. Voici la définition formelle.

\begin{definition}

\label{definition-birational}

\begin{reference}

\cite[(2.2.9)]{EGA1}

\end{reference}

Laisser $X$, $Y$ les régimes. Supposons $X$ et $Y$ ont finiment beaucoup de composantes irréductibles. On dit un morphisme $f : X \to Y$ est {\it birationnel} si

\begin{enumerate}

\item  $f$ induit une bijection entre l'ensemble de points génériques de composantes irréductibles de $X$ et l'ensemble des points génériques des composantes irréductibles de $Y$,
\item pour chaque point général $\eta \in X$ d'un composante irréductible de $X$ l'application locale de l'anneau
$\mathcal{O}_{Y, f(\eta)} \to \mathcal{O}_{X, \eta}$
est un isomorphisme.

\end{enumerate}

\end{definition}

\noindent
Nous verrons ci-dessous que les fibres d'un morphisme birationnel sur les points génériques sont des singletons. De plus, nous verrons que dans la plupart des cas on rencontre en pratique l'existence d'un morphisme birationnel entre des schémas irréductibles $X$ et $Y$ implique que $X$ et $Y$ sont des schémas birationnels.

\begin{lemma}

\label{lemma-birational-dominant}
Laissez $f : X \to Y$ être un morphisme de schémas ayant finiment de nombreux composantes irréductibles. Si $f$ est birationnel alors $f$ est dominant.
\end{lemma}

\begin{proof}
Suit de lemme \ref{lemma-generic-points-in-image-dominant} et la définition.
\end{proof}

\begin{lemma}

\label{lemma-birational-generic-fibres}
Laissez $f : X \to Y$ être un morphisme birationnel de schémas ayant finiment de nombreux composantes irréductibles. Si $y \in Y$ est le point général d'un composante irréductible, alors le changement de base $X \times_Y \Spec(\mathcal{O}_{Y, y}) \to \Spec(\mathcal{O}_{Y, y})$ est un isomorphisme.
\end{lemma}

\begin{proof}

Nous pouvons supposer $Y = \Spec(B)$ est affine et irréductible. $X$ est irréductible aussi. Si nous prouvons le résultat pour n'importe quel non vide ouvert affine $U \subset X$, alors le résultat tient pour $X$
(petit argument omis). Par conséquent, nous pouvons supposer $X$
est affine aussi, disons $X = \Spec(A)$. Laisse $y \in Y$ correspondent au minimum de prime $\mathfrak q \subset B$. Par hypothèse $A$ dispose d'un minimum unique $\mathfrak p$ couché sur $\mathfrak q$ et $B_\mathfrak q \to A_\mathfrak p$ est un isomorphisme. $A_\mathfrak q \to \kappa(\mathfrak p)$ est surjectif, donc $\mathfrak p A_\mathfrak q$ est un idéal maximal. D'un autre côté
$\mathfrak p A_\mathfrak q$ est le premier minimum unique de $A_\mathfrak q$. Nous concluons que $\mathfrak p A_\mathfrak q$ est le premier de
$A_\mathfrak q$ et que $A_\mathfrak q = A_\mathfrak p$Depuis
$A_\mathfrak q = A \otimes_B B_\mathfrak q$ le lemme suit.

\end{proof}

\begin{example}

\label{example-birational-not-iso-over-open}
Voici deux exemples de morphismes birationnels qui ne sont pas des isomorphismes au-dessus d'une ouverture de la cible.

\medskip\noindent
Premier exemple. Laissez $k$ être un champ infini. Laissez $A = k[x]$. Laissez $B = k[x, \{y_{\alpha}\}_{\alpha \in k}]/
((x-\alpha)y_\alpha, y_\alpha y_\beta)$. Il y a une inclusion $A \subset B$ et une rétractation $B \to A$ pour tout $y_\alpha$ égal à zéro. Tant le morphisme $\Spec(A) \to \Spec(B)$ que le morphisme $\Spec(B) \to \Spec(A)$ sont birationnels mais pas un isomorphisme sur n'importe quel ouvert.

\medskip\noindent
Deuxième exemple. Laissez $A$ être un domaine. Laissez $S \subset A$ être un sous-ensemble multiplicatif ne contenant pas $0$. Avec $B = S^{-1}A$ le morphisme $f : \Spec(B) \to \Spec(A)$ est birationnel. S'il existe un $U$ ouvert de $\Spec(A)$ tel que $f^{-1}(U) \to U$ est un isomorphisme, alors il existe un $a \in A$ tel que chaque élément de $S$ devient inversible dans la localisation principale $A_a$. Prenant $A = \mathbf{Z}$ et $S$ l'ensemble des entiers impairs donne un contre-exemple.
\end{example}

\begin{lemma}

\label{lemma-birational-birational}

Laisser $f : X \to Y$ être un morphisme birationnel de schémas ayant finiment de nombreux composantes irréductibles sur un schéma de base $S$Supposons que l'une des conditions suivantes est remplie

\begin{enumerate}

\item  $f$ est localement de type fini et $Y$ réduit,
\item  $f$ est le lieu de présentation finie.

\end{enumerate}

Puis il y a des ouvertures denses $U \subset X$ et $V \subset Y$
de telle manière que $f(U) \subset V$ et $f|_U : U \to V$ est un isomorphisme. En particulier si $X$ et $Y$ sont irréductibles, alors
$X$ et $Y$ sont $S$-Birationnel.

\end{lemma}

\begin{proof}
Il y a une réduction immédiate dans le cas où $X$ et $Y$ sont irréductibles que nous omettons. De plus, après avoir rétréci davantage et nous pouvons supposer que $X$ et $Y$ sont affines, disons $X = \Spec(A)$ et $Y = \Spec(B)$. Par hypothèse $A$, rép.\ $B$ a un minimum unique de $\mathfrak p$, rép.\ $\mathfrak q$, le $\mathfrak p$ primaire se trouve sur $\mathfrak q$, et $B_\mathfrak q = A_\mathfrak p$. Par leme \ref{lemma-birational-generic-fibres} nous avons $B_\mathfrak q = A_\mathfrak q = A_\mathfrak p$.

\medskip\noindent
Supposons que $B \to A$ est de type fini, disons $A = B[x_1, \ldots, x_n]$. Il existe un $b_i \in B$ et $g_i \in B \setminus \mathfrak q$ de telle sorte que $b_i/g_i$ mappe à l'image de $x_i$ dans $A_\mathfrak q$. Par conséquent $b_i  - g_ix_i$ mappe à zéro dans $A_{g_i'}$ pour certains $g_i' \in B \setminus \mathfrak q$. Réglage $g = \prod g_i g'_i$ nous voyons que $B_g \to A_g$ est surjectif. Si d'ailleurs $Y$ est réduit, alors l'application $B_g \to B_\mathfrak q$ est injectable et donc $B_g \to A_g$ est injective aussi. Ceci prouve le cas (1).

\medskip\noindent
démonstration de (2). Par l'argument donné dans le paragraphe précédent, nous pouvons supposer que $B \to A$ est surjectif. Comme $f$ est localement de présentation finie le noyau $J \subset B$ est un idéal fini. Écrivons $J = (b_1, \ldots, b_r)$. Depuis $B_\mathfrak q = A_\mathfrak q$ il existe $g_i \in B \setminus \mathfrak q$ tel que $g_i b_i = 0$. Réglage $g = \prod g_i$ nous voyons que $B_g \to A_g$ est un isomorphisme.
\end{proof}

\begin{lemma}

\label{lemma-criterion-birational-finite-presentation}

Laisser $S$ Être un plan. $X$ et $Y$ les régimes irréductibles localement de présentation finie over $S$. Laisse $x \in X$ et $y \in Y$
les points génériques. Les conditions suivantes sont équivalentes

\begin{enumerate}

\item  $X$ et $Y$ sont $S$- birationnel,
\item il existe des ouvertures non vides de $X$ et $Y$
qui sont $S$-isomorphe, et
\item  $x$ et $y$ carte au même point $s$ de $S$ et
$\mathcal{O}_{X, x}$ et $\mathcal{O}_{Y, y}$ sont isomorphes comme
$\mathcal{O}_{S, s}$- des algèbres.

\end{enumerate}

\end{lemma}

\begin{proof}
Nous avons vu l'équivalence de (1) et (2) dans le lemme \ref{lemma-birational-integral}. Il est immédiat que (2) implique (3). Pour finir, nous supposons (3) tient et nous prouvons (1). Par le lemme \ref{lemma-rational-map-finite-presentation} il y a une application rationnelle $f : U \to Y$ qui envoie $x \in U$ à $y$ et induit l'isomorphisme donné $\mathcal{O}_{Y, y} \cong \mathcal{O}_{X, x}$. Ainsi $f$ est un birationnel morphisme et induit donc un isomorphisme sur non vide s'ouvre par lemme \ref{lemma-birational-birational}. Ceci termine la démonstration.
\end{proof}

\begin{lemma}

\label{lemma-common-open}

Laisser $S$ Être un plan. $X$ et $Y$ les régimes intégrés de type fini sur $S$. Laisse $x \in X$ et $y \in Y$ les points génériques. Les conditions suivantes sont équivalentes

\begin{enumerate}

\item  $X$ et $Y$ sont $S$- birationnel,
\item il existe des ouvertures non vides de $X$ et $Y$ qui sont $S$-isomorphe, et
\item  $x$ et $y$ carte au même point $s \in S$ et
$\kappa(x) \cong \kappa(y)$ comme $\kappa(s)$- des extensions.

\end{enumerate}

\end{lemma}

\begin{proof}

Nous avons vu l'équivalence de (1) et (2) dans le \ref{lemma-birational-integral}. Il est immédiat que (2) implique (3). Pour finir, nous supposons (3) tient et nous prouvons (1). Observez que $\mathcal{O}_{X, x} = \kappa(x)$ et
$\mathcal{O}_{Y, y} = \kappa(y)$ par Algèbre, lemme \ref{algebra-lemma-minimal-prime-reduced-ring}. Par lemme \ref{lemma-rational-map-finite-presentation}
il y a une application rationnelle $f : U \to Y$
qui envoie $x \in U$ à $y$ et induit l'isomorphisme donné $\mathcal{O}_{Y, y} \cong \mathcal{O}_{X, x}$. Ainsi $f$ est un morphisme birationnel et induit donc un isomorphisme sur non vide ouvert par lemme \ref{lemma-birational-birational}. Cela finit la démonstration.

\end{proof}







\section{Morphismes génériquement finis}
\label{section-generically-finite}

\noindent
Dans cette section, nous caractérisons les cartes entre les schémas qui sont localment de type fini et qui sont ``génériquement fini'' dans un sens ou un autre.

\begin{lemma}

\label{lemma-generically-finite}

Laisser $X$, $Y$ les régimes. Laisser $f : X \to Y$ être localement de type fini. $\eta \in Y$ être un point général d'un composante irréductible de $Y$. Les conditions suivantes sont équivalentes:

\begin{enumerate}

\item l'ensemble $f^{-1}(\{\eta\})$ est fini,
\item il existe des ouverts affines $U_i \subset X$, $i = 1, \ldots, n$
et $V \subset Y$ avec $f(U_i) \subset V$,
$\eta \in V$ et $f^{-1}(\{\eta\}) \subset \bigcup U_i$
de telle sorte que chacun $f|_{U_i} : U_i \to V$ C'est fini.

\end{enumerate}

Si $f$ est quasi-séparé, alors ceux-ci sont également équivalents à

\begin{enumerate}

\item[(3)] il existe des ouverts affines $V \subset Y$, $U \subset X$ avec $f(U) \subset V$,
$\eta \in V$ et $f^{-1}(\{\eta\}) \subset U$
de telle manière que $f|_U : U \to V$ C'est fini.

\end{enumerate}

Si $f$ est quasi-compact et quasi-séparé, alors ils sont également équivalents à

\begin{enumerate}

\item[(4)] il existe un ouvert affine $V \subset Y$, $\eta \in V$
de telle manière que $f^{-1}(V) \to V$ C'est fini.

\end{enumerate}

\end{lemma}

\begin{proof}
La question est locale sur la base. Par conséquent, nous pouvons remplacer $Y$ par un quartier affiné de $\eta$, et nous pouvons et faisons supposer tout au long de la démonstration ci-dessous que $Y$ est affine, disons $Y = \Spec(R)$.

\medskip\noindent
Il est clair que (2) implique (1). Supposons que $f^{-1}(\{\eta\}) = \{\xi_1, \ldots, \xi_n\}$ est fini. Choisissez ouverts affines $U_i \subset X$ avec $\xi_i \in U_i$. Par Algèbre, lemme \ref{algebra-lemma-generically-finite} nous voyons qu'après avoir remplacé $Y$ par un principal ouvert dans $Y$ chacun des morphismes $U_i \to Y$ est fini. En d'autres termes (2) tient.

\medskip\noindent

Il est clair que (3) implique (1). $f$ est quasi-séparé et (1). $f^{-1}(\{\eta\}) = \{\xi_1, \ldots, \xi_n\}$. Il n'y a pas de spécialisation parmi les $\xi_i$ par lemme \ref{lemma-finite-fibre}. Depuis chaque $\xi_i$ les cartes au point général $\eta$ d'un composante irréductible de $Y$, il ne peut y avoir une spécialisation non triviale $\xi \leadsto \xi_i$ en $X$
(depuis $\xi$ pour cartographier à $\eta$ Nous concluons chacun $\xi_i$ est un point général d'un composante irréductible de $X$Depuis $Y$ est affine et $f$ quasi-séparé nous voyons $X$ est quasi-séparé. Par Propriétés, lemme \ref{properties-lemma-maximal-points-affine}
nous pouvons trouver un ouvert affine $U \subset X$ contenant chacun $\xi_i$. Par Algèbre, lemme \ref{algebra-lemma-generically-finite} nous voyons qu'après avoir remplacé $Y$ par un principal d'ouverture en
$Y$ les morphismes $U \to Y$ est fini. En d'autres termes (3) tient.

\medskip\noindent

Il est clair que (4) implique tout (1) -- (3) sans autres hypothèses sur $f$. Supposons que $f$ Nous devons montrer que les conditions équivalentes (1) -- (3) impliquent (4). $U$, $V$ comme dans (3). Remplacer $Y$ par $V$Depuis $f$ est quasi compacte et $Y$ est quasi-compact (être affine) nous voyons que $X$ est quasi-compact. $Z = X \setminus U$ est quasi-compact, d'où le morphisme
$f|_Z : Z \to Y$ est quasi-compact. Par construction de $Z$ nous voyons que
$\eta \not \in f(Z)$. D'où par lemme \ref{lemma-quasi-compact-generic-point-not-in-image}
nous voyons qu'il existe un quartier affiné ouvert $V'$ de $\eta$ en $Y$ de telle manière que $f^{-1}(V') \cap Z = \emptyset$. Alors nous avons $f^{-1}(V') \subset U$ et cela signifie que $f^{-1}(V') \to V'$ C'est fini.

\end{proof}

\begin{example}

\label{example-counter-generically-finite}

Laisser $A = \prod_{n \in \mathbf{N}} \mathbf{F}_2$. Chaque élément de $A$ est un idempotent. Ainsi, chaque idéal premier est maximal avec corps résiduel $\mathbf{F}_2$. Ainsi la topologie sur $X = \Spec(A)$ est totalement déconnecté et quasi-compact. $A \to \mathbf{F}_2$ définir les points ouverts de $\Spec(A)$. Il ne peut être que tous les points de $X$ sont ouverts depuis $X$ est quasi-compact. $x \in X$ être un point fermé qui n'est pas ouvert. Alors nous pouvons former un schéma $Y$ qui est deux exemplaires de $X$ collé le long $X \setminus \{x\}$. En d'autres termes, c'est $X$ avec $x$ doublé, comparer Schémas, exemple \ref{schemes-example-affine-space-zero-doubled}Le morphisme
$f : Y \to X$ est quasi-compact, type fini et a des fibres finies mais n'est pas quasi-séparé. $x \in X$ est un point général d'un composante irréductible de $X$ (depuis $X$ est totalement déconnecté). Mais les propriétés (3) et (4) de lemme \ref{lemma-generically-finite} ne pas tenir. La raison est que pour tout quartier ouvert $x \in U \subset X$ l'image inverse
$f^{-1}(U)$ n'est pas affine parce que les fonctions sur $f^{-1}(U)$ ne peut pas séparer les deux points $x$ (prouve omise; c'est un bel exercice). D'où la condition que $f$ est quasi-séparé est nécessaire dans les parties (3) et (4) du lemme.

\end{example}

\begin{remark}

\label{remark-quasi-finite-finite-over-dense-open}
Une alternative à lemme \ref{lemma-generically-finite} est l'affirmation qu'un morphisme quasi-fini est fini sur un ouverture dense de la cible. Ceci sera montré dans Compléments sur les morphismes, lemme \ref{more-morphisms-lemma-quasi-finite-finite-over-dense-open}.
\end{remark}

\begin{lemma}

\label{lemma-quasi-finiteness-over-generic-point}

Laisser $X$, $Y$ les régimes. Laisser $f : X \to Y$ être localement de type fini. $X^0$, rép.\ $Y^0$ désignent l'ensemble de points génériques de composantes irréductibles de $X$, rép.\ $Y$. Laisse $\eta \in Y^0$. Les conditions suivantes sont équivalentes

\begin{enumerate}

\item  $f^{-1}(\{\eta\}) \subset X^0$,
\item  $f$ est quasi-fini à tous les points couchés $\eta$,
\item  $f$ est quasi-fini du tout $\xi \in X^0$ couché sur $\eta$.

\end{enumerate}

\end{lemma}

\begin{proof}
Condition (1) implique qu'il n'y a pas de spécialisation entre les points de la fibre $X_\eta$. Par conséquent (2) détient par lemme \ref{lemma-quasi-finite-at-point-characterize}. L'implication (2) $\Rightarrow$ (3) est immédiate. $\eta$ est un point général de $Y$, les points génériques de $X_\eta$ sont des points génériques de $X$. Par conséquent (3) et lemme \ref{lemma-quasi-finite-at-point-characterize} impliquent les points génériques de $X_\eta$ sont également fermés. Ainsi, tous les points de $X_\eta$ sont génériques et nous voyons que (1) détient.
\end{proof}

\begin{lemma}

\label{lemma-finite-over-dense-open}

Laisser $X$, $Y$ les régimes. Laisser $f : X \to Y$ être localement de type fini. $X^0$, rép.\ $Y^0$ désignent l'ensemble de points génériques de composantes irréductibles de $X$, rép.\ $Y$Supposons que

\begin{enumerate}

\item  $X^0$ et $Y^0$ sont finis et $f^{-1}(Y^0) = X^0$,
\item soit $f$ est quasi-compact ou $f$ est séparé.

\end{enumerate}

Puis il y a une ouverture dense $V \subset Y$
de telle manière que $f^{-1}(V) \to V$ C'est fini.

\end{lemma}

\begin{proof}
Étant donné que $Y$ a finiment beaucoup de composantes irréductibles, nous pouvons trouver une ouverture dense qui est une union disjointe de ses composantes irréductibles. Ainsi, nous pouvons supposer que $Y$ est une affine irréductible avec point général $\eta$. Puis la fibre sur $\eta$ est finie comme $X^0$ est finie.

\medskip\noindent
Supposons que $f$ est séparé et $Y$ affine irréductible. Choisissez $V \subset Y$ et $U \subset X$ comme dans le lemme \ref{lemma-generically-finite} partie (3). Puisque $f|_U : U \to V$ est fini, nous voyons que $U \subset f^{-1}(V)$ est fermé ainsi que ouvert (lemmes \ref{lemma-image-proper-scheme-closed} et \ref{lemma-finite-proper}). Ainsi $f^{-1}(V) = U \amalg W$ pour certains sous-systèmes ouverts $W$ de $X$. Cependant, puisque $U$ contient tous les points génériques de $X$ nous concluons que $W = \emptyset$ comme souhaité.

\medskip\noindent

Présumer $f$ est quasi compacte et $Y$ affine irréductible. $X$ est quasi-compact, il existe donc un sous-système ouvert dense $U \subset X$
qui est séparé (Propriétés, lemme \ref{properties-lemma-quasi-compact-dense-open-separated}Depuis l'ensemble des points génériques $X^0$ est fini, nous voyons que $X^0 \subset U$. Ainsi $\eta \not \in f(X \setminus U)$Depuis $X \setminus U \to Y$ est quasi-compact, nous concluons qu'il y a un ouvert non vide $V \subset Y$
de telle manière que $f^{-1}(V) \subset U$, voir lemme \ref{lemma-quasi-compact-dominant}. Après remplacement $X$ par $f^{-1}(V)$ et $Y$ par $V$ Nous nous réduisons à l'affaire séparée que nous avons traitée dans le paragraphe précédent.

\end{proof}

\begin{lemma}

\label{lemma-birational-isomorphism-over-dense-open}

Laisser $X$, $Y$ les régimes. Laisser $f : X \to Y$ être un morphisme birationnel entre des schémas qui ont finiment de nombreux composantes irréductibles.

\begin{enumerate}

\item soit $f$ est quasi-compact ou $f$ est séparé, et
\item soit $f$ est localement de type fini et $Y$ est réduit ou
$f$ est le lieu de présentation finie.

\end{enumerate}

Puis il y a une ouverture dense $V \subset Y$
de telle manière que $f^{-1}(V) \to V$ est un isomorphisme.

\end{lemma}

\begin{proof}
Par lemme \ref{lemma-finite-over-dense-open} nous pouvons supposer que $f$ est fini. Puisque $Y$ a finiment de nombreux composantes irréductibles, nous pouvons trouver un ouvert dense qui est une union disjointe de ses composantes irréductibles. Ainsi nous pouvons supposer que $Y$ est irréductible. Par lemme \ref{lemma-birational-birational} nous trouvons un ouvert non vide $U \subset X$ tel que $f|_U : U \to Y$ est un ouverte d'immersion. Après avoir enlevé le sous-ensemble fermé (comme $f$ fini) $f(X \setminus U)$ de $Y$ nous voyons que $f$ est un isomorphisme.
\end{proof}

\begin{lemma}

\label{lemma-finite-degree}

Laisser $X$, $Y$ de faire partie intégrante des régimes. $f : X \to Y$ être localement de type fini. Supposons que $f$ est dominante. Les conditions suivantes sont équivalentes:

\begin{enumerate}

\item l'extension $R(Y) \subset R(X)$ a degré de transcendance $0$,
\item l'extension $R(Y) \subset R(X)$ est fini,
\item il existe des ouverts non vides affines $U \subset X$
et $V \subset Y$ de telle manière que $f(U) \subset V$
et $f|_U : U \to V$ est fini, et
\item le point général de $X$ est le seul point de $X$ cartographie au point général de $Y$.

\end{enumerate}

Si $f$ est séparé ou si $f$ est quasi-compact, alors ceux-ci sont également équivalents à

\begin{enumerate}

\item[(5)] il existe un ouvert affine non vide $V \subset Y$ de telle manière que $f^{-1}(V) \to V$ C'est fini.

\end{enumerate}

\end{lemma}

\begin{proof}
Choisissez tous les ouverts affines $\Spec(A) = U \subset X$ et $\Spec(R) = V \subset Y$ de sorte que $f(U) \subset V$. Ensuite, $R$ et $A$ sont des domaines par définition. L'homomorphisme d'anneaux $R \to A$ est de type fini (lemme \ref{lemma-locally-finite-type-characterize}). Par lemme \ref{lemma-dominant-finite-number-irreducible-components} le point généralique de $X$ cartes au point généralique de $Y$ par conséquent $R \to A$ est injectable. Laissez $K = R(Y)$ être le champ de fraction de $R$ et $L = R(X)$ le champ de fraction de $A$. Puis $L/K$ est une extension de champ générée de façon limitée.

\medskip\noindent
Supposons (2) tient.Soit $x_1, \ldots, x_n \in A$ des générateurs de $A$ sur $R$. Par hypothèse, il existe des polynomes non zéro $P_i(X) \in R[X]$ de telle sorte que $P_i(x_i) = 0$.Soit $f_i \in R$ le coefficient de tête de $P_i$. Puis nous concluons que $R_{f_1 \ldots f_n} \to A_{f_1 \ldots f_n}$ est fini, c'est-à-dire (3) tient. Notez que (3) implique (2). Donc maintenant nous voyons que (1), (2) et (3) sont tous équivalents.

\medskip\noindent

Laisser $\eta$ être le point général de $X$, et laisser $\eta' \in Y$ être le point général de $Y$. Supposons que (4). Puis
$\dim_\eta(X_{\eta'}) = 0$ et nous voyons que $R(X) = \kappa(\eta)$ a degré de transcendance $0$ terminé $R(Y) = \kappa(\eta')$ par lemme \ref{lemma-dimension-fibre-at-a-point}. En d'autres termes (1) tient. Supposons que les conditions équivalentes (1), (2) et (3). Supposons que $x \in X$ est une cartographie de points pour $\eta'$. Comme $x$ est une spécialisation de $\eta$, cela donne des inclusions $R(Y) \subset \mathcal{O}_{X, x} \subset R(X)$, ce qui implique $\mathcal{O}_{X, x}$ est un champ, voir Algèbre, lemme \ref{algebra-lemma-integral-over-field}. D'où $x = \eta$. Ainsi, nous voyons que (1) -- (4) sont tous équivalents.

\medskip\noindent
Il est clair que (5) implique (3) sans hypothèses supplémentaires sur $f$. Ce qui reste est de prouver que si $f$ est soit séparé ou quasi-compact, alors les conditions équivalentes (1) -- (4) impliquent (5). Ceci découle de lemme \ref{lemma-finite-over-dense-open}.
\end{proof}

\begin{definition}

\label{definition-degree}
Laissez $X$ et $Y$ être des schémas intégrés. Laissez $f : X \to Y$ être localement de type fini et dominant. Supposons que $[R(X) : R(Y)] < \infty$, ou toute autre des conditions équivalentes (1) -- (4) de lemme \ref{lemma-finite-degree}. Puis l'entier positif $$
\deg(X/Y) = [R(X) : R(Y)]
$$ est appelé le {\it degré de $X$ sur $Y$}.
\end{definition}

\noindent
Il est possible d'étendre cette notion à un morphisme $f : X \to Y$ si (a) $Y$ fait partie intégrante du point général $\eta$, (b) $f$ est localement de type fini, et (c) $f^{-1}(\{\eta\})$ est fini. Dans ce cas, nous pouvons définir $$
\deg(X/Y)
=
\sum\nolimits_{\xi \in X, \ f(\xi) = \eta}
\dim_{R(Y)} (\mathcal{O}_{X, \xi}).
$$ Namely, étant donné que $R(Y) = \kappa(\eta) = \mathcal{O}_{Y, \eta}$ (lemme \ref{lemma-integral-scheme-rational-functions}) les dimensions ci-dessus sont finies par lemme \ref{lemma-generically-finite} ci-dessus. Cependant, pour la plupart des applications, la définition donnée ci-dessus est la bonne.

\begin{lemma}

\label{lemma-degree-composition}
Laissez $X$, $Y$, $Z$ être des schémas intégrés. Laissez $f : X \to Y$ et $g : Y \to Z$ être des morphismes dominants localement de type fini. Supposons que $[R(X) : R(Y)] < \infty$ et $[R(Y) : R(Z)] < \infty$. Puis $$
\deg(X/Z) = \deg(X/Y) \deg(Y/Z).
$$
\end{lemma}

\begin{proof}
Ceci vient de la multiplication des degrés dans les tours des extensions finies des champs, voir Corps, lemme \ref{fields-lemma-multiplicativity-degrees}.
\end{proof}

\begin{remark}

\label{remark-definition-generically-finite}

Laisser $f : X \to Y$ Il existe au moins deux propriétés que nous pourrions utiliser pour définir {\it morphismes généraux fini}. Ils correspondent à savoir si vous voulez que la propriété soit locale sur la source ou locale sur la cible :

\begin{enumerate}

\item (Local sur la cible; suggéré par Ravi Vakil.) Supposons que chaque quasi-compact ouvert de $Y$ a finiment beaucoup de composantes irréductibles (par exemple si $Y$ L'exigence est que l'image inverse de chaque point général est finie, voir lemme \ref{lemma-generically-finite}.
\item (Local sur la source.) L'exigence est qu'il existe un ouvert dense $U \subset X$ de telle manière que $U \to Y$ est localement quasi-finie.

\end{enumerate}

Dans le cas 1), la prescription peut être formulée sans condition auxiliaire sur $Y$La propriété (2) telle que formulée n'implique pas que les fibres sur les points génériques sont finies; cependant, si $f$ est quasi compacte et $Y$ c'est comme dans (1) alors il le fait.

\end{remark}

\begin{definition}

\label{definition-modification}
Laissez $X$ être un schéma intégral. Une {\it modification de $X$} est un morphisme birationnel propre $f : X' \to X$ avec $X'$ intégrale.
\end{definition}

\noindent

Laisser $f : X' \to X$ être une modification comme dans la définition. \ref{lemma-finite-degree} il existe un non vide $U \subset X$ de telle manière que $f^{-1}(U) \to U$ est fini. Par platitude générale (proposition \ref{proposition-generic-flatness}) nous pouvons supposer $f^{-1}(U) \to U$ est plat et de présentation finie. $f^{-1}(U) \to U$ est fini localement libre (lemme \ref{lemma-finite-flat}Depuis $f$ est birationnel, le degré de $X'$ terminé $X$ est $1$. D'où $f^{-1}(U) \to U$ est finie localement sans degré $1$, autrement dit c'est un isomorphisme. Ainsi, on peut {\it redéfinir} une modification comme un morphisme propre $f : X' \to X$ des régimes intégrés tels que $f^{-1}(U) \to U$ est un isomorphisme pour certains ouverts non vides
$U \subset X$.

\begin{definition}

\label{definition-alteration}

\begin{reference}
\cite[définition 2.20]{altérations}
\end{reference}
Laissez $X$ être un schéma intégral. Une altération {\it de $X$} est un morphisme dominant $f : Y \to X$ avec l'intégrale $Y$ de telle sorte que $f^{-1}(U) \to U$ est finie pour certains $U \subset X$ non vides.
\end{definition}
.

\noindent
C'est la définition donnée dans \cite{alterations}, sauf que nous n'exigeons pas ici que $X$ et $Y$ soient noethérien. Arguant comme ci-dessus, nous voyons qu'une altération est un morphisme approprié dominant $f : Y \to X$ des schémas intégrés qui induit une extension finie des champs de fonction, c'est-à-dire de sorte que les conditions équivalentes de lemme \ref{lemma-finite-degree} tiennent.



\section{La formule de dimension}
\label{section-dimension-formula}

\noindent
Pour les morphismes entre les schémas Noétheriens, nous pouvons dire un peu plus sur les dimensions des anneaux locaux. Voici un résultat important (et pas si difficile à prouver). Rappelez-vous que $R(X)$ désigne le champ de fonction d'un schéma intégral $X$.

\begin{lemma}

\label{lemma-dimension-formula}
Laissez $S$ être un schéma. Laissez $f : X \to S$ être un morphisme de schémas. Laissez $x \in X$, et définir $s = f(x)$. Supposons que
\begin{enumerate}
\item $S$ est localment noethérien,
\item $f$ est localment de type fini,
\item $X$ et $S$ intégrale, et
\item $f$ dominante.
\end{enumerate}
Nous avons
\begin{equation}
\label{equation-dimension-formula}
\dim(\mathcal{O}_{X, x})
\leq
\dim(\mathcal{O}_{S, s}) + \text{trdeg}_{R(S)}R(X)
- \text{trdeg}_{\kappa(s)} \kappa(x).
\end{equation}
En outre, l'égalité est maintenue si $S$ est universellement caténaire.
\end{lemma}

\begin{proof}
L'algèbre correspondant est Algèbre, lemme \ref{algebra-lemma-dimension-formula}.
\end{proof}

\begin{lemma}

\label{lemma-dimension-formula-general}
Laissez $S$ être un schéma. Laissez $f : X \to S$ être un morphisme de schémas. Laissez $x \in X$, et définir $s = f(x)$. Supposons que $S$ est localement noethérien et $f$ est localement de type fini, Nous avons
\begin{equation}
\label{equation-dimension-formula-general}
\dim(\mathcal{O}_{X, x})
\leq
\dim(\mathcal{O}_{S, s}) + E - \text{trdeg}_{\kappa(s)} \kappa(x).
\end{equation}
où $E$ est le maximum de $\text{trdeg}_{\kappa(f(\xi))}(\kappa(\xi))$ où $\xi$ fonctionne sur les points génériques de composantes irréductibles de $X$ contenant $x$.
\end{lemma}

\begin{proof}

Laisser $X_1, \ldots, X_n$ être les composantes irréductibles de $X$ contenant $x$
La structure du schéma induite est réduite et correspond aux primes minimales. $\mathfrak q_i$ de
$\mathcal{O}_{X, x}$ et donc il y a finiment beaucoup d'entre eux (Schémas, lemme \ref{schemes-lemma-specialize-points} et d ' Algèbre, lemme \ref{algebra-lemma-Noetherian-irreducible-components}). Puis
$\dim(\mathcal{O}_{X, x}) = \max \dim(\mathcal{O}_{X, x}/\mathfrak q_i)
= \max \dim(\mathcal{O}_{X_i, x})$. Les $\xi$Il apparaît dans la définition de $E$ sont exactement les points génériques $\xi_i \in X_i$. Laisse $Z_i = \overline{\{f(\xi_i)\}} \subset S$
avec la structure de schéma induite réduite. $X_i \to X \to S$ des facteurs $Z_i$
(Schémas, lemme) \ref{schemes-lemma-map-into-reduction}Ainsi, nous pouvons appliquer la forme de dimension (lemme \ref{lemma-dimension-formula}) de voir que
$\dim(\mathcal{O}_{X_i, x}) \leq \dim(\mathcal{O}_{Z_i, x}) +
\text{trdeg}_{\kappa(f(\xi))}(\kappa(\xi)) -
\text{trdeg}_{\kappa(s)} \kappa(x)$. En rassemblant tout, nous obtenons le lemme.

\end{proof}

\noindent
Une application est la construction d'une fonction dimension sur n'importe quel schéma de type fini sur un schéma universel caténaire doté d'une fonction dimension. Pour la définition des fonctions dimension, voir Topologie, définition \ref{topology-definition-dimension-function}.

\begin{lemma}

\label{lemma-dimension-function-propagates}
Laissez $S$ être une localisation noéthérien et schéma universel. Laissez $\delta : S \to \mathbf{Z}$ être une fonction de dimension. Laissez $f : X \to S$ être un morphisme de schémas. Supposons que $f$ localement de type fini. Puis l'application
\begin{align*}
\delta = \delta_{X/S} : X & \longrightarrow \mathbf{Z} \\
x & \longmapsto \delta(f(x)) + \text{trdeg}_{\kappa(f(x))} \kappa(x)
\end{align*}
est une fonction de dimension sur $X$.
\end{lemma}

\begin{proof}
Laissez $f : X \to S$ être localement de type fini. Laissez $x \leadsto y$, $x \not = y$ être une spécialisation dans $X$. Nous devons montrer que $\delta_{X/S}(x) > \delta_{X/S}(y)$ et que $\delta_{X/S}(x) = \delta_{X/S}(y) + 1$ si $y$ est une spécialisation immédiate de $x$.

\medskip\noindent

Choisissez un ouvert affine $V \subset S$ contenant l'image de
$y$ et choisissez un ouvert affine $U \subset X$ cartographie dans $V$
et contenant $y$. Nous pouvons clairement remplacer $X$ par $U$ et
$S$ par $V$. Ainsi, nous pouvons supposer que $X = \Spec(A)$
et $S = \Spec(R)$ et que $f$ est donné par un homomorphisme d'anneaux $R \to A$. La bague $R$ est universellement caténaire (lemme \ref{lemma-universally-catenary-local}) et l'application $R \to A$ est de type fini (lemme \ref{lemma-locally-finite-type-characterize}).

\medskip\noindent

Laisser $\mathfrak q \subset A$ être l'idéal premier correspondant au point $x$ et laisser $\mathfrak p \subset R$ être l'idéal premier correspondant à $f(x)$. La restriction $\delta'$ de $\delta$
à $S' = \Spec(R/\mathfrak p) \subset S$ est une fonction dimension. L'anneau $R/\mathfrak p$ est universellement caténaire. $\delta_{X/S}$ à $X' = \Spec(A/\mathfrak q)$
est clairement égal à la fonction $\delta_{X'/S'}$ construit à l'aide de la fonction dimension $\delta'$. Par conséquent, nous pouvons supposer en plus de ce qui précède que $R \subset A$ sont des domaines, en d'autres termes que $X$ et $S$ sont des régimes intégrés, et que
$x$ est le point général de $X$ et $f(x)$ est le point général de $S$.

\medskip\noindent
Notez que $\mathcal{O}_{X, x} = R(X)$ et que depuis $x \leadsto y$, $x \not = y$, le spectre de $\mathcal{O}_{X, y}$ a au moins deux points (Schémas, lemme \ref{schemes-lemma-specialize-points}) d'où $\dim(\mathcal{O}_{X, y}) > 0$ . Si $y$ est une spécialisation immédiate de $x$, puis $\Spec(\mathcal{O}_{X, y}) = \{x, y\}$ et $\dim(\mathcal{O}_{X, y}) = 1$.

\medskip\noindent
Write $s = f(x)$ et $t = f(y)$. Nous calculons
\begin{align*}
\delta_{X/S}(x) - \delta_{X/S}(y)
& =
\delta(s) + \text{trdeg}_{\kappa(s)} \kappa(x)
- \delta(t) - \text{trdeg}_{\kappa(t)} \kappa(y) \\
& =
\delta(s) - \delta(t) +
\text{trdeg}_{R(S)} R(X) - \text{trdeg}_{\kappa(t)} \kappa(y) \\
& =
\delta(s) - \delta(t) + \dim(\mathcal{O}_{X, y})
- \dim(\mathcal{O}_{S, t})
\end{align*}
où nous utilisons l'égalité dans (\ref{equation-dimension-formula}) dans la dernière étape. Puisque $\delta$ est une fonction de dimension sur le schéma $S$ et $s \in S$ est le point général, la différence $\delta(s) - \delta(t)$ est égale à $\text{codim}(\overline{\{t\}}, S)$ par Topologie, lemme \ref{topology-lemma-dimension-function-catenary}. Ceci est égal à $\dim(\mathcal{O}_{S, t})$ par Propriétés, lemme \ref{properties-lemma-codimension-local-ring}. Par conséquent, nous concluons que $$
\delta_{X/S}(x) - \delta_{X/S}(y) = \dim(\mathcal{O}_{X, y})
$$ et la lemme suit de ce que nous avons dit ci-dessus sur $\dim(\mathcal{O}_{X, y})$.
\end{proof}
.

\noindent
Une autre application de la formule de dimension est que la dimension ne change pas sous ``altérations'' (à définir ultérieurement).

\begin{lemma}

\label{lemma-alteration-dimension}

Laisser $f : X \to Y$ être un morphisme de schémas. Supposons que

\begin{enumerate}

\item  $Y$ est localement noethérien,
\item  $X$ et $Y$ sont des régimes intégrés,
\item  $f$ est dominante, et
\item  $f$ est localement de type fini.

\end{enumerate}

Alors nous avons
$$
\dim(X) \leq \dim(Y) + \text{trdeg}_{R(Y)} R(X).
$$
Si $f$ est fermée\footnote{Par exemple si $f$ est correct, voir définition \ref{definition-proper}.} alors l'égalité se maintient.

\end{lemma}

\begin{proof}
Laissez $f : X \to Y$ être comme dans la lemme.Soit $\xi_0 \leadsto \xi_1 \leadsto \ldots \leadsto \xi_e$ une séquence de spécialisations dans $X$. Posons $x = \xi_e$ et $y = f(x)$. Observez que $e \leq \dim(\mathcal{O}_{X, x})$ comme les spécialisations données se produisent dans le spectre de $\mathcal{O}_{X, x}$, voir Schémas, lemme \ref{schemes-lemma-specialize-points}. Par la formule de dimension, lemme \ref{lemma-dimension-formula}, nous voyons que
\begin{align*}
e & \leq \dim(\mathcal{O}_{X, x}) \\
& \leq \dim(\mathcal{O}_{Y, y}) +
\text{trdeg}_{R(Y)} R(X) - \text{trdeg}_{\kappa(y)} \kappa(x) \\
& \leq \dim(\mathcal{O}_{Y, y}) + \text{trdeg}_{R(Y)} R(X)
\end{align*}
Nous concluons donc que $e \leq \dim(Y) + \text{trdeg}_{R(Y)} R(X)$ comme désiré.

\medskip\noindent

Ensuite, supposez $f$ est également fermée. $\overline{\xi}_0 \leadsto \overline{\xi}_1 \leadsto \ldots
\leadsto \overline{\xi}_d$ est une séquence de spécialisations en $Y$. Nous voulons montrer que $\dim(X) \geq d + r$. Nous pouvons supposer que $\overline{\xi}_0 = \eta$ est le point général de $Y$. La fibre générale $X_\eta$ est un système localement de type fini sur
$\kappa(\eta) = R(Y)$. Il n'est pas vide comme $f$ est dominant. \ref{lemma-ubiquity-Jacobson-schemes} c'est une schéma de Jacobson. \ref{lemma-jacobson-finite-type-points}
on peut trouver un point fermé $\xi_0 \in X_\eta$ et l'extension
$\kappa(\eta) \subset \kappa(\xi_0)$ est une extension finie. Notez que
$\mathcal{O}_{X, \xi_0} = \mathcal{O}_{X_\eta, \xi_0}$ parce que
$\eta$ est le point général de $Y$. Nous voyons donc que
$\dim(\mathcal{O}_{X, \xi_0}) = r$ par lemme \ref{lemma-dimension-formula}
appliqué au régime $X_\eta$ sur le schéma universel caténaire $\Spec(\kappa(\eta))$ (voir lemme \ref{lemma-ubiquity-uc}) et le point $\xi_0$. Cela signifie que nous pouvons trouver
$\xi_{-r} \leadsto \ldots \leadsto \xi_{-1} \leadsto \xi_0$
en $X$. D'autre part, comme $f$ est fermé spécialisations le long $f$, voir Topologie, me laisse \ref{topology-lemma-closed-open-map-specialization}. Ainsi, comme $\xi_0$ se repose sur $\eta = \overline{\xi}_0$ nous pouvons trouver des spécialisations $\xi_0 \leadsto \xi_1 \leadsto \ldots \leadsto \xi_d$
couché sur $\overline{\xi}_0 \leadsto \overline{\xi}_1 \leadsto \ldots
\leadsto \overline{\xi}_d$En d'autres termes, nous avons
$$
\xi_{-r} \leadsto \ldots \leadsto \xi_{-1} \leadsto \xi_0
\leadsto \xi_1 \leadsto \ldots \leadsto \xi_d
$$
ce qui signifie que $\dim(X) \geq d + r$ comme souhaité.

\end{proof}

\begin{lemma}

\label{lemma-alteration-dimension-general}
Laissez $f : X \to Y$ être un morphisme de schémas. Supposons que $Y$ est localement noethérien et $f$ est localement de type fini. Puis $$
\dim(X) \leq \dim(Y) + E
$$ où $E$ est le supremum de $\text{trdeg}_{\kappa(f(\xi))}(\kappa(\xi))$ où $\xi$ passe par les points génériques des composantes irréductibles de $X$.
\end{lemma}

\begin{proof}
Résultat immédiat de \ref{lemma-dimension-formula-general} et Propriétés, \ref{properties-lemma-dimension}.
\end{proof}








\section{Normalisation relative}
\label{section-normalization-X-in-Y}

\noindent
Dans cette section, nous construisons la normalisation {\it d'un schéma dans un autre}.

\begin{lemma}

\label{lemma-integral-closure}
Laissez $X$ être un schéma. Laissez $\mathcal{A}$ être un faisceau quasi-cohérent de $\mathcal{O}_X$-algèbres. Le sous-chapeau $\mathcal{A}' \subset \mathcal{A}$ défini par la règle $$
U \longmapsto \{f \in \mathcal{A}(U) \mid
f_x \in \mathcal{A}_x \text{ integral over } \mathcal{O}_{X, x}
\text{ for all }x \in U\}
$$ est un $\mathcal{O}_X$-algèbre quasi-cohérent, la tige $\mathcal{A}'_x$ est la fermeture intégrale de $\mathcal{O}_{X, x}$ dans $\mathcal{A}_x$, et pour toute ouverture affine $U \subset X$ l'anneau $\mathcal{A}'(U) \subset \mathcal{A}(U)$ est la fermeture intégrale de $\mathcal{O}_X(U)$ dans $\mathcal{A}(U)$.
\end{lemma}

\begin{proof}

Il s'agit d'un sous-chapeau par la nature locale des conditions. $\mathcal{O}_X$-algèbre d'Algèbre, lemme \ref{algebra-lemma-integral-closure-is-ring}. Laisse $U \subset X$ Sois un ouvert affine. $U = \Spec(R)$
et dire $\mathcal{A}$ est le faisceau quasi-cohérent associé au $R$-Algèbre $A$. Ensuite, selon Algèbre, lemme \ref{algebra-lemma-integral-closure-stalks}
la valeur de $\mathcal{A}'$ terminé $U$ est donné par la fermeture intégrale $A'$ de $R$ en $A$. Cela prouve la dernière affirmation du lemme. Pour prouver que $\mathcal{A}'$ est quasi-cohérent, il suffit de montrer que $\mathcal{A}'(D(f)) = A'_f$. Cela résulte du fait que la fermeture intégrale et le déplacement de localisation, voir Algèbre, lemme \ref{algebra-lemma-integral-closure-localize}. Le même fait montre que les tiges sont comme annoncé.

\end{proof}

\begin{definition}

\label{definition-integral-closure}
Laissez $X$ être un schéma. Laissez $\mathcal{A}$ être un faisceau quasi-cohérent de $\mathcal{O}_X$-algèbres. La fermeture intégrale {\it de $\mathcal{O}_X$ dans $\mathcal{A}$} est la quasi-cohérente $\mathcal{O}_X$-subalgèbre $\mathcal{A}' \subset \mathcal{A}$ construite en lemme \ref{lemma-integral-closure} ci-dessus.
\end{definition}

\noindent

Dans le cadre de la définition ci-dessus, nous pouvons considérer le morphisme des spectres relatifs
$$
\xymatrix{
Y = \underline{\Spec}_X(\mathcal{A}) \ar[rr] \ar[rd] & &
X' = \underline{\Spec}_X(\mathcal{A}') \ar[ld] \\
& X &
}
$$
voir lemme \ref{lemma-affine-equivalence-algebras}Le régime $X' \to X$ sera la normalisation de $X$ dans le régime $Y$. Voici un cadre légèrement plus général. Supposons que nous avons un quasi-compact et quasi-morphisme séparé $f : Y \to X$
Dans ce cas-ci, la gerbe de
$\mathcal{O}_X$-Algèbres $f_*\mathcal{O}_Y$ est quasi-cohérent, voir Schémas, lemme \ref{schemes-lemma-push-forward-quasi-coherent}. Prise de la fermeture intégrale $\mathcal{O}' \subset f_*\mathcal{O}_Y$
nous obtenons un faisceau quasi-cohérent de $\mathcal{O}_X$-algèbres dont le spectre relatif est la normalisation $X$ en $Y$. Voici la définition formelle.

\begin{definition}

\label{definition-normalization-X-in-Y}

Laisser $f : Y \to X$ être un quasi-compact et quasi-morphisme séparé de schémas. $\mathcal{O}'$ soit la fermeture intégrale de $\mathcal{O}_X$ en
$f_*\mathcal{O}_Y$. Les {\it la normalisation des $X$ en $Y$} est le régime\footnote{Le régime $X'$ ne doivent pas être normaux, par exemple si
$Y = X$ et $f = \text{id}_X$Alors $X' = X$.}
$$
\nu : X' = \underline{\Spec}_X(\mathcal{O}') \to X
$$
terminé $X$. Il est équipé d'une factorisation naturelle
$$
Y \xrightarrow{f'} X' \xrightarrow{\nu} X
$$
du morphisme initial $f$.

\end{definition}

\noindent
La factorisation est la composition du morphisme canonique $Y \to \underline{\Spec}_X(f_*\mathcal{O}_Y)$ (voir Constructions, lemme \ref{constructions-lemma-canonical-morphism}) et le morphisme des spectres relatifs provenant de l'application d'inclusion $\mathcal{O}' \to f_*\mathcal{O}_Y$. Nous pouvons caractériser la normalisation comme suit.

\begin{lemma}

\label{lemma-characterize-normalization}

Laisser $f : Y \to X$ être un quasi-compact et quasi-morphisme séparé des schémas. $f = \nu \circ f'$où: $\nu : X' \to X$ est la normalisation de $X$ en $Y$ est caractérisée par les deux propriétés suivantes:

\begin{enumerate}

\item le morphisme $\nu$ fait partie intégrante, et
\item pour toute factorisation $f = \pi \circ g$Avec $\pi : Z \to X$
intégrale, il existe un diagramme commutatif
$$
\xymatrix{
Y \ar[d]_{f'} \ar[r]_g & Z \ar[d]^\pi \\
X' \ar[ru]^h \ar[r]^\nu & X
}
$$
pour un morphisme unique $h : X' \to Z$.

\end{enumerate}

En outre, le morphisme $f' : Y \to X'$ est dominant et dans (2) le morphisme $h : X' \to Z$ est la normalisation de $Z$ en $Y$.

\end{lemma}

\begin{proof}

Laisser $\mathcal{O}' \subset f_*\mathcal{O}_Y$ soit la fermeture intégrale de
$\mathcal{O}_X$ comme dans la définition \ref{definition-normalization-X-in-Y}Le morphisme $\nu$ fait partie intégrante de la construction, ce qui prouve (1). $f = \pi \circ g$ avec $\pi : Z \to X$
intégrale comme dans (2). \ref{definition-integral}

$\pi$ est affine, et donc $Z$ est le spectre relatif d'un faisceau quasi-cohérent de $\mathcal{O}_X$-Algèbres $\mathcal{B}$Le morphisme $g : Y \to Z$ correspond à une application de $\mathcal{O}_X$-Algèbres
$\chi : \mathcal{B} \to f_*\mathcal{O}_Y$Depuis $\mathcal{B}(U)$ fait partie intégrante $\mathcal{O}_X(U)$ pour chaque ouvert affine $U \subset X$
(par définition) \ref{definition-integral}) nous voyons de Lemme \ref{lemma-integral-closure}
que $\chi(\mathcal{B}) \subset \mathcal{O}'$. Par la functorialité du spectre relatif \ref{lemma-affine-equivalence-algebras}
cela nous fournit un morphisme unique
$h : X' \to Z$. Nous omettons la vérification que le diagramme se déplace.

\medskip\noindent
Il est clair que (1) et (2) caractérisent la factorisation $f = \nu \circ f'$ puisqu'elle la caractérise comme un objet initial dans une catégorie.

\medskip\noindent
De la propriété universelle en (2) nous voyons que $f'$ ne fait pas de facteur à travers un bon sous-système fermé de $X'$. D'où l'image schématique de $f'$ est $X'$. Puisque $f'$ est quasi-compact (par Schémas, lemme \ref{schemes-lemma-quasi-compact-permanence} et le fait que $\nu$ est séparé comme un morphisme affine) nous voyons que $f'(Y)$ est dense dans $X'$. Par conséquent, $f'$ est dominante.

\medskip\noindent

Observer que $g$ est quasi compacte et quasi séparée par Schémas, lemmes \ref{schemes-lemma-compose-after-separated} et
\ref{schemes-lemma-quasi-compact-permanence}. Ainsi la dernière affirmation du lemme a un sens. Le morphisme $h$ en (2) est intégré par lemme \ref{lemma-finite-permanence}. Compte tenu d'une factorisation $g = \pi' \circ g'$ avec $\pi' : Z' \to Z$
intégrale, nous obtenons une factorisation $f = (\pi \circ \pi') \circ g'$ et nous obtenons un morphisme $h' : X' \to Z'$. L'unicité implique que
$\pi' \circ h' = h$. D'où la caractérisation (1), (2) s'applique au morphisme $h : X' \to Z$ qui donne l'affirmation finale du lemme.

\end{proof}

\begin{lemma}

\label{lemma-functoriality-normalization}

Laisser
$$
\xymatrix{
Y_2 \ar[d]_{f_2} \ar[r] & Y_1 \ar[d]^{f_1} \\
X_2 \ar[r] & X_1
}
$$
être un diagramme commutatif des morphismes des schémas. $f_1$, $f_2$ quasi compacte et quasi séparée. $f_i = \nu_i \circ f_i'$, $i = 1, 2$
les facteurs canoniques, où $\nu_i : X_i' \to X_i$ est la normalisation de $X_i$ en $Y_i$. Puis il existe une flèche unique $X'_2 \to X'_1$ s'insérer dans un diagramme commutatif
$$
\xymatrix{
Y_2 \ar[d]_{f_2'} \ar[r] & Y_1 \ar[d]^{f_1'} \\
X_2' \ar[d]_{\nu_2} \ar[r] & X_1' \ar[d]^{\nu_1} \\
X_2 \ar[r] & X_1
}
$$

\end{lemma}

\begin{proof}
Par lemmes \ref{lemma-characterize-normalization} (1) et \ref{lemma-base-change-finite} le changement de base $X_2 \times_{X_1} X'_1 \to X_2$ est intégral. Notez que $f_2$ facteurs à travers ce morphisme. C'est pourquoi nous obtenons un morphisme unique $X'_2 \to X_2 \times_{X_1} X'_1$ de lemme \ref{lemma-characterize-normalization} (2). Ceci donne la flèche $X'_2 \to X'_1$ fixation dans le diagramme commutatif et l'unicité suit également.
\end{proof}

\begin{lemma}

\label{lemma-normalization-localization}
Laissez $f : Y \to X$ être un séparé quasi-compact et quasi-morphisme de schémas. Laissez $U \subset X$ être un sous-système ouvert et définir $V = f^{-1}(U)$. Ensuite, la normalisation de $U$ dans $V$ est l'image inverse de $U$ dans la normalisation de $X$ dans $Y$.
\end{lemma}

\begin{proof}
Supprimer de la construction.
\end{proof}

\begin{lemma}

\label{lemma-normalization-is-normalization}
Laissez $f : Y \to X$ être un quasi-compact et quasi-morphisme séparé de schémas. Laissez $X'$ être la normalisation de $X$ dans $Y$. Ensuite, la normalisation de $X'$ dans $Y$ est $X'$.
\end{lemma}

\begin{proof}

Si $Y \to X'' \to X'$ est la normalisation de $X'$ en $Y$, alors nous pouvons appliquer lemme \ref{lemma-characterize-normalization}
à la composition $X'' \to X$ pour obtenir un morphisme canonique
$h : X' \to X''$ terminé $X$. Nous omettons la vérification que les morphismes $h$ et $X'' \to X'$ sont mutuellement inverses (en utilisant l'unicité de la factorisation dans le lemme).

\end{proof}

\begin{lemma}

\label{lemma-normalization-in-reduced}
Laissez $f : Y \to X$ être un séparé quasi-compact et quasi-morphisme de schémas. Laissez $X' \to X$ être la normalisation de $X$ dans $Y$. Si $Y$ est réduit, ainsi est $X'$.
\end{lemma}

\begin{proof}
Ceci résulte du fait qu'un anneau réduit est réduit. Certains détails omis.
\end{proof}

\begin{lemma}

\label{lemma-normalization-generic}
Laissez $f : Y \to X$ être un quasi-compact et quasi-morphisme séparé de schémas. Laissez $X' \to X$ être la normalisation de $X$ dans $Y$. Chaque point général d'un composite irréductible de $X'$ est l'image d'un point général d'un composite irréductible de $Y$.
\end{lemma}

\begin{proof}
Par lemme \ref{lemma-normalization-localization} nous pouvons supposer que $X = \Spec(A)$ est affine. Choisissez un recouvrement fini ouvert affine $Y = \bigcup \Spec(B_i)$. Puis $X' = \Spec(A')$ et les morphismes $\Spec(B_i) \to Y \to X'$ définissent conjointement une application $A$-algèbre injectable $A' \to \prod B_i$. Ainsi le lemme suit d'Algèbre, lemme \ref{algebra-lemma-injective-minimal-primes-in-image}.
\end{proof}

\begin{lemma}

\label{lemma-normalization-in-disjoint-union}
Laissez $f : Y \to X$ être un séparé quasi-compact et quasi-morphisme de schémas. Supposons que $Y = Y_1 \amalg Y_2$ est une union disjointe de deux schémas. Écrivez $f_i = f|_{Y_i}$. Laissez $X_i'$ être la normalisation de $X$ dans $Y_i$. Alors $X_1' \amalg X_2'$ est la normalisation de $X$ dans $Y$.
\end{lemma}

\begin{proof}

En termes de fermetures intégrales, cela correspond au fait suivant: $A \to B$ être un homomorphisme d'anneaux. Supposons que $B = B_1 \times B_2$. Laisse $A_i'$ soit la fermeture intégrale de $A$ en $B_i$. Ensuite
$A_1' \times A_2'$ est la fermeture intégrale de $A$ en $B$. La raison pour laquelle cela fonctionne est que les éléments $(1, 0)$ et $(0, 1)$ de $B$
sont idempotents et donc intègrent $A$. Ainsi la fermeture intégrale
$A'$ de $A$ en $B$ est un produit et il n'est pas difficile de voir que les facteurs sont les fermetures intégrales $A'_i$ comme décrit ci-dessus (certains détails ont été omis).

\end{proof}

\begin{lemma}

\label{lemma-normalization-in-universally-closed}
Laissez $f : X \to S$ être un morphisme quasi-compact, quasi-séparé et universel de fermé des schémas. Alors $f_*\mathcal{O}_X$ est intégrale sur $\mathcal{O}_S$. En d'autres termes, la normalisation de $S$ dans $X$ est égale à la factorisation $$
X \longrightarrow \underline{\Spec}_S(f_*\mathcal{O}_X)
\longrightarrow S
$$ de Constructions, lemme \ref{constructions-lemma-canonical-morphism}.
\end{lemma}

\begin{proof}

La question est locale sur $S$, donc nous pouvons supposer $S = \Spec(R)$
C'est affiné. $h \in \Gamma(X, \mathcal{O}_X)$. Nous devons montrer que
$h$ satisfait à une équation monique sur $R$. Pensez à $h$ comme un morphisme comme dans le diagramme commutatif suivant
$$
\xymatrix{
X \ar[rr]_h \ar[rd]_f & & \mathbf{A}^1_S \ar[ld] \\
& S &
}
$$
Laisser $Z \subset \mathbf{A}^1_S$ être l'image schématique de $h$, voir définition \ref{definition-scheme-theoretic-image}Le morphisme $h$ est quasi-compact comme $f$ est quasi compacte et
$\mathbf{A}^1_S \to S$ est séparé, voir Schémas, lemme \ref{schemes-lemma-quasi-compact-permanence}. Par lemme \ref{lemma-quasi-compact-scheme-theoretic-image} le morphisme $X \to Z$ est dominant. \ref{lemma-image-proper-scheme-closed} le morphisme
$X \to Z$ est fermée. $h(X) = Z$ Ainsi, nous pouvons utiliser lemme \ref{lemma-image-proper-is-proper}
de conclure que $Z \to S$ est universellement fémé (et même propre). $Z \subset \mathbf{A}^1_S$, nous voyons que $Z \to S$ est affine et propre, donc intégrée par lemme \ref{lemma-integral-universally-closed}. Écriture $\mathbf{A}^1_S = \Spec(R[T])$ nous concluons que l'idéal $I \subset R[T]$ de $Z$ contient un polynôme monique
$P(T) \in R[T]$. D'où $P(h) = 0$ et nous gagnons.

\end{proof}

\begin{lemma}

\label{lemma-normalization-in-integral}
Laissez $f : Y \to X$ être un morphisme entier. Ensuite, la normalisation de $X$ dans $Y$ est égale à $Y$.
\end{lemma}

\begin{proof}
Par lemme \ref{lemma-integral-universally-closed} c'est un cas spécial de lemme \ref{lemma-normalization-in-universally-closed}.
\end{proof}

\begin{lemma}

\label{lemma-normal-normalization}
Laissez $f : Y \to X$ être un quasi-compact et quasi-morphisme séparé de schémas. Laissez $X'$ être la normalisation de $X$ dans $Y$. Supposons que
\begin{enumerate}
\item $Y$ est un schéma normal,
\item les ouvertures quasi-compact de $Y$ ont finiment de nombreux composantes irréductibles.
\end{enumerate}
Alors $X'$ est une union disjointe de schémas normaux intégrals. De plus, le morphisme $Y \to X'$ est dominant et induit une bijection de composantes irréductibles.
\end{lemma}

\begin{proof}

Laisser $U \subset X$ être un ouvert affine. Considérez l'image inverse $U'$ de $U$ en $X'$. Prêt $V = f^{-1}(U)$. Par lemme \ref{lemma-normalization-localization}
nous $V \to U' \to U$ est la normalisation de $U$ en $V$. Dites
$U = \Spec(A)$. Ensuite $V$ est quasi-compact, et a donc un nombre fini de composantes irréductibles par hypothèse.
$V = \coprod_{i = 1, \ldots n} V_i$ est une union disjointe finie des schémas intégrés normaux par Propriétés, lemme \ref{properties-lemma-normal-locally-finite-nr-irreducibles}. Par lemme \ref{lemma-normalization-in-disjoint-union}
nous voyons que $U' = \coprod_{i = 1, \ldots, n} U_i'$où: $U'_i$ est la normalisation de $U$ en $V_i$. Par Propriétés, me laisse \ref{properties-lemma-normal-integral-sections}
nous voyons que $B_i = \Gamma(V_i, \mathcal{O}_{V_i})$ est un domaine normal. Notez que $U_i' = \Spec(A_i')$où: $A_i' \subset B_i$
est la fermeture intégrale de $A$ en $B_i$, voir lemme \ref{lemma-integral-closure}. Par Algèbre, lemme \ref{algebra-lemma-integral-closure-in-normal}
nous voyons que $A_i' \subset B_i$ est un domaine normal. $U' = \coprod U_i'$ est une union finie de schémas intégrés normaux et est donc normale.

\medskip\noindent

Comme $X'$ a un recouvrement ouvert par les régimes $U'$ nous concluons de Propriétés, lemme \ref{properties-lemma-locally-normal} que $X'$ est normal. D'autre part, $U'$ est une union disjointe finie de schémas irréductibles, d'où chaque quasi-compact ouvert de $X'$ a finiment beaucoup de composantes irréductibles (par un argument topologique que nous omettons). $X'$
est une union disjointe de régimes intégrés normaux par Propriétés, lemme \ref{properties-lemma-normal-locally-finite-nr-irreducibles}. Il ressort clairement de la description de $X'$ au-dessus de cela $Y \to X'$
est dominante et provoque une bijection sur les composantes irréductibles
$V \to U'$ pour chaque ouvert affine $U \subset X$. Le bijection des composantes irréductibles pour le morphisme $Y \to X'$
en découle un argument topologique (omis).

\end{proof}

\begin{lemma}

\label{lemma-nagata-normalization-finite-general}

Laisser $f : X \to S$ être un morphisme. Supposons que

\begin{enumerate}

\item  $S$ est une école de Nagata,
\item  $f$ est quasi compacte et quasi séparée,
\item quasi-compact ouvre de $X$ avoir finiment de nombreux composantes irréductibles,
\item si $x \in X$ est un point général d'un composite irréductible, puis l'extension de champ $\kappa(x)/\kappa(f(x))$ est finiement générée, et
\item  $X$ est réduit.

\end{enumerate}

Ensuite, la normalisation $\nu : S' \to S$ de $S$ en $X$ C'est fini.

\end{lemma}

\begin{proof}

Il y a une réduction immédiate de l'affaire $S = \Spec(R)$
où $R$ est un anneau Nagata par hypothèse (1). Nous devons montrer que la fermeture intégrale $A$ de $R$ en $\Gamma(X, \mathcal{O}_X)$ est fini sur $R$Depuis $f$ est quasi-compact par hypothèse (2) nous pouvons écrire
$X = \bigcup_{i = 1, \ldots, n} U_i$ avec chacun $U_i$ Affinez. $U_i = \Spec(B_i)$. Chacun $B_i$ est réduit par hypothèse (5) et a finiment de nombreux premiers minimums
$\mathfrak q_{i1}, \ldots, \mathfrak q_{im_i}$
par hypothèse (3) et Algèbre, lemme \ref{algebra-lemma-irreducible}Nous avons
$$
\Gamma(X, \mathcal{O}_X) \subset B_1 \times \ldots \times B_n
\subset
\prod\nolimits_{i = 1, \ldots, n}
\prod\nolimits_{j = 1, \ldots, m_i} (B_i)_{\mathfrak q_{ij}}
$$
la deuxième inclusion d'Algèbre, lemme \ref{algebra-lemma-reduced-ring-sub-product-fields}Nous avons $\kappa(\mathfrak q_{ij}) = (B_i)_{\mathfrak q_{ij}}$ par Algèbre, lemme \ref{algebra-lemma-minimal-prime-reduced-ring}. D'où la fermeture intégrale $A$ de $R$ en $\Gamma(X, \mathcal{O}_X)$
est contenu dans le produit des fermetures intégrales $A_{ij}$ de
$R$ en $\kappa(\mathfrak q_{ij})$Depuis $R$ est Noétherian il suffit de montrer que $A_{ij}$ est un fini $R$-module pour chacun $i, j$. Laisse $\mathfrak p_{ij} \subset R$ être l'image de $\mathfrak q_{ij}$. Comme $\kappa(\mathfrak q_{ij})/\kappa(\mathfrak p_{ij})$
est une extension de champ finiment générée par hypothèse (4), nous voyons que $R \to \kappa(\mathfrak q_{ij})$ est essentiellement de type fini. $R \to A_{ij}$ est fini par Algèbre, lemme
\ref{algebra-lemma-nagata-in-reduced-finite-type-finite-integral-closure}.

\end{proof}

\begin{lemma}

\label{lemma-nagata-normalization-finite}
Laissez $f : X \to S$ être un morphisme. Supposons que
\begin{enumerate}
\item $S$ est un schéma de Nagata,
\item $f$ est de type fini,
\item $X$ est réduit.
\end{enumerate}
Puis la normalisation $\nu : S' \to S$ de $S$ dans $X$ est finie.
\end{lemma}

\begin{proof}

Il s'agit d'un cas particulier de lemme. \ref{lemma-nagata-normalization-finite-general}. A savoir, (2) tient comme le type de fini morphisme $f$ est quasi-compact par définition et quasi-séparé par lemme \ref{lemma-finite-type-Noetherian-quasi-separated}. La condition (3) est maintenue parce que $X$ est localement noethérien par lemme \ref{lemma-finite-type-noetherian}. Enfin, la condition (4) tient parce qu'un morphisme de type fini induit des extensions de corps résiduel finies.

\end{proof}

\begin{lemma}

\label{lemma-relative-normalization-normal-codim-1}
Laissez $f : Y \to X$ être un type de fini morphisme des schémas avec $Y$ réduit et $X$ Nagata. Laissez $X'$ être la normalisation de $X$ dans $Y$. Laissez $x' \in X'$ être un point tel que
\begin{enumerate}
\item $\dim(\mathcal{O}_{X', x'}) = 1$, et
\item la fibre de $Y \to X'$ sur $x'$ est vide.
\end{enumerate}
Alors $\mathcal{O}_{X', x'}$ est un anneau de valorisation distinct.
\end{lemma}
X

\begin{proof}
Nous pouvons remplacer $X$ par un quartier affiné de l'image de $x'$. Par conséquent, nous pouvons supposer $X = \Spec(A)$ avec $A$ Nagata. Par lemme \ref{lemma-nagata-normalization-finite} le morphisme $X' \to X$ est fini. Par conséquent, nous pouvons écrire $X' = \Spec(A')$ pour un $A$-algèbre fini $A'$. Par lemme \ref{lemma-normalization-is-normalization} après avoir remplacé $X$ par $X'$ nous réduisons au cas décrit dans le paragraphe suivant.

\medskip\noindent

L'affaire $X = X' = \Spec(A)$ avec $A$ Noétherien. $\mathfrak p \subset A$ être l'idéal premier correspondant à notre point $x'$. Choisissez $g \in \mathfrak p$
non contenus dans un minimum de $A$ (utiliser la première évitement et le fait que $A$ a finiment beaucoup de primes minimales, voir Algèbre, lemmes \ref{algebra-lemma-silly} et
\ref{algebra-lemma-Noetherian-irreducible-components}). Définir $Z = f^{-1}V(g) \subset Y$; il s'agit d'un sous-système fermé de $Y$. Ensuite $f(Z)$ ne contient aucun point général par choix de $g$
et ne contient pas $x'$ parce que $x'$ n'est pas à l'image de $f$. La fermeture de $f(Z)$ est l'ensemble de spécialisations de points de $f(Z)$ par lemme \ref{lemma-reach-points-scheme-theoretic-image}. Ainsi la fermeture de $f(Z)$ ne contient pas $x'$ parce que l'état $\dim(\mathcal{O}_{X', x'}) = 1$ implique uniquement les points génériques de $X = X'$ spécialisé à $x'$. En d'autres termes, après avoir remplacé $X$ par un quartier affiné d'ouverture de $x'$ nous pouvons supposer que $f^{-1}V(g) = \emptyset$. Ainsi $g$ des cartes vers une fonction globale inversible sur $Y$ et nous obtenons une factorisation
$$
A \to A_g \to \Gamma(Y, \mathcal{O}_Y)
$$
Depuis $X = X'$ cela implique que $A$ est égal à la fermeture intégrale de $A$ en $A_g$. Par Algèbre, lemme \ref{algebra-lemma-integral-closure-localize}
nous concluons que $A_\mathfrak p$ est la fermeture intégrale de $A_\mathfrak p$ en $A_\mathfrak p[1/g]$. Par notre choix $g$Depuis $\dim(A_\mathfrak p) = 1$
et depuis $A$ est réduit nous voyons que $A_\mathfrak p[1/g]$
est un produit fini de champs (le produit des corps résiduels des premiers minimums contenus dans $\mathfrak p$). D'où $A_\mathfrak p$ est normale (Algèbre, lemme \ref{algebra-lemma-characterize-reduced-ring-normal}) et la démonstration est complète. Certains détails omis.

\end{proof}






\section{Normalisation}
\label{section-normalization}

\noindent
Ensuite, nous arrivons à la normalisation d'un schéma $X$. Nous ne le définissons/construisons que lorsque $X$ a localement de nombreux composants irréductibles.Soit $X$ un schéma tel que chaque quasi-compact ouvert ait finiment de nombreux composantes irréductibles.Soit $X^{(0)} \subset X$ l'ensemble de points génériques de composantes irréductibles de $X$.Soit
\begin{equation}
\label{equation-generic-points}
f :
Y = \coprod\nolimits_{\eta \in X^{(0)}} \Spec(\kappa(\eta))
\longrightarrow
X
\end{equation}
l'inclusion des points génériques dans $X$ en utilisant les applications canoniques de Schémas, section \ref{schemes-section-points}. Notez que ce morphisme est quasi-compact par hypothèse et quasi-séparé comme $Y$ est séparé (voir Schémas, section \ref{schemes-section-separation-axioms}).

\begin{definition}

\label{definition-normalization}
Laissez $X$ être un schéma tel que chaque quasi-open compacte a finiment beaucoup de composantes irréductibles. Nous définissons la {\it normalisation} de $X$ comme le morphisme $$
\nu : X^\nu \longrightarrow X
$$ qui est la normalisation de $X$ dans le morphisme $f : Y \to X$ (\ref{equation-generic-points}) construit ci-dessus.
\end{definition}

\noindent
Tout schéma localement noethérien a un ensemble local fini de composants irréductibles et la définition s'y applique. Habituellement, la normalisation est définie uniquement pour les schémas réduits. Avec la définition ci-dessus, la normalisation de $X$ est la même que la normalisation de la réduction de $X_{red}$ de $X$.

\begin{lemma}

\label{lemma-normalization-reduced}
Laissez $X$ être un schéma tel que chaque quasi-open compacte a finiment de nombreux composantes irréductibles. Le morphisme de normalisation $\nu$ facteurs par la réduction $X_{red}$ et $X^\nu \to X_{red}$ est la normalisation de $X_{red}$.
\end{lemma}

\begin{proof}
Laissez $f : Y \to X$ être le morphisme (\ref{equation-generic-points}). Nous obtenons une factorisation $Y \to X_{red} \to X$ de $f$ de Schémas, lemme \ref{schemes-lemma-map-into-reduction}. Par lemme \ref{lemma-characterize-normalization} nous obtenons un morphisme canonique $X^\nu \to X_{red}$ et que $X^\nu$ est la normalisation de $X_{red}$ dans $Y$. La lemme qui suit comme $Y \to X_{red}$ est identique au morphisme (\ref{equation-generic-points}) construit pour $X_{red}$.
\end{proof}

\noindent
Si $X$ est réduit, alors la normalisation de $X$ est la même que le spectre relatif de la fermeture intégrale de $\mathcal{O}_X$ dans la bandeau des fonctions méromorphiques $\mathcal{K}_X$ (voir Diviseurs, section \ref{divisors-section-meromorphic-functions}). À savoir, $\mathcal{K}_X = f_*\mathcal{O}_Y$ dans ce cas, voir Diviseurs, lemme \ref{divisors-lemma-reduced-finite-irreducible} et sa preuve.

\begin{lemma}

\label{lemma-description-normalization}

Laisser $X$ soit un schéma réduit de telle sorte que chaque quasi-compact ouvert ait finiment de nombreux composantes irréductibles. $\Spec(A) = U \subset X$
être un ouvert affine.

\begin{enumerate}

\item  $A$ a finiment beaucoup de primes minimales
$\mathfrak q_1, \ldots, \mathfrak q_t$,
\item l'anneau total des fractions $Q(A)$ de $A$ est
$Q(A/\mathfrak q_1) \times \ldots \times Q(A/\mathfrak q_t)$,
\item la fermeture intégrale $A'$ de $A$ en $Q(A)$ est le produit des fermetures intégrales des domaines $A/\mathfrak q_i$
dans les champs $Q(A/\mathfrak q_i)$,
\item  $\nu^{-1}(U)$ est identifié avec le spectre de $A'$ où
$\nu : X^\nu \to X$ est le morphisme de la normalisation.

\end{enumerate}

\end{lemma}

\begin{proof}

Les primes minimales correspondent à des composantes irréductibles (Algèbre, lemme) \ref{algebra-lemma-irreducible}), donc nous avons (1) par hypothèse.
$(0) = \mathfrak q_1 \cap \ldots \cap \mathfrak q_t$ parce que $A$ est réduit (Algèbre, lemme \ref{algebra-lemma-Zariski-topology}). Alors nous avons
$Q(A) = \prod A_{\mathfrak q_i} = \prod \kappa(\mathfrak q_i)$
par Algèbre, lemmes \ref{algebra-lemma-total-ring-fractions-no-embedded-points}
et \ref{algebra-lemma-minimal-prime-reduced-ring}. Cela prouve (2). La partie 3) découle d'Algèbre, lemme \ref{algebra-lemma-characterize-reduced-ring-normal}, ou lemme \ref{lemma-normalization-in-disjoint-union}. La partie 4 tient parce qu'il est clair que $f^{-1}(U) \to U$ est le morphisme
$$
\Spec\left(\prod \kappa(\mathfrak q_i)\right)
\longrightarrow
\Spec(A)
$$
où $f : Y \to X$ est le morphisme (\ref{equation-generic-points}).

\end{proof}

\begin{lemma}

\label{lemma-stalk-normalization}

Laisser $X$ être un schéma tel que chaque quasi-compact ouvert a un nombre fini de composantes irréductibles. $\nu : X^\nu \to X$
soit la normalisation de $X$. Laisse $x \in X$. Ensuite, les éléments suivants sont canoniquement isomorphiques comme $\mathcal{O}_{X, x}$-Algèbres

\begin{enumerate}

\item la tige $(\nu_*\mathcal{O}_{X^\nu})_x$,
\item la fermeture intégrale $\mathcal{O}_{X, x}$ dans le cycle total des fractions de $(\mathcal{O}_{X, x})_{red}$,
\item la fermeture intégrale $\mathcal{O}_{X, x}$ dans le produit des corps résiduels des primes minimales de $\mathcal{O}_{X, x}$
(et il y en a finiment beaucoup).

\end{enumerate}

\end{lemma}

\begin{proof}

Après remplacement $X$ par un quartier affiné d'ouverture de $x$ nous pouvons supposer que $X$ a finiment beaucoup de composantes irréductibles et que $x$ est contenu dans chacun d'eux. Puis la tige $(\nu_*\mathcal{O}_{X^\nu})_x$ est la fermeture intégrale de $A = \mathcal{O}_{X, x}$ dans le produit $L$
des corps résiduels des primes minimales de $A$. Cela résulte de la construction de la normalisation et de la \ref{lemma-integral-closure}. Vous pouvez également utiliser lemme
\ref{lemma-description-normalization}
et le fait que la normalisation se déplace avec la localisation (Algèbre, lemme \ref{algebra-lemma-integral-closure-localize}Depuis $A_{red}$ a finiment beaucoup de primes minimales (car celles-ci correspondent exactement aux points génériques des composantes irréductibles de $X$ de passage $x$) nous voyons que $L$ est l'anneau total des fractions de $A_{red}$
(Algèbre, lemme) \ref{algebra-lemma-total-ring-fractions-no-embedded-points}). Ainsi notre bague est aussi la fermeture intégrale de $A$ dans le cycle total des fractions de $A_{red}$.

\end{proof}

\begin{lemma}

\label{lemma-normalization-normal}

Laisser $X$ être un schéma tel que chaque quasi-compact ouvert a finiment beaucoup de composantes irréductibles.

\begin{enumerate}

\item La normalisation $X^\nu$ est une union disjointe de régimes normaux intégrés.
\item Le morphisme $\nu : X^\nu \to X$ est intégrale, surjective et induit une bijection sur les composantes irréductibles.
\item Pour tout morphisme entier $\alpha : X' \to X$ tel que pour
$U \subset X$ quasi-compact ouvrir l'image inverse $\alpha^{-1}(U)$ a finiment de nombreux composantes irréductibles et
$\alpha|_{\alpha^{-1}(U)} : \alpha^{-1}(U) \to U$ est birationnel\footnote{Cette formulation maladroite est nécessaire car nous n'avons défini que ce qu'elle signifie pour un morphisme à être birationnel si la source et la cible ont finiment de nombreux composantes irréductibles.
$X'_{red} \to X_{red}$ satisfait à la condition.} il existe une factorisation
$X^\nu \to X' \to X$ et $X^\nu \to X'$ est la normalisation de $X'$.
\item Pour tout morphisme $Z \to X$ avec $Z$ un régime normal tel que chaque composante irréductible de $Z$ domine un composante irréductible de $X$ il existe une factorisation unique $Z \to X^\nu \to X$.

\end{enumerate}

\end{lemma}

\begin{proof}

Laisser $f : Y \to X$ être comme dans (\ref{equation-generic-points}) — Le régime $X^\nu$ est une union disjointe de régimes intégrés normaux parce que $Y$ est normal et chaque ouvert affine de $Y$ a finiment beaucoup de composantes irréductibles, voir lemme \ref{lemma-normal-normalization}. Cela prouve (1). On peut aussi déduire (1) de lemmes \ref{lemma-normalization-reduced} et
\ref{lemma-description-normalization}.

\medskip\noindent
Le morphisme $\nu$ est intégré par lemme \ref{lemma-characterize-normalization}. Par lemme \ref{lemma-normal-normalization} le morphisme $Y \to X^\nu$ induit une bijection sur les composantes irréductibles, et par la construction de $Y$ cela implique que $X^\nu \to X$ induit une bijection sur les composantes irréductibles. Par la construction $f : Y \to X$ est dominante, donc aussi $\nu$ est dominante. Comme un morphisme entier est fermé (lem \ref{lemma-integral-universally-closed}) cela implique que $\nu$ est surjectif. Ceci prouve (2).

\medskip\noindent
Supposons que $\alpha : X' \to X$ soit comme dans (3). Il est clair que $X'$ satisfait aux hypothèses sous lesquelles la normalisation est définie.Soit $f' : Y' \to X'$ le morphisme (\ref{equation-generic-points}) construit à partir de $X'$. Comme $\alpha$ est localement birationnel, il est clair que $Y' = Y$ et $f = \alpha \circ f'$. D'où l'existence de la factorisation $X^\nu \to X' \to X$ et $X^\nu \to X'$ est la normalisation de $X'$ par lemme \ref{lemma-characterize-normalization}. Ceci prouve (3).

\medskip\noindent

Laisser $g : Z \to X$ être un morphisme dont le domaine est un schéma normal et tel que chaque composante irréductible domine un composante irréductible de $X$. Par lemme \ref{lemma-normalization-reduced}
Nous avons $X^\nu = X_{red}^\nu$ et par Schémas, lemme \ref{schemes-lemma-map-into-reduction}

$Z \to X$ des facteurs $X_{red}$. Par conséquent, nous pouvons remplacer $X$ par
$X_{red}$ et supposer $X$ est réduit. De plus, comme la factorisation est unique, il suffit de la construire localement sur $Z$. Laisse $W \subset Z$ et $U \subset X$ être ouverts affinés telle que $g(W) \subset U$. Écrire $U = \Spec(A)$ et
$W = \Spec(B)$Avec $g|_W$ données par $\varphi : A \to B$. Nous utiliserons les résultats de lemme \ref{lemma-description-normalization} librement. Laisser $\mathfrak p_1, \ldots, \mathfrak p_t$ être les premiers minima de $A$. Comme $Z$ c'est normal, nous voyons que $B$ est un anneau normal, en particulier réduit. $\mathfrak q \subset B$ Nous avons ça $\varphi^{-1}(\mathfrak q)$
est une prime minimale de $A$. Par conséquent, si $x \in A$ est un diviseur non zéro, c'est-à-dire,
$x \not \in \bigcup \mathfrak p_i$Alors $\varphi(x)$ est un non-zérodiviseur en $B$. Ainsi nous obtenons un homomorphisme canonique d'anneaux $Q(A) \to Q(B)$. Comme $B$ est normal il est égal à sa fermeture intégrale dans $Q(B)$ (voir Algèbre, lemme) \ref{algebra-lemma-normal-ring-integrally-closed}C'est pourquoi nous voyons que la fermeture intégrale $A' \subset Q(A)$ de $A$
des cartes dans $B$ via l'application canonique $Q(A) \to Q(B)$Depuis $\nu^{-1}(U) = \Spec(A')$ cela donne la factorisation canonique $W \to \nu^{-1}(U) \to U$ de $\nu|_W$. Nous omettons la vérification qu'elle est unique.

\end{proof}

\begin{lemma}

\label{lemma-normalization-in-terms-of-components}
Laissez $X$ être un schéma tel que chaque quasi-open compacte a finiment beaucoup de composantes irréductibles. Laissez $Z_i \subset X$, $i \in I$ être les composantes irréductibles de $X$ doté de la structure induite réduite. Laissez $Z_i^\nu \to Z_i$ être la normalisation. Alors $\coprod_{i \in I} Z_i^\nu \to X$ est la normalisation de $X$.
\end{lemma}

\begin{proof}

Nous pouvons supposer $X$ est réduit, voir lemme \ref{lemma-normalization-reduced}. Ensuite, le lemme suit soit de la description locale en lemme \ref{lemma-description-normalization}
ou à partir de lemme \ref{lemma-normalization-normal} partie 3) parce que
$\coprod Z_i \to X$ est intégrale et localement birationnelle (comme $X$ est réduit et possède localement de nombreux composantes irréductibles).

\end{proof}

\begin{lemma}

\label{lemma-normalization-birational}
Laissez $X$ être un schéma réduit avec des composantes irréductibles finis. Puis le morphisme de normalisation $X^\nu \to X$ est birationnel.
\end{lemma}

\begin{proof}

La normalisation induit une bijection de composantes irréductibles par lemme. \ref{lemma-normalization-normal}. Laisse $\eta \in X$ être un point général d'un composante irréductible de $X$ et laisser $\eta^\nu \in X^\nu$
être le point général de la composition correspondante irréductible de $X^\nu$. Ensuite $\eta^\nu \mapsto \eta$ et pour finir la démonstration nous devons montrer que
$\mathcal{O}_{X, \eta} \to \mathcal{O}_{X^\nu, \eta^\nu}$
est un isomorphisme, voir définition \ref{definition-birational}Parce que $X$ et $X^\nu$ sont réduits, nous voyons que les deux anneaux locaux sont égaux à leurs corps résiduels (Algèbre, lemme \ref{algebra-lemma-minimal-prime-reduced-ring}D'autre part, par la construction de la normalisation comme la normalisation de $X$ en $Y = \coprod \Spec(\kappa(\eta))$
nous voyons que nous avons
$\kappa(\eta) \subset \kappa(\eta^\nu) \subset \kappa(\eta)$
et la démonstration est complète.

\end{proof}

\begin{lemma}

\label{lemma-finite-birational-over-normal}

Un morphisme birationnel fini (ou même intégral) $f : X \to Y$
des régimes intégrés avec $Y$ normal est un isomorphisme.

\end{lemma}

\begin{proof}

Laisser $V \subset Y$ être un ouvert affine avec image inverse $U \subset X$ qui est une affinité d'ouverture aussi. $f$ est un morphisme birationnel de schémas intégrés, l'homomorphisme
$\mathcal{O}_Y(V) \to \mathcal{O}_X(U)$ est une application injective de domaines qui induit un isomorphisme des champs de fractions. $Y$ c'est normal, l'anneau $\mathcal{O}_Y(V)$ est entièrement fermé dans le champ de fraction. $f$ est fini (ou intégral) chaque élément de $\mathcal{O}_X(U)$
fait partie intégrante $\mathcal{O}_Y(V)$. Nous concluons que
$\mathcal{O}_Y(V) = \mathcal{O}_X(U)$. Cela prouve que $f$ est un isomorphisme tel que souhaité.

\end{proof}

\begin{lemma}

\label{lemma-normalization-trivial}
Laissez $X$ être un schéma avec localement de nombreux composantes irréductibles. Le morphisme de normalisation $\nu : X^\nu \to X$ est un isomorphisme si et seulement si $X$ est normal.
\end{lemma}

\begin{proof}

Si $\nu$ est un isomorphisme, alors $X$ est normal par lemme \ref{lemma-normalization-normal}. Inversement, supposer $X$ C'est normal. \ref{lemma-normalization-localization} et Propriétés, lemme \ref{properties-lemma-normal-locally-finite-nr-irreducibles}
Nous pouvons supposer $X$ C'est une partie intégrante. \ref{lemma-normalization-normal} le morphisme
$\nu$ fait partie intégrante et $X^\nu$ dispose d'un composante irréductible unique (il s'agit donc d'un schéma intégral). \ref{lemma-normalization-birational} et
\ref{lemma-finite-birational-over-normal}.

\end{proof}

\begin{lemma}

\label{lemma-Japanese-normalization}
Laissez $X$ être un schéma intégré, japonais. La normalisation $\nu : X^\nu \to X$ est un morphisme fini.
\end{lemma}

\begin{proof}

Suit de la définition (Propriétés, définition) \ref{properties-definition-nagata}) et lemme \ref{lemma-description-normalization}. A savoir, dans ce cas, le lemme dit que $\nu^{-1}(\Spec(A))$ est le spectre de la fermeture intégrale de $A$ dans son champ de fractions.

\end{proof}

\begin{lemma}

\label{lemma-nagata-normalization}
Laissez $X$ être un schéma de Nagata. La normalisation $\nu : X^\nu \to X$ est un morphisme fini.
\end{lemma}

\begin{proof}

Notez qu'une schéma de Nagata est localement noethérien, définition \ref{definition-normalization}
Le lemme est maintenant un cas particulier de lemme. \ref{lemma-nagata-normalization-finite-general}
mais nous pouvons aussi le prouver directement comme suit. $X^\nu \to X$ comme composition
$X^\nu \to X_{red} \to X$. Comme $X_{red} \to X$ est une immersion fermée elle est finie. Il suffit donc de prouver la lemme pour une réduction de schéma de Nagata (par lemme \ref{lemma-composition-finite}). Laisser $\Spec(A) = U \subset X$ être un ouvert affine. \ref{lemma-description-normalization} Nous avons
$\nu^{-1}(U) = \Spec(\prod A_i')$ où $A_i'$ est la fermeture intégrale de $A/\mathfrak q_i$ dans son champ de fractions. $A$ est une bague Nagata (voir Propriétés, lemme \ref{properties-lemma-locally-nagata}) chacune des extensions d'anneau
$A/\mathfrak q_i \subset A'_i$ sont finis. $A \to \prod A'_i$
est un homomorphisme fini d'anneaux et nous gagnons.

\end{proof}

\begin{lemma}

\label{lemma-normalization-geom-unibranch-univ-homeo}

Laisser $X$ être un schéma unibranchique géométriquement irréductible. Le morphisme de normalisation $\nu : X^\nu \to X$ est un homéomorphisme universel.

\end{lemma}

\begin{proof}

Nous devons montrer que $\nu$ est intégrale, universellement injectable, et surjectif, voir lemme \ref{lemma-universal-homeomorphism}. Par lemme \ref{lemma-normalization-normal} le morphisme $\nu$
C'est une partie intégrante. $x \in X$ et mise en place $A = \mathcal{O}_{X, x}$Depuis $X$ est irréductible nous voyons que $A$ a une seule prime minimale $\mathfrak p$ et $A_{red} = A/\mathfrak p$. Par lemme \ref{lemma-stalk-normalization} la tige
$A' = (\nu_*\mathcal{O}_{X^\nu})_x$ est la fermeture intégrale de
$A$ dans le champ de fraction de $A_{red}$. Par Compléments sur l'algèbre, définition \ref{more-algebra-definition-unibranch}
nous voyons que $A'$ a un seul premier $\mathfrak m'$ couché sur $\mathfrak m_x \subset A$ et $\kappa(\mathfrak m')/\kappa(x)$
est purement inséparable. $\nu$ est bijectif (d'où surjectif) et universellement injectable par \ref{lemma-universally-injective}.

\end{proof}







\section{Normalisation faible}
\label{section-weak-normalization}

\noindent
Nous ne définirons la normalisation impossible d'un schéma que lorsqu'il a localement de nombreux composantes irréductibles, comme dans le cas de la normalisation.

\begin{lemma}

\label{lemma-relative-weak-normalization-algebra}
Laissez $A \to B$ être un homomorphisme d'anneaux induisant un morphisme dominant $\Spec(B) \to \Spec(A)$ des spectres. Il existe un $A$-subalgèbre $B' \subset B$ tel que
\begin{enumerate}
\item $\Spec(B') \to \Spec(A)$ est un homéomorphisme universel,
\item donné une factorisation $A \to C \to B$ telle que $\Spec(C) \to \Spec(A)$ est un homéomorphisme universel, l'image de $C \to B$ est contenue dans $B'$.
\end{enumerate}

\end{lemma}

\begin{proof}
Nous allons utiliser lemme \ref{lemma-reduction-universal-homeomorphism} sans autre mention. Considérons le diagramme commutatif $$
\xymatrix{
B \ar[r] & B_{red} \\
A \ar[u] \ar[r] & A_{red} \ar[u]
}
$$ Pour toute factorisation $A \to C \to B$ de $A \to B$ comme dans (2), nous voyons que $A_{red} \to C_{red} \to B_{red}$ est une factorisation de $A_{red} \to B_{red}$ comme dans (2). Il s'ensuit que si la lemme tient pour $A_{red} \to B_{red}$ et produit la $A_{red}$-subalgebra $B'_{red} \subset B_{red}$, puis le réglage de $B' \subset B$ égal à l'image inverse de $B'_{red}$ résout la lemme pour $A \to B$. Cela nous réduit au cas discuté dans le paragraphe suivant.

\medskip\noindent

Présumer $A$ et $B$ dans ce cas-ci $A \subset B$ par Algèbre, lemme \ref{algebra-lemma-image-dense-generic-points}. Laisse $A \to C \to B$ être une factorisation comme dans (2). Alors nous pouvons appliquer la proposition \ref{proposition-universal-homeomorphism}
à $A \subset C$ de voir que chaque élément de $C$ est contenu dans une extension $A[c_1, \ldots, c_n] \subset C$ tel que pour
$i = 1, \ldots, n$ Nous avons

\begin{enumerate}

\item  $c_i^2, c_i^3 \in A[c_1, \ldots, c_{i - 1}]$, ou
\item il existe un nombre premier $p$ avec
$pc_i, c_i^p \in A[c_1, \ldots, c_{i - 1}]$.

\end{enumerate}

Ainsi la propriété (2) détient si nous définissons $B' \subset B$ être le sous-ensemble d'éléments $b \in B$ qui sont contenues dans une extension
$
A[b_1, \ldots, b_n] \subset B
$
tel que (*) détient: pour $i = 1, \ldots, n$ Nous avons

\begin{enumerate}

\item  $b_i^2, b_i^3 \in A[b_1, \ldots, b_{i - 1}]$, ou
\item il existe un nombre premier $p$ avec
$pb_i, b_i^p \in A[b_1, \ldots, b_{i - 1}]$.

\end{enumerate}

Il n'y a que deux choses à vérifier : a) $B'$ est un $A$-subalgèbre, et b) $\Spec(B') \to \Spec(A)$ est un homéomorphisme universel. La partie (a) suit parce que donné $n \geq 0$ et $b_1, \ldots, b_n \in B$
satisfaisant (*) et $m \geq 0$ et $b'_1, \ldots, b'_m \in B$
satisfaisant (*), le nombre entier
$n + m$ et $b_1, \ldots, b_n, b'_1, \ldots, b'_m \in B$
Enfin, la partie b) tient par proposition \ref{proposition-universal-homeomorphism} et notre construction de $B'$.

\end{proof}

\begin{lemma}

\label{lemma-relative-weak-normalization-localization}

Laisser $A \to B$ être un homomorphisme d'anneaux induisant un morphisme dominant
$\Spec(B) \to \Spec(A)$ de spectres. Formation du $A$- sous-algèbre
$B' \subset B$ en lemme \ref{lemma-relative-weak-normalization-algebra}
les déplacements avec localisation (voir la démonstration pour les explications).

\end{lemma}

\begin{proof}
Laissez $S \subset A$ être un sous-ensemble multiplicatif. Alors $S^{-1}A \to S^{-1}B$ est un homomorphisme d'anneaux qui induit un morphisme dominant $\Spec(S^{-1}B) \to \Spec(S^{-1}A)$ (voir lemmes \ref{lemma-base-change-dominant} et \ref{lemma-generalizations-lift-flat}). Par conséquent, lemme \ref{lemma-relative-weak-normalization-algebra} produit un $S^{-1}A$-subalgèbre $(S^{-1}B)' \subset S^{-1}B$. L'énoncé signifie que $S^{-1}B' = (S^{-1}B)'$ sous $S^{-1}A$-subalgèbres de $S^{-1}B$.

\medskip\noindent
Pour voir cela est vrai, nous allons utiliser la construction de $B'$ et $(S^{-1}B)'$ dans la démonstration de lemme \ref{lemma-relative-weak-normalization-algebra}. Dans la première étape, nous voyons que $B'$ est l'image inverse de la $A_{red}$-subalgèbre $B'_{red} \subset B_{red}$ construite pour l'homomorphisme d'anneaux $A_{red} \to B_{red}$ et de la même façon pour $(S^{-1}B)'$. Notant que $S^{-1}B_{red} = (S^{-1}B)_{red}$ cela nous réduit à l'affaire discutée dans le paragraphe suivant.

\medskip\noindent
Si $A$ et $B$ sont réduits, nous avons construit $B'$ comme l'union des sous-algèbres $A[b_1, \ldots, b_n]$ de telle sorte que pour $i = 1, \ldots, n$ nous avons
\begin{enumerate}
\item $b_i^2, b_i^3 \in A[b_1, \ldots, b_{i - 1}]$, ou
\item il existe un numéro principal $p$ avec $pb_i, b_i^p \in A[b_1, \ldots, b_{i - 1}]$.
\end{enumerate}
De même pour $(S^{-1}B)' \subset S^{-1}B$. Il est donc clair que l'image de $B' \to B \to S^{-1}B$ est contenue dans $(S^{-1}B)'$. Pour montrer que l'application correspondante $S^{-1}B' \to (S^{-1}B)'$ est surjective, on utilise lemme \ref{lemma-special-elements-and-localization} pour dégager successivement les dénominateurs; nous omettons les détails.
\end{proof}

\begin{lemma}

\label{lemma-relative-seminormalization-algebra}
Laissez $A \to B$ être un homomorphisme d'anneaux induisant un morphisme dominant $\Spec(B) \to \Spec(A)$ des spectres. Il existe un $A$-subalgèbre $B' \subset B$ tel que
\begin{enumerate}
\item $\Spec(B') \to \Spec(A)$ est un homéomorphisme universel induisant des isomorphismes sur les corps résiduels,
\item donné une factorisation $A \to C \to B$ tel que $\Spec(C) \to \Spec(A)$ est un homéomorphisme universel induisant des isomorphismes sur les corps résiduels, l'image de $C \to B$ est contenue dans $B'$.
\end{enumerate}

\end{lemma}

\begin{proof}
Cette démonstration est exactement la même que la démonstration de lemme \ref{lemma-relative-weak-normalization-algebra} sauf que nous utilisons la proposition \ref{proposition-universal-homeomorphism-equal-residue-fields} au lieu de la proposition \ref{proposition-universal-homeomorphism}
\end{proof}

\begin{lemma}

\label{lemma-relative-seminormalization-localization}

Laisser $A \to B$ être un homomorphisme d'anneaux induisant un morphisme dominant
$\Spec(B) \to \Spec(A)$ de spectres. Formation du $A$- sous-algèbre
$B' \subset B$ en lemme \ref{lemma-relative-seminormalization-algebra}
les déplacements avec localisation (voir la démonstration pour les explications).

\end{lemma}

\begin{proof}
La démonstration est la même que la démonstration de lemme \ref{lemma-relative-weak-normalization-localization}.
\end{proof}

\begin{lemma}

\label{lemma-relative-seminormalization}

Laisser $f : Y \to X$ être un quasi-compact, quasi-séparé, et morphisme dominant des schémas.

\begin{enumerate}

\item La catégorie des facteurs $Y \to X' \to X$ où $X' \to X$
est un homéomorphisme universel a un objet initial $Y \to X^{Y/wn} \to X$.
\item La catégorie des facteurs $Y \to X' \to X$ où $X' \to X$
est un homéomorphisme universel inducteur d'isomorphismes sur les corps résiduels a un objet initial $Y \to X^{Y/sn} \to X$.

\end{enumerate}

En outre, la formation de la factorisation $Y \to X^{Y/wn} \to X$ et
$Y \to X^{Y/sn} \to X$ commutations avec changement de base pour ouvrir des sous-régimes de $X$.

\end{lemma}

\begin{proof}
Nous prouverons (1) et omettrons la démonstration de (2); aussi l'affirmation finale découlera de la construction de la factorisation. Nous utiliserons lemme \ref{lemma-universal-homeomorphism} sans autre mention. Premièrement. Soit $Y \to X^{Y/n} \to X$ la normalisation de $X$ dans $Y$, voir définition \ref{definition-normalization-X-in-Y}. Pour $Y \to X' \to X$ comme dans (1), nous obtenons un morphisme unique $X^{Y/n} \to X'$ compatible avec les morphismes donnés, voir lemme \ref{lemma-characterize-normalization}. Ainsi, il suffit de prouver la lemme avec $f$ remplacé par $X^{Y/n} \to X$. Cela nous réduit au cas étudié dans le paragraphe suivant.

\medskip\noindent
Supposons que $f$ est une partie intégrante (le reste de la démonstration fonctionne plus généralement si $f$ est affine).Soit $U = \Spec(A)$ une affine ouverte de $X$ et que $V = f^{-1}(U) = \Spec(B)$ soit l'image inverse dans $Y$. Alors $A \to B$ est un homomorphisme d'anneaux qui induit un morphisme dominant sur les spectres. Par lemme \ref{lemma-relative-weak-normalization-algebra} nous obtenons un $A$-subalgèbre $B' \subset B$ tel que le réglage de $U^{V/wn} = \Spec(B')$ la factorisation $V \to U^{V/wn} \to U$ est initial dans la catégorie des factorisations $V \to U' \to U$ où $U' \to U$ est un homéomorphisme universel.

\medskip\noindent

Si $U_1 \subset U_2 \subset X$ sont ouverts affines, puis paramétrage
$V_i = f^{-1}(U_i)$ nous obtenons un morphisme canonique
$$
\rho_{U_1}^{U_2} : U_1^{V_1/wn} \to U_1 \times_{U_2} U_2^{V_2/wn}
$$
terminé $U_1$ par la propriété universelle de $U_1^{V_1/wn}$. Ces morphismes satisfont une functorialité naturelle que nous laissons au lecteur de formuler et de prouver. $\rho_{U_1}^{U_2}$
est un isomorphisme; cela découle de lemme \ref{lemma-relative-weak-normalization-localization}
à condition que $U_1 \subset U_2$ est un principal ouvert et dans le cas général peut être réduit à ce cas par la nature functoire de ces cartes et Schémas, lemme \ref{schemes-lemma-standard-open-two-affines}
(détails omis). Ainsi par collage relatif (Constructions, lemme \ref{constructions-lemma-relative-glueing}) nous obtenons un morphisme $X^{Y/wn} \to X$ qui se limite $U^{V/wn} \to U$ terminé $U$ Compatible avec le $\rho_{U_1}^{U_2}$Bien sûr, les morphismes $V \to U^{V/wn}$ colle à un morphisme $Y \to X^{Y/wn}$
(voir Constructions, remarque)
\ref{constructions-remark-relative-glueing-functorial}) et nous obtenons notre factorisation $Y \to X^{Y/wn} \to X$ où le second morphisme est un homéomorphisme universel.

\medskip\noindent

Enfin, laissons $Y \to X' \to X$ être une factorisation comme dans (1). Avec $V \to U^{V/wn} \to U \subset X$ comme ci-dessus, nous obtenons une factorisation $V \to U \times_X X' \to U$ où la deuxième flèche est un homéomorphisme universel et nous obtenons un morphisme unique $g_U : U^{V/wn} \to U \times_X X'$ terminé $U$
par la propriété universelle de $U^{V/wn}$. Ces $g_U$ sont compatibles avec les morphismes $\rho_{U_1}^{U_2}$; détails omis. Il y a donc un morphisme unique $g : X^{Y/wn} \to X'$
terminé $X$ d ' accord avec $g_U$ terminé $U$, voir Constructions, remarque \ref{constructions-remark-relative-glueing-functorial}. Cela prouve que $Y \to X^{Y/wn} \to X$ est initiale dans notre catégorie et la démonstration est complète.

\end{proof}

\begin{definition}

\label{definition-relative-seminormalization}
Laissez $f : Y \to X$ être un morphisme quasi-compact, quasi-séparé et dominant des schémas.
\begin{enumerate}
\item La factorisation $Y \to X^{Y/sn} \to X$ construite en lemme \ref{lemma-relative-seminormalization} partie (2) est la {\it semi-normalisation de $X$ dans $Y$}.
\item La factorisation $Y \to X^{Y/wn} \to X$ construite en lemme \ref{lemma-relative-seminormalization} partie (1) est la {\it normalisation de $X$ dans $Y$}.
\end{enumerate}

\end{definition}

\noindent
Voici une façon de réinterpréter la semi-normalisation d'un schéma qui a localement de nombreux composantes irréductibles.

\begin{lemma}

\label{lemma-seminormal-and-normalization}
Laissez $X$ être un schéma tel que chaque quasi-open compacte a finiment beaucoup de composantes irréductibles. Laissez $\nu : X^\nu \to X$ être la normalisation de $X$. Ensuite, la semi-normalisation de $X$ dans $X^\nu$ est la semi-normalisation de $X$. Dans une formule: $X^{sn} = X^{X^\nu/sn}$.
\end{lemma}

\begin{proof}

Laisser $f : Y \to X$ être comme dans (\ref{equation-generic-points}) de sorte que $X^\nu$ est la normalisation de $X$ en $Y$. La semi-normalisation $X^{sn} \to X$ de $X$ est l'objet initial dans la catégorie des homéomorphismes universels $X' \to X$ induisant des isomorphismes sur les corps résiduels. $Y$ est l'union disjointe des spectres des corps résiduels aux points génériques de composantes irréductibles de $X$, nous voyons cela pour n'importe quel $X' \to X$ dans cette catégorie nous obtenons un ascenseur canonique $f' : Y \to X'$ de $f$. Puis par lemme \ref{lemma-characterize-normalization}
nous obtenons un morphisme canonique $X^\nu \to X'$. D'où à son tour un morphisme canonique $X^{X^\nu/sn} \to X'$
par la propriété universelle de $X^{X^\nu/sn}$. De cette façon, nous voyons que $X^{X^\nu/sn}$ satisfait à la même propriété universelle que $X^{sn}$ a et nous concluons.

\end{proof}

\noindent
lemme \ref{lemma-seminormal-and-normalization} motive la définition suivante. Puisque nous n'avons construit la normalisation que dans le cas $X$ a localement de nombreux composantes irréductibles finis, nous nous limiterons également à ce cas pour la normalisation factible.

\begin{definition}

\label{definition-weak-normalization}
Laissez $X$ être un schéma tel que chaque quasi-open compacte a finiment beaucoup de composantes irréductibles. Nous définissons la {\it normalisation faible} de $X$ comme la normalisation faible $$
X^\nu \longrightarrow X^{wn} \longrightarrow X
$$ de $X$ dans la normalisation $X^\nu$ de $X$ (définition \ref{definition-normalization}). Dans une formule: $X^{wn} = X^{X^\nu/wn}$.
\end{definition}

\noindent
Combiné avec lemme \ref{lemma-seminormal-and-normalization} nous voyons que pour un schéma $X$ qui a localement de nombreux composantes irréductibles il y a des morphismes canoniques $$
X^\nu \to X^{wn} \to X^{sn} \to X
$$ Ayant fait cette définition, nous pouvons dire ce que cela signifie pour un schéma d'être faiblement normal (à condition qu'il ait localement de nombreux composants irréductibles).

\begin{definition}

\label{definition-weakly-normal}

Laisser $X$ être un schéma tel que chaque quasi-compact ouvert a finiment beaucoup de composantes irréductibles. $X$ est {\it faiblement normal} si la normalisation est impossible $X^{wn} \to X$ est un isomorphisme (définition \ref{definition-weak-normalization}).

\end{definition}

\noindent
Il s'ensuit immédiatement des définitions que pour un schéma $X$ de telle sorte que chaque quasi-open compacte a finiment beaucoup de composantes irréductibles nous avons $$
X\text{ normal} \Rightarrow
X\text{ weakly normal} \Rightarrow
X\text{ seminormal}
$$ Nous pouvons trouver la signification de la normale faible dans le cas affiné comme suit.

\begin{lemma}

\label{lemma-affine-weakly-normal}

Laisser $X = \Spec(A)$ être un schéma affiné qui a finiment beaucoup de composantes irréductibles. $X$ est faiblement normal si et seulement si

\begin{enumerate}

\item  $A$ est semi-normale (définition \ref{definition-seminormal-ring}),
\item pour un numéro principal $p$ et $z, w \in A$ de telle manière que a) $z$ est un diviseur non zéro, b) $w^p$ est divisible par $z^p$, et c) $pw$ est divisible par $z$Alors $w$ est divisible par $z$.

\end{enumerate}

\end{lemma}

\begin{proof}

Présumer $X$ Comme un schéma faiblement normal est semi-normal, nous voyons que (1) tient (par notre définition des schémas faiblement normaux). $A$ est réduit. $p, z, w$ être comme dans (2). Choisissez $x, y \in A$ de telle manière que $z^p x = w^p$ et $zy = pw$. Ensuite $p^p x = y^p$. L'homomorphisme d'anneaux
$A \to C = A[t]/(t^p - x, pt - y)$ induit un homéomorphisme universel sur les spectres. $X^\nu$ de $X$
est le spectre de la fermeture intégrale $A'$
de $A$ dans le cycle total des fractions de $A$, voir lemme \ref{lemma-description-normalization}. Notez que $a = w/z \in A'$ parce que $a^p = x$. Nous avons donc un $A$- homomorphisme de l'algèbre $A \to C \to A'$
envoi $t$ à $a$. À ce stade, la propriété définissant $X = X^{wn} = X^{X^\nu/wn}$
d'être faiblement normaux nous dit que $C \to A'$ des cartes dans $A$. Ainsi nous trouvons $a \in A$ comme souhaité.

\medskip\noindent

À l ' inverse, supposer les paragraphes 1 et 2. $A'$ Comme dans le paragraphe précédent, nous devons montrer que $X^{X^\nu/wn} = X$. Par construction dans la démonstration de lemme \ref{lemma-relative-weak-normalization-algebra}, le régime
$X^{X^\nu/wn}$ est le spectre du subring de $A'$
qui est l'union des sous-anneaux $A[a_1, \ldots, a_n] \subset A'$
tel que pour $i = 1, \ldots, n$ Nous avons

\begin{enumerate}

\item[(a)]  $a_i^2, a_i^3 \in A[a_1, \ldots, a_{i - 1}]$, ou
\item[(b)] il existe un nombre premier $p$ avec
$pa_i, a_i^p \in A[a_1, \ldots, a_{i - 1}]$.

\end{enumerate}

Ensuite, nous pouvons utiliser (1) et (2) pour voir inductivement que $a_1, \ldots, a_n \in A$; nous omettons les détails. $X = X^{X^\nu/wn}$ et donc
$X$ est faiblement normal.

\end{proof}

\noindent
Voici le lemme obligatoire.

\begin{lemma}

\label{lemma-locally-weakly-normal}

Laisser $X$ soit un schéma tel que chaque quasi-compact ouvert ait finiment de nombreux composantes irréductibles. Les conditions suivantes sont équivalentes:

\begin{enumerate}

\item Le régime $X$ est faiblement normal.
\item Pour chaque ouvert affine $U \subset X$ la bague $\mathcal{O}_X(U)$
remplit les conditions (1) et (2) de l'article \ref{lemma-affine-weakly-normal}.
\item Il existe un recouvrement ouvert affine $X = \bigcup U_i$ tel que chaque bague $\mathcal{O}_X(U_i)$
remplit les conditions (1) et (2) de l'article \ref{lemma-affine-weakly-normal}.
\item Il existe une ouverture de recouvrement $X = \bigcup X_j$
de telle sorte que chaque sous-système ouvert $X_j$ est faiblement normal.

\end{enumerate}

En outre, si $X$ est faiblement normal alors que chaque sous-système ouvert est faiblement normal.

\end{lemma}

\begin{proof}
La condition pour $X$ être faiblement normale est que le morphisme $X^{wn} = X^{X^\nu/wn} \to X$ est un isomorphisme. Depuis la construction de $X^\nu \to X$ navettes avec changement de base pour ouvrir des sous-systèmes et depuis la construction de $X^{X^\nu/wn}$ navettes avec changement de base pour ouvrir des sous-systèmes de $X$ (lemme \ref{lemma-relative-seminormalization}) la lemme est claire.
\end{proof}










\section{Théorème principal de Zariski (version algébrique)}
\label{section-Zariski}

\noindent
C'est la version que vous pouvez prouver en utilisant des méthodes purement algébriques. Avant de pouvoir prouver des versions plus puissantes (pour les affines non-morphismes) nous devons développer plus d'outils. Voir Cohomologie de Schémas, section \ref{coherent-section-applications-formal-functions} et Compléments sur les morphismes, section \ref{more-morphisms-section-application-etale-neighbourhoods}.

\begin{theorem}
[Version algébrique du théorème principal de Zariski]
\label{theorem-main-theorem}
Laissez $f : Y \to X$ être un morphisme affine de schémas. Supposons que $f$ est de type fini. Laissez $X'$ être la normalisation de $X$ dans $Y$. Image: $$
\xymatrix{
Y \ar[rd]_f \ar[rr]_{f'} & & X' \ar[ld]^\nu \\
& X &
}
$$ Ensuite, il existe un sous-système ouvert $U' \subset X'$ tel que
\begin{enumerate}
\item $(f')^{-1}(U') \to U'$ est un isomorphisme, et
\item $(f')^{-1}(U') \subset Y$ est l'ensemble de points où $f$ est quasi-fini.
\end{enumerate}

\end{theorem}

\begin{proof}

Il y a une réduction immédiate dans le cas où $X$ et donc $Y$
Ils sont affinés. $X = \Spec(R)$ et $Y = \Spec(A)$. Ensuite $X' = \Spec(A')$où: $A'$ est la fermeture intégrale de
$R$ en $A$, voir définitions \ref{definition-integral-closure}
et \ref{definition-normalization-X-in-Y}. Par Algèbre, théorème \ref{algebra-theorem-main-theorem}
pour chaque $y \in Y$ où $f$ est quasi-fini, il existe un ouvert $U'_y \subset X'$ de telle manière que $(f')^{-1}(U'_y) \to U'_y$
est un isomorphisme. $U' = \bigcup U'_y$ où $y \in Y$ plages sur tous les points où $f$ est quasi-finie.
$f$ est quasi-fini à tous les points de $(f')^{-1}(U')$Si $y \in (f')^{-1}(U')$ avec image $x \in X$, alors nous voyons que
$Y_x \to X'_x$ est un isomorphisme dans un quartier de $y$. Il n'y a donc pas de raison de $Y_x$ qui se spécialise dans $y$, puisque c'est vrai pour $f'(y)$ en $X'_x$, voir lemme \ref{lemma-integral-fibres}. Par lemme \ref{lemma-quasi-finite-at-point-characterize} partie 3) cela implique $f$ est quasi-fini à $y$.

\end{proof}

\noindent
Nous pouvons utiliser la version algébrique du théorème principal de Zariski pour montrer que l'ensemble de points où un morphisme est quasi-fini est ouvert.

\begin{lemma}

\label{lemma-quasi-finite-points-open}

\begin{slogan}
Le lieu quasi-fini d'un morphisme est ouvert
\end{slogan}
Laissez $f : X \to S$ être un morphisme de schémas. L'ensemble de points de $X$ où $f$ est quasi-fini est un $U \subset X$ ouvert. Le morphisme induit $U \to S$ est localement quasi-fini.
\end{lemma}

\begin{proof}
Supposons que $f$ soit quasi-fini à $x$. Laissez $x \in U = \Spec(A) \subset X$, $V = \Spec(R) \subset S$ être ouverts affines comme dans la définition \ref{definition-quasi-finite}. Par soit le théorème \ref{theorem-main-theorem} ci-dessus ou Algèbre, lemme \ref{algebra-lemma-quasi-finite-open}, l'ensemble des premiers $\mathfrak q$ à laquelle $R \to A$ est quasi-fini est ouvert dans $\Spec(A)$. Puisque tous correspondent à des points de $X$ où $f$ est quasi-fini nous obtenons la première déclaration. La deuxième déclaration est évidente.
\end{proof}

\noindent
Nous améliorerons les morphismes séparés quasi-finis généraux suivants plus tard, voir Compléments sur les morphismes, lemme \ref{more-morphisms-lemma-quasi-finite-separated-pass-through-finite}.

\begin{lemma}

\label{lemma-quasi-finite-affine}
Laissez $f : Y \to X$ être un morphisme de schémas. Supposons que
\begin{enumerate}
\item $X$ et $Y$ sont affines, et
\item $f$ est quasi-finie.
\end{enumerate}
Puis il existe un schéma $$
\xymatrix{
Y \ar[rd]_f \ar[rr]_j & & Z \ar[ld]^\pi \\
& X &
}
$$ avec $Z$ affines, $\pi$ finies et $j$ une ouverture d'immersion.
\end{lemma}

\begin{proof}
Ceci est Algèbre, lemme \ref{algebra-lemma-quasi-finite-open-integral-closure} reformulé dans la langue des schémas.
\end{proof}

\begin{lemma}

\label{lemma-image-nowhere-dense-quasi-finite}
Laissez $f : Y \to X$ être un morphisme quasi-fini de schémas. Laissez $T \subset Y$ être un sous-ensemble fermé nulle part dense de $Y$. Alors $f(T) \subset X$ est un sous-ensemble nulle part dense de $X$.
\end{lemma}

\begin{proof}

Comme dans la démonstration de lemme \ref{lemma-image-nowhere-dense-finite} cela se réduit immédiatement au cas où la base $X$ Dans ce cas, $Y = \bigcup_{i = 1, \ldots, n} Y_i$ est une union finie d'overts affines (comme $f$ est quasi-compact). Depuis chaque $T \cap Y_i$
n'est nulle part dense, et puisque l'union finie des ensembles denses nulle part n'est nulle part dense (voir Topologie, lemme \ref{topology-lemma-nowhere-dense}), il suffit de prouver que l'image $f(T \cap Y_i)$ est nulle part dense en $X$. Cela nous réduit au cas où les deux $X$ et $Y$ A ce stade, nous appliquons \ref{lemma-quasi-finite-affine} ci-dessus pour obtenir un diagramme
$$
\xymatrix{
Y \ar[rd]_f \ar[rr]_j & & Z \ar[ld]^\pi \\
& X &
}
$$
avec $Z$ affine, $\pi$ finis et $j$ une ouverture d'immersion. $\overline{T} = \overline{j(T)} \subset Z$Par Topologie, lemme \ref{topology-lemma-image-nowhere-dense-open}
nous voyons $\overline{T}$ est nulle part dense en $Z$Depuis $f(T) \subset \pi(\overline{T})$
le lemme suit du résultat correspondant dans le cas fini, voir lemme \ref{lemma-image-nowhere-dense-finite}.

\end{proof}




\section{Fibres universellement bornées}
\label{section-universally-bounded}

\noindent

Laisser $X$ être un régime sur un terrain $k$Si $X$ est fini sur $k$Alors $X = \Spec(A)$ où $A$ est un fini $k$- Algèbre. Une autre façon de dire ceci est que $X$ est fini localement libre sur $\Spec(k)$, voir définition \ref{definition-finite-locally-free}. D'où
$X \to \Spec(k)$ a une {\it degré} qui est un entier $d \geq 0$, à savoir $d = \dim_k(A)$. Nous appelons parfois cela le {\it degré} du régime (fini) $X$ terminé $k$.

\begin{definition}

\label{definition-universally-bounded}
Laissez $f : X \to Y$ être un morphisme de schémas.
\begin{enumerate}
\item Nous disons que l'entier $n$ {\it limite les degrés des fibres de $f$} si pour toutes les $y \in Y$ la fibre $X_y$ est un schéma fini sur $\kappa(y)$ dont le degré sur $\kappa(y)$ est $\leq n$.
\item Nous disons que les {\it fibres de $f$ sont universellement liés}\footnote{Ceci est probablement une notation non standard.} s'il existe un entier $n$ qui lie les degrés des fibres de $f$.
\end{enumerate}
.
\end{definition}
ZXZX

\noindent
Notez que, en particulier, le nombre de points dans une fibre est également limité par $n$. (L'inverse ne tient pas, même si toutes les fibres sont des schémas réduits finis.)

\begin{lemma}

\label{lemma-characterize-universally-bounded}

Laisser $f : X \to Y$ être un morphisme de schémas. $n \geq 0$. Les conditions suivantes sont équivalentes:

\begin{enumerate}

\item le nombre entier $n$ limite les degrés des fibres de $f$,
\item pour chaque morphisme $\Spec(k) \to Y$où: $k$ est un champ, le produit fibré $X_k = \Spec(k) \times_Y X$ est fini sur $k$
de degré $\leq n$.

\end{enumerate}

Dans ce cas, les fibres de $f$ sont universellement limités et les régimes
$X_k$ avoir au plus $n$ plus précisément, si
$X_k = \{x_1, \ldots, x_t\}$, alors nous avons
$$
n \geq \sum\nolimits_{i = 1, \ldots, t} [\kappa(x_i) : k]
$$

\end{lemma}

\begin{proof}

Implications (2) $\Rightarrow$ (1) est trivial. L'autre implication tient parce que si l'image de $\Spec(k) \to Y$ est $y$Alors
$X_k = \Spec(k) \times_{\Spec(\kappa(y))} X_y$. Par définition, les fibres de $f$ être universellement limité signifie que certains $n$ Enfin, supposez que $X_k = \Spec(A)$. Ensuite $\dim_k A = n$. D'où $A$ est Artinian, tous les idéaux principaux sont des idéaux maximaux $\mathfrak m_i$, $A$ est le produit des localisations à ces idéaux maximaux. Voir Algèbre, lemmes \ref{algebra-lemma-finite-dimensional-algebra}
et \ref{algebra-lemma-artinian-finite-length}. Ensuite $\mathfrak m_i$
correspond à $x_i$, nous avons
$A_{\mathfrak m_i} = \mathcal{O}_{X_k, x_i}$
et donc il y a une surjection
$A \to  \bigoplus \kappa(\mathfrak m_i) = \bigoplus \kappa(x_i)$
qui implique l'inégalité dans l'énoncé du lemme par l'algèbre linéaire.

\end{proof}

\begin{lemma}

\label{lemma-finite-locally-free-universally-bounded}
Si $f$ est un morphisme finiment localement libre de degré $d$, alors $d$ limite le degré des fibres de $f$.
\end{lemma}

\begin{proof}
Ceci est vrai parce que tout changement de base de $f$ est fini localement exempt de degré $d$ (lemme \ref{lemma-base-change-finite-locally-free}) et donc les fibres de $f$ ont tous un degré $d$.
\end{proof}

\begin{lemma}

\label{lemma-composition-universally-bounded}
Une composition de morphismes avec des fibres universelles est un morphisme avec des fibres universelles. Plus précisément, supposez que $n$ limite les degrés des fibres de $f : X \to Y$ et $m$ limite les degrés de $g : Y \to Z$. Puis $nm$ limite les degrés des fibres de $g \circ f : X \to Z$.
\end{lemma}

\begin{proof}

Laisser $f : X \to Y$ et $g : Y \to Z$ ont des fibres universellement nées. $\deg(X_y/\kappa(y)) \leq n$ pour tous $y \in Y$, et que
$\deg(Y_z/\kappa(z)) \leq m$ pour tous $z \in Z$. Laisse $z \in Z$ Par hypothèse, le régime
$Y_z$ est fini sur $\Spec(\kappa(z))$. En particulier, l'espace topologique sous-jacent de $Y_z$
est un ensemble discret fini. Les fibres du morphisme
$f_z : X_z \to Y_z$ sont les fibres de $f$ aux points correspondants de $Y$, qui sont des ensembles discrets finis par le raisonnement ci-dessus. Nous concluons donc que l'espace topologique sous-jacent de $X_z$ est un ensemble discret fini aussi bien. $X_z$ est un schéma affine (c'est un exercice agréable; il suit aussi par exemple de Propriétés, lemme \ref{properties-lemma-maximal-points-affine}
appliqué à l'ensemble de tous les points de $X_z$). Écrire $X_z = \Spec(A)$,
$Y_z = \Spec(B)$, $k = \kappa(z)$. Ensuite $k \to B \to A$
et nous savons que (a) $\dim_k(B) \leq m$, et b) pour chaque idéal maximal $\mathfrak m \subset B$ Nous avons
$\dim_{\kappa(\mathfrak m)}(A/\mathfrak mA) \leq n$. Nous prétendons que cela implique que $\dim_k(A) \leq nm$. Notez que $B$ est le produit de ses localisations $B_{\mathfrak m}$, par exemple parce que $Y_z$ est une union disjointe de $1$-les régimes de points, ou par Algèbre, lemmes \ref{algebra-lemma-finite-dimensional-algebra} et
\ref{algebra-lemma-artinian-finite-length}. Donc nous voyons que
$\dim_k(B) = \sum_{\mathfrak m} \dim_k(B_{\mathfrak m})$ et
$\dim_k(A) = \sum_{\mathfrak m} \dim_k(A_{\mathfrak m})$ dans les deux cas $\mathfrak m$ passe sur les idéaux maximaux de
$B$ (sans $A$Par ce qui précède, et le lemme de Nakayama (Algèbre, lemme \ref{algebra-lemma-NAK}) nous voyons que chacun $A_{\mathfrak m}$ est un quotient de
$B_{\mathfrak m}^{\oplus n}$ en tant que $B_{\mathfrak m}$- Module.
$\dim_k(A_{\mathfrak m}) \leq n \dim_k(B_{\mathfrak m})$. Tout mettre ensemble nous voyons que
$$
\dim_k(A) = \sum\nolimits_{\mathfrak m} \dim_ta	(A_{\mathfrak m})
\leq \sum\nolimits_{\mathfrak m} n \dim_k(B_{\mathfrak m})
= n \dim_k(B) \leq nm
$$
comme souhaité.

\end{proof}

\begin{lemma}

\label{lemma-base-change-universally-bounded}
Un changement de base d'un morphisme avec des fibres universelles est un morphisme avec des fibres universelles. Plus précisément, si $n$ limite les degrés des fibres de $f : X \to Y$ et $Y' \to Y$ est n'importe quel morphisme, alors les degrés des fibres du changement de base $f' : Y' \times_Y X \to Y'$ est également limité par $n$.
\end{lemma}

\begin{proof}
Ceci est clair du résultat de lemme \ref{lemma-characterize-universally-bounded}.
\end{proof}

\begin{lemma}

\label{lemma-descent-universally-bounded}
Laissez $f : X \to Y$ être un morphisme de schémas.Soit $Y' \to Y$ un morphisme de schémas, et que $f' : X' = X_{Y'} \to Y'$ soit le changement de base de $f$. Si $Y' \to Y$ est surjectif et que $f'$ a des fibres universellement nées, alors $f$ a des fibres universellement nées. Plus précisément, si $n$ limite le degré des fibres de $f'$, alors $n$ limite également les degrés des fibres de $f$.
\end{lemma}

\begin{proof}
Laissez $n \geq 0$ être un nombre entier limitant les degrés des fibres de $f'$. Nous prétendons que $n$ fonctionne aussi pour $f$. À savoir, si $y \in Y$ est un point, alors choisissez un point $y' \in Y'$ situé sur $y$ et observez que $$
X'_{y'} = \Spec(\kappa(y')) \times_{\Spec(\kappa(y))} X_y.
$$ Puisque $X'_{y'}$ est supposé fini de degré $\leq n$ sur $\kappa(y')$, il s'ensuit que $X_y$ est fini de degré $\leq n$ sur $\kappa(y)$. (Certains détails omis.)
\end{proof}
X

\begin{lemma}

\label{lemma-immersion-universally-bounded}

Une immersion a des fibres universelles.

\end{lemma}

\begin{proof}

Le nombre entier $n = 1$ fonctionne dans la définition.

\end{proof}

\begin{lemma}

\label{lemma-etale-universally-bounded}
Laissez $f : X \to Y$ être un morphisme \'etale des schémas. Laissez $n \geq 0$. Les suivants sont équivalents
\begin{enumerate}
\item le nombre entier $n$ limite les degrés des fibres,
\item pour chaque champ $k$ et morphisme $\Spec(k) \to Y$ le changement de base $X_k = \Spec(k) \times_Y X$ a à la plupart des points $n$, et
\item pour chaque $y \in Y$ et chaque fermeture algébrique séparable $\kappa(y) \subset \kappa(y)^{sep}$ le schéma $X_{\kappa(y)^{sep}}$ a à la plupart des points $n$.
\end{enumerate}

\end{lemma}

\begin{proof}
Ceci découle de \ref{lemma-characterize-universally-bounded} lemme et du fait que les fibres $X_y$ sont des unions disjointes des spectres des extensions de champ séparables finies de $\kappa(y)$, voir lemme \ref{lemma-etale-over-field}.
\end{proof}

\noindent
Avoir des fibres universelles est une notion absolue et non une notion relative. C'est pourquoi la condition dans la lemme suivante est que $X$ est quasi-compact, et non que $f$ est quasi-compact.

\begin{lemma}

\label{lemma-locally-quasi-finite-qc-source-universally-bounded}
Laissez $f : X \to Y$ être un morphisme de schémas. Supposons que
\begin{enumerate}
\item $f$ est localement quasi-fini, et
\item $X$ est quasi-compact.
\end{enumerate}
Alors $f$ a des fibres universellement nées.
\end{lemma}

\begin{proof}

Depuis $X$ est quasi-compact, il existe un recouvrement fini ouvert affine
$X = \bigcup_{i = 1, \ldots, n} U_i$ et ouverts affines $V_i \subset Y$,
$i = 1, \ldots, n$ de telle manière que $f(U_i) \subset V_i$. En raison de la nature locale de la « quasi-finité locale » (voir lemme \ref{lemma-quasi-finite-at-point-characterize} partie (4)) nous voyons que les morphismes $f|_{U_i} : U_i \to V_i$ sont localement morphismes quasi-finis des affines, donc quasi-finis, voir lemme \ref{lemma-quasi-finite-locally-quasi-compact}. Pour $y \in Y$ il est clair que $X_y = \bigcup_{y \in V_i} (U_i)_y$
c'est un recouvrement ouvert. Il suffit donc de prouver le lemme pour un morphisme quasi-fini d'affines (à savoir, si $n_i$ travaille pour le morphisme $f|_{U_i} : U_i \to V_i$Alors $\sum n_i$
travaille pour $f$).

\medskip\noindent
Supposons que $f : X \to Y$ est un morphisme quasi-fini d'affines. Par lemme \ref{lemma-quasi-finite-affine} nous pouvons trouver un diagramme $$
\xymatrix{
X \ar[rd]_f \ar[rr]_j & & Z \ar[ld]^\pi \\
& Y &
}
$$ avec $Z$ affine, $\pi$ finie et $j$ une immersion ouverte. Depuis $j$ a des fibres universellement créées (lemme \ref{lemma-immersion-universally-bounded}) cela nous réduit à montrer que $\pi$ a des fibres universellement liées (lemme \ref{lemma-composition-universally-bounded}).

\medskip\noindent
Cela nous réduit à un morphisme de la forme $\Spec(B) \to \Spec(A)$ où $A \to B$ est fini. Écrivons $B$ est généré par $x_1, \ldots, x_n$ sous forme de module $A$, de sorte que $$
A^{\oplus n} \longrightarrow B, \quad
(a_1, \ldots, a_n) \longmapsto \sum a_i x_i
$$ est une application surjective $A$-module. Pour n'importe quel $\mathfrak p \subset A$ prime, cela induit une application surjective $\kappa(\mathfrak p)$-vector espaces $$
\kappa(\mathfrak p)^{\oplus n} \longrightarrow B \otimes_A \kappa(\mathfrak p)
$$ En d'autres termes, l'entier $n$ fonctionne dans la définition d'un morphisme avec des fibres universelles.
\end{proof}

\begin{lemma}

\label{lemma-universally-bounded-permanence}
Considérez un diagramme commutatif des morphismes des schémas $$
\xymatrix{
X \ar[rd]_g \ar[rr]_f & & Y \ar[ld]^h \\
& Z &
}
$$ Si $g$ a des fibres universelles, et $f$ est surjectif et plat, alors $h$ a aussi des fibres universelles. Plus précisément, si $n$ limite le degré des fibres de $g$, alors $n$ limite également le degré des fibres de $h$.
\end{lemma}

\begin{proof}

Présumer $g$ a des fibres universellement nées, et $f$ est surjectif et plat. Dites le degré des fibres de $g$ est limitée par: $n \in \mathbf{N}$. Nous réclamons $n$ travaille aussi pour $h$. Laisse $z \in Z$. Considérez le morphisme des schémas $X_z \to Y_z$. Il est plat et surjectif. Par hypothèse $X_z$ est un schéma fini sur $\kappa(z)$, en particulier c'est le spectre d'un anneau artinien (par Algèbre, lemme \ref{algebra-lemma-finite-dimensional-algebra}). Par lemme \ref{lemma-Artinian-affine} le morphisme $X_z \to Y_z$ est affine en particulier quasi-compact. Il résulte de lemme \ref{lemma-fpqc-quotient-topology}
que $Y_z$ est un discrète fini comme cela tient pour $X_z$. D'où $Y_z$ est un schéma affine (c'est un exercice agréable; il suit aussi par exemple de Propriétés, lemme \ref{properties-lemma-maximal-points-affine}
appliqué à l'ensemble de tous les points de $Y_z$). Écrire $Y_z = \Spec(B)$ et $X_z = \Spec(A)$. Ensuite $A$ est fidèlement plat over $B$Donc... $B \subset A$. D'où $\dim_k(B) \leq \dim_k(A) \leq n$ comme souhaité.

\end{proof}





\section{Divers}
\label{section-misc}

\noindent
Résultats qui ne conviennent pas ailleurs.

\begin{lemma}

\label{lemma-factor-through-point}
Laissez $f : Y \to X$ être un morphisme de schémas. Laissez $x \in X$ être un point. Supposons que $Y$ est réduit et $f(Y)$ est défini-théoriquement contenu dans $\{x\}$. Puis $f$ facteurs à travers le morphisme canonique $x = \Spec(\kappa(x)) \to X$.
\end{lemma}

\begin{proof}
Démonstration omise. Suggestions: travailler affine localement on réduit à une algèbre commutative lemme. Compte tenu d'un homomorphisme d'anneaux $A \to B$ avec $B$ réduit de telle sorte qu'il existe un idéal unique premier $\mathfrak p \subset A$ dans l'image de $\Spec(B) \to \Spec(A)$, puis $A \to B$ facteurs par $\kappa(\mathfrak p)$. C'est un bel exercice.
\end{proof}

\begin{lemma}

\label{lemma-factor-through-set-of-points}
Laissez $f : Y \to X$ être un morphisme de schémas. Laissez $E \subset X$. Supposez que $X$ est localement noethérien, il n'y a pas de spécialisations non triviales parmi les éléments de $E$, $Y$ est réduit, et $f(Y) \subset E$. Puis $f$ facteurs par $\coprod_{x \in E} x \to X$.
\end{lemma}

\begin{proof}

Quand $E$ est un singleton qui suit de lemme \ref{lemma-factor-through-point}Si $E$ est fini, alors $E$
(avec la topologie induite de $X$) est un espace discret fini par notre hypothèse sur les spécialisations. C'est pourquoi ce cas se réduit au cas singleton. En général, il y a une réduction au cas où $X$ et $Y$ sont des plans affinés. Dites $f : Y \to X$ correspond à l'homomorphisme d'anneaux $\varphi : A \to B$. Dénotation $A' \subset B$ l'image de $\varphi$. Laisse $E' \subset \Spec(A') \subset \Spec(A)$ soit l'ensemble des primes minimales de $A'$. Par Algèbre, lemme \ref{algebra-lemma-injective-minimal-primes-in-image}
l'ensemble $E'$ est contenu dans l'image de
$\Spec(B) \to \Spec(A') \subset \Spec(A)$. Nous concluons que $E' \subset E$Depuis $A'$ C'est Noétherien que nous avons $E'$ est fini par Algèbre, lemme
\ref{algebra-lemma-Noetherian-irreducible-components}. Depuis tout autre point à l'image de $\Spec(B) \to \Spec(A)$ est une spécialisation d'un élément de $E'$ et dans $E$, nous concluons que l'image est contenue dans $E'$ (par notre hypothèse sur les spécialisations entre points de $E$). Ainsi nous réduisons au cas où
$E$ est fini que nous avons traité ci-dessus.

\end{proof}




\begin{multicols}{2}[\section{Autres chapitres}]
\noindent
Préliminaires
\begin{enumerate}
\item \hyperref[introduction-section-phantom]{Introduction}
\item \hyperref[conventions-section-phantom]{Conventions}
\item \hyperref[sets-section-phantom]{Théorie des ensembles}
\item \hyperref[categories-section-phantom]{Catégories}
\item \hyperref[topology-section-phantom]{Topologie}
\item \hyperref[sheaves-section-phantom]{Faisceaux sur les espaces}
\item \hyperref[sites-section-phantom]{Sites et faisceaux}
\item \hyperref[stacks-section-phantom]{Champs}
\item \hyperref[fields-section-phantom]{Corps}
\item \hyperref[algebra-section-phantom]{Algèbre commutative}
\item \hyperref[brauer-section-phantom]{Groupes de Brauer}
\item \hyperref[homology-section-phantom]{Algèbre homologique}
\item \hyperref[derived-section-phantom]{Catégories dérivées}
\item \hyperref[simplicial-section-phantom]{Méthodes simpliciales}
\item \hyperref[more-algebra-section-phantom]{Compléments d'algèbre}
\item \hyperref[smoothing-section-phantom]{Lissage des morphismes d'anneaux}
\item \hyperref[modules-section-phantom]{Faisceaux de modules}
\item \hyperref[sites-modules-section-phantom]{Modules sur les sites}
\item \hyperref[injectives-section-phantom]{Objets injectifs}
\item \hyperref[cohomology-section-phantom]{Cohomologie des faisceaux}
\item \hyperref[sites-cohomology-section-phantom]{Cohomologie sur les sites}
\item \hyperref[dga-section-phantom]{Algèbre différentielle graduée}
\item \hyperref[dpa-section-phantom]{Algèbre à puissances divisées}
\item \hyperref[sdga-section-phantom]{Faisceaux différentiels gradués}
\item \hyperref[hypercovering-section-phantom]{Hyperrecouvrements}
\end{enumerate}
Schémas
\begin{enumerate}
\setcounter{enumi}{25}
\item \hyperref[schemes-section-phantom]{Schémas}
\item \hyperref[constructions-section-phantom]{Constructions de schémas}
\item \hyperref[properties-section-phantom]{Propriétés des schémas}
\item \hyperref[morphisms-section-phantom]{Morphismes de schémas}
\item \hyperref[coherent-section-phantom]{Cohomologie des schémas}
\item \hyperref[divisors-section-phantom]{Diviseurs}
\item \hyperref[limits-section-phantom]{Limites de schémas}
\item \hyperref[varieties-section-phantom]{Variétés}
\item \hyperref[topologies-section-phantom]{Topologies sur les schémas}
\item \hyperref[descent-section-phantom]{Descente}
\item \hyperref[perfect-section-phantom]{Catégories dérivées de schémas}
\item \hyperref[more-morphisms-section-phantom]{Compléments sur les morphismes}
\item \hyperref[flat-section-phantom]{Compléments sur la platitude}
\item \hyperref[groupoids-section-phantom]{Schémas en groupoïdes}
\item \hyperref[more-groupoids-section-phantom]{Compléments sur les schémas en groupoïdes}
\item \hyperref[etale-section-phantom]{Morphismes étales de schémas}
\end{enumerate}
Thèmes de théorie des schémas
\begin{enumerate}
\setcounter{enumi}{41}
\item \hyperref[chow-section-phantom]{Homologie de Chow}
\item \hyperref[intersection-section-phantom]{Théorie de l'intersection}
\item \hyperref[pic-section-phantom]{Schémas de Picard des courbes}
\item \hyperref[weil-section-phantom]{Théories cohomologiques de Weil}
\item \hyperref[adequate-section-phantom]{Modules adéquats}
\item \hyperref[dualizing-section-phantom]{Complexes dualisants}
\item \hyperref[duality-section-phantom]{Dualité pour les schémas}
\item \hyperref[discriminant-section-phantom]{Discriminants et différentes}
\item \hyperref[derham-section-phantom]{Cohomologie de de Rham}
\item \hyperref[local-cohomology-section-phantom]{Cohomologie locale}
\item \hyperref[algebraization-section-phantom]{Géométrie algébrique et formelle}
\item \hyperref[curves-section-phantom]{Courbes algébriques}
\item \hyperref[resolve-section-phantom]{Résolution des surfaces}
\item \hyperref[models-section-phantom]{Réduction semi-stable}
\item \hyperref[functors-section-phantom]{Foncteurs et morphismes}
\item \hyperref[equiv-section-phantom]{Catégories dérivées de variétés}
\item \hyperref[pione-section-phantom]{Groupes fondamentaux de schémas}
\item \hyperref[etale-cohomology-section-phantom]{Cohomologie étale}
\item \hyperref[crystalline-section-phantom]{Cohomologie cristalline}
\item \hyperref[proetale-section-phantom]{Cohomologie pro-étale}
\item \hyperref[relative-cycles-section-phantom]{Cycles relatifs}
\item \hyperref[more-etale-section-phantom]{Compléments de cohomologie étale}
\item \hyperref[trace-section-phantom]{La formule des traces}
\end{enumerate}
Espaces algébriques
\begin{enumerate}
\setcounter{enumi}{64}
\item \hyperref[spaces-section-phantom]{Espaces algébriques}
\item \hyperref[spaces-properties-section-phantom]{Propriétés des espaces algébriques}
\item \hyperref[spaces-morphisms-section-phantom]{Morphismes d'espaces algébriques}
\item \hyperref[decent-spaces-section-phantom]{Espaces algébriques décents}
\item \hyperref[spaces-cohomology-section-phantom]{Cohomologie des espaces algébriques}
\item \hyperref[spaces-limits-section-phantom]{Limites d'espaces algébriques}
\item \hyperref[spaces-divisors-section-phantom]{Diviseurs sur les espaces algébriques}
\item \hyperref[spaces-over-fields-section-phantom]{Espaces algébriques sur des corps}
\item \hyperref[spaces-topologies-section-phantom]{Topologies sur les espaces algébriques}
\item \hyperref[spaces-descent-section-phantom]{Descente et espaces algébriques}
\item \hyperref[spaces-perfect-section-phantom]{Catégories dérivées d'espaces}
\item \hyperref[spaces-more-morphisms-section-phantom]{Compléments sur les morphismes d'espaces}
\item \hyperref[spaces-flat-section-phantom]{Platitude sur les espaces algébriques}
\item \hyperref[spaces-groupoids-section-phantom]{Groupoïdes dans les espaces algébriques}
\item \hyperref[spaces-more-groupoids-section-phantom]{Compléments sur les groupoïdes dans les espaces}
\item \hyperref[bootstrap-section-phantom]{Amorçage}
\item \hyperref[spaces-pushouts-section-phantom]{Sommes amalgamées d'espaces algébriques}
\end{enumerate}
Thèmes de géométrie
\begin{enumerate}
\setcounter{enumi}{81}
\item \hyperref[spaces-chow-section-phantom]{Groupes de Chow des espaces}
\item \hyperref[groupoids-quotients-section-phantom]{Quotients de groupoïdes}
\item \hyperref[spaces-more-cohomology-section-phantom]{Compléments sur la cohomologie des espaces}
\item \hyperref[spaces-simplicial-section-phantom]{Espaces simpliciaux}
\item \hyperref[spaces-duality-section-phantom]{Dualité pour les espaces}
\item \hyperref[formal-spaces-section-phantom]{Espaces algébriques formels}
\item \hyperref[restricted-section-phantom]{Algébrisation des espaces formels}
\item \hyperref[spaces-resolve-section-phantom]{Résolution des surfaces revisitée}
\end{enumerate}
Théorie des déformations
\begin{enumerate}
\setcounter{enumi}{89}
\item \hyperref[formal-defos-section-phantom]{Théorie formelle des déformations}
\item \hyperref[defos-section-phantom]{Théorie des déformations}
\item \hyperref[cotangent-section-phantom]{Le complexe cotangent}
\item \hyperref[examples-defos-section-phantom]{Problèmes de déformation}
\end{enumerate}
Champs algébriques
\begin{enumerate}
\setcounter{enumi}{93}
\item \hyperref[algebraic-section-phantom]{Champs algébriques}
\item \hyperref[examples-stacks-section-phantom]{Exemples de champs}
\item \hyperref[stacks-sheaves-section-phantom]{Faisceaux sur les champs algébriques}
\item \hyperref[criteria-section-phantom]{Critères de représentabilité}
\item \hyperref[artin-section-phantom]{Axiomes d'Artin}
\item \hyperref[quot-section-phantom]{Espaces de Quot et de Hilbert}
\item \hyperref[stacks-properties-section-phantom]{Propriétés des champs algébriques}
\item \hyperref[stacks-morphisms-section-phantom]{Morphismes de champs algébriques}
\item \hyperref[stacks-limits-section-phantom]{Limites de champs algébriques}
\item \hyperref[stacks-cohomology-section-phantom]{Cohomologie des champs algébriques}
\item \hyperref[stacks-perfect-section-phantom]{Catégories dérivées de champs}
\item \hyperref[stacks-introduction-section-phantom]{Introduction aux champs algébriques}
\item \hyperref[stacks-more-morphisms-section-phantom]{Compléments sur les morphismes de champs}
\item \hyperref[stacks-geometry-section-phantom]{La géométrie des champs}
\end{enumerate}
Thèmes de théorie des espaces de modules
\begin{enumerate}
\setcounter{enumi}{107}
\item \hyperref[moduli-section-phantom]{Champs de modules}
\item \hyperref[moduli-curves-section-phantom]{Espaces de modules de courbes}
\end{enumerate}
Divers
\begin{enumerate}
\setcounter{enumi}{109}
\item \hyperref[examples-section-phantom]{Exemples}
\item \hyperref[exercises-section-phantom]{Exercices}
\item \hyperref[guide-section-phantom]{Guide bibliographique}
\item \hyperref[desirables-section-phantom]{Desiderata}
\item \hyperref[coding-section-phantom]{Style de programmation}
\item \hyperref[obsolete-section-phantom]{Obsolète}
\item \hyperref[fdl-section-phantom]{Licence de documentation libre GNU}
\item \hyperref[index-section-phantom]{Index généré automatiquement}
\end{enumerate}
\end{multicols}


\bibliography{my}
\bibliographystyle{amsalpha}

\end{document}
